% =====================================================================
%  Harald Høffding, DEN STORE HUMOR.  En psykologisk Studie.  København og
%  Kristiania: Gyldendalske Boghandel, Nordisk Forlag, 1916.  178 pp.
%
%  TRANSCRIPTION -- GENERATED; DO NOT EDIT BY HAND.  Made in three steps, all rerunnable:
%    python3 ebook-method/den-store-humor/draft.py --emit --force
%      this template + the scan's text layer (the Royal Library's OCR, in KB's scan;
%      ~/bibliotek/Høffding, Harald/den-store-humor.pdf, see SCANS.tsv), then
%      patches.PATCHES (corrections read at the image that a word decision cannot express);
%    python3 ebook-method/collate/check.py den-store-humor --italics --apply
%      the decisions in ebook-method/den-store-humor/decisions.tsv, taken against an
%      independent tesseract reading of the images (rules.py writes the ones settled by
%      rule, each marked with its rule; the rest were read at the image), and the
%      italics decisions (--italics must be in the same call as --apply);
%    python3 ebook-method/den-store-humor/draft.py --post
%      patches.POST: corrections written against the collated text (the footnotes,
%      read at the image note by note; specks read as punctuation).
%  Run them chained (&&): each step exits without writing if a patch fails to match.
%  STATE (2026-10-04): collated, OPEN 0 (words and italics); all 108 notes read at
%  the image; pp. 55 and 151 spot-read in full; the translation pass read every
%  page (121 more faults fixed, 18 misprints marked).  See ACCURACY-LEDGER.tsv.
%  No e-book of this book exists.
%
%  CONVENTIONS.  "% --- p. N ---" + \opage{N} at the first word of printed page N
%  (a word broken over the turn is split with %).  Footnotes as printed, 1), 2), ...
%  numbered per page (\footnote[k]).  Italics \textit; letterspacing \emph.  A note
%  running over a page is set whole at its mark.  Running heads, folios and the
%  library's stamps are not transcribed.  Print governs; misprints are kept and
%  marked PRINTED AS IS.
% =====================================================================
\documentclass[12pt]{article}
\usepackage[utf8]{inputenc}
\usepackage[T1]{fontenc}
\usepackage{libertinus}
\usepackage{libertinust1math}
\usepackage[danish]{babel}
\usepackage{textalpha}            % Greek in the notes (α β γ δ ε; πρῶτος …)
\usepackage[protrusion=true,expansion=true]{microtype}
\usepackage{setspace}
\usepackage[margin=1.3in]{geometry}
\usepackage{fancyhdr}
\usepackage{hyperref}
\hypersetup{pdftitle={Den store Humor}, pdfauthor={Harald Høffding}, hidelinks}
\pagestyle{fancy}
\fancyhf{}
\fancyhead[LE,RO]{\thepage}
\fancyhead[C]{\textit{Harald Høffding: Den store Humor (1916)}}
\setstretch{1.3}
\usepackage{marginnote}
\usepackage{graphicx}
\renewcommand{\thefootnote}{\arabic{footnote})}
\newcommand{\tocline}[3]{\noindent\makebox[2.2em][r]{#1}\hspace{0.6em}#2\dotfill\ #3\par}
\newcommand{\entry}{\par\noindent\hangindent=1.5em }
\newcommand{\opage}[1]{\marginnote{\footnotesize\textit{[#1]}}}

\begin{document}

% --- front matter (PDF 8-12), read at the image 2026-10-03.
% p. [1], half-title
\begin{center}\vspace*{6em}
{\small HARALD HØFFDING}\\[0.5em]
{\large DEN STORE HUMOR}
\end{center}
\clearpage
% p. [2], the author's other works (in a dotted frame)
\begin{center}{\large HARALD HØFFDING}\end{center}
{\small\setlength{\parindent}{0pt}\setlength{\parskip}{0.2em}
\everypar{\hangindent=2em}
DEN ANTIKE OPFATTELSE AF MENNESKETS VILLIE. 1870.\par
PHILOSOPHIEN I TYSKLAND EFTER HEGEL. 1872.\par
DEN ENGELSKE PHILOSOPHI I VOR TID. 1874.\par
OM GRUNDLAGET FOR DEN HUMANE ETHIK. 1876.\par
PSYKOLOGI i Omrids paa Grundlag af Erfaring. (1882). Sjette paa ny gennemsete Udgave. 1911.\par
FORMEL LOGIK. Til Brug ved Forelæsninger. (1884). Sjette Udgave. 1913.\par
ETIK. En Fremstilling af de etiske Principer og deres Anvendelse paa de vigtigste Livsforhold. (1887). Fjerde, paa ny gennemsete og delvis ændrede Udgave. 1913.\par
PSYKOLOGISKE UNDERSØGELSER. 1889. (Videnskabernes Selskabs Skrifter).\par
ETISKE UNDERSØGELSER. Om Muligheden af en filosofisk Etik. --- Om Velfærdsprincipet. --- Forholdsloven i Etiken. 1891.\par
SØREN KIERKEGAARD SOM FILOSOF. 1892.\par
KONTINUITETEN I KANTS FILOSOFISKE UDVIKLINGSGANG. 1893. (Videnskabernes Selskabs Skrifter).\par
DEN NYERE FILOSOFIS HISTORIE. En Fremstilling af Filosofiens Historie fra Renaissancens Slutning til vore Dage. To Bind. (1894---95). Anden gennemsete Udgave. 1904.\par
JEAN JACQUES ROUSSEAU og hans Filosofi. (1896). Anden Udgave. 1912.\par
KORT OVERSIGT OVER DEN NYERE FILOSOFIS HISTORIE. (1898). Femte, paany gennemsete Udgave. 1913.\par
MINDRE ARBEJDER. Første Række. 1899.\par
DET PSYKOLOGISKE GRUNDLAG FOR LOGISKE DOMME. 1899. (Videnskabernes Selskabs Skrifter).\par
VED AARHUNDREDESKIFTET. Tale holdt paa Københavns Universitet. 1901.\par
RELIGIONSFILOSOFI (1901). Andet Oplag. 1906.\par
FILOSOFISKE PROBLEMER. 1902. (Universitetsprogram).\par
OM NOGLE RELIGIONSFILOSOFISKE ARBEJDER fra den nyeste Tid. 1903. (Universitetsprogram).\par
MODERNE FILOSOFER. 1904.\par
MINDRE ARBEJDER. Anden Række. 1905.\par
DANSKE FILOSOFER. 1909.\par
DEN MENNESKELIGE TANKE, dens Former og dens Opgaver. 1910.\par
RELIGION OG VIDENSKAB. 1910.\par
UDVALGTE STYKKER af dansk filosofisk Litteratur. Med Indledninger af \textit{Harald Høffding}. 1910.\par
PERSONLIGHETSPRINCIPEN I FILOSOFIN. Stockholm. 1911.\par
MINDRE ARBEJDER. Tredie Række. 1913.\par
HENRI BERGSON'S FILOSOFI. Karakteristik og Kritik. 1914.\par}
\clearpage
% p. [3], title
\begin{center}
HARALD HØFFDING\\[2em]
{\LARGE DEN STORE HUMOR}\\[2em]
{\small EN PSYKOLOGISK STUDIE}\\[14em]
% (a rule)
{\small G Y L D E N D A L S K E\ \ B O G H A N D E L\\
N O R D I S K\ \ F O R L A G\ -\ K J Ø B E N H A V N\\
O G\ \ K R I S T I A N I A\ -\ M D C C C C X V I}
\end{center}
\clearpage
% p. [4], verso.  The Royal Library's stamp (BIBLIOTHECA REGIA HAFNIENSIS) and a
% shelfmark in pencil on the title page are not transcribed.
\begin{center}\vspace*{\fill}{\scriptsize KJØBENHAVN --- FORLAGSTRYKKERIET}\end{center}
\clearpage
% p. [5], INDHOLD
\begin{center}INDHOLD\end{center}
\medskip
\hfill{\scriptsize Pagina}\par
\tocline{I.}{Humor som Totalfølelse}{7}
\tocline{II.}{Latter og Humor}{48}
\tocline{III.}{Ironi og Humor}{59}
\tocline{IV.}{Hovedformer af Humor}{75}
\tocline{V.}{Tragik og Humor}{88}
\tocline{VI.}{Forstaaelse og Humor}{112}
\tocline{VII.}{Handling og Humor}{124}
\tocline{VIII.}{Historiske Betingelser for Humor}{135}
\tocline{IX.}{Humor og Filosofi}{145}
\clearpage

% ===== TEXT (pp. 7--175), REGISTER (pp. 177--178) =====
% --- p. 7 ---
\setcounter{footnote}{0}%
\opage{7}\begin{center}I. HUMOR SOM TOTALFØLELSE\end{center}


1. Ingen sjælelig Tilstand er aldeles enkelt. Saa langt Iagttagelse
kan naa, med Bortseen fra Alt hvad der skyldes Eftertankens
Indflydelse, forefinder den stedse en større eller mindre Mangfoldighed
af Trang eller Drift, Sansning, Lyst og Ulyst, Forestilling,
kort sagt af, hvad der med ét Ord kan kaldes psykiske
Elementer, og den forefinder tillige en løsere eller fastere Sammenhæng
mellem disse Elementer. Enhver sjælelig Tilstand er en
Totalitet, indenfor hvilken en vis Kontinuitet gør sig gældende,
saa at en skarp Sondring mellem de forskellige Elementer altid
bliver mere eller mindre kunstig.

Ogsaa mellem de forskellige sjælelige Tilstande indbyrdes gør
der sig Kontinuitet gældende. Tilstandene gaa ofte umærkeligt over
i hverandre, og hvor Overgangen sker pludseligt, med et Ryk
eller et Slag, bevirker netop denne Pludselighed, at den følgende
Tilstand faar sin Beskaffenhed bestemt ved Modsætningen til den
foregaaende, saa at den ikke kan forstaas uden denne. Der er altid
noget Kunstigt i at holde de forskellige sjælelige Tilstande skarpt
ude fra hverandre.

Naar jeg alligevel som Indledning til den Undersøgelse, der her
skal foretages, skelner mellem Enkelttilstande og Totaltilstande,
da finder jeg Forskellen mellem dem deri, at hine ere bestemte
ved en enkelt Oplevelse, medens disse ere bestemte ved en Række
af Oplevelser. Saa fyldigt et Indhold en Enkelttilstand end kan
have, saa er dette Indhold --- Fornemmelser, Tanker, Drifter, Følelser
--- dog betinget ved det fælles Forhold til en enkelt Oplevelse.
Derimod have flere, maaske mange Oplevelser medvirket
ved Totaltilstandenes Tilblivelse. Der er naturligvis mange Over%
% --- p. 8 ---
\setcounter{footnote}{0}%
\opage{8}gangsformer mellem en enkelt Oplevelse og en Række af Oplevelser.
Aldeles enkelt er kun Oplevelsen af en samtidigt given Mangfoldighed,
f. Ex. Synet af en hvilende Gruppe af Mennesker eller
Ting. Naar det er en successiv Mangfoldighed, der opleves, som
naar man hører en Melodi eller følger Omridset af en Figur,
eller naar en historisk Begivenhed udfolder sig gennem en længere
Tid, bestaar Oplevelsen af en Række af enkelte Oplevelser. Det
kommer da an paa, om de enkelte Episoder i deres Forhold til
den, der oplever dem, fremtræde i indre Sammenhæng; kun da
vil en virkelig Totaltilstand kunne opstaa. Var der absolut ingen
Sammenhæng mellem Oplevelserne, vilde der, saavidt vi kunne
se, tilsidst ingensomhelst sjælelige Tilstande gives.

Allerede Enkelttilstandes Opstaaen frembyder Vanskeligheder
nok for den psykologiske Forsken; men Totaltilstande indeholde
ifølge Sagens Natur endnu større Problemer. Simple Enkelttilstande
kunne maaske fremkaldes experimentalt, og man behøver
i hvert Tilfælde ikke at udstrække Betragtningen ud over visse
forholdsvis simple og overskuelige Forhold. Man kan benytte Analogier
fra Fysiologien, ja maaske fra Mekaniken. Totaltilstande
ere derimod vanskeligere at analysere. De kræve større Kendskab
til Sjælelivets hele Udviklingshistorie hos vedkommende Menneske,
maaske tillige til den Tid og det Folk, det tilhører, ja
endog til Fortidens aandelige Udvikling. En historisk (sammenlignende,
sociologisk) Metode maa anvendes ved Siden af den iagttagende
og analyserende Metode. Intet Under, at Psykologerne
med Forkærlighed have holdt sig til Enkelttilstandene og ofte have
ment, at Totaltilstandene uden videre kunne forstaas ud fra dem,
ligesom mekaniske Helheder kunne forstaas, naar man blot kender
Delene.

Lad os, inden vi gaar videre, anføre nogle Exempler paa Enkelttilstande.


Midt i den Stemning, en smuk Foraarsmorgen har fremkaldt,
faar jeg pludseligt en sørgelig Efterretning. Alt samler sig nu for
mig om Budskabet; dettes Indhold lægger Beslag paa al Eftertanke;
dets Indgriben i Forhold, der ere mig dyrebare, tegner sig
for mig i store Træk; den lyse Stemning afløses af Nedslagenhed
% --- p. 9 ---
\setcounter{footnote}{0}%
\opage{9}og dump Hensynken; Alt, hvad der ellers gav Livet Værd, forsvinder
fra Bevidstheden. Forsøg paa at tænke bort, hvad der er
sket, eller finde Midler til at afværge dets Følger, hæmmes ved
Indsigten i det Uundgaaelige. Den organiske Livsfølelse, hele min
Tilstand faar Hæmningens, Forstemthedens Præg. Ingen Tanke
kan komme op, naar den ikke angaar det Skete. Dagens Arbejde
udøves tvangsmæssigt, uden indre Tilslutning. Efter et enkelt Øjebliks
Forglemmelse staar den mørke Virkelighed igen for mig i
sit hele Omfang. Og der rører sig maaske endog en Stræben efter
ret at fordybe sig i den mørke Virkelighed og holde alt Lys ude,
ja, der kan føles Samvittighedsnag over, at man overhovedet har
kunnet tænke paa andre Ting.

Et andet Exempel. En ny Ide gaar op for en Tænker. Han ser
Muligheden af en helt ny Retning i Forskningen. Med Intuitionens
Styrke, der er saa meget des større, som der forud er gaaet en
Periode af Tvivl, Mismod og uheldige Forsøg, aabner der sig nu
Udblik over nye Egne, og spredte Veje mødes for Blikket i et
fælles Maal. Der er foregaaet en pludselig Koncentration; en syntetisk
Intuition\footnote[1]{Se min Afhandling om \textit{Begrebet Intuition} (Vid. Selsk. Oversigt 1914) p. 81. (Ogsaa \textit{Henri Bergson’s Filosofi} p. 24 f.).} har dannet sig. Ikke blot Tankelivet, hele det
sjælelige Liv kommer i stærk Bevægelse. Ingen anden Interesse
end den intellektuelle kan gøre sig gældende, ingen anden Stræben;
og dog er den højeste Begejstring tilstede. En Oplevelse af
denne Art har Réné Descartes havt den 10. November 1619, medens
han laa i Vinterkvarter i Bayern, da det efter lang Tids Søgen
gik op for ham, hvilken Vej han skulde følge i sin Forsken, idet
Grundtrækkene i hans filosofiske og matematiske Fremtidsarbejde
aabenbarede sig for ham.\footnote[2]{\textit{Den nyere Filosofis Historie.}\textsuperscript{2} I p. 205.} Da Entusiasmen, den pludselige Enkelttilstand,
havde fordelt sig, kom de lange Tiders Arbejde, ja
først atten Aar efter saa hans grundlæggende Værk Lyset. ---
Kunstens Historie frembyder lignende Exempler. Bekendt er saaledes
Mozart’s Beskrivelse af, hvorledes Motiver, han ofte uvilkaarligt
havde nynnet for sig selv, samlede sig til et Tonehele, der
paa én Gang stod for ham (gleich Alles zusammen), saa at han
% --- p. 10 ---
\setcounter{footnote}{0}%
\opage{10}kunde overskue det, som om det var et Maleri eller et Skulpturarbejde.
Rousseau’s Undfangelse af Kulturproblemet frembyder
lignende Træk.

Et tredie Exempel. Efter lang Tids Overvejelse har jeg taget en
Beslutning i en vigtig Sag. Mit hele Væsen, det vil sige, al Tanke,
Følelse og Stræben, samler sig paa ét Punkt, i ét Ønske, som dog
er mere end et Ønske, idet jeg allerede ser mig selv ifærd med at
udføre Handlingen, og idet dens Virkninger, som jeg vedkender
mig, aftegne sig for mig i mere eller mindre klare Omrids. Jeg
mener at mærke, at Noget fra mit Indre, maaske fra mit Inderste,
strømmer ud i den forestillede Fremtid, i hvilken Handlingen skal
foregaa. Hele mit Selv, saaledes som det nu engang er, faar Luft
i Afgørelsen, kommer gennem den til at aande ud.\footnote[1]{Smlgn. om Beslutningens Psykologi min \textit{Psykologi i Omrids paa Grundlag af Erfaring} VII B, 2, (Sjette Udgave p. 409—411).}

Et sidste Exempel henter jeg fra Ludvig Feilbergs rige Samling
af Enkelttilstande. Feilbergs Styrke og Svaghed er, at han er altfor
optagen af Enkelttilstande; men han beskriver dem med Mesterskab.
Totaltilstande er han tilbøjelig til at se ned paa som mindre
ægte. De lade sig heller ikke saa let gribe i Flugten, og kunne
ikke beskrives saa anskueligt, som f. Ex. følgende lille Oplevelse.
„En lille, frisk affalden grøn Kastanie laa paa Stenbroen under
Holbergs Træ ved Siden af nogle affaldne Smaagrene, ældre affaldne
Kastanier og halvvisne Kastanieblomster, der af Vinden
var fejede sammen. Det var overraskende at se, hvor den var
smuk i denne Omgivelse; det var den slet ikke, medens den sad i
Træet. Tvertimod vakte den lille piggede Ting, der mindede om,
at Blomsternes Tid var forbi, dengang nærmest Mishag. Her var
den ganske ejendommeligt hyggelig. Man kunde ikke lade være
at lægge Mærke til den.“\footnote[2]{\textit{Ludvig Feilberg: Samlede Skrifter.} I p. 172.}

2. Disse Exempler paa Enkelttilstande ere med Villie ikke
valgte blandt de simpleste. De have det Særkende, at et mangfoldigt
Indhold samler sig i hvert af dem om en forholdsvis enkelt
ydre eller indre Oplevelse. Selv om vi havde valgt ganske simple
Enkelttilstande, vilde det dog have vist sig, at der stedse forud%
% --- p. 11 ---
\setcounter{footnote}{0}%
\opage{11}sættes Betingelser, der skyldes Menneskets tidligere Erfaringer.
Tusinder have set grønne Kastanier ligge i Fiolstræde uden at
opleve, hvad Feilberg oplevede, og den Prisopgave, der lod Kulturproblemet
opgaa for Rousseau, er bleven set af Mange, uden at
virke, som den virkede paa ham. Ogsaa Enkelttilstande kunne
altsaa have en historisk Karakter. Om Individer ved Forsøg over
den fysiologiske Tid reagere sensorisk eller motorisk, vil, i hvert
Tilfælde i Begyndelsen, bero paa deres Øvelser og Vaner. I Smaat
som i Stort viser det sig, at vor Nutid ikke kan forstaas uden vor
Fortid. Men i ganske særlig Grad gælder dette om Totaltilstande,
som bestemmes ved en lang Række af forskellige Oplevelser. Enhver
af disse Oplevelser har fremkaldt sin Enkelttilstand, og det
er gennem disse Enkelttilstande, at en Totaltilstand kan blive til.
Forudsætningen er, at de ikke forsvinde sporløst, men efterlade
sig Dispositioner, som kunne medvirke ved senere Tilstandes Opstaaen.
Al Genkendelse og al Genkaldelse saavel som al Eftertanke
beror herpaa. Ikke mindst virker denne Ophoben paa Følelseslivets
Omraade. Gennem Skiften af Glæde og Sorg, Haab
og Frygt, Begejstring og Nedstemthed kunne nye, varige Grundstemninger
arbejde sig frem. Under Brydning med nedarvede Dispositioner
udformes der gennem Oplevelsernes Række, hvad vi i
en mere uvilkaarlig, naturbunden Form kalde Temperament, og
ved mer eller mindre bevidst Stræben og Handlen kalde Karakter.
Jeg har i min Psykologi (V B., 5) indført Udtrykket realt Jeg for
at betegne, hvad der saaledes dels er Forudsætning for, dels Virkning
af Oplevelsernes Indflydelse paa de sjælelige Tilstande.

For at forebygge en Misforstaaelse, der selv hos fortræffelige
Psykologer er altfor hyppig, maa det udtrykkeligt fremhæves, at
dette reale Jeg ikke behøver at være Genstand for Bevidsthed hos
Individet selv. Det behøver ikke at afføde nogen „Jeg-Forestilling“.
At Ordet „jeg“ forholdsvis tidligt bruges, betyder Intet; det er
fra først af et blot Navn, et Navn, om hvilket Barnet senere til sin
store Forundring opdager, at det er et Navn, alle Mennesker med
samme Ret gøre Krav paa at have. Der kan være en herskende
Interesse og Stræben --- og dermed et realt Jeg --- uden at der
dannes nogen Forestilling derom, --- uden at det bliver Objekt.
Det reale Jeg kan undergaa endog meget indgribende Ændringer,
% --- p. 12 ---
\setcounter{footnote}{0}%
\opage{12}uden at Individet bliver sig det bevidst. Ændringerne lægge sig
umiddelbart og uvilkaarligt for Dagen i dets Færd og bliver maaske
derigennem aabenbare for Andre, skønt Individet selv ikke
opdager dem.

Ludvig Feilberg gik endog saa vidt, at han helt vilde afskaffe
Ord som „Selv“ eller „Personlighed“, uden forsaavidt de kunde
bruges i dadlende Betydning, nemlig om indbildsk Selvbevidsthed,
raffineret Selvoptagethed, der paa fordringsfuld Maade kredser
om sin egen lille Verden istedetfor at glemme sig selv i „Ligeløb“.
Han var saa optagen af de Blomster, han fandt under sin psykologiske
Botaniseren, at han ikke saa, hvorledes der i og gennem
Enkelttilstandenes Række kan lægge sig en sjælelig Naturorden
for Dagen, som netop danner Grundlaget for, hvad der i Psykologien
menes med Ord som Jeg, Selv eller Personlighed. --- I sine
senere Skrifter var han dog bleven klar over de Misforstaaelser,
som hans Polemik mod „Personlighed“ kunde give Anledning
til.

Det er rigtigt nok, at udtrykkelig Bevidsthed om, hvad der rører
sig hos os, kan være betænkelig, forsaavidt den fører ud over det
Umiddelbare og Uvilkaarlige, som er den frugtbare Naturgrund
for aandeligt Liv. Men ligesom Instinktet kan føre til, at Dyret
gaar i Fælden, saaledes kan Livets uvilkaarlige Udfoldelse meget
let blive af det Onde, og det kan være nødvendigt for Selvopholdelsen,
at vi besinde os paa, hvilken Vej der kan føre os til vort
Maal. Det vil da ofte blive nødvendigt at skyde den uvilkaarlige
Trang tilside. Det er især paa saadanne Vendepunkter, vi blive os
vort reale Jeg bevidst, idet vi rette Opmærksomheden paa vort
indre Væsen, mod vor centrale Stræben, og det kan da betegne
et stort Fremskridt, at Bevidstheden gør sig selv til Genstand.

Med Hensyn til det reale Jeg ere Totaltilstande mere afgørende
end Enkelttilstande. Disse sidste kunne forsvinde uden indgribende
Virkninger, uden hverken i deres Opstaaen eller gennem deres
Indflydelse at staa i nøje Sammenhæng med det i os, der er det
egentlige Centrale. Psykologien antager vel, at alle Tilstande og
Oplevelser afsætte Virkninger og give Bidrag til den indre Naturordens
Udvikling, og altsaa direkte eller indirekte, positivt eller
% --- p. 13 ---
\setcounter{footnote}{0}%
\opage{13}negativt virke til at frembringe Totaltilstande.\footnote[1]{Den saakaldte psykoanalytiske Skole i Sindssygevidenskaben (Freud og hans Disciple) lægger i den nyeste Tid stor Vægt paa, at Oplevelser af indgribende Betydning stadigt vedblive at øve deres Indflydelse paa Sindet og derved kunne blive Aarsag til Sindssygdomme. Under Hypnose eller i Drømme kan Patienten komme til at røbe saadanne Oplevelser, medens han maaske ellers uvilkaarligt eller vilkaarligt søger at trænge dem tilbage. Især en vis Gruppe af Sindssygdomme (de „hysteriske“) forklares som Virkninger af sjælelige Saar (som traumatiske Nevroser), og det bliver Lægens Opgave at opdage disse Saar og at helbrede dem ved at drage dem frem for den klare Bevidsthed og gøre dem til Genstand for Forstaaelse.} Men ved denne
Antagelse stilles der egentligt kun et stort Ideal for psykologisk
Forsken, et Ideal, der stadigt maa haves for Øje, men hvis Virkeliggørelse
stedse vil være ufuldstændig. Det vil desuden ikke altid
være en eneste Totaltilstand, der kommer til at raade. Flere Totaltilstande
kunne afløse hverandre i samme Menneskes Liv, maaske
endog i den Grad, at selve Begrebet realt Jeg i det enkelte Tilfælde
kan blive problematisk. Hos meget udprægede Personligheder
vil der dog være en Totaltilstand, der giver alt Andet i Sjælelivet
dets Ejendommelighed. Den behøver ikke uafbrudt at være
tilstede; men den udøver stadigt sine Eftervirkninger og spiller i
hvert Tilfælde en indirekte Rolle. Og i de for Personligheden afgørende
Øjeblikke vil det være den, der fører Ordet. Det er til
den, Personligheden svinger tilbage, naar det gælder Samling,
Koncentration. Den udtrykker det indre Menneske, hvad enten
dette er tilgængeligt for Andre eller ikke. Hos nogle Naturer, som
man har kaldt de „én Gang Fødte“ („de Enfødte“), vil hele Livet
være baaret af samme Totaltilstand, selv om der opstaar Stemninger
og Krusninger, som ikke helt igennem betinges ved den,
og selv om der, naar der skal handles under bestemte Forhold,
kan opstaa Sindsbevægelser, der ikke synes at have Noget at gøre
med den. De Formaal, et Menneske stiller sig, bestemmes gennemgaaende
ved den hos ham herskende Totaltilstand, det raadende
Sindelag. Men naar der skal arbejdes og kæmpes for disse
Formaal, kan der paa Grund af Hindringer, der frembyde sig, eller
Midler, der skulle skaffes, udløses Tanke- og Sindsbevægelser,
% --- p. 14 ---
\setcounter{footnote}{0}%
\opage{14}som i og for sig ikke høre Totaltilstanden til. Menneskekærligheden
kan saaledes føre til med Mod og Harme at bekæmpe Modstand;
Fædrelandskærlighed kan udløse krigerske Stemninger;
intellektuel eller kunstnerisk Begejstring kan føre til Bitterhed
overfor aandelig Sløvhed og Brutalitet. Den berømte Kriminalist
Anselm v. Feuerbach har særligt lagt Vægt paa denne tilsyneladende
Modsætning mellem Sindelag og Sindsbevægelse som
nødvendig til Forstaaelse af Forbryderkarakterer. (Se Psykologi
VI E, 5).

Hos de „to Gange Fødte“ („de Tvefødte“) erhverves den raadende
Totaltilstand gennem et Brud eller en Krise, hvorved forud
tilstedeværende Tendenser trænges tilbage. Kontinuiteten er her
ikke saa tydeligt fremtrædende som hos de Enfødte, skønt ofte
en dybere gaaende Undersøgelse kan paavise, at det nye og det
gamle Sindelag have et fælles Præg trods den helt forskellige
Retning, i hvilken Sjælelivets Udvikling foregaar indenfor dem.
Medens de Enfødtes Udvikling bestaar i en Udfoldelse, betinges de
Tvefødtes Sjæleliv ved Overvindelse af indre Disharmonier.\footnote[1]{De to Benævnelser skrive sig fra Francis Newman (Kardinalens humanistiske Broder). De to Typer beskrives, under forskellige Navne, ofte i den nyere Religionspsykologi. Smlgn. min \textit{Religionsfilosofi}\textsuperscript{2} p. 258 -261 og \textit{Mindre Arbejder.} Anden Række. p. 124 f.; 134 f.}

Der er allerede i det Foregaaende antydet Exempler paa Totaltilstande,
og jeg skal blot tilføje nogle faa Exempler endnu. Ærefrygt,
Vemod, Resignation, Trofasthed, Erkendelsesglæde ere Totaltilstande,
naar de forudsætte en Række Livserfaringer, Oplevelser
af mere eller mindre sammensat og intensiv Art. De ere da forskellige
fra momentan Opblussen og fra instinktiv Stræben, fra Beslutninger,
der fremkaldes ved forbigaaende Situationer, saavel
som fra Stemninger og Tilskyndelser, der forsvinde saa hurtigt
som de komme. Men de kunne under deres Brydning med de
praktiske Forhold tage saadanne mere elementære psykiske Funktioner
i deres Tjeneste. Overhovedet er det raadende Totaltilstande,
der ligge til Grund ved ethvert Forsøg paa Karakteristik
af Personligheder, ved enhver Angivelse af de Egenskaber, der
betragtes som ejendommelige for et Menneske.

Allerede hos Platon findes en træffende Skildring af, hvorledes
% --- p. 15 ---
\setcounter{footnote}{0}%
% CHECK p. 15: notes paired with the marks in order (the layer misread a number)
\opage{15}Totaltilstanden bliver til. Han sammenligner i „Staten“s fjerde
Bog (p. 429 D E) de forskellige Oplevelser med de forskellige
Lag af Farver, som maa anvendes paa et Stof, før den bestemte
Farve, man ønsker, at Stoffet skal have, kan anvendes. „Du véd
dog, at naar Farverne ville farve Uld purpurrød, udvælge de først
af de mange Farver én Art, nemlig den hvide, og derpaa bearbejde
de den ved mangfoldige Tilberedelser, forat den paa bedste
Maade kan optage Purpurfarven, og da først farve de. Hvad der
er farvet paa denne Maade er ægte, og ingen Vaskning kan berøve
det sin Farve!“ Den Egenskab eller Totaltilstand, hvis Betingelse
Platon vil angive, er Tapperhed, ikke i Betydning af uvilkaarligt
Mod, men i Betydning af Energi til under alle Forhold at
fastholde Overbevisningen om, hvad der med Rette er at frygte,
og hvad ikke. Det er Karakterfasthedens eller Troskabens Totaltilstand,
han vil beskrive. Han vil anvise den Vej, Opdragelsen
maa gaa for at kunne føre til en saadan Tilstand. Saa systematisk
som han vil gaa frem ved Forberedelsen af de Karakteregenskaber,
som han vil fremelske, gaar det virkelige Liv naturligvis ikke. Til
Gengæld er det en rigere Mangfoldighed af Karakteregenskaber,
Livets Gang kan føre til hos de forskellige Mennesker. Skæbnen
„grunder“ os med en Række af Farvestrøg, før den Totaltilstand
bliver til, som betinger vor Karakters ejendommelige Farvetone.

3. Ad forskellige Veje kunne Oplevelser komme i Berøring
med hverandre, saa at de kunne indgaa Forbindelser og der under
disses Indflydelse kan dannes nye Totaltilstande indenfor Sjælelivet.


Man plejer at skelne mellem tre Arter af Elementer, Sider eller
Beskaffenheder, som Eftertanken kan skelne i enhver sjælelig
Tilstand, nemlig Erkendelse (Fornemmelse og Forestilling), Følelse
af Lyst og Ulyst, Stræben eller Villen. Benytte vi denne Tredeling,
som der ingen Grund er til at forlade\footnote[1]{Se min Afhandling om \textit{Begrebet Villie} (Mindre Arbejder. Tredie Række).}, vil der være tre
Grupper af Love eller Tendenser, som kunne antages at spille
en Rolle ved Totaltilstandes Opstaaen. Om de ogsaa gøre sig gæl%
% --- p. 16 ---
\setcounter{footnote}{0}%
\opage{16}gende ved de allersimpleste Enkelttilstandes Opstaaen, er et meget
vanskeligt Spørgsmaal. Ti vi tør ikke uden videre gaa ud fra, at
de Inddelinger eller Distinktioner, Eftertanken kan og maa foretage,
naar den gør Sjælelivet til Genstand for Undersøgelse, ogsaa
kunne gælde for Sjælelivets primitive Former. Psykologisk Analyse
danner ikke nogen fuldstændig Analogi til kemisk Analyse. Et
sammensat Stof kan opløses i Dele, der kunne bestaa hver for
sig, uafhængigt af Sammensætningen. Men hvad vi kunne skelne
i de faktisk foreliggende psykiske Tilstande er ikke Dele, der
kunne have Existens afset fra den Helhed, i hvilken vi mene at
opdage dem. Der existerer ingen ren Erkenden, Følen eller Villen,
stedse kun Komplexer af disse psykiske Elementer. Og disse
Komplexer vise sig at være inderligere sammenføjede, jo længere
vi gaa tilbage i Sjælelivets Historie. Saavidt vi kende denne, for
det enkelte Menneskes og for Slægtens Vedkommende, er det et
Hovedtræk, at de senere Stadier frembyde udprægede Forskelligheder,
der ikke fremtræde paa tidligere Stadier, og som vi ikke
have Ret til at forudsætte der. Hin Tredeling gælder da kun for
Sjælelivets senere Stadier, hos den Voxne i Modsætning til Barnet,
hos Mennesket i Modsætning til Dyret, hos civiliserede Mennesker
i Modsætning til Barbarer og Vilde. En skematisk og deduktiv
Behandling af Psykologien, der ud fra de Elementer, der
kunne skelnes indenfor det højere udviklede Sjæleliv, vil kombinere
de Tilstande, som fremtræde med en Helhedskarakter, fører
derfor til illusoriske Resultater.\footnote[1]{Se nærmere herom min \textit{Psykologi} Kap. 4 og \textit{Den menneskelige Tanke} (1909) p. 9—19. --- Paa Grundlag af sammenlignende Psykologi har \textit{Edward Thorndike} energisk indskærpet denne Sandhed. „Den gamle Opfattelse af det menneskelige Bevidsthedsliv,“ siger han \textit{(Animal Inintelligence.} New York 1911. p. 154), „er den, at det er bygget op af elementære Fornemmelser, --- at der først opstaa meget smaa Bevidsthedsfragmenter (bits of consciousness), og at disse saa gradvis bygges op til komplexe Væv.“ I Modsætning dertil hævder han, at „Udviklingen ikke sker fra det Smaa og Simple til det Store og Komplexe, men fra umiddelbar Forbindelse til middelbar Forbindelse, indenfor hvilken en Hob isolerede Elementer optræde, --- fra ren Erfaring eller udifferentierede Tilstande til Opstaaen af Forskelligheder paa d. e. Side, Generalisationer og Abstraktioner paa d. a. Side.“} % PRINTED AS IS: „In-intelligence“ (the syllable doubled over the line break)

% --- p. 17 ---
\setcounter{footnote}{0}%
\opage{17}For Erkendelseselementernes Vedkommende er Forestillingernes
Association det vigtigste Fænomen. De fleste Psykologer antage,
at der til Grund for al Forestillingsforbindelse ligger en Trang
eller Tendens til at udfylde Forestillinger, der optræde forholdsvis
isolerede, og indlemme dem i en Helhed. Hvad der forhen har
været erfaret i Sammenhæng, men nu kun erfares isoleret, søges
udfyldt til den tidligere Helhed, og denne kan yderligere udfyldes
og udvides, naar nye Erfaringer vise Muligheden af endnu mere
omfattende Helheder. Der kan desuden spores en Tendens til at
fortsætte en begyndt Forestillingsrække efter samme Regel, som
hidtil har ført fra Led til Led, selv om der ikke foreligge umiddelbare
Erfaringer om flere Led end de hidtil givne. Saadanne Overgange
fra givne Forestillinger til beslægtede Forestillinger ville
især indtræde, naar Sindet er livligt, og naar der ved de givne
Forestillinger er udløst mere Energi, end der kan forbruges
til at fastholde dem; da bevæger Sindet sig kredsende uden
om dem.

Naar saa igen forskellige Forestillingsrækker sammenholdes,
kan der opstaa nye Forbindelser eller Helhedsdannelser derved,
at et Led i den ene Række viser sig at kunne indsættes (substitueres)
istedetfor et Led i den anden Række. Ogsaa en enkelt
Række kan ændres, saa en ny Helhed dannes, nemlig naar Led
indenfor Rækken kunne skydes ud saaledes, at Rækkens Lov vedbliver
at gælde for de resterende Led. Ad disse Veje opstaa Slutninger,
der ere logiske Helhedsdannelser.\footnote[1]{Se nærmere herom \textit{Den menneskelige Tanke.} p. 168— 173.}

Under saadanne Forestillingsprocessers Forløb forholde Følelseselementerne
sig paa forskellig Vis.\footnote[2]{Smlgn. med Hensyn til det Følgende Afsnittet om Følelsen (Kap. VI) i min Psykologi.} Ofte drages de gennem Forestillingsforbindelser
over fra deres oprindelige Aarsager eller Genstande
til andre. Derved bliver Værdiforskydning mulig, da det,
der først havde Værdi blot som Middel, kan faa Værdi som Formaal,
eller det, der stod som Formaal, for Fremtiden kan komme til
at staa som Middel for et fjernere liggende Formaal. Men det er
ogsaa muligt, at Følelseselementerne kunne hindre eller hæmme
Forestillingsforbindelser, der rent objektivt eller intellektuelt set
% --- p. 18 ---
\setcounter{footnote}{0}%
\opage{18}vilde ligge nær. En Forestilling, der én Gang er knyttet til stærk
Følelse, isoleres ved denne, og hvis Association overhovedet er
mulig, ville kun saadanne Forestillinger have Mulighed for at gøre
sig gældende, som umiddelbart eller middelbart begunstige Fastholdelsen
af den givne Forestilling. Desuden kan der utvivlsomt
øves en umiddelbar Indflydelse af Følelseselementer, som ere
knyttede til en Forestillingskreds, paa Følelseselementer, der ere
knyttede til en anden Forestillingskreds. Ved Samtidighed eller
umiddelbar Rækkefølge af to Forestillingskredse, kunne ikke blot
selve Forestillingerne, men ogsaa de til dem knyttede Følelseselementer
komme i umiddelbar Berøring. Og her frembyder sig
da flere Muligheder. Det kan være, at Modsætningen er saa stærk,
at der opstaar Spænding og maaske endog trues med Opløsning
af Bevidsthedens Enhed. Men det kan ogsaa være, at den ene
Stemning netop baner Vej for den anden ved at udtømme den
Kraft, der staar til Raadighed i en vis Retning, saa at den modsatte
Retning begunstiges, saalænge der overhovedet er Kraft tilovers.
I saa Fald begunstiger den ene Stemning den anden ved
Kontrastvirkning. Det kan fremdeles være, at Følelsen faar en saa
dyb Rod i Sindet, at den ikke ophører, selv om den Anledning, der
har fremkaldt den, hører op, men breder sig over hele Tilstanden
og fortrænger de til andre Oplevelser knyttede Følelser, om saadanne
overhovedet have kunnet gøre sig gældende.

Men foruden disse Muligheder er der endnu en Mulighed, som
er af særlig Interesse for den følgende Undersøgelse. Følelseselementer,
der ere knyttede til forskellige Oplevelser, kunne indgaa
Forbindelser med hverandre, saa at der opstaa Følelseskomplexer.
Erfaring viser, hvad der senere i en speciel Sammenhæng
skal nærmere udvikles, at der er to Hovedformer for saadanne
Følelsestotaliteter eller Totalfølelser. Forbindelsen af de forskellige
Følelseselementer kan have Karakteren af en Sammensmeltning,
idet der formedelst Forbindelsen opstaar en ny Følelse med
en fra de forudgaaende Følelser forskellig Kvalitet. Kun ved at
udforske den nye Følelses Udviklingshistorie kan man opdage,
hvilke Elementer der have medvirket ved denne Dannelse; for umiddelbar Iagttagelse staar den som simpel og udelt. En anden
Art af Totalfølelse har Karakteren af Organisation, idet de for%
% --- p. 19 ---
\setcounter{footnote}{0}%
\opage{19}skellige Følelseselementers Selvstændighed ikke ophæves, men
de ordne sig indbyrdes efter deres Forhold til og Betydning for
en herskende Interesse, en Grundværdi, der kan finde Udtryk i
en ledende Tanke.

Der findes analoge Helhedsformer indenfor de andre Sider af
Sjælelivet. En Analogi til Sammensmeltning frembyder indenfor
Erkendelsens Psykologi f. Ex. den Maade, hvorpaa der til en
Række af meget hurtigt paa hverandre følgende Sanseindtryk kun
svarer en eneste Fornemmelse, medens der ved langsommere
Rækkefølge vilde opstaa en Række særskilte Sansefornemmelser.
Ved umiddelbar Genkendelse sker der en Sammensmeltning af
den øjeblikkelige Fornemmelse med Eftervirkningen af tidligere
Fornemmelser af samme Art. Forestillinger, der hyppigt ere fremkomne
i umiddelbar Sammenhæng, kunne komme til at danne en
eneste Forestilling; saaledes ved Overgangen fra diskursiv Undersøgelse
til intuitiv Viden. En Analogi til Organisationen af en
Gruppe af Følelser danner Sanseanskuelse, i hvilken en Mangfoldighed
af Fornemmelsesindhold ere ordnede i Tid og Rum, og det
Samme gælder om Erindrings- og Fantasibilleder. Paa Villieslivets
Omraade have vi en Analogi til Sammensmeltningen i Instinkt
og Drift med den umiddelbare Koncentration, og en Analogi
til Organisationen i den Maade, paa hvilken Tanken om et
Formaal, der skal virkeliggøres, bestemmer Rangordningen mellem
de Midler, der skulle anvendes, eller de Handlinger, der
skulle udføres; de høre alle med til Beslutningen, omsluttes af
denne og udgøre en leddelt Helhed indenfor den.

Alle Helhedsformer paa Erkendelsens og Følelsens Omraade
ville tilsidst være betingede og bestemte ved den Stræben efter
Bestaaen og Udvikling, der gør sig gældende hos ethvert levende
Væsen, altsaa ved en Villen i dette Ords videste Betydning. Hvad
der skal leve, maa optræde som en mere eller mindre sluttet Helhed.
Lyst og Ulyst ere tilsidst kun Symptomer paa, om denne Stræben
begunstiges eller hæmmes; Sansefornemmelser tjene til Udløsning
af Reaktioner mod Omverdenen; og Forestillinger og Forestillingsforbindelser
tjene til at fastholde og udforme Formaal og
Midler, hvor uvilkaarlig Trang og Evne ikke slaa til. Spinoza beholder
Ret i sin Paastand, at vi ikke tilstræbe, ville eller ønske
% --- p. 20 ---
\setcounter{footnote}{0}%
\opage{20}Noget, fordi vi anse det for godt, men vi anse det for godt, fordi
vi tilstræbe, ville og ønske det.

Det er da intet Tilfælde, at der danner sig Totaltilstande. Hvor
langt Helhedsdannelsen skrider frem hos de enkelte Mennesker,
og hos det enkelte Menneske indenfor en enkelt Livsperiode, vil
bero paa de specielle indre og ydre Forhold. Fuldt totaliseret bliver
Sjælelivet neppe nogensinde. Dels vil der være Enkelttilstande,
som staa forholdsvis isolerede, dels ville forskellige Totaltilstande
kunne brydes mod hverandre. Derfor er der et psykisk Arbejde
at gøre, saalænge Livet varer.

4. Der har vel til alle Tider kunnet danne sig Totaltilstande i
den menneskelige Bevidsthed. Men de ville lettere optræde, naar
Oplevelserne blive rigere og mangfoldigere. Derfor ville de være
mere fremtrædende i nyere Tid end i ældre Tider, og de ville
være mere udprægede, frembyde flere Nuancer og flere individuelle
Forskelligheder end i tidligere Tider, hvor Livet førtes indenfor
en snevrere Horizont, hvor selve Verdensbilledet var begrænset
og overskueligt, og hvor heller ikke den historiske Erfaring
aabnede større Vidder og talrige Muligheder for Blikket. Imidlertid
maa man vel vogte sig for at gøre Modsætningen i denne Henseende
større, end berettiget er. Hvad der, naar Forholdene medføre
Modsætninger og Forskelligheder, gør sig gældende i fuld
Tydelighed, kan have været tilstede i mere udifferentieret Form
paa tidligere Trin, ligesom man i Fosteret kan se Antydninger af
de senere fuldt udviklede Organer. Hvis Tendensen til Dannelse af
Totaltilstande har sin Grund i selve Sjælelivets Natur og Vilkaar,
maa den antages at være tilstede paa alle Sjælelivets Udviklingstrin.
Det vilde være en vigtig Opgave for sammenlignende psykologisk
Forsken at undersøge dette nærmere.

Herved maa det nu ikke overses, at Totaltilstandens Optræden
ikke nødvendigvis ledsages af Evne til at beskrive og analysere
dem. Komplexe Tilstande karakteriseres af en ubehjælpsom Psykologi
ofte ved et enkelt Træk, der menes at indbefatte det Hele
i sig. Og da nu den intellektuelle Side af Sjælelivet ligger klarest
for Dagen, vil den primitive Psykologi tale mest om Tanke,
Fornuft og Viden, som da for den ikke blot betyde en enkelt Side
% --- p. 21 ---
\setcounter{footnote}{0}%
\opage{21}af Sjælelivet, men omfatte Alt, ogsaa hvad der for en mere indtrængende
Analyse staar som særegne Sider af Sjælelivet. Sokrates’
Personlighed og Virken bliver f. Ex. uforstaaelig, naar man
ikke fastholder, at han ved den Viden, hvori han saa den rette
Handlens Grundlag, mente „noget Stærkt, Ledende og Beherskende“,
„stærkt nok til at hjælpe Mennesket gennem Livet“
(Platon’s Protogoras p. 357 B C). En Sondring mellem Viden og
Villen, mellem Hoved og Hjerte, gjorde sig først senere gældende;
i hvert Tilfælde manglede Evnen til at opfatte og karakterisere
Kløvnings- og Brydningsforhold. Af Totaltilstandene ville igen
de, der ere opstaaede ved Sammensmeltning være særligt vanskelige
at beskrive, medens Organisationsformen lettere lader sig
analysere.

Totaltilstande ville som Regel ikke have den faste og ligevægtige
Karakter som Enkelttilstande. De Modsætninger, som ere
overvundne og forbundne i dem, kunne bryde frem paany. De
svare jo til flere Oplevelser, ikke til en enkelt, og Koncentrationen
bliver derfor ikke let fuldstændig. Der er ofte et Spændingsforhold
mellem de forskellige Elementer, som kan føre til Opløsning,
eller dog til Omdannelse til nye Former, maaske gennem
Udskilning af Elementer, der ikke forligedes med de andre. Et Menneske
kan maaske kun bevare sin indre Sammenhæng ved saavidt
muligt at glemme en enkelt bestemt Oplevelse, trænge den tilbage
i det Ubevidste. Lear trængte voldsomt Erindringen om sine
Døtres Utaknemmelighed bort for at den ikke skulde gøre ham
afsindig. O, that way madness lies; let me shun that; no more
of that! (King Lear III, 4, 21). Det er ikke blot virkelige Oplevelser,
der kunne volde Kløvning eller Opløsning. Ogsaa Mulighederne
spille, jo mangfoldigere Forestillingerne ere, en stor Rolle, snart
lokkende, snart ængstende, snart begge Dele. „Kun den,“ siger
Kierkegaard,\footnote[1]{\textit{Samlede Skrifter}. IV p. 422. --- I \textit{Begrebet Angst} (hvor den anførte Udtalelse findes) og i sine efterladte Papirer har Kierkegaard lagt den største Vægt paa ængstende Muligheders lokkende Magt. De volde Svimmelhed, virke ved den blotte Bevidsthed om, at de ere til.} „der blev dannet ved Muligheder, dannes efter sin
Uendelighed.“ Saalænge Livet varer, høre de Muligheder, der for
Alvor fremstille sig for os, ogsaa med til Oplevelserne, og de
% --- p. 22 ---
\setcounter{footnote}{0}%
\opage{22}kunne være vanskeligere at faa Bugt med end de faste og begrænsede
Oplevelser, vi kalde de virkelige. Det er af saadanne Muligheder,
store Digtere skabe deres Dramaer. Men i os Alle kan der
opføres indre Dramaer med Knuder og Løsninger, med Sammenstød
og Afgørelser, Dramer, der maaske slet ikke anes af vore
Omgivelser, men som dog kunne afsætte deres dybe Spor i vor
Karakters fortsatte Udvikling.

De Skikkelser og Skæbner, Digterne fremstille for os, høre ogsaa
med til vore Oplevelser. Selve det, at jeg maa indrømme, at
Livet rummer noget Saadant som det, der skildres, hvad enten det
er stort eller lidet, godt eller ondt, kan blive af afgørende Betydning
for mig. Det er ikke nødvendigt, at jeg selv skal have oplevet
noget Lignende, naar det blot staar som Noget, der kan
have sin Plads i Virkeligheden. Man har ment, at Bevidstheden
om, at vi selv ikke ere med i det, vi læse eller se, men
sidde i fuld Sikkerhed som Læsere eller Tilskuere, skulde være
et væsentligt Element i den æstetiske Følelse. Men det er i hvert
Tilfælde ikke Alle, der have den rolige Bevidsthed om sig selv
som Læsere eller Tilskuere ved Siden af Opfattelsen af Digterens
Fremstilling. Efter min personlige Erfaring vil en saadan Bevidsthed
umuliggøre den fulde Optagethed og Betagethed. Og da vi
lære os selv at kende, ikke blot gennem den Maade, paa hvilken
Livets virkelige Begivenheder paavirke os, men ogsaa gennem
den Maade, paa hvilken vi stemmes ved de Muligheder, vor egen
Fantasi eller Digternes Fremstillinger lægge frem for os, er det
med fuld Ret, naar vi her tale om Oplevelser. Derfor kan det
naturligvis alligevel have sin Betydning at fastholde Forskellen
mellem de Oplevelser, der skyldes Muligheder, og dem, der skyldes
Virkeligheder. Hine indeholde unegteligt flere Fristelser til
Selvbedrag, til at anse sig som Helt uden at være det, eller til i
Angst at tænke for ringe om sig selv, maaske kaste sig selv helt
hen.

I Henhold til denne Betragtning kan jeg ikke saa skarpt som
Alfr. Lehmann\footnote[1]{\textit{Die Hauptgesetze des menschlichen Gefühlslebens.} Leipzig 1914. p. 312—316. --- Det er beklageligt, at denne nye, udvidede og ændrede Udgave af „Følelseslivets Hovedlove“ ikke foreligger paa Dansk.} skelne mellem autopatiske, sympatiske og æste%
% --- p. 23 ---
\setcounter{footnote}{0}%
\opage{23}tiske Følelser, saaledes at disse sidste skulle være karakteriserede
ved den rolige Bevidsthed om at staa helt udenfor, hvad Fantasi
eller Kunst lader fremtræde for os. Der gives utvivlsomt Tilstande,
i hvilke Forskellene mellem disse tre Arter af Følelse ere forsvundne,
idet vort reale eller centrale Jeg netop kan bestaa i, udgøres
af den levende Medfølelse med virkelige eller forestillede
personlige Væsener og først senere vender tilbage til sin isolerede
Existens, saaledes at denne dog vil være undergaaet en Forandring
gennem Selvforglemmelsens Tilstande. Vort reale Jeg kan svinge
mellem Selvforglemmelse og Selvbevidsthed, og Udslag i Retning
af Selvforglemmelse er ofte mest frugtbart. Lehmann ender
ganske vist med at anerkende Overgange mellem de tre Arter af
Følelser, men denne Anerkendelse er ikke tilstrækkelig. Hovedsagen
er, om han har Ret i, at Forestillingen om mig selv som
forskellig fra alt Andet og alle Andre, altsaa Jegforestillingen, har
den Betydning og den Allestedsnærværelse, han tillægger den.
Som allerede ovenfor bemærket, kan der gives et realt Jeg uden
Jegforestilling. Et Andet er, at vi i Psykologien, naar der spørges
om, hvorpaa vi grunde Berettigelsen til overhovedet at tale om
et Jeg, bl. A. henvise til det reale Jeg i den Betydning, i hvilken
vi ovenfor have taget dette Ord. ---

Det Spørgsmaal, om det er et Gode, at Forholdene paa Sjælelivets
Omraade blive mere ustadige ved fremskridende Udvikling,
skal ikke drøftes her. Store Farer og store Goder følges ofte ad.
Jo større en Fylde der er at holde sammen paa, des mere Energi
behøves der, for at Koncentration skal blive mulig. Det kan blive
et Spørgsmaal om at være eller ikke være, Koncentrationens Mulighed
stiller.

I denne Henseende staa de to Arter af Helhedsdannelse, der
ere antydede i det Foregaaende, og som jeg i det Følgende skal
vende tilbage til, ikke ens. Ved Sammensmeltning indgaa Elementerne
en nøjere Forbindelse end ved Organisation. Der opstaar
ved Sammensmeltning en ny Egenskab, en ny Umiddelbarhed;
de medvirkende Elementer have tabt deres Selvstændighed
og dermed Evnen til at gøre Oprør. Den dramatiske Karakter, det
indre Liv kan antage formedelst Elementernes indbyrdes Modsætning,
forudsætter, at en Sammensmeltning ikke har fundet Sted.

% --- p. 24 ---
\setcounter{footnote}{0}%
\opage{24}I Undersøgelsen af Totaltilstande har Psykologien sine højeste
og uafviselige Opgaver. Med Forkærlighed sysle Psykologerne
gerne med Enkelttilstande, og endnu hellere med de Elementer,
der kunne udskilles af dem. Det bliver da Spørgsmaalet, om Totaltilstande,
komplexe psykiske Fænomener uden videre kunne
udledes af de elementære Synspunkter eller Love, som kunne
findes ved Analyse eller Experiment. Naar dette ikke kan lade
sig gøre, taber man ofte Interessen for Totaltilstandene og lader,
som om de ikke existere, eller overlader til Poeter og Prædikanter
at beskæftige sig med dem. Det synes dog klart, at vi her som
overalt maa gaa frem ad to Veje, naar vi ville vinde Erkendelse,
nemlig ikke blot forsøge at trænge frem fra det Simple til det
Komplexe, men ogsaa omvendt iagttage de komplexe Fænomener
i deres Ejendommelighed og i deres historiske Sammenhæng,
fæstne dem ved Beskrivelse og muligvis gennem Analyse naa tilbage
til Synspunkter, der ere vundne ved Undersøgelse af Sjælelivets
mere elementære Fænomener. Maaske vil det saa vise sig,
at det Elementære ikke er saa simpelt endda. Det er en Fordom,
som paa mange Omraader atter og atter stikker Hovedet frem, at
„i Begyndelsen“ vare Elementerne, og at de desuden ere „det
Egentlige“. Men ved alle de Begyndelser, vi kunne iagttage, er
det Tilstande, ikke Elementer, som foreligge, og det Tilstande med
et mere eller mindre afgjort Helhedspræg. De Processer, gennem
hvilke Totaltilstande blive til, kunne bedst studeres paa mere fremskredne
Udviklingstrin, hvor Betingelserne bedre kendes, og hvor
Bevidstheden er mere vaagen. Og det er da ikke blot Totaltilstandene,
der kunne belyses ved Analogi fra Enkelttilstandene, men
ogsaa omvendt. Ligesaalidt som en Totaltilstand kan opfattes som
en Sum af Enkelttilstande, ligesaalidt kan en Enkelttilstand opfattes
som en en Sum af Elementer. Sjælelivets historiske Karakter, dets
Afhængighed af bestemte Tidsforhold, fremtræder efter Sagens
Natur mest ved Totaltilstandene. Den historiske Metode maa derfor,
som allerede bemærket, supplere de andre psykologiske Metoder,
og Psykologien maa benytte, hvad Kultur- og Literaturhistorie
kunne lære os.

Det er en Undersøgelse af denne Art, som er Hensigten med
denne Afhandling. Men da den Totaltilstand, som skal undersøges,
% --- p. 25 ---
\setcounter{footnote}{0}%
\opage{25}faar sin Hovedejendommelighed ved de Følelseselementer, den
indeholder, maa vi først dvæle lidt ved Følelseslivet og se, hvorledes
det hidtil Udviklede finder Anvendelse der.

5. Som Enkeltfølelse betegner jeg en sjælelig Enkelttilstand,
i hvilken Følelseselementerne ere overvejende og afgørende, og
som ikke blot er bestemt ved en enkelt Oplevelse, men tillige er
forholdsvis enkelt i sin Natur. En saadan Enkeltfølelse er f. Ex.
Velbehaget ved en Farve eller en Duft, Glæden ved en Tanke
eller et Billede, den uvilkaarlige Billigelse, Synet af en brav Handling
kan fremkalde, Tilfredsstillelsen ved selv at have handlet
med Dygtighed. Glæde og Sorg, Haab og Frygt, Vrede og Kærlighed
kunne optræde som Enkeltfølelser.

Saadanne Enkeltfølelser kunne indgaa Forbindelser, af hvilke
der kan fremgaa Totalfølelser. Totalfølelse forudsætter Oplevelser
af forskellig Art, som staa i indbyrdes Sammenhæng. Dog kan en
objektivt set enkelt Oplevelse have saa forskellige Virkninger ved
de forskellige Egenskaber, den frembyder, at den subjektivt set
kommer til at staa som flere forskellige samtidige Oplevelser. Totalfølelser
kunne fremkaldes i Enkelttilstande (se de i § 1 anførte
Exempler), men kun naar de først have udviklet sig som
Totaltilstande gennem en Række af Oplevelser. Der er her en
Analogi til partiel Genkendelse, ved hvilken vi kende en sammensat
Genstand paa en enkelt Del eller Egenskab, naar vi tidligere
have havt Fornemmelser, der svarede til hele Gruppen af
Dele eller Egenskaber. (Sml. Psykologi V B., 2).

Allerede Enkeltfølelserne have en forskellig Karakter hos hvert
Individ, hvert Folk, hver Stand og hver Tid. Men endnu større
Forskelligheder optræder ved Totalfølelserne. ---

Forskellen mellem Enkeltfølelse og Totalfølelse har jeg allerede
fremhævet i min Psykologi (VI B, C og E), uden dog at bruge
disse Betegnelser. Jeg finder den i Literaturen først antydet hos
Platon paa det ovenfor anførte Sted og, som det senere skal omtales,
gør den sig paa tydelig og karakteristisk Maade gældende
hos Renaissancens mest fremragende Aander. I nyere Tid er den
først med Klarhed beskreven af F. C. Sibbern (Psykologie 1856
p. 140. 380), som baade lægger Vægt paa Følelsernes sporadiske
% --- p. 26 ---
\setcounter{footnote}{0}%
\opage{26}Optræden i Overgangstilstande og Kriser og paa Muligheden af, at
der kan opstaa, hvad han kalder en „blandet Følelse“, i hvilken
modstridende Følelser have indgaaet saa inderlig en Forbindelse,
at de virke sammen til at danne én eneste Totalfølelse. Dette
Fænomen vil jeg hellere kalde Sammensmeltning, især da „blandede
Følelser“ i daglig Tale bruges omtrent i Betydning af Ubehag.
Men der gives en anden Art af Totalfølelse, som har Karakteren
af Organisation i den Betydning, i hvilken vi ovenfor have
brugt dette Ord. Forskellige Følelser kunne nemlig virke sammen
under Ledelse af en enkelt Følelse, der bestemmer den Værdi
eller det Formaal, i hvis Tjeneste de staa. Denne Art af Totalfølelse
blev godt beskreven af Th. Ribot i hans Bog „Essai sur les
Passions“ (1907), hvor Ordet „emotion“ bruges om Enkeltfølelse,
Ordet „passion“ om Totalfølelse. Og Alexander Shand har paa
en mesterlig Maade skildret, hvad jeg kalder Organisation, under
Navn af Systematisation, først i sin Afhandling „Character of the
Emotions“ (Mind 1896), senere i sit Værk „The Foundations of
Character (1914) % PRINTED AS IS: no closing quote or stop

Vi betragte nu nærmere de to Former, i hvilke Totalfølelser
kunne fremtræde.

6. Sammensmeltning er den inderligste Forbindelse, forskellige
Følelser kunne indgaa, og der gives mange Mellemformer eller
Trin, før denne Grad naas. Sindet kan bølge frem og tilbage, idet
forskellige Sindsbevægelser afløse hverandre. Sokrates’ Disciple
snart lo, snart græd under den sidste Samtale med deres Mester.
Situationen var jo den, at hans Død var nær forestaaende, men at
de dog endnu sammen med ham øvede deres sædvanlige Tankearbejde.
Der var, siger Platon (Faidon p. 59 A), „en forunderlig
Stemning“, en Blanding af Sorg og Glæde. At Graad og Latter
skiftede, viser, at indenfor Blandingen har snart det ene, snart
det andet Element havt Overhaand, og det er ikke kommet til
nogen samlet, ny Følelse. Det var her tydeligt to forskellige Aarsager,
Tanken om Mesterens Bortgang og den filosofiske Samtale,
der hver satte sin Følelse i Bevægelse. Naar Disciplene senere
tænkte paa deres Mester, have de forskellige Sindsbevægelser
% --- p. 27 ---
\setcounter{footnote}{0}%
\opage{27}kunnet udligne sig med hinanden, saa at Forestillingen om Tabet
og Forestillingen om den uforlignelige Personlighed have kunnet
virke sammen til at fremkalde en Totalfølelse af vemodsfuld Beundring;
en Sammensmeltning har kunnet fuldbyrdes.

En Sammensmeltning kan ofte kun finde Sted i enkelte særligt
koncentrerede Øjeblikke, medens Enkeltfølelserne i andre Øjeblikke
skifte eller brydes med hinanden. Vemod f. Ex. er da ikke
den raadende Følelse; den faar kun under gunstige Betingelser
Herredømmet, medens ellers Sorgen over Tabet og Glæden ved
det Tabtes Værdi afløse hinanden. Et Menneskes Karakter --- det
reale Jeg --- kommer overhovedet ikke udelt frem i hver eneste
Tilstand. Det er paa vort Livs Højdepunkter, at vi virkeligt ere
os selv, d. v. s. beherskes af den Totalfølelse, der har dannet sig
under Indflydelse af vore Oplevelser. Til andre Tider raade partielle,
mere eller mindre periferiske Jeger. Og Forholdet mellem
det Centrale og det Periferiske kan være meget forskelligt hos
forskellige Naturer. Ofte kan kun et meget nøje Kendskab finde
de Traade, der forbinde det Periferiske og Partielle med det Centrale
og Totale.

Ikke altid gør der sig fra først af en bestemt Forskel gældende
for os mellem Aarsager, der paavirke os i forskellig Retning.
Deres Følelsesvirkninger kunne fra første Færd forbinde sig, saa
at Totalfølelsen opstaar umiddelbart, uden at være forberedt gennem
Svingninger mellem forskellige, mere eller mindre modsatte
Følelser. Saaledes kan et Billede ved første Syn vække en Totalfølelse,
hvis Elementer først bagefter kunne spores. Exempler herpaa
afgive Velasquez’ Billeder af de Dverge, han traf ved det
spanske Hof. Disse Smaafyre staa som komiske ved deres utvetydige
Selvfølelse i Kontrast med deres Lidenhed, men det Komiske
kompliceres ved et melankolsk Udtryk i Øjet, der indfører et
Element af Sympati i Beskuerens Stemning, som derved kan blive
af den Art, vi i det Følgende skulle beskæftige os med, --- Humor,
ligesom der maaske har ligget Humor til Grund hos Kunstneren,
da han gik til sit Værk. Julius Lange bemærker om Velasquez’
Dvergportræter, at saa stærkt karakteriserede Dvergeskikkelserne
end ere, saa viser Mesterens noble Skønhedssans sig dog i hans
% --- p. 28 ---
\setcounter{footnote}{0}%
\opage{28}Blik for det Punkt, hvor Legemets Vanskæbne har gjort Sjælen
Uret.\footnote[1]{\textit{Udvalgte Skrifter} II p. 142.} Det kan naturligvis ikke afgøres med fuld Sikkerhed, hvorvidt
de to Elementer have været forbundne til Totalfølelse i Kunstnerens
Sjæl. Ogsaa for Beskuerens Vedkommende kan der maaske
være Uenighed om, hvorvidt Totalfølelsen opstaar umiddelbart.
Det kunde jo være, at Beskueren først betragtede Skikkelsen og
derefter dvælede ved Blikket. For mit Vedkommende forholdt det
sig saaledes, at jeg strax fik Øje paa Blikket som Centrum i dets
hele Skikkelse. Først bagefter stavede jeg mig gennem Billedets
forskellige Dele. Jeg maa tilføje, at jeg ikke har set Originalerne.

Den tydeligste og skarpeste Opfattelse af Modsætninger faa vi
ganske vist ved successiv Iagttagelse af dem. Alt Andet lige, virker
successiv Forskel stærkere end simultan Forskel. (Psykologi
V A, 5). Men William James\footnote[2]{\textit{Principles of Psychology} I p. 495.} har neppe Ret, naar han gaar et
langt Skridt videre og paastaar, at to Led eller Elementer ikke
kunne opfattes som forskellige ved samtidig Opfattelse, naar vi
ikke successivt have opfattet dem hvert for sig og derved faaet
den ejendommelige Overgangsfornemmelse af Forskel (the transitional
sensation of difference). Der sker, hævder han, ingen Analyse
af et Bevidsthedsfænomen, hvor komplex det saa er, hvis
ikke nogle af dets Elementer ere undergaaede Forandring. --- Hertil
vil jeg først bemærke, at der er flere Overgangsformer mellem
den umiddelbare Anskuen og den Analyse, James taler om. Selv
om der sker Forandring i en Del af et Anskuelsesbillede (f. Ex.
ved at et enkelt Træk betragtes nøjagtigere eller belyses stærkere
udefra), opløses Billedet derfor ikke. Virkningen bliver maaske
kun, at Billedet ligesom tilspidses i en enkelt Retning eller faar et
nyt Tyngdepunkt. Anskuelsen artikuleres, leddeles, men opløses
ikke i sin samtidige Omfatten af Billedets Indhold.\footnote[3]{Smlg. om artikulerende Anskuen som Mellemform mellem Anskuelse og Dom \textit{Det psykologiske Grundlag for logiske Domme.} (Vid. Selsk. Skr. 6. Række). p. 18—21. (Kortere i \textit{Den menneskelige Tanke.} p. 46; 53—56).} --- Dernæst
er det i hvert Tilfælde hvad Følelsesvirkningen angaar, ikke selve
Forskelsopfattelsen, der er afgørende, men netop den Omstæn%
% --- p. 29 ---
\setcounter{footnote}{0}%
\opage{29}dighed, at de forskellige Elementer virke sammen trods deres Forskellighed.
Til den komplexe Paavirkning kan meget vel umiddelbart
svare en komplex Følelsesvirkning. Og det kan saa netop
være Forundringen over denne Følelsesvirkning, der gør, at vi
ved Hjælp af successiv Iagttagelse af Billedets enkelte Dele, altsaa
ved Analyse, søge at fremdrage de samvirkende Elementer
hvert for sig. Den egentlige Forskelsopfattelse bliver da sekundær,
men derfor er Følelsesvirkningen ikke sekundær i Forhold til de
enkelte Elementers Følelsesvirkninger. --- Det maa naturligvis indrømmes,
at det er vanskeligt at skelne mellem Samtidighed og
meget hurtig Succession. Og den Kendsgerning, at vi kun i meget
korte Øjeblikke kunne samle vor Opmærksomhed om et Emne,
altsaa at Opmærksomheden vibrerer, peger ogsaa i Retning af, at
successiv Opfattelse meget hurtigt indtræder. ---

Som Exempler paa Sammensmeltning nævner Sibbern to Grupper
af Fænomener.

For begge Grupper gælder det, at der efter et Spændingsforhold
mellem to Følelser opstaar en Totalfølelse, saa at de to simplere
Følelser ikke mere mærkes hver for sig. Istedetfor de Følelser,
som brødes med hinanden eller afløste hinanden, er der indtraadt
en Følelse af ny Kvalitet, som dog vilde falde bort, naar en af de
forbundne Modsætninger faldt bort.

Den ene Gruppe af Exempler betegner Sibbern som den subjektive
Tilfredsstillelses Følelser. De skyldes Fyldestgørelse af en
Trang, Udfyldning af et Savn, Lindring af Ubehag, Overvindelse
af Modstand. Ogsaa Følelsen ved det Spændende og Hemmelighedsfulde,
ved Fare og ved Kamp hører herhen.

Mod de førstnævnte Exempler paa denne Gruppe indvender
Alfr. Lehmann, at da Trangen bortfalder ved Fyldestgørelse, sker
der her ingen Sammensmeltning. Sibbern kunde have svaret hertil,
at enhver Følelse har en Tendens til at holde sig, efterat den
oprindelige Anledning er falden bort. Den breder sig ved en Slags
Expansion, ud over de følgende Tilstande og kan derved komme
til at smelte sammen med disse Tilstandes Følelseselementer.
(Psykologi VI B, 2; E, 4). Det at have følt en Trang, er ikke
nødvendigvis et rent Fortidsfænomen. Sprogforskerne lære da ogsaa,
at et Perfectum absolutum foruden Tidsmomentet ogsaa inde%
% --- p. 30 ---
\setcounter{footnote}{0}%
\opage{30}holder noget Andet, nemlig Forestillingen om Resultatet, om „det
Efterlevende, Permansive.“\footnote[1]{\textit{Otto Jespersen: Tid og Tempus.} (Vid. Selsk. Oversigt 1914). p. 320 f.} Og Sibbern taler udtrykkeligt om
den ejendommelige Glæde ved \emph{at være} undsluppen fra en Fare
eller Nød, eller \emph{at have} overstaaet den. Der er altsaa Noget, man
stadigt har, naar en Fare er overstaaet. Saaledes taler han ogsaa
(i anden Del af Gabrielis’ Breve) om „indoptagen Usalighed“ som
forskellig fra blot overstaaet Usalighed. ---

Den anden Gruppe af Sibberns Exempler indeholde Følelser
som Vemod, Ærefrygt, Følelse ved det Ophøjede. At de kunne
have Karakteren af Sammensmeltning, kan der ikke rejses Tvivl
om. Hertil høre ogsaa, som det siden vil vise sig, mange Former
for Humor.

7. Den anden Form for Følelsesforbindelser fremtræder, naar
en Række af Oplevelser udløse Følelser, der vel staa i nøje indbyrdes
Sammenhæng, men ikke helt tabe deres Ejendommelighed
hver for sig, hvorimod de danne en Helhed, hvis indre Ordning
bestemmes ved en enkelt Værdi, et ledende Formaal som organiserende
Kraft. En saadan Ordning kan opstaa, naar der gør sig en
Grundværdi gældende, der til sin fulde Virkeliggørelse kræver en
Række af Midler og Veje, saavel som en mere eller mindre bevidst
Stræben. Der vil da under Grundværdiens Brydning med
Virkelighedens Forhold udløses en Række forskellige Følelser,
der saa at sige ligge paa Vejen til det egentlige Maal. Der kan
være Mennesker, hvis hele Liv bestemmes ved en saadan Grundværdi
og dens Forhold til den Virkelighed, indenfor hvilken den
skal gennemføres. For andre Mennesker er det kun visse Livsomraader,
f. Eks. Kunst eller Videnskab eller Samfundsliv, hvor
en Grundværdi overhovedet existerer for dem. For andre igen
har hver Livsperiode sin Grundværdi, Ungdommen sin egen, den
senere Alder sin; eller det er maaske kun i en enkelt Livsperiode,
at der overhovedet kan tales om en Grundværdi. Ethvert alvorligt
Forhold til en Sag, til et Menneske, til en Bestræbelse (Kunst,
Videnskab, Menneskeomsorg, social og politisk Frihedsstræben),
vil naturligt afføde en Organisation af Følelseslivet, saa at de 
en%
% --- p. 31 ---
\setcounter{footnote}{0}%
\opage{31}kelte Følelser faa deres bestemte Plads i den sjælelige Husholdning
efter deres Betydning for hint Forholds Hævdelse. De Enkeltfølelser,
Livets almindelige Gang medføre, --- Glæde og Sorg,
Haab og Frygt, Velvillie og Vrede, --- ville begunstiges eller
hæmmes ved Arbejdet i Grundværdiens Tjeneste, og selve dette
Arbejde vil give Anledning til en Række Enkeltfølelsers Opstaaen.

Der vil ved den Sammenhæng, i hvilken Enkeltfølelserne træde
ind, naar der sker en Organisation om en raadende Interesse, bevirkes
Ændringer, Nuancer i Enkeltfølelsernes Beskaffenhed.
Dels vil den Retning, i hvilken den hele Stræben gaar, faa Indflydelse
paa de enkelte Følelsers Art og Grad, dels vil den indbyrdes
Vexelvirkning mellem dem, naar de mødes under Bestræbelsernes
Forløb, give dem en særegen Karakter, de under andre
Forhold ikke vilde faa. Vreden vil f. Ex. have en anden Karakter,
naar den har et egoistisk Udspring, end naar den skyldes Krænkelse
af en altruistisk Følelse, og fremdeles optræder den anderledes,
naar den udløses sammen med eller lige efter Frygt, end
naar den er udløst sammen med eller lige efter Glæde, --- saadan
Frygt eller Glæde, som Arbejdet i Grundværdiens Tjeneste frembringer.


Alexander Shand, der i sit ovenfor nævnte Skrift „Foundations
of Character“ paa en udmærket Maade har beskrevet Enkeltfølelsers-
(emotions’) Organisation (eller som han kalder det, Systematisering)
til Totalfølelser (sentiments), synes dog at nægte
den Indflydelse, det maa have paa en Enkeltfølelses Beskaffenhed,
at den indgaar i en Organisation. Den uinteresserede Karakter af
en Følelse beror, siger han (p. 43— 50), ikke paa selve Følelsen,
men paa den Aarsag, der fremkalder den, eller paa det System, til
hvilket den hører; Vrede og Frygt kunne være ligesaa uinteresserede
som Medlidenhed og Sympati, naar den Aarsag, der fremkalder
dem, ikke hører ind under Individets egen Selvopholdelsestrang
og den Systematisering, som den Trangs Herredømme bevirker.
Ganske tydelig er denne Udtalelse ikke, men den synes
at gaa ud paa, at Følelsen selv ikke faar forskellig Karakter ved
at tjene forskellige Formaal. Karakteren af uinteresseret eller selvisk
skulde da kun angaa Formaalet, ikke de Enkeltfølelser, Arbejdet
for Formaalet fremkalder. Men faar Frygt ikke en forskellig
% --- p. 32 ---
\setcounter{footnote}{0}%
\opage{32}Nuance, naar man er bange for sin egen Person, og naar man er
bange for Andres Vel eller for en Sag, man er optagen af? Baade
Art og Grad ville sikkert være forskellige. Selve det forskellige
Formaal vil gøre, at der i de forskellige Tilfælde indgaa forskellige
Elementer i Frygtfølelsen; ogsaa her kan der ske Sammensmeltning
eller Organisation, og derved kan ny Klangfarve opstaa for
Enkeltfølelsernes Vedkommende. Der kan ikke være et aldeles
ligegyldigt Forhold mellem Følelsen selv og den Sammenhæng, i
hvilken den forekommer. Mangehaande Kontrast- og Udligningsforhold
ville her gøre sig gældende.

Det hænger sikkert sammen med Shand’s Opfattelse af dette
Punkt, at han ikke ret anerkender den første, ovenfor omtalte
Hovedform for Totalfølelse. Hvad han kalder en Følelses „System“,
er Indbegrebet af de Tanker, Følelsen er knyttet til, ---
den Stemning (feeling) og de Tendenser, der gøre sig gældende
ved den, --- de organiske Processer, der ere knyttede til den, ---
og den Adfærd eller Handlemaade, der følger af den. Men disse
Elementer skulle ikke undergaa nogen Forandring ved at indtræde
i „Systemet“, Teorien om fuld Sammensmeltning til en kvalitativt
ny Følelse vilde ifølge Shand (p. 55) ikke kunne forklare,
hvorledes en Totalfølelse under forskellige Forhold --- Genstandenes
Nærværelse eller Fraværelse, Opgaaen i Forventning eller
i Erindring --- kan optræde paa højst forskellig Maade. En sammensmeltet
Følelse (compound emotion) vilde ikke optræde paa samme
Maade, kun „systematiserede“ Følelser (complex emotions, ---
emotions combining into a system) kunne have den Plasticitet, som
Erfaring godtgør for visse Følelsers Vedkommende (p. 55; 172).
Shand har Ret i, at der praktisk set kan være stor Forskel paa de
to Arter af Totalfølelse; men det udelukker ikke, at der her foreligger
to ejendommelige Former for Totalfølelse, og det er det,
det her drejer sig om. Iøvrigt kunne de hver have sine Fortrin.
Sammensmeltning gør en mere koncentreret Stemning og Færd
mulig, medens Organisationsformen har Plads for større Variationer
i de enkelte Tilfælde. Den større Plasticitet i Organisationen
kan dog meget vel forenes med, at „Systemets“ Elementer gensidigt
bestemme hverandres Nuance eller Klangfarve. Ethvert Element
stiller visse Fordringer, men der stilles ogsaa Fordringer til
% --- p. 33 ---
\setcounter{footnote}{0}%
\opage{33}det. Det staar baade i et negativt og i et positivt Forhold til den
Helhed, i hvilken det indgaar. Den Vrede, en Faders Omsorg for
hans Søns Opdragelse har fremkaldt, faar sin Begrænsning ved
Formaalet. Hos koleriske Fædre kan der maaske kræves en Kamp
for at holde Vreden indenfor denne Begrænsning. Hos Forbrydere
kan der (sml. § 2) under deres Handlings Udøvelse undertiden
røre sig en Grusomhed, der ikke er betinget ved det egentlige
Formaal for deres Adfærd. Ligeledes hos Soldater i Krigen. Ofte
er det Frygt, der gør grusom, idet den gør det umuligt at stanse,
naar man engang er begyndt at slaa løs for at sikre sig. De Elementer,
der organiseres, kunne have en stadig Tendens til at
bryde ud og røre sig ubehersket og ubegrænset af noget Maal undtagen
det ene: at rase ud! Men Klangfarven ændres saa i samme
Øjeblik, som de bryde ud. Negationen eller Begrænsningen falder
da bort. ---

Ikke alle Enkeltfølelser kunne indgaa som Led i Totalfølelse,
i den ene eller den anden Form. I saa Fald vilde det menneskelige
Sjæleliv have en mindre kaotisk og disharmonisk Karakter, end
det har. Der er baade Glæde og Sorg, Harme og Beundring, Haab
og Frygt, Tro og Tvivl, der ikke indgaa som Led i nogen fast
sjælelig Sammenhæng.

I Ingemann’s prægtige Digt „Der strømmer en Flod mod det
evige Hav“ ligger den optimistiske Tanke til Grund, at alle Sjælelivets
Elementer kunne harmoniseres. Det siges om hin Flod:

\begin{quote}
Hver Sorg den forvandler til Vemodslyst,\\
hvert Suk til smeltende \emph{Sang} --- --- ---\\
Selv Stenen, der knuged med Bjergevægt,\\
opløses til gyldene Gran.
\end{quote}

Ja, dette kan ske. Men det sker ikke altid. Der er Suk, som
ikke forvandles til Sang, og der er Stene, i hvilke der intet Guld
findes. Og da enhver Totalfølelse har sine Farer, idet Følgen kan
blive Afslutning, Uimodtagelighed for de nye Oplevelser, saa er
det maaske godt, at ikke alt sjæleligt Kaos er harmoniseret.

To store Aander have udtalt sig om, hvorledes Modsætninger
kunne holde sig i Sindet og dele dette, uden at der sker enten
Sammensmeltning eller Organisation. De bruge begge samme Bil%
% --- p. 34 ---
\setcounter{footnote}{0}%
\opage{34}lede til at belyse deres Erfaringer. Saaledes siger Carlyle (Sartor
Resartus II, 9): „I Forgaarden danses der til Guitarmusik, men
indenfor lyder af og til en sagte Klagen.“ Og Kierkegaard skriver
i sine Dagbøger (Papirer IV p. 35): „Jeg sidder og lytter til Tonerne
i mit Indre, Musikens glade Vink og Orglets dybe Alvor.“

Den psykiske Proces, der hos Mange uvilkaarligt fører til harmonisk
Forbindelse, er hos Naturer som disse to Mænd umulig, og
hos dem er det ofte kun ved en voldsom Villiesanstrengelse, at
der tilvejebringes en Enhed, idet de kaotiske Strømme tvinges
ind i ét Leje, hvor man dog stadigt kan mærke Forskellen mellem
de forskellige Kilders Vand.

8. Tendensen til Totalfølelsers Opstaaen hænger (som allerede
antydet i § 2) baade sammen med Sjælelivets inderste Natur og
med de Fordringer, Kampen for Tilværelsen stiller. Der gør sig
paa alle Trin og i alle Former af Sjæleliv en Trang og Evne gældende,
som gaar ud paa Sammenfatten, Syntese. Der virker her,
som Kant har udtrykt det, en skjult Kunst i os, uden hvilken vi
ikke vilde have Bevidsthed, men som vi kun sjældent blive os
bevidst. Det er den almindelige Psykologis Opgave at paavise
dette i det Enkelte; det var i hvert Tilfælde den Opgave, jeg satte
mig i min i Aaret 1882 udkomne Fremstilling af Psykologien.
Sjælelivet bestaar i og under et stadigt Arbejde paa at sammenfatte
og ordne et Kaos. Ogsaa paa Følelseslivets Omraade øves et saadant
Arbejde uvilkaarligt, og de Processer, der føre til Dannelse
af Totalfølelse, ere interessante Vidnesbyrd derom. Der gives
maaske, som vi have set, Naturer, hos hvem forskellige Interesser
eller Stemningstendenser kunne ligge jævnsides, næsten
selvstændige, eller saaledes, at de skiftevis faa Herredømmet.
Men der vil, afset fra den Indflydelse, ydre Paavirkninger kunne
øve, gøre sig en naturlig Vexeldrift gældende, idet en stærk Svingning
i den ene Retning afføder Trang til Svingning i den anden
Retning, og herigennem viser det sig da allerede, at de enkelte
Elementer ere underlagte de Krav, Hævdelsen af den hele sjælelige
Husholdning stiller. Men der vil desuden fordetmeste være
en Interesse, en Værdi, som i det lange Løb normerer de andre
% --- p. 35 ---
\setcounter{footnote}{0}%
\opage{35}Værdiers Optræden og bestemmer dem mere end den selv bestemmes
ved dem. Hvilket det Element er, der bestemmer Sammensmeltningens
Klangfarve eller Organisationens Grundværdi,
det vil ikke altid med Sikkerhed kunne erkendes af Individet selv
gennem Selviagttagelse, men det vil fremgaa af dets Adfærd, den
Maade, hvorpaa det tager Livet og griber ind i (eller ikke griber
ind) i Livet. Det er, som jeg har fremhævet i min Psykologi, ofte
først under en dybtgaaende, med fuld Energi foretagen Overvejelse
i en bestemt Situation, som kræver Handling, at det viser
sig, hvad vort reale Selv, vor Grundværdi egentligt er. Men det
kan meget vel være, at det reale Selv, der tidligere raadede i os,
under en Overvejelse, der stiller os overfor Livets Realiteter, ændres
paa afgørende Maade, saa at hvad vi opdage i saadanne Øjeblikke,
netop ikke er det, der hidtil laa til Grund i os. Naar man
fordyber sig i denne Betragtning, faar man et lille Indblik i, hvilken
uendelig Vexelvirkning der gaar for sig i det Inderste af vort
Væsen. Kunde en simpel Intuition en Gang for alle hjælpe os til
Klarhed her, vilde det være lettere at leve. Men Livet vilde da til
Gengæld miste sin Fylde og sin Spænding.

Naar der saaledes i selve Sjælelivets Natur, som Erfaringen
viser os det, ligger en Tendens, paa én Gang en Trang og en Evne,
til Samling og Koncentration, saa stemmer dette med de Fordringer,
Livsforholdene stille. Maalestokken for et levende Væsens
Udviklingstrin bestaar for en væsentlig Del i den Grad, i hvilken
det kan optræde som Helhed overfor sin Omverden. Planten formaar
dette i højere Grad end Stenen, og Dyret i højere Grad end
Planten. Indenfor Dyreverdenen hænger Nervesystemets Betydning
under Kampen for Tilværelsen sammen med, at det gør hel
og sluttet Optræden mulig. De Dyr, vi kalde de højere, have et
mere koncentreret Nervesystem end de, vi kalde de lavere. Og til
denne fysiologiske Koncentration svarer netop Sjælelivets Virken
som Syntese. I en Handlingssituation kan det gælde at samle alle
Erfaringer og Erindringer, al Stemning og Stræben og koncentrere
dem i en eneste Retning. Naar der er dannet en Totalfølelse, er
der allerede paa Forhaand sket en Mobilisering af Sindets Kræfter
i en bestemt Retning, og dette kan blive af den største Betydning
for Handlingslivet, --- medens det paa den anden Side
% --- p. 36 ---
\setcounter{footnote}{0}%
\opage{36}ogsaa kan føre til tragiske Konflikter, dersom Mobiliseringen er
sket i en anden Retning end den, i hvilken der skal handles. ---

I Sammenhæng med denne Betragtning er det af stor Betydning,
hvad Alexander Shand paaviser, at enhver Totalfølelse har
Tendens til at blive Grundlag for en Etik. Koncentration om en
Værdi medfører Stræben efter at tilvejebringe Alt, hvad der kræves
til at fastholde og udvikle denne Værdi. Og Dommene om
Godt og Ondt ville afhænge af, om denne Stræben begunstiges
eller hæmmes. Først og fremmest gælder det om, at selve Totalfølelsen
bevares i sin Frihed og Styrke, ti da ville de Egenskaber
(„Dyder“), som hin Stræben kræver, uvilkaarligt blive til. Skorter
det paa Friskhed og Styrke, maa der opstilles Idealer og fastslaas
Pligter. Dyder, Idealer og Pligter variere med Totalfølelsen. Saaledes
kræves der af Elskov eller Venskab som Totalfølelse Troskab
og Ærlighed, af Forældrekærlighed Taalmodighed, Selvbeherskelse
og Visdom, af intellektuel Interesse Udholdenhed og
Upartiskhed o. s. v. Hver Totalfølelse har sin Samvittighed. Den
saaledes udviklede personlige Etik staar i stadig Vexelvirkning
med den af sociale Krav bestemte Etik. (Foundations of Character
p. 111---120; 172 f.).

Denne Opfattelse af det etiske Problem gaar i samme Retning
som den, jeg forlængst har gjort gældende i min 1887 udkomne
Etik (Kap. III), og som jeg stadigt fastholder. Jeg har særligt vist,
at der er en Gruppe af etiske Begreber, som alle etiske Systemer
benytte, skønt de stille dem i forskellig Rangordning efter den
Grundværdi, de bygge paa. Saadanne Begreber ere Begreberne
Goder, Dyder, Pligter og Rettigheder (Etik VIII, 1). Enhver
etisk Opfattelse, ligegyldigt hvilken Totalfølelse den forudsætter,
faar Anvendelse for disse Begreber i en bestemt indbyrdes Orden.
Disse Begreber kunne derfor siges at udgøre Indholdet af en formal
Etik, der maa kunne fremstilles i Sammenhæng paa Grundlag
af en sammenlignende Undersøgelse af de forskellige etiske Systemer.
Herom har jeg udtalt mig nærmere i „Den menneskelige
Tanke“ (p. 350— 359).

Da der faktisk kan bygges paa forskellige Grundværdier, er der
her en Grænse for Etikens Rationalitet, en Grænse, som er given
ved ethvert etisk Systems sidste Forudsætning. Den sociale eller
% --- p. 37 ---
\setcounter{footnote}{0}%
\opage{37}sociologiske Etik, der bygger paa Samfundets Bestaaen som
Grundværdi, hvis Indhold er Formulering af de sociale Krav, og
som ofte betragter sig selv som den egentlige, rette Etik, er kun
et af de mange, psykologisk og historisk mulige etiske Systemer.

Medens jeg er kommen til denne Opfattelse gennem en erkendelsesteoretisk
Undersøgelse af Etikens rationelle Berettigelse
(smlg. Etik III, 2, 13), er Shand kommen til den ad rent psykologisk
Vej, idet han saa, at saadanne „Systematiseringer“ af Følelseslivet,
som han har undersøgt, uvilkaarligt maatte blive
Grundlag for de etiske Domme, der fældedes af de Mennesker,
hos hvem der var foregaaet „Systematisering“ i en eller anden
Retning.\footnote[1]{I en Afhandling om \textit{La Morale et la Science} har \textit{Henri Poincaré} (Derniéres Pensées 1913) udtalt en lignende Opfattelse af Etikens videnskabelige Stilling som den ovenfor fremstillede.} % PRINTED AS IS: „Derniéres“

9. Som allerede tidligere bemærket har der vel til alle Tider
dannet sig Totalfølelser. Men de træde ikke lige stærkt frem til
alle Tider, og Forholdene ere ikke altid lige gunstige for dem. Saalænge
Livet udfolder sig indenfor en begrænset Ramme, der staar
som uforanderlig, vil Følelseslivet have en forholdsvis enkelt Karakter.
Der vil være et fælles Præg, bestemt ved Fællesskabet
indenfor Familie, Stamme og Stat. Hvad der søger ud over denne
Ramme, alle Bestræbelser, for hvilke den er for snæver, alle
Modsætninger, der true med at sprænge den, hæmmes og trænges
tilbage. Betydningsfulde Totalfølelsers Opstaaen forudsætte ikke
blot stor Fylde af Erfaringer og Udblik til store Horizonter, men
ogsaa en Selvstændighed hos det enkelte Individ, der gør det muligt,
at der kan foregaa specielle, maaske enestaaende Krystallisationer
af de Følelser, der fremkaldes under mangfoldige Oplevelsers
Indflydelse.

Holde vi os til de for os vigtigste Kulturperioder, den antike og
den moderne, ville vi for den førstes Vedkommende ikke kunne
vente at finde mange karakteristiske Exempler paa Totalfølelse.
Den antike Oldtid beherskedes af Grundforudsætningen om Uforanderlighed
som Udtryk for Tilværelsens inderste Natur og tillige
som Karaktermærke for den højeste Værdi. Det Foranderlige stod
% --- p. 38 ---
\setcounter{footnote}{0}%
\opage{38}for græsk Tænkning tilsidst som Illusion, og det en værdiløs Illusion.
Selv i den rigeste Fremstilling af menneskelig Stræben, Oldtiden
har givet os, den, der foreligger i Platons „Symposion“, selv
der staar denne Stræben med al den Uro og al den Begejstring,
den medfører, som hjemmehørende i en lavere Sfære i Sammenligning
med den uforanderlige Væren, som er Tankens højeste
Genstand.\footnote[1]{Se herom min Afhandling om Platons Symposion i \textit{Mindre Arbejder}. Tredie Række. p. 88—93.} Følelser af Lyst og Ulyst forudsætte nu netop Forandring,
og det mere eller mindre stærk og pludselig Forandring,
og de naa deres mest intensive Former, hvor store Modsætninger
mødes eller afløse hinanden. Lyst og Ulyst ere ikke mindst knyttede
til Overgange og Kriser, i hvilke modsatte Bestræbelser
kæmpe om Herredømmet hos de Enkelte og i de Enkeltes indbyrdes Forhold. Kampen gælder tilsidst, hvad der skal være
Grundværdi og derved anvise alle andre Værdier deres Plads.
Mennesket drages til modsatte Sider, og at der er en Enhed i dets
Væsen lægger sig under en saadan Brydning kun for Dagen ved,
at Striden og Splittelsen føles som Smerte. Smerte og Sorg, Anger
og Tvivl forudsætte en indre Enhed; det er netop Deltheden, der
føles som Smerte. Lidelsen er det sidste Baand, der holder Sjælelivet
sammen. „Intet Glædeligt,“ lader Sibbern sin Gabrielis sige,
„intet Glædeligt havde jeg, der kunde give min Existens Kontinuitet;
saa blev da Smerte og Fortvivlelse det, der holdt min Tilværelse
sammen og bevarede mig fra at give min Sjæl hen i
Søndersplittelse.“ Derfor lærer man i Smerte og Sorg sig selv
saa godt at kende. Kunde et af Modsætningsleddene kastes bort
og glemmes, vilde Lidelsen falde bort i samme Øjeblik. Det er
Trofasthed imod sig selv og mod de Værdier, man har lært at
kende, der gør, at man bestaar i Lidelsens Skærsild, eller ogsaa
bevirker et stort Haab, at man trods Alt ikke havner i Selvopgivelse.


Koret i den græske Tragedie følger vel Heltens Lidelse og Sorg,
men vender dog stedse tilbage til Ønsket om et Liv, der ikke stedes
i saa svar en Strid og til Beslutningen om ikke at glemme de
Dødeliges Begrænsning. Og denne Tankegang udtrykte Grundsætningen
i den antike Livsanskuelse. Platon var vel klar over,
% --- p. 39 ---
\setcounter{footnote}{0}%
\opage{39}at Livet var baade Komedie og Tragedie, og at de hørte naturligt
sammen. Men denne Tanke blev dog ikke af afgørende Betydning
for ham. Maaske indtager Sokrates her en Særstilling. Det vil
senere blive nærmere undersøgt.

Med alt det Forbehold, der maa tages ved Betragtninger af
denne Art, tør man sige, at det hører den nyere Tid til at være
fuldt forstaaende overfor Livets Mangfoldighed og Modsætningsrigdom.
De første psykologiske Skildringer, der vidne om en saadan
Forstaaelse, fremkom omtrent samtidigt i Giordano Bruno’s
vidunderlige Værk: „De gl’ heroici furori“ og i Shakespeares Tragedier.


Hvad Bruno kalder heroisk Lidenskab, betinges ved det højere
Sjælelivs sammensatte Karakter. Kun hos de Dumme, der ikke
kunne se ud over det umiddelbart Foreliggende og ikke have
nogen omfattende Erfaring, ikke kende Tilværelsens Mangfoldighed,
kun hos dem kunne Følelse og Trang have en enkelt og simpel
Karakter. Hos dem, der have rig og mangesidig Erfaring, er
Følelseslivet sammensat af Lyst og Smerte, Sødme og Bitterhed,
og det saaledes, at disse modsatte Elementer gensidigt bestemme
hinanden. Uvidenheden er Moder til Nydelsen i det dyriske Paradis,
og den, der forøger sin Viden, forøger sin Smerte. Men den
heroiske Sjæl føler Lyst ved selve Smerten, fordi denne er Tegn
paa, at der higes mod en uendelig og ophøjet Genstand, som maa
og kan være Maal for stedse fornyet Stræben. Den heroiske Sjæl
vil hellere segne i Kampen eller dog føle sin Mangel og Ufuldkommenhed
overfor det store Maal end opnaa Fuldkommenhed
ved at sætte sig et lavere Maal.

Sammenligning med Platon’s „Symposion“ viser Forskellen
mellem antik og moderne Tankegang. Platon vil føre ud over den
stadige Higen, ud over den Forandringens Verden, i hvilken der
ene kan være Tale om Stræben. Sindsbevægelserne skulle afløses
af intellektuel Klarhed, og Sindet skal lægge sig til Ro i det Uforanderlige.
Men for Bruno er Modsætningernes Spil og den Spænding,
det medfører, netop Udtryk for sand Realitet, fordi det er
Livets Rigdom og Verdens Uendelighed, der medfører de stadige
Forandringer og de store Modsætninger. Ligesaa lidt som det ydre
Univers for ham havde nogen Grænse, ligesaa lidt kunde Idea%
% --- p. 40 ---
\setcounter{footnote}{0}%
\opage{40}lernes og Værdiernes Verden have Grænse. Den store Grundtanke,
for hvilken han blev Martyr, har gennemtrængt hans inderste
Sjæl og bestemt hans Livsanskuelse. --- Den Totalfølelse,
Bruno har beskrevet, er udsprungen af Trangen til Erkendelse og
af de Livserfaringer, der betinges ved Alt, hvad der enten begunstiger
eller hæmmer denne Trangs Tilfredsstillelse. Det var
den fremtrængende, selvstændige Tankes Skæbner, der betingede,
hvad han fik ud af Livet. Sin mest intensive aandelige Livsfølelse
havde han, hvor han saa de Skranker, det gamle Verdensbillede
havde sat for Aanden og Tanken, vige for større Horizonter. Som
det senere vil blive vist (se Kap. IX), er denne Art Totalfølelse
den hyppigste hos store Filosofer.

Hvad Shakespeare angaar, vil der senere ofte blive Lejlighed
til at fremdrage ham i Anledning af den Totalfølelse, denne Afhandling
særligt skal beskæftige sig med. Her skal blot henpeges
paa, at hans Karakterer, saavel som de Skæbner, han skildrer,
ikke ere byggede paa ganske enkelte Motiver, men, hvor han naar
højest, hver for sig ere som et helt Orkester, i hvilket mange Motiver,
Modsætninger, Opløsninger og Harmonier klinge sammen
til en mægtig Totalitet. Derved er han bleven en Fundgrube for
Psykologien i højere Grad end nogen anden Digter.

10. Forskellen mellem Enkeltfølelse og Totalfølelse lægger sig;
paa en ret interessant Maade for Dagen i Ordet Humors Historie,
og det er nødvendigt at dvæle lidt herved for at forstaa den specielle
Betydning af Ordet, i hvilken det skal tages i den følgende
Undersøgelse.\footnote[1]{Foruden mine egne Studier har jeg her havt megen Nytte af en Afhandling (Addison’s Humour, its matter and its form) af den unge, tidligt bortgangne Literaturforsker \textit{Mary Suddard} i hendes \textit{Studies and Essays} (1912).}

Ordets Betydning er først fysisk: Fugtighed, Væde, Væske. Derefter
bruges det i den gamle Fysiologi om de Væsker i Legemet,
der bestemme Temperamentets Beskaffenhed. De bekendte fire
gamle Temperamenter svarede hvert til sin Væskes Fremhersken.
Fra Fysiologien er Ordet gaaet over til Psykologien for at betegne en
vis sjælelig Tilstand, og den psykologiske Sprogbrug viser saa
% --- p. 41 ---
\setcounter{footnote}{0}%
\opage{41}igen en Række af Overgange, som vidne om Forandringer i Følelseslivet
i nyere Tid, eller dog om en klarere Opfattelse af Følelseslivets
Fænomener. Det er naturligvis ikke Navnet, det kommer
an paa. Altfor ofte tro Psykologer endnu at kunne forklare os,
hvad et vist Navn paa en Følelse eller et andet sjæleligt Fænomen
(Opmærksomhed, Villie, Egoisme, Religion o. s. v.) „egentligt“
betyder. Det Afgørende er at faa en Følelse beskreven og karakteriseret.
Saa kan man bagefter se at finde det Navn til den, som
er mest hensigtsmæssigt, naturligvis med alt muligt Hensyn til
den overleverede Sprogbrug.

Hvad nu Ordet Humor angaar, ses det tydeligt, hvorledes den
fysiologiske Sprogbrug har bestemt den psykologiske.\footnote[1]{Overgangen fra den fysiologiske til den psykologiske Betydning synes at være foregaaet i Frankrig i Begyndelsen af det 16. Aarhundrede. Ifølge \textit{Kr. Nyrop (Grammaire Historique de la Langue Française.} IV. 1913. p. 228) anføres den psykologiske Betydning af Henri Etienne som en Nyhed.} Saadanne
Karakterytringer, i hvilke man særligt mente at kunne spore Virkningen
af en enkelt af Legemets „Væsker“, kaldte man i England
„humorous“. Man tænkte derved særligt paa Egenskaber som
stædig Ensidighed eller Lunefuldhed eller Excentricitet. Hos to
yngre Samtidige af Shakespeare fremtræder denne Sprogbrug.
Robert Burton (Anatomy of Melancholy. 1621. I, 3, 1, 2) beskriver
Mennesker, der ere „humorous“ saaledes: de le snart af Hjertens
Grund, snart græde de uden ydre Grund, snart sukke de og
ere bedrøvede og eftertænksomme, snart hilde de sig i urimelige
Forestillinger. Ben Jonson giver i Prologen til sin Komedie „Every
man out of his humour“ følgende Forklaring: „Naar en enkelt
Egenskab i den Grad raader hos et Menneske, at den faar alle hans
Tanker og Evner til at gaa en og samme Vej, da maa et saadant
Menneske siges at have Præget af Humor (to be a humour).“ Ordet
„humourist“ bruges hos Ben Jonson om den, der er „humorous“
(Act. I. Sc. 3). Jonson fremfører sin Forklaring af Ordet som en
Protest mod den Misbrug, at enhver Nar eller Laps kaldte sig „humorous“.
Han vilde begrænse Betydningen til en udpræget, temperamentfuld
Egenskabs Fremhersken.

Naar Hertug Frederik i „As you like it“ siges at være „humo-
% --- p. 42 ---
\setcounter{footnote}{0}%
\opage{42}rous“, maa efter Sammenhængen den Egenskab, der tænkes paa
være Stædighed. Naar Hamlet ved Efterretningen om Skuespillernes
Ankomst bl. A. lover, at „the humorous man“ skal faa Lov
at udføre sin Rolle i Fred, tænkes der sandsynligvis paa, hvad vi
vilde kalde en Karakterskuespiller, særligt paa en Fremstiller af
paafaldende, extravagante Karakterer. Dowden tænker (i sin Udgave
af Hamlet, Noten til II, 2, 340) paa Karakterer som Jacques
og Mercutio.

Hvad enten man ved Ordet humorous nærmest tænkte paa Vægelsind\footnote[1]{\textit{Holberg} forudsætter hos „den Vægelsindede“ en pludselig Skiften af Stemninger, uden at den foregaaende faar nogen Indflydelse paa den følgende Karakter. Dette maa vel forstaas saaledes, at Lucretia sættes i Stemning ved den mindste Anledning, og at den saaledes vakte Stemning kommer til at farve hele Tilstanden. Paa en interessant Maade har \textit{Johanne Louise Heiberg (Et Liv genoplevet i Erindringen}\textsuperscript{3} II p. 27—31) hævdet Nødvendigheden af mimisk Motivering af Overgangen fra et „Sind“ til et andet og henvist til de smaa Antydninger af saadanne Overgange, Texten ved nærmere Undersøgelse frembyder. --- Efter Burton’s Beskrivelse vilde Lucretia være at kalde „humorous“: „sometimes profoundly laughing, and then again weeping without a cause“ (Anatomy of Melancholy I, 3, 1, 2). --- En af successiv Dobbeltbevidsthed lidende Patient, som \textit{Morton Prince} har iagttaget, og som efter sin Helbredelse selv har beskrevet sin Sygdom \textit{(B. C. A.: My Life as a Dissociated Personality.} Boston 1909) kunde pludseligt og uden ydre Anledning faa en helt anden Karakter, der ytrede sig i Miner, Optræden, Tale o. s. v., og det saaledes, at den ene Tilstand ikke havde nogen Erindring om den anden.}
eller paa Ensidighed, betyder det en Enkeltfølelses Herredømme.
Det har endnu ikke faaet den Betydning af en sammensat,
gennem modsatte Elementers Brydning opstaaet Følelse, som forberedes
i Løbet af det attende og bestemt fastslaas i det nittende
Aarhundrede.\footnote[2]{Se nærmere herom nedenfor § 47—50.}

Det er Shakespeares Kunst, der for en væsentlig Del har banet
Vejen for denne nye Betydning og for Forstaaelsen af det ejendommelige
sjælelige Fænomen, der ligger til Grund for den. Udadtil
fremtræder denne Forstaaelse i den geniale Maade, paa hvilken
han har indarbejdet det ældre Skuespils Clownvæsen i en stor tragisk
Handlings Helhed. Livets modsatte Strømme forenes i
Shakespeares Kunst til en mægtig Flod. Spøg og Alvor, Latter og
% --- p. 43 ---
\setcounter{footnote}{0}%
\opage{43}Graad, Højhed og Lavhed, --- Alt yder sit Bidrag til den Totalstemning,
som Værkerne efterlader, og som sandsynligvis har været
tilstede hos Digteren selv som Frugt af hans Livserfaring. Han
havde ikke kunnet løse sin kunstneriske Opgave, naar han ikke
i sit Liv havde løst de tilsvarende praktiske Opgaver. Digteren behøver
naturligvis ikke selv at have oplevet de Modsætninger, han
skildrer. Men han maa have set dem som Muligheder. I de Karakterer
og Skæbner, han iagttog, i de Stemninger og Tilskyndelser,
han har erfaret hos sig selv, i de Tilskikkelser, der ere overgaaede
ham, maa for ham have aftegnet sig Antydninger, Kløvningslinier,
Muligheder, overfor hvilke det kunde blive et alvorligt Arbejde at
hævde Sindets Kraft og Harmoni. Hvad der er naat gennem dette
Arbejde, har han saa formet i en Verden af Skikkelser og Skæbner,
der for Tilskuer eller Læser staar som et af de største Billeder af
Livet, der nogensinde er givet.

Kun ud fra en dyb Forstaaelse af Menneskenaturen i dens mangfoldige
Former og af Menneskelivet i dets skiftende Skæbner kan
der dvæles saaledes ved det Ringe, Lave, Urimelige, at det trods
Alt ses som Element i det Hele, og kun paa et saadant Grundlag
kan Spøgen have sin Gænge, skønt den store Alvor danner det
sidste Grundlag.

I Humoren --- i den specielle Betydning, i hvilken vi her have
med den at gøre --- ytrer sig paa én Gang et stort Blik paa Livet
og en ærlig Erfaring om Livets Smaaligheder og Ulykker. Totalfølelsen kan ogsaa her fremtræde enten som Følge af Sammensmeltning
eller som Følge af Organisation. Hist ere Modsætningerne
bundne, ere kun potentielt tilstede; her træde de mere frit op, men
forbindes ved det fælles Forhold til en Grundværdi.

Shakespeares Humor er, som Dowden har sagt, et Udtryk for
hans Troskab mod Virkeligheden. Han har ikke mistet Troen paa
det Store, fordi Livet frembyder saa megen Lavhed og saa megen
Lidelse; heller ikke har han i sin Forvisning om Værdien af det
Store mistet Sansen for det Smaa og det Smertelige.

Den store spanske Digter, hvis „Don Quixote“ udkom paa samme
Tid som Shakespeares „Hamlet“, betragtes af Nogle som Mesteren
for al humoristisk Poesi, fordi hans Spøg altid er Alvor, og
% --- p. 44 ---
\setcounter{footnote}{0}%
\opage{44}hans Alvor stedse træder op i Spøgens Form.\footnote[1]{Saaledes \textit{C. H. Weisse: Kleine Schriften zur Aesthetik} 1867.} I sin Fremstilling
af den vandrende Ridder og hans Kamp med Vejrmøller, med en
Sancho Pansa som Vaabendrager og Chorus, har Cervantes stedse
Forstaaelse og Medfølelse overfor det gode Hjerte og den oprigtige
Villie, skønt det først er efter Illusionernes Sammenbrud, at
disse Egenskaber ret træde frem hos hans Helt.

De her antydede Elementer i, hvad jeg vil kalde den store Humor,
skulle i de følgende Afsnit nærmere belyses i deres forskellige
Nuancer og i deres forskellige indbyrdes Forhold. Den Alvor,
der ligger bag Spøgen, kan skyldes Sympati og Forstaaelse, eller
Vemod og Længsel, eller Højsind og Overlegenhed. Vi ville komme
til at dvæle noget nærmere ved disse Typer. ---

Der gives ogsaa en lille Humor, den mest populære Form, som
er Et med mere eller mindre godmodig Spøg. Godmodigheden
kan have mange Grader, gennem hvilke Humoren kan gaa over
til Ironi, Satire eller Haan. Det er ofte denne lille Humor, der
tænkes paa, naar Ordet Humor eller Humorist nævnes. Den er
en Enkeltfølelse, ikke en Totalfølelse. Det er en enkelt Oplevelse,
der tages op til Behandling i let Spøge-Form, i Regelen med en
Understrøm af Forstaaelse og Sympati, og saaledes at den Spøgende
meget vel kan drage sin egen Person og sin egen Færd med
ind blandt sine Emner. --- Sprogbrugen tillader dog ogsaa at tale
om Humor som Evne til at fremkalde Latter, ganske afset fra det
Sindelag, der ligger bag ved. Det er den lille bitte Humor.

Den lille Humor kan blive en ren Form, der kan anvendes ud
fra højst forskellige Motiver. Den har ikke det dybe Grundlag som
den store Humor, der er en Livsanskuelse eller rettere et Sindelag
overfor Livet, en Totaltilstand, til hvilken alle Oplevelser og
Bestræbelser have ydet deres Bidrag. Den lille Humor kan benyttes
af Haan og Uvillie som en blot Form; den i Spøgen skjulte Alvor
er da forstilt. Ogsaa Ironikeren kan, som det senere vil vise
sig, bruge den lille Humor. Ofte kan der bag denne ligge en Trang
til at lege og gennem Legen frigøre sig for, hvad der trykker Sindet.
I det enkelte Tilfælde kan det være vanskeligt at afgøre, om
% --- p. 45 ---
\setcounter{footnote}{0}%
\opage{45}det er den lille eller den store Humor, man har for sig, og i første
Tilfælde, i hvilken Tjeneste den staar. I det Følgende vil den lille
Humor kun undtagelsesvis komme paa Tale, da det væsentligt er
Humor som Livsanskuelse eller Livstype, Livsstade, der er Emnet
for denne Afhandling.\footnote[1]{For at prøve, hvilken Betydning af Ordet Humor, der nutildags maatte ligge literaturkyndige Mænd nærmest, bad jeg tre Kolleger, der hver fra sin Side beskæftigede sig med Literaturforsken, at sige mig aldeles improviseret, hvilken Forestilling Ordet nærmest vakte hos dem. Det ene Svar gik nærmest i Retning af den lille Humor; det andet opfattede Humor som komisk Sans paa Grundlag af Sympati og Forstaaelse; det tredie karakteriserede Humor som en Livsanskuelse, der betinges ved, at Livet i Tiden ses med Evigheden som Baggrund. Alle tre Former, og flere til, kommer jeg i det følgende til at omtale. Det første Svar gav Ordets populære Betydning, det andet den senere engelske Sprogbrug, det tredie tog Humor i Romantikens og Kierkegaards Betydning.}

11. Paa Forhaand maa den Misforstaaelse afværges, at det
skulde være Meningen at godtgøre en bestemt Livsanskuelses Eneberettigelse.
Det vilde være en umulig Ting. Forsaavidt er det rigtigt,
hvad der ofte siges, i Tide og i Utide (især i Utide), at om
Livsanskuelser kan der ikke disputeres. Ved Livsanskuelse menes
ikke en Teori eller en Opfattelse, som man blot \emph{har}, men et Stade,
man staar paa, eller et Liv, man faktisk \emph{lever}. Den er ikke et
Syn, man ser, men et Øje, som ser. Man behøver ikke at være
sig sin Livsanskuelse bevidst. For Filosofien, særlig for Psykologien
kunne Livsanskuelser kun være Genstande for Beskrivelse
og Undersøgelse. Filosofiens Forhold til Livsanskuelser er
ikke udtømt med den Bemærkning, at de ere forskellige eller modstridende.
Den har den Opgave positivt at analysere dem, prøve
deres Konsekvens og finde deres Oprindelse og deres Udviklingsgang.
Der er for Filosofien her ikke blot en psykologisk, men ogsaa
en erkendelsesteoretisk og en etisk Undersøgelse at gennemføre.
Det er af Vigtighed, at Filosofien ikke skyder denne Opgave
fra sig. Ikke blot indskrænker den derved sit Omraade, men naar
Livsanskuelser ikke tages op til objektiv Behandling, liste de sig
let frem ad Omveje og Bagveje og øve en Indflydelse, der kan være
des større, jo mere ubevidst den gaar for sig. Filosofien har stadigt
det sokratiske „Kend dig selv!“ til Motto. Og at kende sig selv
% --- p. 46 ---
\setcounter{footnote}{0}%
\opage{46}betyder at være sig sine Forudsætninger, sit Stade bevidst i saa
høj Grad, som det overhovedet er muligt.

Opgaven for nærværende Afhandling er den at undersøge den
psykologiske Mulighed for et paa den store Humor grundet Livsstade.
Naar Psykologien beskæftiger sig med Sjælelivets mere konkrete
og komplicerede Tilstande, maa den være klar over, at den
har med Typer, individuelle Former, særegne sjælelige Helhedsdannelser
at gøre. Hvad den raader over af almene Synspunkter
eller Hypoteser, kan anvendes til, i deres Sammenspil at kaste Lys
over saadanne Dannelser; men det vil være umuligt at bevise Nødvendigheden
af, at disse Dannelser overhovedet opstaa og bestaa.
Man kan kun undersøge Betingelserne for en sjælelig Helhed, naar
den virkeligt foreligger. En saadan Undersøgelse kan afgive gavnlige
og ofte højst nødvendige Korrektiver overfor Tilbøjeligheden
til at anvende en Gang for alle opstillede Skemaer paa det virkelige
Sjælelivs Omraade. ---

Det er maaske heller ikke overflødigt at bemærke, at nærværende
Afhandling er psykologisk, ikke æstetisk. Begrebet Humor
spiller ganske vist en Rolle i Æstetikens Historie, og fra den poetiske
Literatur kan der hentes gode Exempler til at belyse det. Men
saadanne Exempler skulle i denne Afhandling kun benyttes, hvor
der er Sandsynlighed for, at det er Livserfaring, der ligger til Grund
for den literære Fremstilling. Allerede Henry Home, der i sine
„Elements of Criticism“ (1762) ch. XII gaar ud fra den ældre, ovenfor
omtalte Betydning af Humor, skelner imellem Forfattere, der
benytte humoristisk Form, og saadanne, der ere „Humorister i Karakter“,
og hvis Værker derfor uvilkaarligt faa Humorens Præg.
„Swift og Fontaine vare,“ siger han, „Humorister i Karakter, og
deres Skrifter ere fulde af Humor. Addison var ikke Humorist i
Karakter, og dog raader der i hans Prosaskrifter en saa fin Humor“.\footnote[1]{Til denne Opfattelse af Addison kommer ogsaa Mary Suddard i sin ovenfor nævnte Afhandling. Humoren var for ham væsentligt en Fremstillingsform, tildels et Sikkerhedsmiddel. --- Om Swift med Rette er at kalde Humorist, i snevrere Betydning, skal senere drøftes.}% the print has no reference mark for this note (read at the image): set at the sentence it comments
En nyere fransk Literaturforsker har bemærket, at medens
man i „klassiske“ Værker kan undgaa at tænke paa Forfatteren,
røbe de „humoristiske“ Værker nødvendigvis Forfatterens Tem%
% --- p. 47 ---
\setcounter{footnote}{0}%
\opage{47}perament og Karakter. (Baldensperger: Leçon d'ouverture). Dette
kan kun gælde, hvor Humor er Andet og Mere end en Form. Literaturen
udspringer af Livet, og hvis det Livsstade, der kaldes den
store Humor, har nogen Betydning, er det ikke, fordi det har fundet
Udtryk i Literaturen, men omvendt: det giver Literaturen forhøjet
Værd, at et betydningsfuldt Livsstade har kunnet finde Udtryk
gennem den. Og det er muligt at lære dette Livsstade at kende
ad andre Veje end Beskæftigelse med Literaturen. ---

Vi ville i det Følgende komme til at give Bidrag til den store
Humors Historie. Begrebsmæssigt er den først bleven karakteriseret
af tyske Æstetikere i Begyndelsen af det nittende Aarhundrede,
og her er i første Linie Solger at nævne (i hans Dialog „Erwin“
1815 og hans efterladte Skrifter). For Solger fremtræder i
Humoren paa én Gang Følelsen af Tilværelsens Herlighed og af
dens Ufuldkommenhed, hvorved der betinges en indre Sammenhæng
mellem det Tragiske og det Komiske, saa at Intet staar som
værdifuldt, uden at det ved sin Begrænsning kan blive komisk,
Intet som komisk, uden at erkendes i sin Værdi! (Erwin. Ny Udg.
1911. p. 350---354).

Før Solger maa nævnes Jean Paul (Vorschule der Aesthetik.
1804), senere især Zeising (Aesthetische Forschungen. 1855), Lazarus
(Leben der Seele. 1855), George Eliot (Essays 1883) og
Th. Lipps: Komik und Humor (1898).

I dansk Literatur er Begrebet Humor smukt udviklet efter dets
æstetiske Side af Sibbern og Wilkens i deres Værker over Æstetik.
Men jeg vil i denne Undersøgelse især holde mig til Kierkegaards
Behandling af Emnet, i mine Øjne et Bidrag af højeste
Rang. Om hans Forhold til Solger skal der senere tales. Kierkegaard
beskriver Humor (i Betydning af den store Humor) som et
af sine „Stadier“ (Uvidenskabelig Efterskrift. 1846. --- Samlede
Skrifter VII). Det har været mit Ønske at føre videre, hvad Kierkegaard
har antydet i faa Træk, men tillige at godtgøre det Uberettigede
i de Skranker, han ud fra sine Forudsætninger sætter for
dette Stades Betydning, baade hvad Omfang, hvad Indhold og hvad
Energi angaar.

% --- p. 48 ---
\setcounter{footnote}{0}%
\opage{48}\begin{center}II. LATTER OG HUMOR\end{center}


12. I Anledning af den ældre Betydning af Ordet Humor gjorde
allerede Henry Home opmærksom paa, at en Karakter kun kaldes
„humorous“, naar den virkede komisk. Ogsaa efter den senere
Betydning af Ordet har det Noget med Latter at gøre, eller i hvert
Tilfælde med Smilet. Lad os begynde Undersøgelsen af Humor
med at betragte dens Forhold til Latteren.

Latteren kan undersøges fra to Sider. Der kan spørges om Beskaffenheden
af de ydre Oplevelser, der vække Latter, og der kan
spørges om den indre Stemning, der faar Udtryk i Latter. I Regelen
har man lagt mest Vægt paa det første Spørgsmaal, og det er især
det, Teorier om det Latterlige søge at besvare. Men for vor Undersøgelse
er det andet Spørgsmaal af størst Betydning. Som Indledning
maa vi dog berøre det første, især da de to Spørgsmaal tilsidst
ikke kunne holdes ude fra hinanden.

Der er her i nyere Tid optraadt to Teorier, som kunde synes at
staa i skarp Modsætning til hinanden. --- Den ene lægger Vægt paa
at naar Noget er latterligt, er det, fordi det først fremtræder med
en vis Tilforladelighed, Selvfølgelighed, Betydelighed, Overlegenhed,
men saa pludseligt aabenbarer sig i sin Intethed og sin Afmagt.
Derved fremkaldes der hos den, for hvem det træder op,
først en Følelse af Tryghed eller Beundring eller Betagethed og
Spænding, og derpaa pludseligt en Skuffelse eller en Hoveren, eller
dog en Lettelse og en Løsning af Spændingen. Hvad der først stod
saa selvsikkert og imponerende, faar nu paa én Gang den modsatte
Karakter. Det Latterlige beror efter denne Opfattelse paa en Kontrastvirkning
i nedadgaaende Retning. I korte, kraftige Træk er
denne Teori allerede fremsat af Thomas Hobbes. I min Psykologi
% --- p. 49 ---
\setcounter{footnote}{0}%
\opage{49}(VI E, 9) har jeg sluttet mig til den, kun med særlig Betoning af
den forskellige Karakter, Latteren faar efter den Leendes Grundstemning.
James Sully (Essay on Laughter. 1902) kalder denne
Teori for Degradationsteorien. --- Den anden Teori lægger Vægt
paa, at det Latterlige altid frembyder en Urimelighed, en Uoverensstemmelse,
en Selvmodsigelse. Den er i dansk Literatur fremsat
af Kierkegaard. Sully kalder den Inkongruensteorien.

Baade Sully og Ribot (Psychologie des Sentiments p. 344)
mene, at de to Teorier passe paa hver sin Gruppe af Tilfælde, og
at det er umuligt at finde en Formel, der kan passe paa begge
Grupper. Utvivlsomt fremhæve de to Teorier hver for sig vigtige
Synspunkter. Og de føre begge til at forstaa Latterens Betydning
i Livet, f. Ex. til at forstaa, hvorledes man har kunnet opfatte Latteren
som Sandhedskriterium (test of truth), som Shaftesbury og
Henry Home have gjort. Gennem Latterens Skærsild udskilles
alt Kunstigt og udefra Paahængt ved et Emne; det kommer til
at fremtræde i sin hele Nøgenhed, hvorved enten dets Værdiløshed
eller den Værdi, det maaske trods Alt har, kan fremtræde
saa meget des stærkere. Holberg opfattede sin Virksomhed som
komisk Digter som en Bestræbelse for at skille Skyggen fra
Legemet.

I Virkeligheden ere dog de to Teorier ikke hinanden modsatte.
Inkongruensteorien er aabenbart en særegen Form for Degradationsteorien.
Modsigelse er jo en speciel Form for Modsætning.
At blive aabenbar i sin indre Uoverensstemmelse og Selvmodsigelse
er jo dog en „Degradation“. Rent logisk er det Selvmodsigende
= 0. Men psykologisk virker Opdagelsen af Selvmodsigelse
som en Kontrast i nedadgaaende Retning. Det, der fremtraadte
som noget Fornuftigt og Vederhæftigt, viser sig nu i sin
Intethed. Derved fremkommer det Choc, som er Betingelse for
Udløsning af Latter. Mellem alle Tilfælde af det Latterlige er der
paa dette Punkt en Analogi, som grunder sig paa en fælles Betingelse.
Spændingens Art og Grad kan være højst forskellig. En
opblæst Person, der pludseligt afsløres som indholdsløs og tankeløs,
ligesom en oppustet Blære, der stikkes Hul i, staar naturligvis
for os paa en anden Maade end en dybsindig Tænker, der i Optagethed
af sit Emne begaar en Regne- eller Tænkefejl paa et
% --- p. 50 ---
\setcounter{footnote}{0}%
\opage{50}underordnet Punkt, selv om dette underordnede Punkt var et nødvendigt
Led i Tankegangen. Man kan naturligvis formulere saadanne
Kontraster som Modsigelser: Fylde-Tomhed; Dybsind-Dumhed.
Men Modsigelsen er her tillige Modsætning, og kun
derved bliver Fænomenet latterligt. Naar Sancho Pansa med Iver
klamrer sig til en Grøftekant, som han i Mørket anser for Randen
af en Afgrund, kan det naturligvis godt formuleres som en Modsigelse,
en „Inkongruens“. Sancho’s store Angst passer unægteligt
ikke til Situationen. Men det er ikke Urimeligheden i sig
selv, der virker komisk; det er dens Modsætning til det Minimum
af Fornuft, vi stedse forudsætte. Se vi et Menneske i Dødsangst,
stole vi foreløbigt paa, at der er Grund til det. Vi tage foreløbigt
Alt for gode Varer; og naar Angsten pludseligt ophæves, er der
Mulighed for Latter. Der var forud for Opløsningen ogsaa
Mulighed for Graad; det Tragiske og det Komiske kunne ligge
hinanden nær. Latteren forudsætter, at de, som le, stadigt ere
trygge og i god Behold paa Fornuftens og Magtens Side, ogsaa
efterat „Degradationen“ er indtraadt. Der maa være Noget, som
ikke opløses ved den nedadgaaende Kontrast, --- i hvert Tilfælde
det Stade, hvorpaa den Leende befinder sig. Ved det Tragiske ere
Forholdene mere indviklede, som vi senere ville se.

Følelsen ved det Latterlige er et udpræget Exempel paa en
Forholdsfølelse, d. v. s. en Følelse, der ikke blot, som al Følelse,
betinges ved sit Forhold til andre Følelser, men som i sit inderste
Væsen bestemmes ved, at den Oplevelse, der fremkalder den,
bestaar i en Kontrast. Ved Latterfølelsen er denne Kontrast nedadgaaende;
ved Følelsen af det Ophøjede er den opadgaaende. I
Latteren føle vi os store overfor det Smaa; ved det Ophøjede føle
vi os smaa overfor det Store.

Den Kontrast, der her spiller en Rolle behøver ikke at være
en Kontrast mellem to bestemt udformede Bevidsthedstilstande.
Der kan i en og samme Oplevelse være indeholdt en Kontrast,
som pludseligt gør sig gældende for os, uden at der kan iagttages
nogen Succession. Ovenfor er der allerede omtalt et Exempel fra
et andet Omraade herpaa, nemlig den Stemning, Velasquez’ Dvergbilleder
fremkalde. Naar --- for at tage et ofte benyttet Exempel
--- Synet af et lille Barn med Mandshat virker komisk, er Kon%
% --- p. 51 ---
\setcounter{footnote}{0}%
\opage{51}trasten ikke mellem Barnehoved og Mandshat, men mellem Barnet
med sin sædvanlige Hovedbedækning og Barnet med Mandshatten,
og Forestillingen om den sædvanlige Hovedbedækning
behøver slet ikke at komme udtrykkeligt frem; den kan virke
som en stiltiende, tilvant Forudsætning. Der kan meget vel opstaa
Kontrastvirkninger til Noget, der ikke selv er oppe over Bevidsthedens
Tærskel. Der gives stille, trygge Forventninger, som ere
saa indgroede, at de kunne danne Betingelse for Kontrastvirkning,
en Betingelse, vi maaske først blive opmærksomme paa, naar vi
forundres over Kontrastfænomenet og søge at forklare os det. Det
gaar her, som naar vi undre os over, at Noget staar for os som
bekendt, og vi saa undersøge, hvor og naar vi have set eller
hørt det.

13. Latter er som det Pust, der fjerner et Fnug. Forholdet mellem
Pust og Fnug kan være højst forskelligt. Der kan her opstilles
en hel Skala.

a. I den barnlige Latter, den „kaade Overgivenhed og Fryd for
Ingenting“, behøver der slet ikke at være noget Fnug, nogen ydre
Anledning. Overstrømmende Livskraft og Livsfølelse fører til at
forholde sig legende og spøgende til Alt. Latteren er her det
sunde og kraftige Sinds uvilkaarlige Ytring. Der er hverken nogen
Bevidsthed om indre Tryghed og Overlegenhed over Forholdene
eller om ydre Anledning. Dette Uvilkaarlighedsmoment forsvinder
aldrig ganske fra Latteren, saalænge Latter overhovedet er
mulig\footnote[1]{\textit{Sully} (\textit{Essay on Laughter} p. 289—293; 427—432) er bekymret for den sunde, spontane Latters Fremtid paa Grund af Pessimisme, Træthed og Cynisme. The reservoir which in the past supplied the stream of national gaiety, has certainly fallen and threatens even to dry up. Gammel Englands Munterhed udviklede sig efter den fælles Kamp imod Rom og mod Spanien. Men nu træde Klassemodsætninger i Forgrunden. --- Maaske ville en ny Tids Kampe kunne virke genfødende gennem at udløse og samle Kræfterne.}. Al Latter forudsætter fast Grund under Fødderne, hvor
begrænset eller illusorisk denne Grund end maaske kan vise sig
at være for en rent objektiv Betragtning.

b. Nærmest i Slægt med denne spontane Latter er den, der
skyldes Lettelse af et Tryk eller Befrielse for en Fare, som har
% --- p. 52 ---
\setcounter{footnote}{0}%
\opage{52}truet det ydre eller det indre Liv. Her er der udtrykkelig Bevidsthed
om Fnugget, om Grunden til Latteren.

Ved et Dampskibs Afgang staar en Dame paa Kajen med en lille
Dreng for at tage Afsked med Bortrejsende. Hun holder nu Drengen
ud over Bolværket og hæver og sænker ham som Hilsen til de
Bortdragende, idet hun raaber: „Gustav siger Farvel, Farvel,
Farvel!“ Hun taber ham i Vandet. Almindelig Bestyrtelse og Fortvivlelse.
Men en Matros springer ud og faar fat i Drengen, og
medens han endnu er nede i Vandet, løfter han Drengen op og
raaber: „Gustav siger Goddag, Goddag, Goddag!“ En befriende
og jublende Latterbølge afløste Spændingen og Angsten. Tillige
laa der til Grund en vis Hoveren over den taabelige Moders
Adfærd og en sympatisk Glæde over, at et Menneskeliv var reddet.
Herved hører dette Exempel tillige ind under de to følgende
Typer.

c. En endnu større Rolle end ved Befrielsens Latter spiller
Latterens Genstand, hvor Latteren er Udtryk for Haan. Genstanden
betyder her en Modstand, der syntes eller mente sig at være
overmægtig, men som overvindes og falder til Jorden. Medens
ved den spontane Latter Genstanden ingen Betydning har, og
medens den ved Befrielsens Latter er Noget man bevæger sig
bort fra og maaske snart glemmer, saa spiller Genstanden ved
Haan en positiv Rolle. Der er en Trang til at fastholde Forestillingen
eller maaske endog Synet af Genstanden for ret at kunne
nyde sin egen Magtfølelse. Her gør sig ikke blot en Reaktion
gældende mod Ens egen forudgaaende Spænding, men ogsaa mod
selve Genstanden. Foragten faar atter og atter Luft overfor den.
Fnugget er her blevet et Bjerg, man har væltet.

d. I Humoren, den fjerde Type for Latterfølelse, spiller Genstanden,
ligesom i Haanen, en væsentlig Rolle, kun at Forholdet
til den ikke er antipatisk, men mere eller mindre sympatisk. Der
kan her igen skelnes forskellige Former, forskellige Arter af Humor,
hvis Forskellighed beror paa den Betydning, Genstanden har.

α. Hvor Overlegenheden overfor en Modstander er stor og dybt
grundet, vil Arbejdet med at overvinde Modstanden staa som en
Leg, og der vil ikke være nogen Trang til atter og atter at dvæle
ved den. Den virkelige Overlegenhed hviler i sig selv, og dens
% --- p. 53 ---
\setcounter{footnote}{0}%
\opage{53}Selvfølelse faar ikke Karakteren af Foragt overfor Andet og Andre.
Trangen til at føle Haan eller Foragt og give dem Udtryk er netop
et Tegn paa, at man er afhængig af den ydre Genstand og kun
kan føle sig selv ret i Modsætning til den. Selv om der ikke føles
positiv Sympati for den, kan den dog meget vel anerkendes i sin
ejendommelige Værdi, og Smilet eller Latteren kan da ogsaa betinges
ved Forstaaelse. Der kan bygges Bro mellem den Leende
og Latterens Genstand.

β. Da den virkelige Overlegenhed raader over mere Energi,
end der behøves til at overvinde Modstanden og sikre sin egen
Bestaaen, er der her Mulighed for et Højsind, som ikke blot forstaar
Genstanden, men søger at skærme og frede om den, alt medens
dens Ringhed og Afmagt gør sig gældende. Højsindet bærer
med et Smil ikke blot sin egen Byrde, men Andres Byrde, idet
det føjer sin Kraft til den ringe Kraft, der har fremkaldt Smilet.
Anerkendelsen af det Ringes og Afmægtiges Værdi er her mere
positiv end, hvor den blotte og bare Overlegenhed raader. Der er
her et indre Baand, ikke blot en ydre Bro mellem Subjekt og
Objekt.

γ. Det er ikke blot Overlegenheden, der kan hvile i sig selv.
Hvor Sindet er optaget af et Værdifuldt, der er Genstand for
Erindring eller Haab, kan det indre Liv være saa rigt, at ydre
Genstande og Forhold overfor det staa som ringe og ubetydelige
uden derfor at miskendes i deres ejendommelige Værdi. Humoren
kan have Vemod eller Længsel til Grundlag. Vemodet er rettet
mod Fortiden, Længselen mod Fremtiden. Overfor begge stræbe
nye Tilskikkelser forgæves at vinde Magt; men derfor haanes de
ikke og betragtes ikke med Ligegyldighed. I Kraft af Følelsens
Expansion (sml. § 6) vil Individets egen Stemning kunne komme
til at optage de Andres Interesse og Stræben under sin Indflydelse,
især naar disse Andre selv ere henviste til Erindringens
eller Haabets Verdener. Vemod og Længsel støde ikke bort, men
kunne føre til dyb Forstaaelse af Alt, som savner og længes.

δ. En særegen Karakter faar Humor, naar en Forstaaelse af
Livets og Tilværelsens store Love danner Baggrunden. Denne
Forstaaelse behøver ikke at være strengt videnskabelig. Et Menneske,
der bruger sin Eftertanke og har omfattende Erfaring, vil
% --- p. 54 ---
\setcounter{footnote}{0}%
\opage{54}naturligt naa til den Overbevisning, at der er en stor Tingenes
Orden, af hvilken alle Begivenheder, selv de, der gribe mest ind
i Menneskers Ve og Vel, ere afhængige. Og denne Overbevisning
vil afgive en Baggrund, en Horizont, i Henhold til hvilken det,
der optager Menneskers Interesse, kan staa som ringe og forsvindende.
Overlegenheden er her betinget ved Forstaaelse. Men
netop fordi det Ringe og Forsvindende ligesaa vel som det Store
hører ind under Tilværelsens lovmæssige Sammenhæng, netop
derfor kan det betragtes i sin Ejendommelighed og med Forstaaelsens
Sympati.

ε. Inderligst bliver Forholdet paa Latterens Omraade mellem
Subjekt og Objekt, naar Subjektet nærer en positiv Sympati for
Objektet. Forud for Smilet eller Latteren kan der være givet et
Baand, som ikke brister, fordi Objektet ses i sin Ringhed eller
sin Selvmodsigelse. Et saadant Baand er der mellem Forældre og
Børn, mellem Søskende, mellem Venner og overalt, hvor Menneskekærlighed
og hvor Sans for menneskeligt Sjæleliv raader.
Overlegenheden er her Kærlighedens Overlegenhed, og Kærligheden
er den højsindede Kærlighed (Etik IX, 1; XII, 2), der er
i Slægt med Retfærdighed. Forstaaelsen er her Sympatiens Forstaaelse
(medens den foregaaende Form for Humor kunde karakteriseres
som Forstaaelsens Sympati). Jo stærkere Baandet er, des
friere vil Smil og Latter kunne røre sig. De Kræfter, der ellers
kunne virke frastødende, --- Viddet og den komiske Sans, --- kunne
her komme til at vidne om Baandets Fasthed og maaske knytte
det fastere. Paa denne Art Humor passer Thackeray’s og George
Eliot’s Definition af Humor som Enhed af Vid og Sympati (wit
and love).

Alle disse Former for Humor kunne optræde som Øjeblikstilstande
(Sindsbevægelser eller Krusninger) eller som varigere
Tilstande (Stemninger) eller som Udtryk for det reale Jeg (Sindelag).
Denne sidste Form er, hvad jeg i denne Afhandling kalder
den store Humor.

I min Psykologi (VI E, 9) har jeg kun taget Hensyn til den
Form for Humor, der kan betegnes som Enhed af Vid og Kærlighed.
James Sully (Essay on Laughter. p. 306) indvendte mod min
Beskrivelse af Humor, at levende Sympati vil hindre Latter. Dette
% --- p. 55 ---
\setcounter{footnote}{0}%
\opage{55}kan ikke nægtes. Men om det sker, beror ikke blot paa Sympatiens
Grad, men ogsaa paa Graden og Arten af den Kontrast, der
er tilstede. Der kan jo være Modsætninger saa skarpe og bitre,
at de volde Saar, som ikke kunne læges. Der kan raade en overvældende
Skæbne, overfor hvilken al den Magt, som kan opbydes,
er Afmagt. Til denne Mulighed kommer jeg tilbage i en særegen
Undersøgelse, der er af afgørende Betydning for min hele Opfattelse,
nemlig Undersøgelsen af Forholdet mellem Humor og
Tragik (Kap. V).

Foreløbigt fremgaar det af den givne Karakteristik af de forskellige
Lattertyper, at hvad et Menneske ler af, beror paa, hvad
Slags Menneske han er. Et Menneskes Karakter aabenbarer sig
gennem hans Forhold til det Latterlige. Det have gode Iagttagere
til alle Tider set. Det vil ikke være uden Interesse at anføre
nogle Udtalelser herom. 

Platon forlangte,.at hvad man fandt latterligt, skulde ikke bestemmes % PRINTED AS IS: „forlangte,.at“
ved, hvad Sanserne vise os, men ved hvad Tanken sætter
højest (Staten V, p. 452 D).

Thomas Hobbes siger (Leviathan cap. 6): „Megen Latter er
Tegn paa en lille Aand (pusillanimitas). Ejendommeligt for en stor
Aand er det derimod at bringe Hjælp og at befri Andre for Foragt,
men at sammenligne sig selv med de største Mennesker alene.“

„Gennem mit hele Liv,“ siger Lichtenberg (Vermischte Schriften,
I p. 173), „har jeg fundet, at et Menneskes Karakter ikke
kan erkendes bedre end af en Spøg, han tager ilde op.“ 

Goethe siger (i sine „Maximen und Reflexionen“): „Menneskene
kundgør, ikke deres Karakter klarere end ved, hvad de finde
latterligt. Det sanselige Menneske ler ofte, hvor der ikke er
Noget at le af. Hvad der saa gør Indtryk paa ham, saa lægger
hans indre Velbehag sig for Dagen. Forstandsmennesket finder
næsten Alt latterligt, Fornuftmennesket næsten Intet.“ --- Og i et
af sine Epigrammer siger han (hvad der bekræfter Lichtenberg’s
Ytring og er af Betydning for Humorens Psykologi og Etik): „Den,
der ikke kan have sig selv til Bedste, hører sikkert selv ikke til
de Bedste.“

Carlyle siger (Sartor Resartus I, 4): „Intet Menneske, der én
Gang har let af sit fulde Hjerte, kan være helt uforbederligt slet.
% --- p. 56 ---
\setcounter{footnote}{0}%
\opage{56}Hvor Meget ligger der ikke i Latter --- den Chiffernøgle, hvormed
vi kunne dechiffrere hele Mennesket!\dots{}. Det Menneske,
der ikke kan le, er ikke blot istand til at øve Forræderi, List og
Rov; men hele hans Liv er allerede Forræderi og Bedrag.“

14. I det Foregaaende er Humor bleven beskrevet som et enkelt
Led i den hele Række af Latterformer. Og fordetmeste behandles
Humoren under Latterens Psykologi; saaledes ogsaa i
den Fremstilling af Psykologien, jeg har givet. Men den adskiller
sig fra de andre Latterformer ved sit Grundlag, ved i højere Grad
end de at have Karakteren af Totalfølelse. Den forudsætter mere
end de andre Former Livserfaring og Eftertanke, en Skærsild, i
hvilken Livets Modsætninger og Tilværelsens stridende Kræfter
have øvet deres Indflydelse paa Sjælelivets inderste Kærne. Det
Store og det Smaa, det Ophøjede og det Lave, det Lyse og det
Mørke, Alt har afsat sit Spor, og Eftertanken har besindet sig paa
Livets almene Love.

Hvad enten Grundlaget er den Overlegenhed, der udelukker
Foragt, eller det Højsind, der som en naturlig Ting bærer Andres
Byrde, eller det Vemod, der er Frugt af rige, men smertelige Erfaringer,
eller den Længsel, der er rettet mod et stort Maal, eller
den Forstaaelse, der ser Alt indordnet i stor Sammenhæng, eller
den Sympati, der føler Livets Puls i det Ringeste som i det Største,
--- saa er der en Mulighed for Smil eller Latter overfor Alt,
der i Virkeligheden eller tilsyneladende optræder som Modsætning
til, hvad der optager Sindet. Livets Modsigelser, Disharmonier,
Skuffelser og Sorger fortone sig overfor det Grundlag og
den Baggrund, som bestemmer de enkelte Oplevelsers Virkning.
En Hovedsag for den Latterform, der har Humorens Karakter, er
det, at Forholdet til Livet er positivt, eller at den Leende eller
Smilende her i hvert Tilfælde tager sig selv med. Overfor den
store Maalestok der anlægges, forsvinder tilsidst den skarpe Modsætning
mellem det Store og det Smaa. I den store Humor høre
vi selv, med hele vor Personlighed og hele vor Stræben, med til
det, som finder sin Retfærdiggørelse ved Digterens Ord
\begin{quote}
„det Ringeste forsvarer\\
sin Plads i Livets Krans.“
\end{quote}

% --- p. 57 ---
\setcounter{footnote}{0}%
\opage{57}Et Exempel paa denne Side af Humorstadet kan hentes fra
Kaalund’s Digt „Naturens Graad“. Naturen længtes efter sine
Børn, der nu havde lukket sig inde i de store Byer. De strømmede
da ud til hende i det Frie. Men den store Moder harmedes
over at se disse pyntede Herrer og Damer, og Harmen fik Luft i
en Tordenregn, som jog dem tilbage til Byen i vild Flugt og under
alskens Kalamiteter. Men Digteren bekender: „jeg selv var med,
jeg rulled i en Drosche!“

Der er i den store Humor en Dobbeltbevægelse. Den er Udtryk
for en Selvhævdelse, der bygger paa Selvforstaaelse. Ordet
om Mesterskabet, der viser sig i Begrænsning, gælder her i fuldeste
Betydning. Ti kun den kan hævde sig selv, som klart véd,
hvor han staar, --- hvad han formaar, og hvad ikke, --- hvad han
er værd opadtil og nedadtil. Den Egenskab, Aristoteles kaldte
Megalopsyki, og som bestod i at være sig sin egen virkelige Værdi
bevidst, udtrykker den rette Selvhævdelse fra den positive Side.

Skrankerne følge da af sig selv, idet de bestemmes indefra, ikke
defra. De ligge ganske simpelt der, hvor Bevidstheden om eget % PRINTED AS IS: „defra“ (sense: udefra)
Værd hører op. Da Romantikerne i sin Tid forsøgte at stille Tieck
op som Goethe’s Jevnbyrdige, udtalte Goethe, at det var, som
om han vilde sammenligne sig med Shakespeare. Goethe var klar
over sin Grænse opadtil og nedadtil. Hans Ord om, at den, der
ikke kan have sig selv til Bedste, sikkert ikke hører til de Bedste, kan
bruges som Motto for den store Humor. Det at have sig
selv til Bedste er ikke Et med Selvhaan. Humoristen haaner
ligesaa lidt sig selv som han haaner Andre. Han holder sig selv
saavel som dem op imod den store Baggrund, Livet har stillet
ham overfor, og Forholdet til den bestemmer baade den Plads,
han anviser sig selv, og den Plads, han anviser dem. Overalt
søger han Kærnen blandt Skallerne. Og da Skallerne ikke altid
røbe Kærnen, vil han være tilbøjelig til at mene, at det er vor
Uvidenhed eller vor Mangel paa sympatisk Forstaaelse (eller paa
Forstaaelsens Sympati), der gør, at vi ofte mene at finde Skaller
uden Kærner. Det er jo ikke alle Nødder, der kunne knækkes,
saa at man kan faa at se, hvad der er inde i dem. Som Goethe
skrev i en ung Students Dagbog: „Vorherre har skabt Nødderne,
men han har ikke knækket dem.“ Der kan netop blive Brug for
% --- p. 58 ---
\setcounter{footnote}{0}%
\opage{58}Humor overfor dem, der mener at have knækket al Verdens
Nødder, skønt de i Virkeligheden kun have knækket deres egne
Tænder. Den store Humor vil være forbunden med stadig Søgen
og staar i Modsætning til al dogmatisk Visdom, hvad enten den
optræder i den sunde Menneskeforstands, Videnskabens eller Religionens Navn.

Det var i Romantikens Tid, at man fik Øje for, at der i et
Sindelag som det beskrevne laa Muligheden for en Livstype, altsaa
for Noget, som var Andet og Mere end en speciel Form for
æstetisk Følelse. Berettigelsen heraf trænger endnu til at hævdes
og begrundes. Herman Lotze har saaledes negtet, at Humor skulde
kunne anerkendes som sammenfattende Livsstemning. Den kan i
hans Øjne kun være en vis Art af komisk Stemning, og det en
Art, der har sin Fare, idet den let virker opløsende overfor alt
Ophøjet. (Geschichte der Aesthetik in Deutschland, p. 389). Og
uagtet Kierkegaard beskriver den store Humor med Sympati og
Forstaaelse, kan den for ham dog ikke betegne et højeste Stade;
det forbeholder han for den religiøse Maade at tage Livet paa.

Der er da al Grund til en nærmere Undersøgelse af dette
Emne.

% --- p. 59 ---
\setcounter{footnote}{0}%
\opage{59}\begin{center}III. IRONI OG HUMOR\end{center}


15. Hvad enten man betragter Humor blot som en særegen
Art af Latterfølelse eller som en Livstype, gives der ingen saa
karakteristisk Modsætning til den som Haan. Men der gives flere
Mellemformer mellem disse to Yderpunkter, og vi ville her betragte
nogle af de vigtigste.

Haan forudsætter, som vi have set, et Modsætningsforhold, som
ikke er tilstede ved den rent uvilkaarlige Latter. Haan er paa antipatisk
Maade optagen af sin Genstand, medfører en Trang til atter
og atter at beskæftige sig med den for at fremhæve og udmale
dens Afmagt, dens Ubetydelighed, dens Dumhed eller dens Nederdrægtighed.
Den er en afgjort reaktiv Følelse, skønt Nietzsche
mente, at den netop egnede sig for høj Overlegenhed. Vel begrundet
Overlegenhed fører, som ovenfor udtalt, snarere til Humor
end til Haan; den hviler i den Grad i sig selv, optagen af
sin rige Selvudfoldelse, at der ingen Trang kan opstaa til Sideblik
eller Nedblik, og hvor den staar i Forhold til andre Væsener,
har den i hvert Tilfælde ikke nogen Grund til Antipati. Det er et
Tegn paa Usikkerhed og paa Mangel af Opgaaen i Livsudfoldelsen,
naar det atter og atter er nødvendigt at dvæle ved Andres
Ringhed. Det var det Tragiske hos Nietzsche, at hans stadige Trang
til at udtale sine Antipatier ikke gav ham Ro og Kraft til positivt
at udvikle, hvad der dog mest laa ham paa Hjerte, Livets store
Magt og Selvbekræftelse. Atter og atter afbryder han Udarbejdelsen
af sit Hovedværk for at lade Haan og Antipati faa Ordet.
Værket blev da heller aldrig færdigt. Men mest grelt fremtræder
Modsigelsen hos ham, naar Afstandsfølelsen (das Pathos der Distanz)
overfor den store Hob skal være et væsentligt Moment i
store Menneskers Sindelag. Han saa ikke den skarpe Strid, der
% --- p. 60 ---
\setcounter{footnote}{0}%
\opage{60}er mellem den Pessimisme, „Afstandens Pathos“ nødvendigvis
forudsætter (idet „den store Hob“ kun skal kunne være Midler,
Hindringer eller Kopier!), og den absolute Optimisme, hans Evangelium
skulde forkynde\footnote[1]{Se min Karakteristik af Nietzsche i \textit{Moderne Filosofer}.}.

Hos store Satirikere kan man studere dette Forhold, og hos
Ingen bedre end hos Swift. Hans Optagethed og Betagethed af,
hvad han forkastede og bekæmpede, træder saa grelt frem, at hele
Skikkelsen ofte er bleven formørket derved for Eftertiden, ligesom
Swift’s eget Sind formørkedes derved. Han søger stedse nye
Udtryk for den Foragt, han følte for Menneskenes Mængde og
deres Færd, saaledes som han under sin Deltagelse (eller Forsøg
paa Deltagelse) i det politiske Liv mente at have lært den at
kende. Han svælger i Billeder og Allegorier, der kunne udtrykke
denne Følelse. Han lader sin Gulliver rejse til Lilliput for i dette
Tommelidenfolk at kunne latterliggøre Europæernes Kultur. Og
naar Gulliver saa kommer til Kæmpefolket Brobdingnag, bliver
der Lejlighed til gennem den Beretning, Gulliver der aflægger
om Livet i Europa, at skildre Europæerne som „den mest fordærvelige
Race af afskyeligt Utøj, Naturen nogensinde har fundet
sig i at lade krybe omkring paa Jordens Overflade“ (Gullivers
Travels. Part II. Chap. 6). Højdepunktet naas dog senere i Beskrivelsen
af det Folk, der kaldes „Houyhnhnms“, et idealt Folk, der
naturligvis ikke har Menneskeskikkelse, men Hesteskikkelse. I
dette Land bor der, foruden Hovedfolket, en hæslig, dyrisk, lavsindet
Race, der har Menneskeskikkelse og lever Menneskeliv. I
disse „Yahoo’er“ har Væmmelsen ved det Menneskelige fundet
sit voldsomste Udtryk.

Den store Satiriker har i disse Skildringer af Lavhed og Nederdrægtighed
ikke blot fundet Udtryk for sin moralske Indignation,
men ogsaa for krænket Selvfølelse, for indre Lidelser og for tragiske
Livsforhold. Det er, for at bruge Leslie Stephen’s Udtryk i
hans Bog om Swift (p. 184), en saaret Sjæl, der taler til os, ikke
blot en Profet, der revser Usselheden.

Maaske har enhver Profet en saaret Sjæl. Men paa Eftertiden
vil han dog virke mest ved sin Evne til at gøre det Store og det
Sande gældende. Det er jo kun den nyere Tids profetiske Skik%
% --- p. 61 ---
\setcounter{footnote}{0}%
\opage{61}kelser, hvis indre Liv vi kende. I de gamle Tider havde man ikke
Sans for individuel Psykologi, og i hvert Tilfælde var det kun,
forsaavidt Profeten i Harme og Begejstring udfoldede et stort
Sinds Kræfter, at han gjorde Indtryk paa den nærmeste Eftertid,
fra hvilken det Billede stammer, vi have af hans Skikkelse. Naar
han da var gaaet over til at være et af Slægtens store Forbilleder,
falde de rent individuelle Træk let bort, især de, der skyldes
Brydningen mellem indre, modstridende Drifter og Tendenser.
Hvis de havde været helt optagne af indre Strid, havde de vel
heller ikke egnet sig til Forbillede for Andre, men vilde have
havt nok at gøre med sig selv. Det er en stor, ofte uløselig Opgave
at finde den psykiske Sammenhæng mellem den vældige
Polemik og det, der er foregaaet i Profetens Sjæl.

16. Til Satirens Midler kan den ironiske Form høre. Den optræder
da med tilsyneladende Anerkendelse og Alvor, skønt dens
egentlige Stræben er at angribe og tilintetgøre. En indirekte Angrebsform
afløser da den direkte. Ironien er her kun Middel, betegner
ikke en særlig Livstype. Det er, hvad man kunde kalde
den lille Ironi, medens den store Ironi ikke blot er Form og Middel,
men et Livsstade. Forskellen vil bl. a. bestaa i, at medens det
ved Ironi som Form eller Talefigur maa tilstræbes, at de Andre
forstaa, at det er Ironi, at Alvoren ikke er alvorligt ment, forholder
det sig anderledes ved den store Ironi. Men før vi undersøge
den, maa vi dvæle noget ved Ironibegrebets Historie.

I græsk Sprogbrug og græsk Tænkning er dette Begreb undergaaet
mærkelige Forandringer, der paa en interessant Maade ere
paaviste af Ribbeck (Über den Begriff des εἴρων. Rheinisches Museum.
N. F. XXXI p. 386).

Ordet forekommer først i Aristofanes’ „Skyer“ som et Skældsord,
idet Ironikeren her stilles sammen med Løgneren og Praleren.
For Grækerne, der baade forudsatte og krævede, at det Indre
skulde være udtrykt i det Ydre, maatte Ironien, i hvilken det
Indre staar i Modsætning til det Ydre som Spøg til Alvor, være en
Uting. Naar Sokrates kaldtes Ironiker, var det et Udtryk for Folks
Uvillie og Forbitrelse over den Maade, paa hvilken han lokkede
dem, han talte med, paa Glatis ved tilsyneladende at anerkende
% --- p. 62 ---
\setcounter{footnote}{0}%
\opage{62}dem og gaa ind paa deres Meninger. Det er, som Ribbeck siger,
en Skæbnens Ironi, at netop den mest sandhedselskende af alle
Ateniensere skulde sættes i Klasse med Løgnere og Pralere. Men
netop Sokrates’ Personlighed, saaledes som Platon har skildret
den, bevirkede da ogsaa en Forandring i Opfattelsen af Begrebet
Ironi. Hos Aristoteles stilles Ironikeren i Modsætning til Praleren,
skønt han, ligesom denne, betegner en Mangel i Modsætning til
den Sandhedskærlige (Eth. Nic. IV, 13). Ironikerens Optræden
forklares ved hans Uvillie mod at slaa stort paa. Men det fremhæves
her endnu ikke, at den store Forskel mellem Ironikeren
paa den ene Side, Hykleren og Praleren paa den anden Side, bestaar
deri, at medens Praleren gør det ydre Skin til sit Formaal,
fordi han netop vil gøre sig større end han er, og medens Hykleren
bruger det ydre Skin som Middel til at opnaa selviske Formaal,
saa er for Ironikeren Skinnet et Middel til at bevare sit indre Liv
uforstyrret, eller maaske endog til gennem Indgaaen paa Andres
Forudsætninger at faa dem til at vinde i Klarhed. Paa denne
Maade opfatte vi nu Sokrates’ Ironi, ledede af den Anvisning,
Platon har givet i Skildringen af sin Mester. Et enkelt Sted giver
Aristoteles en Antydning i denne Retning, naar han om den Mand,
der har en stor Sjæl (ὁ μεγαλόψυχος), og som ikke blot agter sig
selv Noget værd, men ogsaa virkeligt er det værd, siger, at han
optræder ironisk overfor Mængden. (Eth. Nic. IV, 9). Han har jo
Noget i sit Indre, han ikke uden videre kan meddele.

Senere Grækere stode igen fremmede overfor den Modsætning
mellem det Indre og det Ydre, der lagde sig for Dagen i den sokratiske
Ironi. Epikuræerne opfattede Ironien som Udtryk for
Overmod og Praleri\footnote[1]{\textit{Theophrasti Characteres et Philodemi de vitiis liber decimus.} ed. Ussing. p. 4. 57. 175.}, og Ciceros Lærer Epikuræeren Zenon
kaldte endog Socrates en attisk Spasmager (scurra Atticus)\footnote[2]{Cicero: \textit{De natura deorum}. I, 34.}. ---

Den sokratiske Ironi er stor Ironi, hvis Sokrates’ Stade overhovedet
er udtømt ved Begrebet Ironi, hvad der siden skal overvejes
nærmere. Ved den store Ironi falder det polemiske Forhold
til Omgivelserne bort. Der kan, som allerede antydet, være to
Opgaver, der løses ved denne Ironi. Ironi kan være Tilflugt for
% --- p. 63 ---
\setcounter{footnote}{0}%
\opage{63}et Sind, der vil drage sig ud af Forholdet til det Ydre, til andre
Mennesker, til Livets givne Former, ledet af Trangen til at hævde
sig selv og arbejde paa sit indre Livs Udfoldelse, en Trang, der
naturligvis maa blive des stærkere, jo rigere et Indhold der er
ved at arbejde sig frem i Personligheden. En stor indre Fylde kan
gøre Isolation nødvendig, og ved i det Ydre at gaa ind paa og
anerkende Livets givne Former, Menneskenes Meninger og Sæder
og de raadende Autoriteter sparer man al Kraft til indre Arbejde
og indre Væxt. Men Ironien kan ikke blot tjene Selvhævdelsen.
Den kan ogsaa, og paa samme Tid, være en Metode, et pædagogisk
Middel til at muliggøre Forstaaelse af andre Personligheder,
en Forstaaelse, der igen er Betingelse for Hjælp, Opdragelse og
Medarbejde. Gennem den tilsyneladende Indgaaen paa Andres
Forudsætninger og Antagelser, gennem Anerkendelse af den Plads
og det Stade, de staa paa, faar man dem lettere til at aabne sig,
til at udfolde Konsekvenserne af deres Tanker, og dermed er
allerede det Samarbejde, i hvilket al virkelig Hjælp og al sand
Opdragelse bestaar, muliggjort. Medens Satiren støder Genstanden
bort, eller kun fastholder den for atter og atter at støde den bort,
kan Ironien være en Vej, ad hvilken det bliver muligt at trænge
ind i Genstanden og oprette et Fællesskab med den. Selv hvor
Ironien tjener Selvhævdelsen, er den ikke en Frastøden, men en
Unddragen, en Vej til at holde Sindets Helligdom fri for Omverdenens
Indtrængen.

Naar Ironikeren vækker Vrede ved sin Adfærd, kan det være,
fordi Menneskene ikke holde af, at man vil opdrage dem, føre
dem ud over det, i hvilket de trygt hvilede. Saasnart de opdage,
at dette var Enden paa den tilsyneladende Anerkendelse af deres
Stade, krænkes deres Selvfølelse. Men Grunden kan ogsaa være,
at Mennesker nødigt se Nogen unddrage sig Fællesskabet med
dem. Det volder Utryghed, blot det, at der er Nogen, som ikke
vil være med. Man vil have fælles Bekendelse. Ikke blot paa det
religiøse Omraade gives der Dogmatisme. Har jeg Ret, hvorledes
kan saa en Anden have Ret til at gaa andre Veje end jeg? Og især
naar jeg ikke engang véd, hvilke andre Veje han gaar, bliver jeg
maaske uhyggeligt tilmode.

I begge sine Former, baade som Selvhævdelse og som Opdra%
% --- p. 64 ---
\setcounter{footnote}{0}%
\opage{64}gelse, betegner den Art Ironi, som her er beskreven efter Antydninger
hos den historiske Sokrates\footnote[1]{Se nærmere nedenfor § 44.}, men i videre Udførelse, et
stort Vendepunkt i Aandslivet. Den er den første Form, i hvilken
Personlighedsprincipet, Forstaaelsen af og Agtelsen for individuel
menneskelig Personlighed\footnote[2]{Se om dette Princip i dets speciellere Udredning mine Helsingforsforelæsninger, udkomne under Navn af \textit{Personlighetsprincipen i Filosofin.} Stockholm 1911.}, Ens egen og Andres, bryder
frem. Og Sokrates staar her som Banebryderen. I hans
Ironi have, saavidt man kan se, begge de nævnte Former gjort
sig gældende.

Kierkegaard konstruerede i sit geniale Ungdomsskrift (Om Begrebet
Ironi, med særligt Hensyn til Sokrates. 1841) en absolut
Ironi, hvorved han mente en Ironi af rent negativ Art. Det var
den lille Ironi, han gjorde absolut, istedetfor at gøre den stor. Han
betragtede Ironien blot som opløsende, som blottet for ethvert
positivt Indhold. Han fraskriver den baade den indre Fylde og
det sympatiske Element. Den i den ironiske Alvor skjulte Spøg
skal være rettet mod Alt mellem Himmel og Jord. Alt opløses i
Intet --- med Undtagelse af det Selviske. Ironikeren tager Alt forfængeligt
--- undtagen sig selv. (Skrifter XIII p. 332). Sokrates
opløser det græske Aandsliv (ib. p. 338), uden at sætte Nogetsomhelst
istedet. Men --- mener Kierkegaard --- dette var berettiget
af Sokrates, ti den græske Kulturs Tid var væsentligt forbi. I
nyere Tid derimod er, mener han, en saadan Ironi (som den f. Ex.
optræder hos Fr. Schlegel) aldeles uberettiget: ti nu har „Subjektiviteten“
(Personlighedsprincipet) sejret. Den har gennemført sine
Krav i Stat og Kirke, i det sociale og det religiøse Liv, og den
aabenbarer sig stedse gennem disse Former.

Kierkegaard gaar i sit nævnte Ungdomsskrift ud fra sin Tids,
under Hegels Indflydelse udviklede spekulative Teologi, der hos
os repræsenteredes af Heiberg og Martensen, til hvilke han ogsaa
henviser. Derved førtes han, paa dette Stadium af sin Udvikling,
til at miskende den positive Karakter af den sokratiske Ironi.
Endnu senere hen (1846) kalder han (Papirer VII, 1. p. 27) Ironikeren
for en Spasmager. Det minder om Epikuræernes scurra
% --- p. 65 ---
\setcounter{footnote}{0}%
\opage{65}Atticus. Iøvrigt har han i senere Skrifter korrigeret sin Opfattelse,
saaledes som det senere vil blive vist.

Ikke blot Sokrates, ogsaa Fr. Schlegel gjorde Kierkegaard Uret
i sit Ungdomsskrift. Forholdene skulle ifølge Kierkegaard (og
Hegel) have været helt andre, da Schlegel optraadte, end paa Sokrates’
Tid. „Subjektiviteten“ havde jo nu sejret i Stat og Kirke!
--- Hertil maa siges, at det dog, da Schlegel optraadte, næppe har
været ufornødent at hævde den Ret, Aandens frie Liv har overfor
de i Institutioner og Dogmer krystalliserede Overleveringer. Der
er jo her Noget, som ikke kan afgøres en Gang for alle. Der staar
en evig Strid mellem Frihed og Norm. Udvorteshed, Trivilialitet
og Aandløshed ere netop ofte til Huse i de Former, Aandslivets
tidligere friske Opsving har lagt Grunden til. Man slaar sig til Ro
med at fejre de store, aandfulde Begyndelser og mener sig, trods
al Stilstand, i Slægt med dem. Der synes særligt i Tyskland ved
det nittende Aarhundredes Begyndelse, da Schlegel optraadte, at
have været Meget af denne Art at bekæmpe. Saaledes bemærker
Lotze i sin „Geschichte der Aesthetik in Deutschland“ (p. 371):
„den da foreliggende Stemning hos det tyske Folk kunde i æstetisk
Henseende snarere kaldes tom end rig. I dorsk Overlevering
og sneverhjertet Sæd og Skik havde Modtageligheden for det
Skønne tabt sig i den Grad, at det kunde staa som Opgave først
og fremmest at genoprette Driftens uhildede Liv ved at rejse sig
mod utallige Skranker og prøve og bestride utallige Fordomme.“
En af Tidens bedste Mænd, den store Idealist Fichte erklærede
endog i et Brev, at Tiden var karakteriseret ved „absolut Henraadnen
af alle Ideer.“ Den Tid, da Fichte og Schlegel optraadte,
har neppe havt noget Fortrin for den Tid, da Sokrates optraadte.
Tvertimod var Sokrates’ Optræden knyttet til en af Verdenshistoriens
højeste Kulturperioder.

Og trods alle Indvendinger, der kunne gøres mod Fr. Schlegel
som Person og som Skribent, var der hos ham en virkelig Trang
til Frihed og Ærlighed som første Betingelse for et værdigt Liv.
Dette var det, der vandt ham Venskabsforhold til Mænd som
Fichte og Schleiermacher. Og istedetfor med Kierkegaard blot at
støtte sig til en Analyse af Schlegels Roman „Lucinde“, bør man
% --- p. 66 ---
\setcounter{footnote}{0}%
\opage{66}høre, hvorledes Schlegel selv har gjort Rede for sit Begreb om
Ironi i sine Afhandlinger i Tidsskriftet „Athenäum“.\footnote[1]{Jeg har allerede bemærket dette i \textit{Den nyere Filosofis Historie,} II Note 38. --- De i det følgende anførte Udtalelser af Schlegel findes i \textit{Athenäum} (1800). I,2 og III, 1.}

Det er ifølge Schlegel Umuligheden af fuldkommen Selvmeddelelse,
der gør sig gældende i Ironien. Det Kosteligste, Mennesket
har i Eje, selve den indre Tilfredsstillelse, har tilsidst sin Rod i
noget Dunkelt, som dog er det, der bærer det Hele, og som vilde
miste sin Kraft i det Øjeblik, da man vilde prøve paa at opløse
det i ren Forstaaelse. I Ironien fremtræder Bevidstheden om en
evig Bevægelighed (Agilität), et uendeligt Kaos, som aldrig ganske
kan systematiseres. Kun tilsyneladende kan Personligheden
udtømmes i et enkelt Forhold eller i et enkelt Arbejde. Der er i
hver Personlighed, ligesom i enhver af Leibniz’s Monader et Univers,
en uendelig Fylde, der gør det baade muligt og nødvendigt
at hensætte sig, ikke blot med Forstaaelse og Fantasi, men med
hele Sjælen, paa højst forskellige Omraader, skønt intet af dem
kan gøre hel Fyldest.

Disse Udtalelser vise, at Schlegel selv opfattede sit Stade
som positivt, og at han havde Ret deri. Men det fremgaar tillige
som sandsynligt, at et saadant Stade ikke kunde være definitivt.
Det var endnu for kaotisk. Der var en Fylde, som trængte sig
frem, men det stod som vilkaarligt, hvilken Del af denne Fylde,
man særligt slog ind paa, og det vilkaarlige Spil, den lunefulde
Skiften havde i sig selv aabenbart en stor Tiltrækning. Og Fr.
Schlegel formaaede ikke at udforme, hvad der trængte sig frem.
Han manglede Sokrates’ og Sokratikernes Energi til at arbejde
paa Personlighedens frie Mark, hvor alle gode Planter gro. Hans
Sans for det Store og Uendelige var parret med Mangel paa
Evne til stadigt og trofast Arbejde og til derigennem at blive
Herre over det kaotiske Indhold. Derfor endte han i Dilettanteri.\footnote[2]{Smlg. \textit{Dilthey’s} fine Karakteristik af Fr. Schlegel i hans \textit{Leben Schleiermaches.} Anden Bog. Kap. 4. 12. 14.} % PRINTED AS IS: „Schleiermaches“

Der kom hos ham selv en Reaktion mod den lunefulde Skiften,
og der rørte sig en Trang til bestemte Former for Aandslivet.
Derved førtes han til Katolicismen. Allerede i „Athenäum
% --- p. 67 ---
\setcounter{footnote}{0}%
% CHECK p. 67: notes paired with the marks in order (the layer misread a number)
\opage{67}længtes han efter Orienten. „Zunächst rede ich nur mit denen,
die schon nach dem Orient sehen“ (III, 1). Det var fra først af
„det evige Orient“, han søgte; men saa mente han at finde det i
den gamle Kirkes vældige Krystallisation.

17. Ogsaa Shakespeare har man villet gøre til Ironiker. A. W.
Schlegel (Fr. Schlegel’s Broder) har konstrueret en Shakespearesk
Ironi, som dels skulde vise sig i Karakteristiken af de enkelte
Personer i hans Dramer, nemlig i den Maade, paa hvilken store
og smaa Motiver brydes hos dem, dels i Dramaets Handling
idethele, gennem den Upartiskhed, hvormed Skæbnens Gang
skildres. Karakterer skildres i hele deres Nøgenhed, ikke i entusiastisk
Beundrings Dragt, og Skæbnen afsløres i hele sin Virkelighed,
uden at Digterens Hjerte røres derved.\footnote[1]{\textit{Vorlesungen über dramatische Kunst und Literatur}\textsuperscript{4}. II. p. 197---199.}

Herimod indvender en af de ypperste nyere Shakespeareforskere\footnote[2]{\textit{Dowden: Shakespeare. His Mind and Art.}\textsuperscript{5} (1880). p. 345.}
med Rette, at Shakespeare’s Upartiskhed overfor Personer
og Skæbner ikke betyder Ligegyldighed og Fjernhed, men netop
udspringer af hans Interesse for Emnerne, af den faste Villie til
at lade enhver af deres Sider komme til sin Ret; hans Begejstring
svækkes ikke ved de Skranker, menneskelige Handlinger og
Skæbner ere underlagte.

Der gives en kunstnerisk Ironi, som hænger nøje sammen
med Kunstens store Opgave, den at give konkrete og individuelle
Billeder af Karakterer og Skæbner, ikke Abstraktioner og
Utopier. Den Lov, som udspringer af denne Opgave, maa Alt i
Kunsten underligge. Livet er saa rigt og arbejder sig baade indenfor
den enkelte Personlighed og indenfor Slægten frem fra
saa mange Sider, at der maa foregaa en Begrænsning, en Tilslibelse,
dersom Alt skal komme til sin Ret. Baade hvad de
individuelle Karakterer angaar, og hvad de Tilskikkelser, hvori
de stedes, angaar, gør sig nødvendigvis et negativt Element gældende,
idet det enkelte i Træk eller den enkelte Begivenhed underordnes
den hele Sammenhæng og ikke faar Lov at udfolde sig
paa den Maade, det kunde ske, dersom de havde været de eneste.
% --- p. 68 ---
\setcounter{footnote}{0}%
\opage{68}Dertil kommer, at Kunstnerens Opgave er at give de Følelser og
Billeder, der beherske ham, et saadant Udtryk og en saadan
Form, at hans Værk kan faa Værdi og vinde Anerkendelse ogsaa
ud over den Kreds, der nærmest omgiver ham og maaske fuldt
ud deler hans Følelser og ere fortrolige med hans Billeder. Derfor
maa det, der bevæger sig i Digterens Indre, udformes i klare
og bestemte Former, have Objektivitetens udprægede Karakter.\footnote[1]{Interessante Bemærkninger om dette Punkt hos \textit{Yrjo Hirn: Konstens Ursprung.} Stockholm 1902 p. 95—102.}

Men det vilde ikke være rigtigt at lægge Hovedvægten paa
det negative, begrænsede Moment i den kunstneriske Udformning.
Det Væsentlige er paa den ene Side den glødende Masse,
som skal formes, og paa den anden Side de plastiske Skikkelser,
som fremkomme ved Udformningen. Kampen, den ofte bitre Kamp,
og Arbejdet, det ofte strenge Arbejde, ligge mellem disse to
positive Momenter og kendes vel oftest kun af den enkelte
Kunstner selv. Den kunstneriske Evne er Evne til at individualisere
det, der uvilkaarligt trænger sig frem i Følelse og Fantasi.
I Kunsten som overalt fremgaar det Individualiserede og fuldt
Udprægede af et mere ubestemt Grundlag, en uorganiseret
Fylde. Forskellen mellem Kunstneren og os Andre bestaar væsentligt
i Evnen til at omsætte det, der dunkelt forestilles og
føles, i individuelle Skikkelser og Begivenheder. Og denne Evne
betyder noget Positivt, ikke blot en Begrænsning. Er Værket
fuldendt, Støbningen lykkedes, saa er der ingen Modsætning mellem
den uvilkaarlige Fylde og de udformede Skikkelser; hin er
kommen til sin fulde Ret i disse.

18. Den Art af Ironi, der har Karakteren af positiv Indgaaen
paa Andres Forudsætninger, staar nær ved Humor. Der er en
underjordisk Forbindelse mellem Humoristen og hans Genstand,
ligesom der er ved den pædagogiske Ironi. Forskellen mellem
Ironi og Humor er ofte udtrykt saaledes, at i Ironi er der Spøg
bag Alvor, i Humor Alvor bag Spøg. Alvor er overalt betinget
ved Forholdet til en Værdi: Alvoren sætter et Formaal og sætter
Villien i Bevægelse i Retning af dette Formaal. I Ironien gaas
% --- p. 69 ---
\setcounter{footnote}{0}%
\opage{69}der tilsyneladende ind paa Værdier og Bestræbelser, men Ironikeren
foretager kun denne Bevægelse for at bevare sit Indre frit,
eller --- i den pædagogisk Ironi --- for at udløse psykisk Energi
hos den, der vedkender sig visse Værdier og Formaal. At der er
Spøg bag Alvoren, behøver blot at betyde, at Ironikeren har sin
Alvor i andre Retninger end hans Omgivelser. Derfor kan Ironien
være en ydre Form for Humor, saa at vi faa følgende Skema
(hvor Tegnet $<$ betyder „Middel“ eller „Udtryk for“):

(Alvor $<$ Spøg) $=$ Ironi $<$ Humor $=$ (Spøg $<$ Alvor), eller
simplere

\begin{center}Alvor $<$ Spøg $<$ Alvor.\end{center}


Den Spøg, i hvilken Humoristen skjuler sin Alvor, bestaar da
netop i Anvendelse af Ironi. Eller simplere: bag den Spøg, der i
Ironien er skjult i Alvor, ligger igen Alvor skjult.\footnote[1]{Paa tilsvarende Maade kan den lille Humor træde i Ironiens Tjeneste, saa vi faa Skemaet: Spøg $<$ Alvor $<$ Spøg. Den Alvor, der i Ironien er ydre Form for Spøg, er da udtrykt humoristisk, d. v. s. indklædt i Spøg. Formelt kunne psykiske Processer som de her antydede stedse føres videre: Alvor $<$ Spøg $<$ Alvor $<$ Spøg \dots{} Men hvor der har udviklet sig en bestemt Livstype, et realt Selv, vil der tilsidst være Værdier og Formaal, der ikke kunne tages spøgende af Individet selv.} Paa denne
Maade gaar den store Ironi, som Form for Selvhævdelse eller for
opdragende Omsorg, over til Humor.

Man kunde spørge, hvorfor Alvor overhovedet skal udtrykkes
i Spøg. Fr. Schlegel (om hvem Novalis med Rette sagde, at han
egentligt var Humorist, ikke Ironiker) har allerede givet Svaret.
Svaret gaar ud paa, at der kan være et saa inderligt og saa omfattende
aandeligt Livsindhold, at enhver positiv Form bliver
utilfredsstillende, inkommensurabel. Og dette vil navnligt være
Tilfældet, hvor Livserfaringen har stillet os Ansigt til Ansigt med
store Modsætninger, hvis Brod vi have følt. Den Værdi, i hvis
Tjeneste Alvoren staar, vil vise sig at staa i Modsætning til andre
Værdier og peger derved ud over sig selv. Der kommer her,
ligesom ved den kunstneriske Udformning, et negativt Moment
frem. Selv de højeste Værdier maa udformes og begrænses gennem
Forholdet til andre Værdier. Denne Erfaring gør, at Humo-
% --- p. 70 ---
\setcounter{footnote}{0}%
\opage{70}risten helst udtaler sig spøgende og indirekte; derved værner han
om sine indre Opgaver, forbeholder sig det fortsatte Arbejde med
dem. Dertil kommer saa Værdiens Forhold til Virkeligheden, og
med det en ny Række af Opgaver. Virkeligheden kan, som Erfaring
viser, stride mod vore Værdier, eller forholde sig aldeles
ligegyldigt overfor dem. Derfor behøve vi ikke at opgive den
Tro paa Værdiernes Gyldighed, der viser sig i fortsat Arbejde
paa at virkeliggøre dem; men vi blive mere forbeholdne
med at bekende, og den spøgende eller indirekte Form tilbyder
sig da naturligt. Ved at bruge den lukke vi os ikke absolut inde;
den Erfarne og Kyndige vil kunne tyde Alvor ud af Spøgens
Text. Selve Spøgen er dog ikke blot en Tilflugt, en Udvej. Gennem
den godtgør Sindet sin Overlegenhed over de Modsigelser
og Modsætninger, det har at kæmpe med under den inderlige
Fastholden af det, der gør Livet værdt at leve. I alle Humorens
forskellige Former (som vi i næste Kapitel ville vende tilbage til)
træder dette Træk frem.

Af størst Betydning er her Modsætningen mellem Værdiernes
Verden og Virkelighedens Verden.

Stiller man sig paa et rent Erfaringspunkt, staar Virkelighedens
Verden som langt mere omfattende end Værdiernes Verden. Det
Værdifulde er kun en Del, maaske en forsvindende lille Del af
Virkelighedens Verden. I stor Stil traadte dette frem, da det
gamle Verdensbillede med dets hyggelige Overskuelighed indenfor
en afsluttende Ramme opløstes, og det gik op for Bevidstheden,
at det er umuligt at fastslaa Grænser for Verden. Engle kunne
nu ikke mere stige op og ned fra Værdiernes Himmel til Virkelighedens
Jord, og der kan ikke mere paavises noget absolut Centrum,
hverken hvad Værdi eller hvad Virkelighed angaar. Som en lille
Oase staar nu Værdiernes Rige i Modsætning til Virkelighedens
store Ørken. Eller lad mig bruge et andet Billede, der er hentet
fra en Dagbogsoptegnelse.

„Det er koldt og raat idag. En skarp Vind farer gennem de
bladløse Træer. Men denne samme Vind frembringer i Æolsharpen
paa det Hus, jeg bor i, stemningsfulde Toner, en lang,
dyb, beroligende Klang. Saaledes kan Værdiernes Rige aabenbare
sig midt i Virkelighedens raakolde Verden.“

% --- p. 71 ---
\setcounter{footnote}{0}%
\opage{71}Der er her to Momenter, som begge Billeder angive. Oasen
ligger inde i selve Ørkenen; den samme Vind, der volder Uhyggen,
fremkalder ogsaa den trøstende Klang. Altsaa Værdi, er
baade Modsætning til Virkeligheden og en Del af denne. Ved
denne Dobbelthed bliver en Livsstemning som Humor mulig, ---
et Sindelag, der ikke mister Sansen for det Værdifulde og Troen
paa dets Gyldighed, fordi ikke Alt kan udledes og forklares udfra
det, eller fordi alt Andet i Verden synes ligegyldigt overfor det.
Humoren vil fæste sig ved, at det er i Kraft af samme Love, at
der existerer noget Værdifuldt i Verden, og at den ligegyldige 
eller fjendske Virkelighed bliver til. Vinden rusker i Træerne,
men klinger gennem Æolsharpen. Det staar som en Sejr, at der
dog paa nogle Steder i Verden er sket et Gennembrud af noget
Stort og Godt, selv om det er en ren Fantasi, at Alt skulde være
lagt til Rette for, at saadanne Gennembrud kunne være mulige.
Der kan komme humoristisk Stemning frem overfor Forestillingen
om de uhyre Massers overvældende Magt, som dog ikke formaar
at hindre, at Værdiernes Klang høres. Ti det er kun det materielt
eller kvantitativt Store, der fremtræder i Verdensmasserne;
i Tankens og idethele i Værdiernes Verden ytrer sig en Storhed
af anden Art, en Inderlighedens Storhed, med hvilken al Udstrækning
i Rummet og i Tiden er inkommensurabel. Der kan
opstaa en Følelse af Trods, naar Modsætningerne træde frem
i deres skarpeste Form; men jo mere Sindet fordyber sig i selve
Værdierne og i de Aarsagsrækker indenfor Virkelighedens Verden,
der ere værdifrembringende, des mere vil den store Humor
komme til at raade. Og det er jo tillige et saadant Sindelag, der
er Grundlag for energisk Arbejde paa at frembringe og fastholde
Værdier.

En særlig Anledning til Humor overfor Tingenes Gang i Verden
fremkommer ved den Erfaring, at det Værdifulde ofte frembringes
af Aarsager, der syntes at virke i helt anden Retning,
saa at de, for at tale mytologisk, mod deres Villie eller uden deres
Villie komme til at udrette, hvad der for os staar som skønt
og godt. Det er det Samme, der, naar Talen er om menneskelige
Følelser og Handlinger, kaldes Værdiforskydning. (Smlgn. Psykologi
VI C; Etik XIII, 4). Naturformer og Naturscener, som til
% --- p. 72 ---
\setcounter{footnote}{0}%
\opage{72}tale vor æstetiske Følelse, have dog ikke udviklet sig for at frembringe
æstetiske Værdier, men efter bestemte Aarsagslove. Menneskelige
Samfunds- og Kulturformer, som vi tillægge etisk
Værdi, kunne dog have udviklet sig under Indflydelse af den
blinde Trang til Selvopholdelse, Ernæring og Forplantning. Det
er ofte mangehaande og højst forskelligartede Motiver, der i Historien
bevirke de store og værdifulde Resultater. Det gælder i
det Smaa, ved Menneskers gode Raad og venlige Haandsrækning
er\footnote[1]{Smlgn. f. Ex. det Komplex af Motiver, som (i \textit{George Eliot’s The Mill on the Floss.} Book I. Chap. 3). førte Mr. Riley til at raade Mr. Tulliver til at sætte sin Søn i en bestemt Kostskole. Det er et Stykke Psykologi, der ikke kan resumeres.}, som det gælder ved verdenshistoriske Afgørelser. Hvad der
kom ind i Verden som et blot Middel, maaske paatvunget Middel,
fik senere Værdi for sig selv og kunde opstilles som Formaal.
Saaledes have f. Ex. de store Erobrere ofte virket i Humanitetens
Tjeneste ved at føre ud over nationale og konfessionelle Grænser,
skønt deres eget Formaal var Tilfredsstillelse af Magtlyst og Ærgerrighed.
Og selv hvor det var et i sig selv værdifuldt Formaal,
der stilledes, viste det sig ofte kun at være et Middel, en Station
paa Vejen til et større Maal.

Man kan sørge over, at der til de store Formaal ikke altid svarer
store og ædle Midler, men at det ofte er meget elementære
eller endog rent egoistiske Aarsager, der føre til Virkeliggørelse
af de store Værdier. Men man kan ogsaa se det som betydningsfuldt,
at selv saadanne Aarsager komme til at tjene det Værdifulde
i Verden. Det er da ikke Værdierne, som fornedres ved at
virkeliggøres gennem elementære Aarsager, men det er omvendt
disse, der adles ved, at de Love, hvorefter de virke, indeholde
Mulighed for værdifulde Helheders Opstaaen. Der er da Plads for
Smil saavel som for Graad, og begge Stemninger kunne forenes
til Humor, med forskellig Klangfarve, alt efter Forholdet mellem
de to Enkeltfølelser.

En lille Dagbogsoptegnelse kan ogsaa her bruges, dels som
Exempel, dels som Symbol. Den er bleven til ved en Iagttagelse,
der gjordes paa Dampfærgen over Store Belt. --- „Jeg stod og saa
paa Maagerne, der fulgte Skibet, og beundrede deres hurtige og
% --- p. 73 ---
\setcounter{footnote}{0}%
\opage{73}dristige Sving, de formfuldendte Bevægelser, de yndefulde Buer,
deres Flugt beskrev. Jeg nød det pragtfulde Syn af de hvidgraa
Fugle i det straalende Solskin. Og saa kom jeg til at tænke paa
Aarsagen til al denne Herlighed, --- hvad det var, der gjorde det
hele Syn muligt: Ernæringstrangen, Madstræbet, Kampen for
Føden!“ --- Saavidt mindes, tabte Synet ikke sin Glans ved
denne Reflexion; men Stemningen blev noget mere sammensat.
Maaske fortrængtes dog den første Glæde et Øjeblik, men vendte
saa styrket tilbage.

Ved alle ideelle Værdier spille Værdiforskydninger en vigtig
Rolle. For Kunstens Vedkommende har Yrjo Hirn i en meget
interessant Fremstilling vist, hvorledes en hel Række praktiske
(„utilistiske“) Faktorer have fremmet de forskellige Kunstformers
Udvikling og ført til Skønhedssansens Væxt. Saadanne Faktorer
ere: Trang til at øve magisk Indflydelse eller erotisk Fortryllelse,
til at opmuntre til Krig eller Arbejde, og ikke mindst til at
meddele sine Tanker til Andre.\footnote[1]{\textit{Konstens Ursprung} p. 173—294.} Noget Tilsvarende gælder, som
jeg har vist i min Etik, for de etiske Ideers og Følelsers Vedkommende.
---

Den store Humor ser Livet baade som stort og som smaat,
baade som værdigunstigt og som værdifjendsk, baade som tragisk
og som komisk. Solger har først indgaaende beskrevet Humor
som Enhed af komisk og tragisk Følelse, af Ironi og Begejstring.
Han er noget vaklende i sin Anvendelse af Begreberne
Ironi og Humor. Klarest beskriver han i sin Dialog „Erwin“
Humoren saaledes: „Det Tragiske og det Komiske ligge her
endnu usondrede ved Siden af hinanden\dots{} Alt er i Humoren i
én Strøm, og overalt gaa, ligesom i den sædvanlige Verden,
Modsætningerne over i hinanden. Intet er i Humoren latterligt
og komisk, som ikke har et Element af Værdighed eller ikke
kan vække til Vemod, --- Intet er ophøjet og tragisk, som ikke
under sin Udformning i Tidens lave Verden kunde synke ned i
det Betydningsløse eller det Latterlige.“ (Erwin. Berlin 1907. p.
354). Kun i Længselen kan Tanken om det Højeste som hævet
over Forfængelighed bo; men Humoristen maa anse denne Læng%
% --- p. 74 ---
\setcounter{footnote}{0}%
\opage{74}sel for langt mere værd end den begrænsede Form, i hvilken
Fantasien formaar at forme det Højeste (Erwin. p. 358).

Dette Solgerske Standpunkt lader Kierkegaard i sin Afhandling
om „Begrebet Ironi“ ikke komme til dets Ret. Han saa paa det
med teologisk-hegelianske Øjne. I det nævnte Skrift (og i Optegnelser,
der ere samtidige med det) opfattede Kierkegaard Forholdet
mellem Ironi og Humor som Modsætninger. Ironien skal
være egoistisk, og den skal være stedt i stadig Kamp mod Verden.
Humoren derimod skal være befriet fra Verden; den staar som
Udtryk for en religiøs Livsalvor og for Bevidstheden om at høre
ind under en stor Helhed. Ironien skal være aristokratisk og
polemisk, men Humoren er forsonet med Verden og hviler i en
Tro, som fastholdes trods alle Absurditeter. (Skrifter XIII p. 35
---119.)

Senere derimod, efter det definitive Brud med den spekulative
Teologi, betragter Kierkegaaard ikke mere Humor som et religiøst % PRINTED AS IS: „Kierkegaaard“
Stade. Den staar for ham nu som det sidste humane Stade
før det religiøse Stade, --- som Udtryk for det højeste rent Menneskeliges
Grænse. Det karakteriseres nærmere som væsentligt
betinget ved vemodigt smilende Erindring. Det er rettet tilbage,
eller rettere, Tiden spiller ikke længere nogen Rolle for det. Religionen
derimod ser fremad. Paa det religiøse Stade er Livet fra
først til sidst, i Smaat og Stort bestemt ved Gudsforestillingen, i
Forhold til hvilken alt Andet taber sin Værdi, saa at det Højeste
bliver Selvtilintetgørelse for Gud.\footnote[1]{De herhenhørende Skrifter ere: \textit{Studier paa Livets Vej} og \textit{Uvidenskabelig Efterskrift.} --- Smlgn. om Kierkegaard’s Humorstade \textit{Danske Filosofer}. p. 160—162.} --- En Prøvelse af denne
Kierkegaards definitive Opfattelse af Humor vil blive anstillet
nedenfor, i Kapitlet om Tragik og Humor.

% --- p. 75 ---
\setcounter{footnote}{0}%
\opage{75}\begin{center}IV. HOVEDFORMER AF HUMOR\end{center}


19. Fra højst forskellig Side kunne de Enkeltfølelser komme
som senere i Forening frembringe en Totalfølelse. Især er her
som paa saa mange andre Punkter, Forholdet mellem Selvhævdelse
og Hengivelse af Betydning. Humor faar saaledes en anden
Klangfarve, naar den udvikler sig paa Grundlag af Overlegenhed
og Sindshøjhed eller paa Grundlag af Vemod og Længsel, end
naar dens væsentlige Momenter ere Sympati og Forstaaelse. Desuden
kommer det igen indenfor disse forskellige Former an paa,
hvad der er erfaret af Sorg og Glæde, af Savn og Nydelse, ---
hvad Erfaring og Eftertanke have aabnet Blikket for, --- hvor
stor Forundring Livet har fremkaldt, --- hvilke Skuffelser det har
voldet, --- og hvilken Forstaaelse, der er vundet af det Hele.
Her, som paa saa mange andre Punkter, fremtræder den af
Sibbern saa stærkt fremhævede sporadiske Karakter af Sjælelivet
paa dets tidligere Udviklingsstadier. Naar nu baade Grundlaget
for Oplevelserne og selve Oplevelserne ere forskellige, maa
der naturligvis forekomme overordentligt mange Former for Humor.
Der vil desuden være stor Forskel paa, hvor let og hurtigt
Sammensmeltningen eller Organisationen foregaar. Ofte vil der
her være svære Kriser at gennemgaa, indre Strid at gennemkæmpe,
før den naar en vis Fasthed og Sluttethed i en af Humorens
Hovedformer. Saadanne Gæringsperioder hører til de psykologiske
Erfaringer, Sibbern bedst har beskrevet.\footnote[1]{\textit{Danske Filosofer} p. 101. --- \textit{Sibbern} er først kommen til sin Idé om sporadisk Udvikling ad psykologisk Vej. Senere har han fundet analoge Forhold i Naturen og i Samfundsforholdene, saa at hin Idé kunde indgaa som et væsentligt Led i hans Verdensanskuelse. Da han meget ofte bruger Udviklingen indenfor Fosterlivet som Eksempel, antager jeg, at han er paavirket af den franske Embryolog \textit{Serres,} der (1829), i Modsætning til den af Harvey, Wolf og Haller antagne Udvikling fra Centrum til Periferi, hævdede at Fosterudviklingen skrider frem „par des noyaux latéraux qui se soudent ensuite“, Noget, han særligt paaviste for Knokkeldannelsens Vedkommende. \textit{Cuvier: Histoire des Progrès des Sciences Naturelles.} Bruxelles 1837—38. II p. 431—434.}

% --- p. 76 ---
\setcounter{footnote}{0}%
\opage{76}Saa stærkt fremtrædende Sjælelivets sporadiske Karakter end
er, saa er det dog en endnu mere fremtrædende Ejendommelighed
ved det, at ethvert psykisk Element har Tendens til at indgaa
Forbindelse med andre Elementer, saa at en Helhed kan
dannes. Det er Kant’s udødelige Fortjeneste --- større end den,
han selv lagde mest Vægt paa --- at han har hævdet Trang og
Evne til sammenfattende Virken som alt Sjælelivs inderste Ejendommelighed.
Denne Trang og Evne fandt han i Sanseanskuelse
og i Forestillingsassociation, og særligt i de specielle Former,
videnskabelig Forsken forudsætter. Tillige fandt han i den Forklaringen
af den menneskelige Aands stadige Tragten efter at
hige ud over Erfaringens og Videnskabens Omraade, idet den jo
ikke fra først af kunde kende sine egne Grænser. Med denne Betragtningsmaade
har Kant stillet Aandsvidenskaben dens egentlige
Opgave.\footnote[1]{\textit{Den nyere Filosofis Historie.} II. p. 27—29; 42—46. --- Smlg. ovenfor § 2 og § 8.}

Hvad der skal udtrykke Sjælelivets Ejendommelighed paa alle
Omraader, maa naturligvis være af formel Art. Hvad Kant kaldte
„den skjulte Kunst i Menneskets Indre“, er ikke bundet til noget
bestemt Indhold, netop fordi det sætter sit Præg paa alt Indhold,
der optages i Sjælen. Kant selv anvendte overvejende sit Syntesebegreb
paa Tankelivets Omraade. Men hvis han har Ret i sin
Opfattelse, maa noget Tilsvarende ogsaa gælde for de andre Sider
af Sjælelivet. Der maa saaledes være en vis Tendens, en Trang
og Evne til at samle Livets Glæder og Sorger i en Totalfølelse,
som man maaske kunde kalde en syntetisk Følelse, ligesom der
paa Villiens Omraade maa være en Tendens til at samle alle Bestræbelserne
i én stor Stræben.

Hvad man, ofte noget fordringsfuldt, kalder psykologiske Love,
angiver de forskellige Former og Betingelser for den sammenfattende
Virken. Hvad de specielle Helhedsdannelser angaar, vil
% --- p. 77 ---
\setcounter{footnote}{0}%
\opage{77}man, som allerede tidligere bemærket, ikke kunne vente, at de
skulde kunne udledes af det Kendskab, vi have til de sporadiske
Ytringer af Fornemmelse, Forestilling, Lyst og Ulyst, uvilkaarlig
Stræben, Drift og Ønske. Der er saa mange Muligheder for Samvirken
og Modvirken, for Strid og Forening, for Forbindelser af
højst forskellig Grad og Art, at vi ere henviste til at studere Helhederne,
efterat de faktisk have dannet sig, og saa bagefter finde,
hvorledes den sporadiske Udvikling kan have ført til saadanne
Resultater. Vi ere her, paa det psykologiske Omraade, stillede
paa lignende Maade som overfor alle Helhedsdannelser i Naturen.\footnote[1]{Smlgn. herom: \textit{Den menneskelige Tanke} p. 218—237; 338—345. I en særlig Afhandling haaber jeg at belyse dette Punkt nærmere fra et erkendelsesteoretisk Synspunkt.}


Den store Humor, som jeg her forsøger at beskrive den, har
en ganske særegen Mulighed for at være en afsluttende Totalfølelse.
Vi mangle endnu indgaaende psykologiske Undersøgelser
af de forskellige mulige Totalfølelser, der kunde tendere til at faa
en afsluttende Karakter for en Livsperiode eller maaske for et
helt Livs Vedkommende. Beskrivelse og Analyse af en enkelt
Totalfølelse kan ikke afgive Bevis for, at den skulde være den
betydningsfuldeste eller den ædleste. Jeg vil blot henlede Opmærksomheden
paa, hvad jeg kaldte den store Humor, og undersøge
de Betingelser, den har for at være en afsluttende Totalfølelse.
Den bliver til gennem Oplevelse af Livets Modsætninger,
gennem Enkeltfølelsers Brydning, under Medvirkning af Eftertanke
og af Stræben efter Formaal, der fastholdes med Energi og
Troskab gennem Tiderne. Og den har Mulighed for at holde
Modtageligheden aaben for, hvad nye Oplevelser kunne bringe.

Den psykologiske Mulighed af en saadan Totalfølelse, de Farer,
der true den, og de Betingelser, den er underlagt, ville blive
undersøgte i følgende Afsnit af denne Afhandling. Foreløbigt gaas
nærmere ind paa nogle allerede (se § 13) omtalte Hovedformer,
i hvilke den store Humor kan optræde.

20. Grænsefænomerne fremtræde her i saadanne Tilfælde, hvor
det er Overlegenhed eller Sindshøjhed, der danne Grundlaget.
% --- p. 78 ---
\setcounter{footnote}{0}%
\opage{78}Overfor et saadant Grundlag kunne selv meget stærke og indtrængende
Oplevelser prelle af uden blivende Virkning. De rokke
ikke ved det faste Midtpunkt i Sjælelivet og udløse derfor hverken
Haan eller Frygt. Men netop paa Grund af den indre Fasthed
er der Mulighed for et spøgende Forhold overfor Oplevelserne,
de ugunstige saavel som de gunstige, saavel som for
Sympati med Mennesker. Jo friere Mennesket staar i sin egen
Selvhævdelse, jo mere han magter den, des mere kan han tage
Del i Andres Kamp for Selvhævdelse. Der er her en sikker
Magtfølelse tilstede, som ikke medfører Trang til at gøre Overlegenheden
gældende overfor det Smaa, men som netop kan afføde
Trang til at faa Afløb for Energien gennem hjælpende Arbejde.
Den sande Overlegenhed trænger ikke til at sammenligne
sig med dem, der staa lavere, men forudsætter netop, at man har
vundet en grundig Selvvurdering gennem Sammenligning med
de Største. Ofte gaar Vejen gennem Kamp og Brydning. Saalænge
ingen Modstand mærkes, ingen Hindring stanser, ingen Forventning
skuffes, vil ingen grundig Totalfølelse danne sig. Kraften
sættes ikke paa afgørende Prøve. Som ved al Erfaring gælder
det ogsaa her, at lutter positive Tilfælde ingen sikker Slutning
tillade. Ligesom Logiken kræver, at vi skulle skaffe os Besked,
ikke blot om de Tilfælde, i hvilke et Fænomen indtræder, men
ogsaa om dem, i hvilke det ikke indtræder, saaledes vindes der
ogsaa kun Livserfaring, naar der ikke blot arbejdes med Strømmen,
men ogsaa mod Strømmen. Kun derved faas den Grundfarve,
som ifølge Platon var nødvendig, for at Karakterens egentlige
Farve skulde blive ægte. Selv hvor Personlighedens Udvikling
foregaar kontinuerligt, kommer der dog kun Tilforladelighed og
Fasthed i den, naar den er prøvet gennem Skuffelse og Savn, og
naar den har lært at holde et stort Maal fast trods de Hindringer,
som møde, og trods den Fjernhed, i hvilken Opnaaelsen ligger.
En Storm, siger La Rochefoucauld, faar et Baal til at blusse
stærkere, medens den slukker en Praas. Kan en Følelse holde sig,
skønt dens Genstand er i det Fjerne og ingen sanselig Nærværelse
eller noget anskueligt Billede støtter den, vil den blive saa
meget inderligere og fæste des dybere Rod. En stor Sorg kan
udløse Energi, der ellers ikke vilde komme til at virke. Der
% --- p. 79 ---
\setcounter{footnote}{0}%
\opage{79}hører allerede Energi til at sørge, til at udholde og gennemleve
Spændingen mellem Følelsen af det Tabtes Værdi og Visheden
om, at det nu virkeligt er tabt. Den letteste Maade at befries for
denne Spænding vilde være at skyde et af Leddene ud, enten
glemme, hvad der er tabt, eller indbilde sig, at det Tabte ingen
Værdi har. Der hører Overlegenhed og Sindshøjhed til ikke at
gribe nogen af disse Udveje. Eller det forholder sig snarere omvendt:
gennem den Tapperhed, der ikke griber nogen af disse
Udveje, gaar Vejen frem til Overlegenhed og Sindshøjhed.

Den stoiske Sindshøjhed gik ikke denne Vej til Ende.\footnote[1]{Smlgn. om Stoicismen \textit{Mindre Arbejder.} Første Række (Afhandlingen om hedenske Sandhedssøgere) og Tredie Række (Afhandlingerne om Epiktet og Markus Aurelius).}
Den købtes dyrt med Angst for Sindets uvilkaarlige Bevægelser.
Ethvert Baand, der knytter os til Nogen eller Noget, medfører
Muligheden af at blive dragne ind i Glæder og Sorger, der
komme uafhængigt af vor Villie. Derfor skal man, hævdede Stoikerne,
hvis man vil bevare sin aandelige Frihed og Ligevægt,
holde sig til, hvad der staar i vor Magt og ikke agte Noget for
godt eller ondt, uden forsaavidt det begunstiger eller hæmmer
vor Selvbeherskelse. Stoikeren staar ved Livets Skuespil som en
kølig Tilskuer, der opmærksomt følger Spillets Gang, men vel
vogter sig for at lade sig rive hen deraf. Fremfor Alt frygter han
for at komme til „at sukke i sit Indre“. Den indre Helligdom
skal holdes aldeles uberørt af, hvad der sker i Verden.

For Stoicismen er derfor, mere end for nogen anden Retning i
græsk Livsopfattelse, Humor en fremmed Følelse. Stoicismen
søger jo netop at fortrænge og udelukke de Elementer, ved hvis
Forbindelse Humor kunde opstaa. Den stænger Livets Modsætninger
ude istedetfor at arbejde sig igennem dem. Humor er kun
mulig, naar Selvhævdelsen ikke behøver at være saa ængstelig
og krampagtig, at den afskærer sig fra at gribes af, hvad der opleves
under Livets Gang. Humor forudsætter, at der kan udløses
saa megen Energi til Raadighed for Selvhævdelsen, at man kan
tage de Stød som Verden giver, og være Vidne til den Dumhed
og Dorskhed, der breder sig i Verden, uden at rokkes i sin indre
Tillid og Fasthed, ja, endog saa at man har et Smil tilovers derved.

% --- p. 80 ---
\setcounter{footnote}{0}%
\opage{80}Naar man har lært, at den højeste Lykke bestaar i Stræben og
Arbejde, vil man finde den naturligste Forklaring af alt moralsk
Ondt deri, at det udspringer af Træghed, af Mangel paa Evne til
at udvikle sig ud over det Stade, man nu engang er kommen til
at staa paa eller hvile i. Hvad der mangler, er tilstrækkelig
stærk Paavirkning og Tilskyndelse til at udløse de bundne Kræfter.
Ud fra en saadan Betragtning stiller Kampen mellem Godt
og Ondt sig anderledes, end naar man med Teologerne udleder
det Onde af et ubegribeligt Fald, der med et Slag skulde have
bragt den højeste Synd, Trods mod det klart bevidste Gode og
Rette, ind i Verden. Og naar man saa tillige gør den Erfaring, at
netop Kampen mod det Onde kan udløse ideale Kræfter til Bekæmpelse
af det, fremkalde store og skønne Bestræbelser, der
ellers ikke vilde være komne op, saa bliver Følelsen overfor
Livet og Tilværelsen mere sammensat. Den kan minde om den
Stemning, der i en gammel Paaskesang fører til at udbryde O,
felix culpa! Ingen Formel og intet Dogme slaar her til, her hvor
de største og af hinanden afhængige Modsætninger forenes i,
hvad man kunde kalde en kosmisk Livsfølelse. Intetsteds viser
det sig mere end her, at vore dybeste og mest omfattende Livserfaringer
kun kunne udtrykkes i Poesiens symbolske Form.

21. Hvor der er Overlegenhed og Sindshøjhed, er der en fast
Grund, og Spørgsmaalet bliver, hvor stor Bevægelighed og Modtagelighed,
der trods denne Fasthed er Mulighed for. Hvor Humor
har udviklet sig med Vemod og Længsel som Grundlag,
bliver Spørgsmaalet det modsatte, nemlig om der trods Bevægeligheden
og Modtageligheden kan være den Fasthed, som ingen
Totalfølelse kan undvære.

Vemodet er tilbageskuende, lever i Erindringen. Længselen er
fremadskuende, lever i Haab og i Higen. Der er et Element af
Sorg i dem begge, hist fordi en Værdi er gaaet ud af Livet, her
fordi Afstanden fra Maalet er saa stor. Men der er Glæde i dem
begge, fordi der leves og aandes i Noget, som har Værdi, hvad
enten det saa hører Fortiden eller Fremtiden til. Man har med
Rette hævdet,\footnote[1]{\textit{Alexander Shand:} Foundations of Character. p. 395.} at den egentlige, fjendtlige Modsætning til Glæde
% --- p. 81 ---
\setcounter{footnote}{0}%
\opage{81}ikke er Sorg, men Afsky og Modbydelighed. Glæde og Sorg
kunne meget vel indgaa i en Sammensmeltning og en Organisation.
Ogsaa Vemod og Længsel kunne gaa sammen; man kan
længes tilbage til Fortiden, naar den ret bliver levende i Erindringen,
og der kan være Vemod i Længselen, naar Tanken om,
hvor fjernt og vanskeligt Maalet er, fremhersker. Forudsætningen
for begge disse Følelser er, at der er tilstrækkelig psykisk Energi
til ikke blot at leve i det sanseligt Nærværende, men i en vis
Uafhængighed af dette. I sløve og nedstemte Tider kan der, som
sidste Opblussen af psykisk Energi, røre sig en Længsel, maaske
en vemodig Længsel, efter den Tid, da man endnu kunde længes.\footnote[1]{Smlgn. Kapitlet om „Sehnsucht“ i \textit{Holzapfel’s Panideal} (Leipzig 1901).}
Denne Betydning som Udtryk for psykisk Energi miste disse
Følelser, hvor Forholdet til det Nærværende helt falder bort, og
Fortid og Fremtid aldeles ikke staa i Sammenhæng med Nutiden.
Livet bliver da et Drømme- eller Fantasiliv. Denne Karakter fik
i høj Grad den romantiske Humor. Ifølge Steffens var det Bedste
i den Tid, han levede i, Sansen for de gamle, forsvundne Tider.\footnote[2]{Se \textit{Danske Filosofer} p. 43.}
Enhver Livsanskuelse, der har sit Ideal i Fortiden, har her en
stor Fare at kæmpe med.

Det var ved at vække en stor og levende Forventning, en
inderlig Længsel, at den oprindelige Kristendom virkede saa
stærkt i den første Tid. Kun saadanne Stemninger, der direkte
eller indirekte kunne afledes af denne store Længsel, bære med
historisk Ret Navn af kristelige. Denne Længsel var rettet mod
et stort Maal; derfor viste den sig saa voldsomt og saa inderligt.
Men den var ikke rettet mod noget fjernt Maal. I den allernærmeste
Fremtid, som den oprindelige Kristendoms Mennesker
kunde haabe at opleve, ventedes et stort Gennembrud af himmelske
Kræfter, som skulde hidføre en helt ny, ideal Tingenes
Orden („Himmeriges Rige“ i dets bestemte, historiske Betydning).
Længselen blev her en verdenshistorisk Kraft. Og der er maaske
intet bedre Bevis paa Længselens Betydning end det, at selv
efter at hin Forventning er blegnet og dens Maal skudt ud i
det Blaa, i det Ubestemte,— hvilket skete, da Kirken havde indrettet
sig hyggeligt i det Nærværende, --- selv da virke de Stem%
% --- p. 82 ---
\setcounter{footnote}{0}%
\opage{82}ninger, den har fremkaldt, og de Forestillinger, den har givet
Anledning til, som væsentlige Elementer i den helt nye Religion,
moderne Kristendom faktisk er.\footnote[1]{Smlgn. Afsnittet „Urkristendom og moderne Kristendom“ i min \textit{Religionsfilosofi} (IV. C).} ---

Hvor Erindring eller Forventning om store Ting raade, vil
Mangt og Meget af, hvad der ellers kunde have staaet som betydeligt,
indgribende og skæbnesvangert, nu kun fremkalde et
Smil. Hvad der bevæger den Halvbefarne saaledes, at Ligevægten
gaar tabt, det fortoner sig nu overfor Sindets Optagethed
af det Store. Og naar det alligevel forsøger at udøve sin Magt
over Sindet, lokkes Smilet frem, ikke Satirens eller Haanens
beske Smil, eller Ironiens skjulte Smil, men Humorens stille
og lyse Smil. Der er jo Adgang til en Verden, som er langt større
end den, de nye Erfaringer tilhøre. Dog er hin Verden ikke skinsyg
paa denne. Den har Plads aaben for den, skønt maaske ikke
den fremragende Plads, de nye Oplevelser i og for sig vilde gøre
Krav paa.

Hvor Forholdet til det Nærværende helt mangler, eller kun
er negativt, afføde Vemod og Længsel ikke Humor. Det er saaledes
vistnok med Rette blevet bemærket om Shelley, at han
manglede Humor.\footnote[2]{Browning: \textit{Essay on Shelley.} London 1903. p. 60—64. --- \textit{A. C. Bradley: Shelley’s View of Poetry.} (Oxford Lectures on Poetry. London. 1909. p. 163—167.)} Hans overstrømmende Stemningsliv skaffede
sig Luft i vældige eller æteriske Fantasier, der udtrykte hans
Harme eller hans Længsel, og som have Præg af det rene Lys
eller det rene Mørke, men ikke lade Lyset bryde sig i noget
Prisme. Hele den virkelige Verden stod for ham, i dens Endelighed
og Begrænsning, som en Misforstaaelse, en Illusion.

Grundlaget for Humor findes ofte, snarere end hos Fantasimennesker,
hos Mennesker, hvis Liv er gaaet hen med praktiske
Hverv og under mangehaande Omskiftelser, der have bragt Erfaring
og Forstaaelse med sig. Bedrifter og Skæbner behøve
ikke at være i den store Stil for at aabne Sindet for Livets
Hemmeligheder. George Eliot taler i et Brev om „den Poesi og
Patos, den Tragedie og Komedie, der kan gemmes i én Menneske%
% --- p. 83 ---
\setcounter{footnote}{0}%
\opage{83}sjæl, der ser ud gennem matte graa Øjne og udtaler sig gennem
en Røst af ganske sædvanlig Klang.“ Man kan paa sin Vej møde
Mennesker af denne Type, men benytter altfor sjelden Lejligheden
til at lære dem nærmere at kende. En af Bipersonerne i
Fru Gyllembourg’s „To Tidsaldre“ er et smukt Exempel. Det
hedder her om Skipperenken Madam Lyng: „Hun havde gaaet
Meget igennem i Livet og var et af de aldrig noksom paaskønnede
Mennesker, hos hvem Hjertets Godhed og Blødhed træde
istedetfor Opdragelse og Dannelse, og hvem Erindringen om de
egne Lidelser og Forurettelser lærer Finhed i Omgang tiltrods
for enhver udvortes Stilling i Livet.“ --- Denne Karakteristik
indeholder væsentlige Betingelser for Humor, selv om denne
Egenskab ikke direkte tillægges den brave Kone. Det er især
„Finheden“, der maa lægges Mærke til. Den er Evne til at stille
sig hensynsfuldt og forstaaende overfor, hvad Livet medfører,
uden at de ydre Forskelle, det frembyder, formaa at imponere.
Mangt og Meget kan her tages med et Smil, som hos den Uprøvede
vilde vække Uro, Smerte eller Haan. Sorgen har havt en
saa at sige vaccinerende Kraft. Derved kunde den jevne Kone
være saa Meget for Novellens ulykkelige Heltinde og tage hendes
Skæbne i sin Haand. Hvis hun ikke i Virkeligheden har været
Humorist, saa har det været, fordi hendes Sindelag først og
fremmest afsatte sig i Handling (se Kap. VII), ikke forgrenede
sig i Stemninger. ---

Vemod og Længsel kunne ikke blot føre til, at ydre Begivenheder
og Andres Adfærd tages med Humor, men ogsaa vore egne
Fejl kunne vi tage paa samme Maade, især naar vi kunne se
dem som vore Dyders Fejl, som Led eller Dele af vor egen Personlighed.
Naar der saa fra Andres Side rettes Angreb, ville vi maaske
(som Viggo Stuckenberg i „Bekendelse“, overfor en Vens Omvendelsesforsøg)
hævde vore Synders positive Værdi, fordi vi
uden dem ikke vilde være, hvad vi ere. Derfor behøve vi ikke
at undervurdere Angerens Skærsild. Vor Personlighed, som den
nu er, med dens Vemod og dens Længsel, har vel netop gennemgaaet
denne Skærsild; men i Erindringen fortoner den enkelte
Handling sig som Led i hele Personlighedens Udvikling, maaske
endog som en felix culpa, og i Længselen fremad er der ikke
% --- p. 84 ---
\setcounter{footnote}{0}%
\opage{84}Tid til at dvæle ved, hvad der nu ikke mere har praktisk Betydning.


22. De i det Foregaaende skildrede Humortyper staa Selvhævdelsen
nærmere, end de staa Hengivelsen. Overlegenhed og Sindshøjhed,
Vemod og Længsel kunne dog danne Grundlag for en
Spøgen med egen og Andres Færd og Skæbne, der er forbunden
med positiv Værdsættelse af deres Handlen og Liden. I Modsætning
til den Udviklingsgang, der fra den oprindelige Selvopholdelsestrang
fører i Retning af Egoisme (Psykologi VI C, 1), ere
disse Typer Følger af en Udvikling, der har udelukket Hildethed
i det egne isolerede Jeg og dets Interesser. Og det er i Kraft af
stor Energi, at denne Frigørelse er naaet i Livserfaringens Skole.
Der kan da danne sig én Totalfølelse overfor Livets Modsætninger
og Omskiftelser, der ikke blot er betinget ved, at disse have
grebet ind i Menneskets egen isolerede Existens.

Denne Frigørelse fra Modsætningen mellem egen og fremmed
Skæbne bliver endnu fuldstændigere, hvor Humoren udvikler sig
paa Grundlag af Sympati (i Betydning af umiddelbar Medglæde
eller Medsorg). Her er der fra første Færd en Forbindelse
mellem Mennesket og det, der sætter dets Følelse i Bevægelse.
I Modsætning til den Enheds- eller Helhedsfølelse, der her én
Gang for alle forbinde Subjekt og Objekt, kan Meget, der ellers vilde
staa som Genstand for Haan og Uvillie, tages paa spøgende
Maade, fordi det intet Øjeblik glemmes, at der bag det ligger en
Værdi, som maaske endog indirekte lægger sig for Dagen i selve
det Latterlige. Det er ogsaa her det dybe Grundlag, der gør Humor
mulig. Modsætningen til Haan er her den størst mulige, og
dog forsvinder Muligheden for Spøg ikke; det er netop den, der
gør Humor til en Totalfølelse.

Som jeg har vist i min Etik (IX, 1), vilde det være urigtigt,
uden videre at stille Sympati i Modsætning til Selvhævdelse.
Hengivelse kan være en Fortsættelse af Selvhævdelse. Den behøver
ikke at være en rent passiv Tilstand, men kan være Udtryk
for en Kraft, der formaar at omfatte Andet og Mere, end
hvad Selvopholdelsen kræver. Her kan det store Ord gælde, som
er sagt om Kærlighed, at den bærer Alt, taaler Alt, overvinder
% --- p. 85 ---
\setcounter{footnote}{0}%
\opage{85}Alt. Sindshøjhed behøver ikke at have Kærlighedens Præg, men
i Kærligheden, som den kan være, er der Sindshøjhed. Det er
Kærlighed af denne Art, der kan være Grundlag for den store
Humor. Den kan have udviklet sig ad forskellige Veje: gennem
Motivforskydning ud fra Selvhævdelsen, --- gennem Udvikling
paa Grundlag af et Naturinstinkt, --- eller i Sorgens og Smertens
Skole.\footnote[1]{\textit{King Lear} IV, 6, 225 siger Gloster til Edgar: Now, good sir, what are You? --- Og Edgar svarer: A most poor man, made tame to fortune’s blows; who, by the art of known and feeling sorrows, am pregnant to good pity.}

Den aktive, energiske Karakter af den Sympati, som her forudsættes,
vil ogsaa lægge sig for Dagen deri, at den større Kreds
af Værdier og Interesser, til hvilken Mennesket knytter sig i
Sympatien, aabner Mulighed for Sorger og Bekymringer, Mennesket
ellers var bleven fri for. Mennesket vilde kunne have
sluppet nemmere gennem Livet, hvis det ikke havde udvidet
Kredsen for, hvad der kunde sætte dets Følelse i Bevægelse.

Der maa her ikke blot tænkes paa Forholdet til andre Mennesker,
men ogsaa paa Forholdet til Livet eller Tilværelsen idethele.
Der kan gives en Medfølelse, som strækker sig til Naturen,
til alt Liv, der udfolder sig og stræber frem. Mennesket kan knytte
sin Interesse til Noget, der rækker langt ud over det individuelle
Livs Omraade, og hvis Fremgang eller Tilbagegang bestemmer
den Værdi, han tillægger Virkeligheden som den er. Den Humor,
som danner sig paa dette Grundlag, bliver en kosmisk Livsfølelse,
en Grundstemning, der er hævet over Modsætningen mellem Optimisme
og Pessimisme, idet den ikke lukker Øjnene for Disharmonier
og Ulykker, men ikke over Lidelsen i Verden og Undergangen
af saa meget Stort og Skønt mister Enhedsfølelsen med
den Værdistrøm, der stadigt baner sig sin Vej trods alle Hindringer.
Paa dette Punkt er det egentligt først, at Humor bliver til
en Livsanskuelse.

23. Denne sidste Betragtning leder over til en Forudsætning,
som i højere eller lavere Grad er fælles for alle de Former af
Humor, der er skildrede i det Foregaaende. Ved Beskrivelsen af
% --- p. 86 ---
\setcounter{footnote}{0}%
\opage{86}disse Typer er der især lagt Vægt paa, hvad jeg kalder Grundlaget,
d. v. s. det Følelsesliv, der danner Forudsætningen. Men de
intellektuelle Forudsætninger kunne ikke udelukkes fra Beskrivelsen,
hvad enten disse Forudsætninger blot skyldes Menneskets
personlige Erfaringer eller tillige skyldes den videnskabelige Indsigt
i Naturens og Historiens Gang og Love, som Mennesket har
kunnet vinde. Naar disse intellektuelle Forudsætninger fremtræde
med særlig Vægt, danne de en Baggrund, der virker sammen
med Grundlaget til at bestemme den hele Types Klangfarve.
Baggrunden er intellektuel og objektiv, Grundlaget emotionelt og
subjektivt. Baggrunden øver en mere indirekte Indflydelse end
Grundlaget, ligesom Luft og Lys kun faa Indflydelse paa Livet,
naar dette allerede er til i Spiren. Paa samme objektive Baggrund
kunne højst forskellige Stader eller Livsanskuelser udfolde sig,
alt efter det forskellige Grundlag, der er tilstede. Baggrund og
Grundlag ville ikke være aldeles uafhængige af hinanden. Det
kan afhænge af Grundlaget, hvor fast og omfattende en Baggrund,
der udvikler sig, og den intellektuelle Baggrund kan paa
den anden Side hæmme eller fremme Grundlagets Udvikling.

Hvad det saa er, hvortil vi kunne føle os knyttede, var det
end det Største og mest Omfattende, som tænkes kunde, saa er
det dog med i Tilværelsens store Bølgegang. Saa langt vi gaa
tilbage i vort Inderste, og saa højt vi stige i vore Tanker, saa
komme vi dog ikke ud over den store, lovmæssige Sammenhæng,
der viser sig at raade i alle Tilværelsens Egne. Naar
denne Indsigt en Gang er vunden, vil den siden farve alle vore
Oplevelser. Derved hæmmes ikke nødvendigvis Lyst og Evne til
dristig Fremstræben. Erfaring lærer os, at uden store Formaal
og stort Haab udrettes intet Stort. Uden dem staar Livet stille
eller synker ned i passiv Vegeteren. Men vi vide ogsaa, at selv
de største Formaal kun ere Led i Værdiernes Række, der ligesaa
lidt kan tænkes afsluttet som Aarsagernes Række. Det for os sidste
Maal, der danner den Værdihorizont, vort Syn stanser ved, er
dog maaske Begyndelsen til en ny Værdirække. Men det er i
Begrænsningen, Mesterskabet viser sig, og selv om Træerne ikke
kunne voxe op i Himmelen, kan deres Væxt dog være rank og
skøn.

% --- p. 87 ---
\setcounter{footnote}{0}%
\opage{87}Skranke og Modstand møde ikke blot den store Stræben; de
optræde ogsaa der, hvor Kraften knap slaar til for at holde Hovedet
oppe. Lidelsen finder ogsaa Vej til dem, der føre det stille
Liv og ikke udæske Skæbnen. Som det regner over Retfærdige
og Uretfærdige, saaledes rammer Smerten baade de Stille og de
Stræbende. Men ligesom den store Humor ikke mister Smilet,
fordi alvorlig Stræben maa brydes med Skranker, saaledes gaar
dens Smil heller ikke tabt, fordi Lidelsen ogsaa trænger ind paa
den stille Lykkes Vej. Livets Disharmonier fortone sig tilsidst
overfor den store Kraft, hvormed Livet føres. Humoristen flyr
ikke til Ørkenen, fordi ikke alle Blomsterdrømme modnes.

Følgen af det Opgør med Livet, den store Humor forudsætter,
kan blive en ny Ungdom. For den første Ungdom ligger Humor
ikke. Ungdommen lader sin Kraft strømme ud uden at tænke paa
Grænserne, som jo ogsaa kun kunne kendes paa første Haand
ved egen Erfaring. Den er utaalmodig og forarget over enhver
Stansning og enhver Modstand. Den nye Ungdom forudsætter, at
man har formet faste Planer for Kraftanvendelsen, --- at man har
gjort sig det klart, at man lever i en endelig Verden, --- men at
man dog med Frejdighed --- maaske med det, Sibbern har kaldt
en frejdig Livsironi --- stedse tager Livet op paany. Ogsaa den
Kendsgerning, at Livet stadigt byder paa Gentagelser, finder man
en Plads til i sin Livsbetragtning. Den samme Energi, der trods
alle Skranker fastholder de vundne Livsværdier, formaar ogsaa at
bevare Friskheden og Inderligheden trods Gentagelsens sløvende
og udjevnende Tendens.

% --- p. 88 ---
\setcounter{footnote}{0}%
\opage{88}\begin{center}V. TRAGIK OG HUMOR\end{center}


24. Som allerede ovenfor bemærket, kan ingen Totalfølelse godtgøres
at betegne en fuldstændig Afslutning af Følelseslivet. Saalænge
nye Oplevelser komme, vil der ogsaa være Mulighed for
Ændring af den Grundstemning, i hvilken det Oplevede har afsat
sine Virkninger. En fuldstændig Helhedsdannelse, enten som
Sammensmeltning eller som Organisation, er et Ideal. Der kan
ligesaa lidt naas fuldstændig Afslutning paa Følelseslivets som paa
Tankelivets Omraade. Vi staa midt inde i en Udvikling, indenfor
hvilken der stadigt er Mulighed for uventede Erfaringer.

Den store Humor kan derfor ikke proklameres som en Gang
for alle afsluttende Livsstemning. Dog har den mange Betingelser
for at nærme sig dette Ideal. Dens Ejendommelighed er, at Livets
modsatte Erfaringer og Stemninger i den kunne komme til deres
Ret. Den forudsætter, at Sindet har været aabent for alle Sider af
Livet og ikke i Blindhed eller i Fejghed har undgaaet, hvad der
ikke behagede. Den kender baade Tider, hvor der maa leves trods
et Underskud, og Tider, hvor der leves ud af en Fylde, der synes
uudtømmelig; men den gaar ud af disse Svingninger med et Overskud,
der gør Overlegenhed overfor Skæbnens Gang mulig, et
Overskud, der ikke, som Barnets, blot er en Naturgave, men for
en væsentlig Del er selverhvervet, er opsparet Arbejde. Hvad der
i Fortiden er virket og lidt, arbejder man sammen med det Nye,
der kommer. De forskellige Oplevelser staa ikke paa lige Trin.
Det er de mest alvorlige Livserfaringer, om hvilke alle andre Livserfaringer
samle sig uden at formaa at bryde dem, men samle sig
netop saaledes, at de ved Modsætningen til dem fremtræde som Spøg.
I Forvisning om Værdiernes Gyldighed kan der smiles over Virkelighedens
Forsøg paa at være Alt. Der kan smiles ad de stadige
% --- p. 89 ---
\setcounter{footnote}{0}%
\opage{89}Værdiforskydninger, over at man har faaet noget helt Andet ud af
Livet, end man havde tænkt. Wilhelm Meister’s „Lehrjahre“ ender
jo med, at den unge Mand sammenlignes med Saul, Kis’ Søn, der
gik ud for at lede efter sin Faders Æsler og endte med at blive
salvet til Konge og at regnes blandt Profeter. Der er maaske
ogsaa gjort den Erfaring, at de ydre Former for Livet, vi saa
møjsommeligt bygge op til Værn for det, til et Leje for dets Strøm,
netop kun blev Hindringer for dets frie Løb. Kort sagt, Verden er
overvunden i Alvor og kan derfor nu overvindes i Spøg.

Humor stræber at forene Livets Tragedie med dets Komedie.
De to Elementer forenes i den, snart saaledes at de forenes til
Et, snart saaledes at de udgøre en ordnet Helhed. Naar Sproget
har dannet Udtrykket „Tragikomik“, har det derved især betegnet
det Komiskes Overvægt, og derfor passer Ordet ikke paa Humor,
hvor Spøgen bevæger sig paa Alvorens Grund. Ordet har desuden
en afgjort haanende eller nedsættende Betydning. Eller maaske
har Frater Taciturnus Ret, og Ordet betegner Noget, man hverken
kan le eller græde over, fordi det hverken er komisk eller
tragisk. Ordet „komitragisk“ vilde passe bedre paa Humor\footnote[1]{\textit{Kierkegaard: Samlede Skrifter.} VI p. 392.}.

Kierkegaard har hævdet, at den tragiske Kunst trænger mere
til en historisk Baggrund end Komedien\footnote[2]{\textit{Samlede Skrifter.} VI p. 402 f.}. Forsaavidt dette er
rigtigt, maa det vel komme af, at Alvoren først ret gør sig gældende,
naar vi staa overfor Virkeligheden i dens hele Strenghed,
medens det Komiske bedre kan tumle sig i en opdigtet Verden.
Men i den store Komik kan Virkelighedens Baggrund ikke mangle;
den angaar jo Livet, ligesaavel som Tragedien gør. Bergson har
endog anvist det Komiske en ganske særlig Virkelighedsbetydning
ved at fremhæve dets sociale Karakter og dets sociale Funktion\footnote[3]{\textit{Henri Bergson’s Filosofi}, p. 46—51.}.
Kierkegaards Ytring\footnote[4]{Naar Kierkegaard henviser til Lessing, maa det bemærkes, at \textit{Lessing} (\textit{Hamburgische Dramaturgie.} N. 87—92) i Modsætning til Diderot og Hurd hævder, at baade i Komedie og Tragedie skulle Karaktererne være individuelt udprægede, hvad enten de hentes fra Historien, eller ikke.} maa i hvert Tilfælde ikke forstaas, som om
det individuelle og private Liv, Livet i den snevre Kreds ikke
% --- p. 90 ---
\setcounter{footnote}{0}%
\opage{90}skulde have sine Tragedier. Det er den saakaldte realistiske
Literaturs Fortjeneste at hævde dette gennem bestemte Exempler\footnote[1]{Smlgn. \textit{Etik}. XXX, 3.}.
Enhver Lidelse kan faa denne Karakter, hvor den overgaar en betydelig
Personlighed.

Jeg anfører et Exempel fra Virkeligheden, der kan siges at betegne
en Mellemform mellem Tragik og Humor.

Den unge Schweizerinde Adèle Kamm laa fængslet til Sygelejet
af en uhelbredelig Sygdom og hindredes derved i at arbejde for
de Formaal, hun havde sat sig. Hun følte med Angst, at hun var
lukket ude fra Livet. Men hun naaede gennem Resignation til en
ny Livsglæde. Da hun ikke kunde „realisere sit Ideal“, sagde hun,
vilde hun „idealisere sin Realitet“. Hun fik den Tanke at organisere
et skriftlig Samkvem mellem unge Piger, der ligesom hun
vare fængslede til et haabløst Sygeleje og derved udelukkede fra
Deltagelse i det sædvanlige Liv. Organisationen blev meget omfattende
og førte ud over alle konfessionelle Skranker. Da de
Sunde aldrig ret kunne forstaa de Syge, blev dette Fælleskab af
stor Betydning. De unge Patienter meddelte hverandre deres
Stemninger og Erfaringer og udvexlede deres Tanker om, hvad
der holdt dem oppe. Derved støttede de hverandre i deres Kamp
for endnu at faa Noget ud af Livet. Ensomhedsfølelsen forsvandt,
og Adèle Kamm følte sig nu som den, der var med i den store
Kamp mellem Godt og Ondt. Saaledes bevarede hun sig „joyeuse
dans l’affliction“. „Lykken,“ sagde hun, „er kosteligere, naar den
er blandet med Fortvivlelse; Freden er fuldstændigere, naar den
følger efter Storm, Lyset mere straalende efter Mørket.“ Støttet
til sine egne Erfaringer morede hun sig undertiden med at drille
Læger og Præster, naar de kom med deres Anvisninger\footnote[2]{\textit{Paul Seippel: Adèle Kamm.} Lausanne 1912.}.

Det Tragiske i Lidelsen er, at den hindrer Arbejdet for, hvad
der først synes at skulle bringe Enhed og Mening ind i Livet.
Det Komiske kan fremkomme ved, at den alligevel kan vise sin
Afmagt overfor et nyt aandeligt Opsving, den selv indirekte har
fremkaldt. Og anden Gang ligger det Komiske i Andres forgæves
Forsøg paa at behandle Sjæl og Legeme efter overleverede Formler.



% --- p. 91 ---
\setcounter{footnote}{0}%
\opage{91}At der er Smerte og Lidelse i Verden, er en Kendsgerning,
ingen alvorlig Livsanskuelse lukker Øjnene for. Det er en Kendsgerning,
der stiller os overfor store Problemer. Er der nogen
Mening i, at der lides? Selv meget fromme Mennesker kunne
negte, at der er en saadan Mening, og erklære, at naar de tro paa
Gud, er det i den Forvisning, at Aarsag til det Onde kan han ikke
være\footnote[1]{Se f. Ex. \textit{Souvenirs de Marie Flournoy-Burnier.} Genève 1910. (Ikke i Boghandelen). Smertelige Oplevelser lode denne højt begavede og inderligt religiøse Kvinde finde sin eneste Trøst i den Tanke, at al Lidelse skyldtes Naturen, som Gud ikke havde skabt. Med den Konsekvens, der snarere er mulig for en Kalvinist end for en Luteraner, hævder hun: „Si Dieu était véritablement dans le monde, il serait partout le vainqueur.“ --- Hun har ikke vidst, at andre Teister vare naaede til den samme Overbevisning, saa højst forskellige Personligheder som Platon, Holberg, Rousseau, Voltaire, Stuart Mill og William James.}. Det er for at bruge det filosofiske Sprog, Modsætningen
mellem Værdi og Virkelighed, der her træder frem, og særligt
den Vanskelighed, at al Virkelighed ikke kan forstaas ud fra
Værdien. Virkelighedens Rige er langt mere omfattende end
Værdiernes. Kun de mest elementære Værdier, de, der ere umiddelbart
knyttede til Individers og Slægters Bestaaen, --- og knap
nok engang de, --- kunne forstaas rent biologisk. Og dog, som
Tingene nu engang ere, betegner Smerten et højere Trin i Udvikling,
naar den sammenlignes med en Tilstand, der kun er
smertefri, fordi den er begrænset. Der er ofte knyttet Smerte til
Overgangen fra en Livsform til en anden, og den kan være
Symptom paa, at fuld Harmoni mellem Livets Tendenser indbyrdes,
eller mellem Livet og dets Omgivelser endnu ikke er
naat. At det absolut skulde være nødvendigt, at saadanne Overgangs-
og Krisetilstande skulde være forbundne med Smerte, kan
ikke godtgøres. Det maa tages som et blindt Faktum. Men som
Forholdene nu engang ere, er Lidelse et Tegn paa Liv, et Vidnesbyrd
om, at Kampen mellem Værdi og Virkelighed, eller rettere
mellem værdifuld Virkelighed og værdiløs Virkelighed stedse er
staaende. Man vilde kunne slippe for megen Lidelse ved at sænke
sit Niveau og indskrænke sine Fordringer til Livet til et Minimum.

Det er denne faktiske Sammenhæng, den store Humor ser
under Øjne. Den indlader sig ikke paa Spekulationer for at ret%
% --- p. 92 ---
\setcounter{footnote}{0}%
\opage{92}færdiggøre Lidelserne i Verden, ligesaa lidt som den anerkender
teologiske Forklaringer af Lidelsernes Kendsgerning. Men den
tager denne Kendsgerning med op i sin Livsanskuelse, i sit Stemningsopgør.
Det er ikke noget ringe Arbejde at indordne ogsaa
det Bitre og det Mørke i sit Helhedsbillede af Tilværelsen; men
dette Arbejde kan dog lade Kraft tilovers, saa at det Smaa og
Urimelige kan mødes med det Smil, det fortjener.

25. Den store Humor er en Livskunst. Den fører ud over al
æstetisk Leg ved sin klare Indsigt paa én Gang i Livets Disharmonier
og i Livets Værdier. Den dvæler ikke, som Kunst i snevrere
Betydning, ved enkelte Billeder af Livet. Digtekunst og
Billedkunst kunne kun give Exempler paa, hvorledes det kan gaa
i Livet, men kunne ikke afgøre, om disse Exempler betyde Regel
eller Undtagelse. Det gælder baade, naar der skildres Idyller, og
naar der skildres Tragedier. Naturligvis kan man gennem Undersøgelser
af de Emner, en Kunstner vælger, saavel som af den
Maade, han behandler dem paa, lære hans egen Livsstemning at
kende. Her er Kunstens aabne Sted, hvor den paa den inderligste
Maade hænger sammen med Livet. Der ligger Livskunst bag
Digtekunst og Billedkunst, en Livskunst, i hvilken Livets Erfaringer
ere samlede i en praktisk Intuition. Den betyder intet
videnskabeligt Opgør, skønt den videnskabelige Erkendelse, der
maatte staa til Raadighed, kan være indvævet i den. Den store
Humor peger ud over enhver begrænset Erfaring, der vil bestemme
hele Følelseslivet, saavel som ud over ethvert Billede,
Fantasien formaar at forme. Overfor saadanne Begrænsninger anvender
den Spøgen som befriende Middel. Men den er selv ikke
Spøg. Naar Undergangen af noget Værdifuldt oplevedes, --- selvom
det, der gik under, hidtil stod som den højeste Værdi, --- finder
den store Humor ny Bekræftelse paa, at ingen enkelt Værdi kan
omfatte al Værdi. Ved og i selve Undergangen peger den enkelte
Værdi ud over sig selv.

Naar Goethe stod afvisende overfor Begrebet Humor, var det
vistnok, fordi han forargedes over romantiske Digteres burleske
og vilkaarlige Leg med Alt, en Leg, de maaske selv betragtede
som Udslag af Humor; men det var da ikke den store Humor, der
% --- p. 93 ---
\setcounter{footnote}{0}%
\opage{93}laa til Grund. Goethe mente, at Humor ikke havde noget Hold
eller nogen Lov i sig selv og derfor let tidligere eller senere udartede
til Tungsind eller ondt Lune. (Brev til Zelter. 30. Oktober
1808). Han erklærer endog, at den, hvem det er bitter Alvor med
Livet, ikke kan være Humorist; --- Humoristen er, mener han,
mere optagen af sin egen Stemning end af sin Genstand. (Unterhaltungen
mit Kanzler von Müller. 6. Juni 1824)\footnote[1]{Alligevel beundrede Goethe Sterne. Han fandt hans Humor uefterlignelig. Men han tilføjer: „Ikke enhver Humor befrier Sjælen.“ (Sprüche in Prosa).}. --- Hvad Goethe
taler om, kan aabenbart ikke være Humor som Totalfølelse, tilvirket
gennem Oplevelser af aktiv og passiv Art.

I Modsætning til det Begreb om Humor, som Goethe forudsætter,
staar Solger’s og Kierkegaard’s Begreb om Humor som en
paa Alvor som Grundlag og med Forstaaelse som Baggrund udviklet
Livsanskuelse eller Livsstemning, der ikke lader sig det
Ophøjede berøve, fordi det har Sammenhæng med det Lave og
Ringe, men formaar at forene Smilet ved dette med Begejstringen
for hint. Paa en meget klar Maade har Kierkegaard udtrykt sin
Opfattelse ved at give Humor Plads i „Stadiernes“ Række efter
det etiske Stade og foran det religiøse. Det var for ham Hedenskabets
Kulmination, at det havde Aandsstyrke til at se det Komiske
og det Tragiske forbundne til Enhed. Humor var for ham
det sidste Livsstade før Troen. Selvfølgeligt tænkte han kun paa,
hvad der i det Foregaaende er kaldt den store Humor, ikke paa,
hvad han betegnede som den umodne Humor, der kun er Friskfyreri,
et æstetisk Raffinement, der overspringer det Etiske\footnote[2]{\textit{Skrifter}, VI, p. 393, 405, 414. --- VII, p. 249 f.}.

Der gives en „æstetisk“ Maade at tage Livet paa, som ikke har
Karakter af stor Humor. Man stiller sig da udenfor Livet som ren
Betragter og Tilskuer, maaske som lidenskabelig Tilskuer, idet det
kan blive en Lidenskab at se paa Tingene uden at træde i personligt
Forhold til dem og uden at agte paa, om de gribe ind i
andre personlige Væseners Ve og Vel. Man føler maaske netop
en Trang til at trække sig ud af ethvert praktisk Forhold til Tingene
for at leve i Fantasien. Denne Art af æstetisk Følelse er
beslægtet med intellektuel Interesse, kun at den mere gaar ud
% --- p. 94 ---
\setcounter{footnote}{0}%
\opage{94}paa at leve i Billeder, medens den intellektuelle Trang søger Begreber
og Love. Spørgsmaalet er, om en saadan æstetisk-intellektuel
Interesse kan udfylde hele Personlighedens Liv uden at indsnevre
og udtørre det. Goethe’s Protest mod, hvad han forstod ved
Humor, motiveres ved dette Spørgsmaal. Baade i videnskabelig
og i kunstnerisk Virken er det kun en Del af Personligheden, der
er i Funktion. Men hvorledes forholder Tanke og Fantasi sig til
den Glæde og Sorg, den Begejstring og Harme, den Ophøjethed
og Nedstemthed, som vækkes ved den egne personlige Skæbne
eller ved de Erfaringer, som gøres om Kaarene for de Interesser,
der gøre Livet værdt at leve? Her kræves en omfattende Syntese,
selve Personlighedens Syntese, indenfor hvilken de intellektuelle
og æstetiske Synteser skulle finde deres Plads\footnote[1]{Smlgn. min Afhandling \textit{Aandelig Kultur.} (Mindre Arbejder. Tredie Række).}. Den Totalfølelse,
som bliver herskende hos et Menneske, selv hos den største
Videnskabsmand eller Kunstner, bestemmes nok saa meget ved
det, der oplevedes udenfor den særegne intellektuelle eller æstetiske
Sfære, som ved det, der oplevedes indenfor disse\footnote[2]{I Kap. IX vende vi tilbage til dette Punkt.}. Oplevelser
af mere individuel Art ville virke med til at bestemme den
raadende Totalfølelse, og desuden have baade Videnskaben og
Kunsten deres aabne Steder, hvor de peger tilbage paa andre
Sider af det personlige Liv og indtage deres Plads indenfor dette
som Helhed. Det er Mennesket, ikke blot Videnskabsmanden eller
Kunstneren, der faar Udtryk i Livstypen eller Stadet. Den store
Humor er en af de mulige Livstyper (Stader, Personlighedssynteser).

Det er ikke i sig selv nødvendigt, at der er en stor Afstand
mellem Mennesket og Videnskabsmanden eller Kunstneren. Hos
de Største kunne de staa hinanden nær. Men denne Nærhed er
ikke selvfølgelig og kan forudsætte stort Arbejde. Det er i en vis
Forstand menneskeligt, „altfor menneskeligt“, at indrette sig saaledes,
at Aandsarbejdet bliver Noget for sig, som man skyder tilside,
naar man har lukket Døren til sit Arbejdsrum for at deltage
i Livet og tage, hvad det bringer. Men naar den Dør er lukket,
staar det nye Arbejde tilbage at væve Arbejdet og dets Frugter
ind i Livet som Helhed. Ellers biver det let en ynkelig Livstype, % PRINTED AS IS: „biver“ (sense: bliver)
% --- p. 95 ---
\setcounter{footnote}{0}%
\opage{95}der kommer til at svare til et i og for sig maaske betydeligt
Aandsarbejde. Man slipper ikke for at have et Stade, men det kan
blive et Stade derefter.

Nogle Psykologer mene, at Lystfølelsen ved det Tragiske nødvendigvis
forudsætter, at man ikke glemmer, at det Hele ikke er
Alvor. Man er sig bevidst, at man ikke er med, men sidder trygt
paa sin Plads i Parkettet og ser paa, hvad der foregaar paa Scenen.
Paa et Billede af den franske Maler Vibert, „Le Récit du Missionaire“,
ser man en Missionær fremvise sine Ar og fortælle om,
hvad han har lidt for sin Tro; nogle fine og fornemme Gejstlige
lytte med stor Interesse til hans Fortælling; det Hele er for dem
aabenbart en æstetisk Nydelse, som fængsler dem, medens de
sidde blødt og mageligt i deres Sofa. Man faar ikke Indtryk af, at
Fængslingen angaar deres hele Sind. En fuldstændig og bevidst
Tryghedsfølelse maa udelukke den tragiske Virkning. En tydelig
Forestilling om, at det Hele ikke er Alvor, vil hindre den Indleven
i Handlingens Gang, den Deltagelse og Spænding, uden
hvilke hverken Katastrofe eller Løsning faa nogen Betydning.
Selv naar vi i Erindringen fordybe os i vore egne tidligere
Ulykker og Farer, maa vi (deri har Spinoza aabenbart Ret)\footnote[1]{\textit{Ethica} III, 19; 47 Schol.} opleve
dem paany, som om de vare nærværende, for at saa Befrielsen
i Erindringen kan fremkalde en Lystfølelse.

I hvert Tilfælde --- for den store Humor, som vi her beskrive,
er det Alvor med, hvad der optager Sindet. Hvis det er et Skuespil,
saa foregaar det ikke paa „de Brædder, der betyde Verden“,
men i selve den virkelige Verden, den Skueplads, hvor vi selv
spille med, hvad enten vi ville det eller ikke, og hvor vi hverken
kunne overse Spillets Gang eller kende dets Slutning.

26. Enhver Totalfølelse kan trues med Opløsning, enten ved at
væsentlige Elementer falde bort, eller ved at nye Oplevelser kræve
større Energi eller ændret Organisation. Ogsaa den store Humor
maa ofte kæmpe for Livet. Den lyse og energiske Harmoni kan
f. Ex. trues af Hypokondriens klamme Taager, som fremkaldes af
indre organiske Dispositioner eller ydre forstemmende Oplevelser.
En nedtrykt Stemning begunstiger eller hidkalder alle Forestil%
% --- p. 96 ---
\setcounter{footnote}{0}%
\opage{96}linger om Ufuldkommenhed, Ulykke, Ringhed, Utilstrækkelighed.
Intet i os eller udenfor os synes at have Værdi. Selv om der
ellers staar en Tankeverden til vor Raadighed, er det i den hypokondre
Tilstand, som om den aldrig havde existeret, eller som om
det var en Indbildning, at vi levede i og for den. En graa Ensformighed
hersker, og den frugtbare Modsætning mellem det Store,
vi saa op til, og det Smaa, vi smilede ad, denne Modsætning, der
er Betingelsen for Humor, er udvisket. Snart er det det Store, der
ikke kan gøres nærværende, snart er det det Smaa, der ikke kan
ses i dets Værdi. Det daglige Livs Opgaver og Begivenheder staa
som ligegyldige, graa, fortrædelige, hindrende. Selv om Tanken
om det Store ikke er forsvunden, bringes de smaa Ting dog ikke
i Forhold til det, hverken i Kontrastens komiske eller tragiske
Forhold eller i Sympatiens og Forstaaelsens inderlige Forhold, end
sige i alle disse forskellige Forhold paa én Gang. Skulde der ogsaa
falde et Lysskær ind i denne Taage, saa bringer det dog ingen
Glæde.

Paa sit Højdepunkt bliver Hypokondrien til den acedia, den
Utilgængelighed for kraftigt Aandsliv, som i Middelalderen regnedes
til Dødssynderne. Dante satte acediosi i Helvede, hvor de
rode i Dyndet, og lader dem sige (Helvede VII, 121 --- 124):

\begin{quote}
Lidet glade\\
var vi, hvor Solen skinned mildt paa Enge,\\
ti Mismod ulmed i vor Sjæl den lade;\\
nu maa vi drøveligt i Dyndet hænge!
\end{quote}

De gamle Mystikere og Helgener kendte vel denne Tilstand og
betragtede den som Djævelens Værk. Selv om de ellers ønskede
sig Lidelse for derigennem at vise deres Tapperhed og Tro, saa
var denne Art Lidelse dog kun af det Onde, fordi den var ufrugtbar;
den bragte dem, som de udtrykte sig i deres mystisk-erotiske
Sprog, ikke til „at føde“. Den gør Handling umulig\footnote[1]{\textit{Joly: Psychologie des Saints.} Paris 1898, p. 187 f. --- \textit{Vie de Ste. Terèse, écrite par elle-même.} Chap. XXX.}. Boccacio
beskriver den i sin Kommentar til Dante saaledes: „Et saadant
Menneske er ikke istand til at tage noget Initiativ, og nødes han
til det, kan han dog Intet fuldbyrde. Hans Liv flyder hen, som om
% --- p. 97 ---
\setcounter{footnote}{0}%
\opage{97}han ikke levede det; hans Tanker blive stedse mørkere og mørkere;
han skyr Samkvem med Andre\dots{} Først føler han Ligegyldighed,
saa Uvilie, og tilsidst Væmmelse ved alt Godt og Skønt“\footnote[1]{Citeret efter \textit{A. Herrlin: Renaissancens Etik,} p. 50. (Trykt som Manuskript).}.

Ikke blot kan der i saadanne Tilstande intet Initiativ tages, men
saadanne Tilstande ere netop saa pinlige paa Grund af deres ubestemte
Karakter, idet en bestemt Aarsag ikke fremtræder, saa at
man kunde forsøge at gribe ind. Og det kan gaa Hypokondristen,
hvis han i sine gode Tider er Humorist, som det gik Harlekin,
der spurgte en Læge til Raads for sin Melankoli og fik det Raad
--- at gaa hen og se Harlekin. Spøgen, der har hjulpet saa mangen
god Gang til at komme om en skarp Pynt, hjælper nu ikke, af
den gode Grund, at den ikke kan naa op imod Strømmen. Hypokondristen
vilde maaske kunne helbredes, naar han kunde faa et
frit Øjeblik, i hvilken hans Hypokondri selv kunde blive tagen
humoristisk --- ses som en mørk Skygge i Livets Landskab, der
fortoner sig, naar man ser den fra Højden, men overvælder, naar
man er nede midt i den. Men Forudsætningen herfor er, at Højden
er klar, og --- at man kan holde sig paa Højden. ---

Dette dagligdags Exempel peger endda ikke paa de største
Skygger. De fremkomme først, naar det Tragiske truer med for
stedse at sønderrive Livets Sammenhæng. Hvorledes staar den
store Humor overfor denne Mulighed? Kan Livets Tragik ikke
blive et Vidnesbyrd om dens Afmagt? --- Før vi tage dette alvorlige
Spørgsmaal op, ville vi se, hvorledes den største af alle tragiske
Digtere staar overfor det.

27. Shakespeare’s Sind har været aabent for alle Livets Sider,
for alle de Modsætninger, det rummer, og han forstod at give sine
Erfaringer Udtryk i individuelle Skikkelser og Skæbner. Have vi
ikke selv havt Lejlighed til at gøre Erfaringer, kunne vi hos ham
komme til at staa overfor dem. Han vidste, som Spinoza vidste
det, at Naturen er rig nok til at lade de mest forskellige og modstridende
Individer og Skæbner fremstaa. Som Cordelia og Goneril
have samme Forældre, saaledes udspringe Godhed og Ondskab
ifølge de Love, som Tilværelsen nu engang er underlagt. Og Lykke
% --- p. 98 ---
\setcounter{footnote}{0}%
\opage{98}og Ulykke gaa deres Gang gennem Verden ifølge de samme Love.
Ondskaben kan blive saa stor, at man fristes til at tabe Troen
paa Godheden, og Ulykken kan synes saa overvældende, at Alt
synes at være i Opløsning. Der kan være Anledning til at udbryde
med Cordelia, da hun hører om Søstrenes Adfærd overfor den
gamle Fader: Let pity not be believed! (King Lear IV, 3, 30), og
med Kent og Edgar, da den gamle sindssyge Konge kommer bærende
med Cordelias Lig: Is this the promis’d end, or image of that
horror? (King Lear V, 3, 203). At Shakespeares Hjerte er med i
alt dette, mærkes gennem den hele Stemning og gennem enkelte
gribende Ord. Hans Skildringer ere som Krystallisationer, der have
dannet sig under Indflydelse af den dybe Bevægelse, hans Oplevelser
have voldet i hans Sind. Vi staa her ved Kunstens aabne
Sted, hvor vi fra det historisk foreliggende Digterværk komme til
at se ind i Digterens eget Sjæleliv, saaledes som dette stemtes
ved, hvad der vederfores ham, eller hvad han saa i sin Kreds.
Shakespeareforskere mene, at i den Periode af hans Liv, i hvilken
hans mægtigste Tragedier fremkom, og som blev hans Produktions
Højdepunkt, har der for ham eller i hans Kreds fundet rystende
Begivenheder Sted. Det er gaaet ham som Dante: han har
været i Helvede\footnote[1]{Veronas Kvinder pegede paa Dante, naar han gik forbi, og sagde: „Der gaar han, som har været i Helvede.“ Den danske Oversætter af \textit{Vita nuova}, Joh. Dam, føjer til: „Dette var ikke Ironi; han havde været i Helvede, og Alle vidste det“.}. Af hvilken Art disse Begivenheder have været,
derom kan man maaske danne sig en Forestilling gennem de to
tragiske Hovedmotiver, der fremtræde i hans Værk\footnote[2]{\textit{A. C. Bradley: Oxford Lectures on Poetry}, p. 305.}. Det ene
Motiv er skønne og ædle Naturers Forhold til en haard Omverden.
Det andet er voldsomme og dæmoniske Naturers Forhold
til en Omverden, de søge at bøje efter deres Villie. Hvis
den Karakteristik af Shakespeare’s Personlighed, som A. C. Bradley
har givet i sin Afhandling „Shakespeare the man“ (Oxford
Lectures on Poetry) er rigtig, saa har han selv hørt til de aabne
og fine Naturer, som Verdens Stød saa let ramme, og som ogsaa
stærkt paavirkes af, hvad de opleve af Sammenstød og mørk
Skæbne omkring sig. Men han har ikke ligget under for det
% --- p. 99 ---
\setcounter{footnote}{0}%
\opage{99}Overvældende, og maaske har hans Kunst her været hans bedste
Hjælp. Han har i hvert Tilfælde formaaet i sin Fremstilling at
lade Karaktererne, de onde som de gode, udfolde sig i deres
naturlige Konsekvenser, og at lade Skæbnerne, de ulykkelige som
de lykkelige, gaa deres stille og store Gang. Han er ikke bleven
hindret i atter og atter at aabenbare os, hvorledes det Smaa hænger
sammen med det Store, det Groteske med det Tragiske, det
Latterlige med det Ophøjede. Heller ikke har han lagt Skjul paa,
at det Tragiske ikke blot kan fremkomme gennem Karakterens
Sammenstød med Verden, men tillige, som i „Hamlet“, kan udspringe
af Brydninger i Sindets indre Verden mellem stridende
Kræfter. Fremfor Alt have de Rædsler, han har set under Øjne,
eller som hans Fantasi har vist ham Muligheden af, ikke hindret
ham i at vise os det inderlige og dybe Sinds Adel i selve Ulykken
og Undergangen. Som det er blevet sagt\footnote[1]{\textit{A. C. Bradley: Shakespearian Tragedy}, p. 324.}: hvad der hænder et
Væsen som Cordelia, er ikke det Afgørende; det Afgørende er, at
hun er til! Filosofisk kunne vi udtrykke dette saaledes: Shakespeare
mister ikke Troen paa, at der er Værdi i Verden, trods
det, at det Værdifulde saa tit gaar under.

I Shakespeare’s sidste Storværk, „The Tempest“, faar Fredsfølelsen
og Vemodet Overhaand. Det Tragiske mildnes, og det
Komiske forsvinder. Men endnu lige tilsidst, da Prospero mener
at have overvundet alt Ondt, saa han roligt hengiver sig til festlige
Syner, helt glemmende Menneskedyret Kaliban og hans Sammensværgelse,
gennemrystes han ved den pludseligt opdukkende
Tanke om de Disharmonier, der endnu blive tilbage. Han har ikke
kunnet gøre Menneskedyret til Menneske. Og som de festlige
Syner, der saa pludseligt maatte forjages, vil jo hele den Verden,
vi kende, med dens Harmonier og dens Disharmonier, engang opløses
i et Kaos! Men fra Tanken om den sejge Ondskab og om
det Herliges Flygtighed frelses saa hans Sjæl gennem Højsindet
overfor dem, der havde handlet ondt imod ham.

{\small\textbf{Anm.} Der er her paa flere Steder i Afhandlingen, i Tilslutning til forskellige Shakespeareforskere, gaaet ud fra den Forudsætning, at der fra den store Digters Værker maa kunne sluttes tilbage til, hvad han selv har
% --- p. 100 ---
\setcounter{footnote}{0}%
\opage{100}oplevet eller set paa nærmere eller fjernere Hold. Fra denne Forudsætning
gaar ogsaa \textit{Georg Brandes} ud i sin Bog om Shakespeare, hvis
Hovedfortjeneste iøvrigt ligger i de fine Analyser af de enkelte Værkers
Indhold. Derimod negtes denne Forudsætning bestemt af Shakespeare’s
lærdeste og mest kritiske Biograf, \textit{Sidney Lee.} „De biografiske Efterretninger,“
siger han (A \textit{Life of William Shakespeare.} New edition. 1915,
p. 417), „yde ingen Støtte for Antagelsen af en længere personlig Erfaring
om tragisk Lidelse. Ej heller støttes denne taagede Teori af den hele
Retning i hans literære Arbejde\dots{} At søge Nøglen til Shakespeare’s
triumferende Erobring af Tragediens Tinder i hans personlige Erfaring,
som jo nødvendigvis maa have været begrænset, det er at undervurdere
hans skabende Evne og dens magiske Kraft.“ --- Men Spørgsmaalet er jo
netop, om høj kunstnerisk Skaben er psykologisk mulig, uden at der ligger
Livserfaringer til Grund. Og hvis dette maa benegtes, vil jo Antagelsen
af, at Digteren havde bygget paa saadanne, ingenlunde føre til at tænke
ringere om hans Kunst. Snarere maatte man beundre, at Kunsten har
kunnet forme de alvorlige Emner paa saa fuldkommen en Maade.\par}

28. Gives der en absolut Tragik? Gives der Skæbner, ved
hvilke Smerte og Undergang blive det sidste Ord? Hvis det er
saa, staa vi her ved en Grænse for den store Humor. Og er der
en Grænse, maa Humoristen selv anerkende den, og den maa
øge hans Vemod og gøre hans Smil mindre lyst, om der overhovedet
endnu bliver Mulighed for Smil.

I et Brev til Zelter siger Goethe: „Jeg er ikke født til tragisk
Digter, ti min Natur er harmoniserende („conciliant“). Derfor
kan det rent tragiske Tilfælde ikke interessere mig, da det egentligt
i sin Grund maa være umuligt at harmonisere.“ I sine Samtaler
med Kanzler v. Müller kommer han tilbage til denne Tanke.
„Alt Tragisk beror paa en Modsætning, der ikke kan udlignes.
Saasnart der indtræder Udligning, eller den ses som mulig, forsvinder
det Tragiske.“ „Tragisk kalder jeg en Situation, af hvilken
ingen Udgang er mulig, ingen Forsoning („Komposition“) tænkelig.“
(Unterredungen mit Kanzler v. Müller. 6. Juni 1824 og 2. Juli
1830). Han kritiserer Hegel’s Opfattelse af det Tragiske som beroende
paa en Konflikt mellem to hver for sig berettigede Magter.
Hegel’s Yndlingsexempel var „Antigone“, hvis Hovedemne efter
hans Opfattelse var Konflikten mellem Familie og Stat. Saadanne
Sammenstød finde jo Sted indenfor en og samme etisk Helhed;
Familie og Stat ere begge Led i Menneskehedens Liv, og idet
% --- p. 101 ---
\setcounter{footnote}{0}%
\opage{101}denne mere omfattende Helhed netop hævdes gennem ensidige
Bestræbelsers Undergang, fører Tragedien til klarere Forstaaelse
af Livet. Det egentligt Tragiske er altsaa for Hegel ikke Katastrofen,
men Konflikten, og Konflikten har stadigt den nødvendige
Funktion at hævde de forskellige etiske Magters gensidige Begrænsning.


Hegel’s lyse og energiske Optimisme vilde, dersom den kunde
fastholdes, udvide Grænserne for den store Humor som Livsanskuelse.
Men Goethe har desværre Ret overfor Hegel. Der kan
gives, og der gives absolut Tragik. For Hegel er de enkelte Individers
Skæbne ligegyldig, naar blot de store Ideer stadigt godtgøre
deres Magt. Men selv om et Menneskes Værd beror paa den
Sag, han tjener, skulde det saa være ligegyldigt, om han selv,
hvis Sjæl er fyldt af denne Sag, bestaar eller gaar under? Den
tragiske Lidelse kan ikke være ligegyldig, hvor stor Vægt man
saa lægger paa den ofte (ikke altid) frugtbare Konflikt mellem forskellige
Værdier, der ikke have fundet deres Begrænsning. Ideerne
(Værdierne) ere nu engang repræsenterede i Erfaringens Verden
gennem personlige Væsener. Og, som vi have set (§ 27), er
det netop den Maade, paa hvilken den tragiske Helt i selve
Undergangen aabenbarer og hævder sit ophøjede Sindelag, der
betinger Tragediens egentlige Virken paa vort Sind.

Ved det berømte, utallige Gange kommenterede Sted i Aristoteles’
Poetik, hvor han finder Tragediens Virkning at bestaa i, at
den vækker Rædsel og Medlidenhed paa en saadan Maade, at disse
Følelser derved „renses“, har man især været uenig om, hvad der
her forstaas ved „Rensning“. Det er maaske rigtigst, med Erwin
Rhode (Psyche\textsuperscript{2} II, p. 49), at tænke paa en Analogi til den Maade,
hvorpaa Musiken paa en Gang kan vække stærke Følelser og ---
netop gennem Vækkelsen og Udladningen --- bevirke, at de ikke
paa pinlig og hæmmende Maade herske over Sindet. Af en nyere
Tids Musik kunde man her tænke paa Beethoven’s femte Symfoni,
„Skæbnesymfonien“, hvor man først hører Skæbnens tunge
Slag paa Porten og, ved at gennemtrænges af den derved vakte
Stemning, forberedes til at paavirkes af det paa en Gang frejdige
og vemodige Indtryk af Symfoniens sidste Del. Men hvad Aristoteles
i hvert Tilfælde ikke siger Noget om, er, hvilken Ka%
% --- p. 102 ---
\setcounter{footnote}{0}%
\opage{102}rakter „Rædsel og Medlidenhed“ antage, naar de ere blevne „rensede“.
A. W. Benn har derfor kaldt ham en slet Psykolog\footnote[1]{\textit{Aristoteles’ Theory of Tragic Emotion.} Mind. 1914, p. 89.}. Men
Sagen er maaske, at Grækerne ikke havde den Forstaaelse af
Sindets Metamorfoser og Følelsernes Forskydninger, som senere
er vunden. Hvad Aristoteles ikke har kunnet udtrykke paa sit
psykologiske Sprog, kan netop have været det, Benn selv hævder,
at Medlidenheden overfor heroiske Karakterer gaar over til Beundring,
og overfor milde og opofrende Karakterer til Kærlighed.
Og Rædselen maa antages at gennemgaa en tilsvarende Forandring\footnote[2]{Benn fortolker Aristoteles saaledes, at Medlidenheden gaar over til Rædsel, og at denne saa beroliges ved, at den kunstneriske Stemning (plot-interest) faar Overhaand ved, at man véd sig selv i Sikkerhed. Jeg kan ikke finde Noget herom i Texten. Det er baade Rædsel og Medlidenhed, der „renses“.}.
Rædsel og Medlidenhed forsvinde ikke, men „renses“ ved at blive
Elementer i en mere kompleks Følelse\footnote[3]{I den nyeste Tid har den af \textit{Freud} grundlagte „Psychoanalyse“ hævdet, at visse Former for Sindssygdom kunne helbredes, ved at en smertelig Oplevelse, i hvilken den egentlige Grund til Sygdommen laa, drages klart og bestemt frem for Patientens Bevidsthed, som derved „renses“. Metoden kaldes da ogsaa Katarsis, samme Udtryk, som Aristoteles bruger.}.

29. I Tragiken har Humoren en af sine Grænser. Humoristen
bøjer sig ikke som den tragiske Helt sammen om Saaret. Hans
Saar ere ikke dødelige og behøve ikke at forjage Smilet. Den
tragiske Helt kan vise Storhed selv i Undergangen, men Smerten
og Sorgen udelukke Smilets Mulighed.

Humoristen véd, at Livet har sine Tragedier ligesaa vel som
sine Komedier. Hans Livserfaring, baade den, der skyldes personlig
Skæbne, og den, der skyldes Oplevelse af Andres Skæbner,
kan være rig og omfattende. Og hans Stræben gaar, bevidst og
ubevidst, ud paa at lade de store Modsætninger faa deres berettigede
Indflydelse paa den totale Stemnings Klangfarve. Men han
maa indrømme, at hvor det Tragiske er af en saadan Natur, at der
ingen Kraft eller Plads er til at lade Sindets lyse Elementer faa
Indflydelse, der er ingen Organisation eller Sammensmeltning
mulig. Den kosmiske Harmoni, som Humoristen bevidst eller ube%
% --- p. 103 ---
\setcounter{footnote}{0}%
\opage{103}vidst tror paa, styrter sammen, hvor der foreligger en absolut
tragisk Skæbne, --- hvor Lidelse og Undergang ere Et. Kun den
tragiske Helt, ikke Humoristen som saadan, hævder i en saadan
Katastrofe Menneskehedens Adel.

Humoristen er villig til at indrømme den tragiske Helt den 
højeste Plads. Men han viger heller ikke for nogen Anden. Og
især viger han ikke for dem, der tilbringe Livet med at tale om,
hvad tragiske Helte have øvet og lidt, uden selv at have noget
Tragisk i deres Existens. Desuden gør Humoristen Indsigelse
imod, at det Tragiske paa Trods af Erfaring erklæres at udtrykke
Livets inderste Væsen. Det er kun Romantik. Den tragiske Helt
selv vil ikke have Tid og Sans til slige Paastande; han føler sig
netop staaende ude paa Grænsen af alle menneskelige Forhold
og al menneskelig Evne; al hans Energi gaar, uagtet den overgaar
den, Humoristen som saadan raader over, med til at staa,
hvor han staar. Den højeste Patos gør tavs, eller dog ordknap, saa
snakkesalige end de kunne være, der kun kende denne Patos paa
tredie eller fjerde Haand.

Livstragedien, der ligger hinsides den store Humor, kommer,
som et Jordskælv kommer. Den bebudes maaske, ligesom dette,
ved forudgaaende Rystelser. Men det beror ikke paa os selv, om
vi opleve noget Saadant. Maaske kunde vi ved at undlade en
Handling, vort Livsideal kræver, undgaa en tragisk Skæbne. Men
da ligger den etiske Vægt paa Undladelsen af denne Handling, og
dette Punkt hører hen under den Drøftelse af Forholdet mellem
Handling og Humor, som et senere Kapitel vil bringe. Ligesom
Handling betegner Grænsen for Humor fra den aktive Side, betegner
den absolute Tragik denne Grænse fra den passive Side.
Humoristen mener ærligt og redeligt at have søgt og fulgt den
rette Vej uden at have svigtet nogen Opgave, der kunde have
hidført tragiske Konflikter. Han indrømmer, at denne hans Vej har
ført gennem milde Egne, og at den Spænding, hans Stade, som
ethvert alvorligt Livsstade, har medført, er at betragte som Harmoni
i Sammenligning med den tragiske Smerte. Han tør maaske
tilføje, at han ikke vilde vige tilbage, dersom tragisk Skæbne ikke
kunde undgaas, uden at han tog Skade paa sin Sjæl. Hans Beundring
for de tragiske Helte er grundet paa Forstaaelse. Dels
% --- p. 104 ---
\setcounter{footnote}{0}%
\opage{104}kan han meget vel sætte sig ind i Situationer, der for hans eget
Vedkommende ikke komme ud over Mulighedernes Verden, dels
erkender han, at Tragedien aabenbarer en Grænse for Harmonien
mellem Værdi og Virkelighed. Men han giver ikke Afkald paa at
være Helt, fordi han ikke er tragisk. Han fører sin Kamp for at
hævde sig selv gennem Fastholden af sin Overbevisning om Livets
Sammenhæng og Værdi trods alle Disharmonier i det Enkelte.
Denne Kamp kan føre ham lige til Grænsen af det Tragiske. Skal
denne Grænse overskrides, da bliver Spørgsmaalet, om han kan
undergaa den Metamorfose, uden hvilken den tragiske Helt ikke
bliver til.

I Gunnar Gunnarson’s storslaaede Bog, „Borgslægtens Historie“,
hedder det om en ung Mand, der vilde opgive Alt for at drage
om som Asket og Profet (IV, p. 117): „Det faldt ham intet Øjeblik
ind, at han allerede havde lidt et Nederlag, --- at dette Liv,
han havde besluttet at føre, ingen Rod ejede i hans egen Sjæl,
ingen Forudsætning havde i hans tidligere Liv. Han vidste ikke,
at Lidelsen er et Skæbnens Gode, der kun bliver faa Mennesker
skænket i saa fuldt Maal, at de helt blive istand til at se bort fra
eget Ve og Vel og kun leve og tænke for Andre og føle Tilfredsstillelse
derved. Han vidste ikke, at kun hos Smertens Udkaarne
dør al Egenkærlighed.“

Bedre kan Forholdet mellem Tragiken paa den ene Side og det
harmonisk-menneskelige Livsstade, Humoren repræsenterer, paa
den anden Side, ikke udtrykkes. Den Lidelse, der bliver tragisk,
er ikke vilkaarligt frembragt eller opsøgt. Den kommer, naar
Skæbnen banker paa Porten, uden at denne Banken (som hos
Beethoven) er Indledning til en Symfoni.

30. For Kierkegaard var Humor et Stade, der har det religiøse
Stade over sig. Og al sand Religiøsitet var for ham af tragisk
Karakter. Hos det Menneske, der skal være i Tilstand af Religion,
maa efter Kierkegaards Opfattelse Gudsforestillingen virke omdannende
paa hele Livsførelsen. Intet menneskeligt Forhold kan
have Værdi i sig selv, men har kun Værdi, naar det bevidst
underordnes og omdannes ud fra det Evigheds- og Uendelighedssynspunkt,
Gudsforestillingen fører til at anlægge. Af Modsæt%
% --- p. 105 ---
\setcounter{footnote}{0}%
\opage{105}ningen mellem den uendelige Fordring og den menneskelige Endelighed
udspringer en Lidelse, der ikke kan harmoniseres. I
Kristendommen --- vel at mærke, den oprindelige Kristendom,
Urkristendommen --- føres hele Livet med Henblik paa det Hinsidige,
eller rettere paa det eneste Sted i Tiden og i Rummet, hvor
det Guddommelige er fremtraadt. Der gives ifølge Urkristendommen
kun én Værdi, den, der blev sat ind i Verden med Gudmennesket.
Alt Andet er værdiforladt eller værdifjendsk. Urkristendommen
er en Sørgemarsch. For den sande Kristen er Livet
Lidelse, en Lidelse, der ikke blot som i al Religiøsitet fremkaldes
ved Gudsforestillingens absolute Krav, men derved, at der kun
findes Værdi paa et eneste Punkt i Tilværelsen, medens alt Andet
ligger i Mørke.

Jeg gaar her ikke nærmere ind paa, om Kierkegaard har Ret i
denne Karakteristik af Urkristendommen. Jeg skal blot gentage,
hvad jeg flere Steder har udtalt (i min Bog om Kierkegaard og i
min Religionsfilosofi), at Urkristendommen ikke var en Sørgemarsch,
men en Sejrsmarsch, idet Sindet var helt opfyldt af Forventningen
om Kristi nære Genkomst, den egentlige Stiftelse af
„Himmeriges Rige“. Da i Modsætning til denne store, levende
Forventning alt Andet maatte blegne, alle andre Interesser og
Opgaver forsvinde, kan Lidelsen ikke, som Kierkegaard mente,
have været raadende. Overfor de store Billeder, i hvilke Forventningen fandt Udtryk, kunde end ikke Smerten ved Løsrivelse fra
de jordiske og menneskelige Forhold gøre sig gældende. ---

Den Tragedie, Kristendommen ifølge Kierkegaard var, overgik
for ham langt Alt, hvad selv den største Tragiker kunde digte.
„O min Ven,“ siger han („Sygdommen til Døden“. Skrifter XI,
p. 236 f.), „hvad har du vel forsøgt i Livet? Anstreng din Hjerne,
riv enhver Bedækning tilside og blot Følelsernes Indvolde i dit
Bryst, sløjf enhver Befæstning, som adskiller dig fra den, om
hvem du læser, --- og læs saa Shakespeare, og du skal gyse for
Kollisioner. Men de egentlige religiøse Kollisioner synes selv
Shakespeare at have gyst tilbage fra. Maaske lade de sig ogsaa
kun udtrykke i Gudernes Sprog.“ ---

Naar Kierkegaard kalder det religiøse Stade, saaledes som han
opfatter det, tragisk, saa er det en Tragik, som det efter hans Op%
% --- p. 106 ---
\setcounter{footnote}{0}%
\opage{106}fattelse er Menneskets Pligt at gaa ind i. Det er villet Tragik.
Den kommer ikke over Mennesket, men Mennesket skal komme
til den. Den skyldes ikke Sammenstød, der kun falde i Enkeltes,
ikke i Alles Lod. Men den er den højeste Form for menneskeligt
Aandsliv.

I de to Pseudonymer Frater Taciturnus (Stadier paa Livets Vej)
og Johannes Climacus (Uvidenskabelig Efterskrift) har Kierkegaard
fremstillet Repræsentanter for, hvad han forstod ved Humor.
De staa begge med Hatten i Haanden overfor det religiøse Stade.
Dette er konsekvent fra Kierkegaards eget Stade, efter hvilket tilsidst
Lidelse blev den afgørende Maalestok for Aandslivets Højde\footnote[1]{Smlgn. \textit{Danske Filosofer,} p. 165—67.}.
Men det stemmer ikke med Begrebet Humor, saaledes som dette
i det Foregaaende er udviklet. Efter dette er et tragisk Forhold
en Skæbne, der kan hidføre en alvorlig Prøve, men ikke Noget,
det er Enhvers Pligt at opleve.

For Humoren er al Værdi ikke indeholdt paa noget enkelt Sted
i Tid og Rum. Der er en stadig Mulighed for at opdage og opleve
nye Værdier i Verden, til samme Tid som der er klar Bevidsthed
om, at hvis Sætningen om Værdiens Bestaaen gælder, saa er det
i hvert Tilfælde ikke de empiriske Værdier, der bestaa. Gives der
evige Værdier, saa kundgøre de sig gennem Lovene for de empiriske
Værdiers Opstaaen og Forgaaen, om vi kunde finde disse
Love. Evighed og Tid ere for Humoristen ikke absolute Modsætninger,
saa at der nødvendigvis skulde være et smertefuldt Spændingsforhold
mellem dem. Evigheden er ikke Noget, der ventes
eller ligger langt ude i Tiden. Den er et stedse Nærværende, betegner
den store Lovmæssighed, der gør al Udvikling i Tiden forstaaelig.
Vi forudsætte den stadigt, hvad enten vi se tilbage eller
gaa fremad. Kierkegaards Sætning, at vi forstaa baglængs, men
leve forlængs, er ikke ganske rigtig. Ti hin store Lovmæssighed
er af Betydning for vor Fremstræben, ikke mindre end for vort
Tilbageblik; gennem Erindringen om Fortiden bestemmes, hvad vi
kunne gøre i Nutiden for at forberede Fremtiden. Humoren ser
ikke Virkeligheden som et Frafald fra Evigheden, men ser netop
enhver Virkelighed som en enkelt Strækning af en evig Strøm.
Kierkegaard’s Humorister se paa romantisk Vis Erindringens Ver%
% --- p. 107 ---
\setcounter{footnote}{0}%
\opage{107}den som en Drømmeverden, over hvilken man kan glemme den
stedse henrindende Tidsexistens. Men der er et inderligt Baand og
en stadig Vexelvirkning mellem Erindring og Liv. Erindringens Verden
befolkes stadig fra de Oplevelser, det handlende Liv medfører,
og Handlingens Liv besjæles af, hvad der bevares i Erindringen.
Ved Opdagelser af nye Værdier virker Erindringen stedse med.

Muligheden af Tragik skal ikke bestemme hele Livet eller
Synet paa Livet. Naar Humor betyder Troskab mod Livet, fordi
den giver det Ringe og det Hæmmende dets Plads ligesaavel som
det Store og det Hjælpende, saa fører den ogsaa til at gaa den
tragiske Skæbe imøde, hvis denne virkeligt er en Følge af, at man % PRINTED AS IS: „Skæbe“ (sense: Skæbne)
har været tro mod Virkeligheden og dens Opgaver. Her falder
da Modsætningen mellem Humor og Tragik bort, og der viser sig
Muligheden af en energisk gennemført kontinuerlig Overgang fra
det ene Stade til det andet. Der er da intet Spring, eller, om der
er et Spring, saa er Kontinuiteten ikke brudt, naar Springet virkeligt
fører over paa den anden Side. Et Spring betyder jo dog ikke
uden videre et Fald. Det frie Menneske, hvis Tanke ikke sysler
med Døden, men med Livet, vil dog gaa i Døden for ikke at fornægte,
hvad der ene gør Livet værdt at leve og at tænke over.

Humoristen kan godt blive bestyrtet, ligesom Prospero, da han,
fordybet i sine Syner, havde glemt Kaliban og hans Sammensværgelse.
(Hvor Mange er det ikke i Aaret 1914 gaaet som
Prospero!) Men bag den store Humor ligger stedse Troen paa
Muligheden af nye Værdiers Udløsning. Hvis alle Værdikilder i
Verden vare udtømte, vilde der ikke længere være Mulighed for
den Totalfølelse, vi kalder Humor, --- ja, ikke for nogensomhelst
Totalfølelse. Det Samme vilde være Tilfældet, dersom Tilværelsen
skulde vise sig at være saa tragikomisk, at der var bundne (potentielle)
Værdier nok, men ingen Mulighed for at udløse dem. Dette
vilde Humoristen ikke kunne tage humoristisk. Han maatte da
indrømme, at Tragedien fik det sidste Ord. Ja, hvis det saa endda
var Tragedie, med al den Inderlighed og Højhed, den kan fremkalde!
Men hvis Tilværelsens Værdikilder hørte op at rinde, vilde
der indtræde en fremskridende Svækkelse og Sløvelse af det
menneskelige Aandsliv, og den sidste Akt af Menneskelivets
Drama vilde just ikke faa det Ophøjedes Præg.

% --- p. 108 ---
\setcounter{footnote}{0}%
\opage{108}Den store Humor staar og falder med den Forudsætning, at
hvad der forgaar, kun er enkelte, i Erfaringen givne Værdier.
Deres Forgaaen er det tragiske Element i de Oplevelser, den
store Humor fremgaar af. Men dette Element bliver ikke eneraadende,
saalænge der opdages nye Værdikilder. I selve den
Sindshøjhed, der kan aabenbare sig, naar en Værdi forgaar, fremtræder
der en ny Værdi. Og derefter tages saa Livet op paany.
Ligesom Kent, da han erklæres fredløs, „gaar sin gamle Vej paa
nye Steder“, saaledes tages Kampen for Tilværelsens Værdier op
efter et Nederlag og føres maaske i nye Former. Den døende
Hamlet giver i sit sidste Aandedræt sin Stemme ved Kongevalget
til den unge kække Fortinbras, hvis Handlekraft og Vovelyst vil
føre ham til andre Skæbner end dem, der faldt i den tungsindige
Drømmers Lod.

31. Om Kierkegaard har Ret i at sætte den skarpe Modsætning
mellem Humor og Religion, kommer naturligvis an paa, hvorledes
man definerer Begrebet Religion. Har Kierkegaard Ret i sin Definition,
staar Humoristen afvisende overfor Religionen og vægrer
sig ved at anerkende den som et højere Livsstade. Men det maa
stadigt anerkendes som en stor Fortjeneste af Kierkegaard, at han
saa bestemt og positivt anerkender og begrunder Existensen af et
selvstændigt, højt og ædelt menneskeligt Aandsliv udenfor (efter
hans Opfattelse forud for) positiv Religions Omraade. Denne Anerkendelse
hang nøje sammen med hans Livssag, den at lade det
Kristelige fremtræde i sin strenge Renhed. Megen Forstaaelse
kan man just ikke sige, at denne Side af hans Tankegang har
været Genstand for efter hans Bortgang, hverken i Kirken eller
udenfor Kirken. Navnligt synes Niveauet i denne Henseende at
være sænket indenfor Kirken, hvis Prædikanter have en stigende
Tilbøjelighed til at mene, at Alt, hvad Ordene Aand og Sjæl kunne
betyde, hører under dem. Det var dog deres naturlige Pligt og
Opgave at gøre opmærksom paa det selvstændige og høje Aandsliv,
der er levet og kan leves udenfor Kirken, for at undgaa Forvexling
og Sammenblanding af Kristendom og Humanitet. Men
det er naturligvis lettere at præke mod „Fornuften“ og den onde
Villie end at sætte sig ind i Livsstader, der ere byggede op paa
% --- p. 109 ---
\setcounter{footnote}{0}%
\opage{109}helt andet Grundlag end deres eget. Og dog ville Mange have
gjort den Erfaring, at jo mere de have frigjort sig fra dogmatiske
Forudsætninger, og tillige fra de antipatiske og polemiske Stemninger,
som Arbejdet paa Frigørelsen let avler, des mere overbevises
de om, at der gives et selvstændigt menneskeligt Aandsliv,
der har udviklet sig historisk under Indflydelse fra mange
Sider, fra Religion og fra Videnskab, men først og fremmest fra
og i Livet selv, et Menneskeliv, der igen kan fremtræde i meget
forskellige Typer. En af disse Typer er det særlige Emne for
denne Afhandling.

Hvad nu Humoristen angaar, vil det neppe ligge ham meget
paa Sinde, om man vil kalde ham religiøs eller ikke. Han vil
maaske forholde sig humoristisk overfor de forskellige Forsøg paa
at sige, hvad Religion „egentligt“ er. Dog vil han (det maa jeg
selvfølgeligt mene) ikke have Grund til Humor overfor mit religionsfilosofiske
Forsøg, da den af mig foreslaaede Definition udtrykkeligt
opstilles dels som en religionshistorisk Hypotese, dels
som en psykologisk Hypotese. Han vil maaske endog føle et vist
Aandsslægtskab overfor min Opfattelse af Religion som Tro paa
Værdiens Bestaaen. Ti den store Humor bliver kun mulig ved en
Fastholden af, at Værdien i Tilværelsen ikke forsvinder, fordi
enkelte i Erfaringen givne Værdier forsvinde.

Det er et religionshistorisk Faktum, at Menneskers Tro bestemmes
ved, hvad de i deres Liv have lært at tillægge Værdi.
Guderne staa som de Magter, der hævde disse Værdier, og Himmelen
er Indbegrebet af de Værdier, man kender. Naturligvis kan
man i sine Undersøgelser af de historisk givne Religioner holde
sig til de udvortes Kultusformer og Dogmer uden at spørge om,
hvilke psykologiske Forudsætninger de hvile paa. Man kan være
en meget lærd Mand uden at være Psykolog. Men det er i hvert
Tilfælde et Forsøg værdt, om man ikke skulde kunne trænge ind til
det psykologiske Grundlag. Det er selvfølgeligt, at jo mere vi gaa
tilbage til primitive Religionsformer, des mindre Oplysning vil
man kunne faa om den Maade, paa hvilken de organiserede Kultusformer
ere blevne til. Selv en saa elementær Religion som den
centralaustraliske kendes kun som et færdigt Hele. Og dog maa
den have havt sin Udviklingshistorie ligesom alt Andet i Verden,
% --- p. 110 ---
\setcounter{footnote}{0}%
\opage{110}og ethvert Forsøg paa at gøre sig Rede for den fører tilbage fra
Sociologi til Psykologi\footnote[1]{Se min Afhandling \textit{Sociologi og Religionsfilosofi.} (Kritik af Durkheim’s Bog „Les formes elementaires de la vie religieuse“) i Tidsskriftet „Vor Tid“. 1914—15.}. Ved de Religioner, vi kalde de højere,
er det psykologiske Grundlag tydeligt fremtrædende, idet de guddommelige
Væsener, paa hvilke der tros, opfattes som Sandhedens
og Godhedens Forkæmpere. Hvis der til Grund for Religion ligger
en Angst (en Sætning, der ikke kan opstilles som almindelig), saa
er det en Angst for at miste, hvad der --- paa det givne Trin ---
giver Livet Værdi. Sætningen om Værdiens Bestaaen ligger til
Grund for al Religion paa lignende Maade, som f. Ex. Aarsagssætningen
ligger til Grund for den praktiske Menneskeforstand,
naar denne søger Midler til at opnaa sine Ønsker, eller naar den
ved paafaldende Begivenheder uvilkaarligt ser sig om efter, hvad
der kan være gaaet forud. Og ligesaa lidt som den praktiske
Menneskeforstand behøver at kende den abstrakte Aarsagssætning
eller at vide, hvad der ligger i den, ligesaalidt behøver den religiøse
Bevidsthed at kende den abstrakte Sætning om Værdiens
Bestaaen, dens Konsekvenser og de Problemer, den stiller.

Hvorvidt Trangen til at fastholde Troen paa Værdiens Bestaaen
vil holde sig, naar den ikke støttes ved at være iklædt
Kultuslivets og Dogmets Former, det kan kun Erfaring afgøre.
Det er en psykologisk Opgave at undersøge Betingelserne for en
saadan Trangs Vedvaren, ligesom det f. Ex. er en psykologisk Opgave
at undersøge Betingelserne for den store Humors Opstaaen
og Bestaaen.

Det er Spændingsforholdet mellem Værdi og Virkelighed, der
betinger de mest indgribende menneskelige Oplevelser og derved
igen bestemmer Totalfølelsers Bestaaen. Forsaavidt Humoristen
interesserer sig for Fortidens aandelige Liv, vil han kunne bruge
Sætningen om Værdiens Bestaaen som Arbejdshypotese. Ved ethvert
Dogme og enhver Kultus vil han spørge, hvilke Værdier
man her har følt eller føler Trang til at hævde eller se hævdede.
I Dyrkelsen af Kurotrofos, Demeter, Isis og Madonna sporer han
saaledes Trangen til at hævde Moderkærligheden som guddommelig.
I den Række af forskellige Forestillinger, Grækerne til for%
% --- p. 111 ---
\setcounter{footnote}{0}%
\opage{111}skellige Tider gjorde sig om Zeus, vil han spore en stigende
Tendens til at betragte Retfærdighed som den højeste Egenskab.
Men ikke blot til Forstaaelse af Fortiden, ogsaa ved Handling i
Nutiden behøves hin Arbejdshypotese. Ved al Stræben efter et
Formaal maa det forudsættes, at Tilværelsen indeholder Mulighed
for Udløsning af Værdier. Ellers falder det Haab, der maa bære
al Handlen, bort. I en af alle Værdimuligheder forladt Verden, er
Handling umulig.

Denne Betragtning danner en naturlig Overgang til at undersøge
den store Humors Forhold til Forstaaelse og til Handling.

% --- p. 112 ---
\setcounter{footnote}{0}%
\opage{112}\begin{center}VI. FORSTAAELSE OG HUMOR\end{center}


32. Den Baggrund, Humoren forudsætter, betinges for en
væsentlig Del ved den Forstaaelse, Mennesket har af sig selv,
af sine Livsvilkaar og af den Tilværelse, hvis Love han er underlagt.
Den Forstaaelse, der her er Tale om, kan fremtræde i saare
mange Grader og Former. Men skal den faa Betydning, maa den
tilfredsstille to Fordringer.

Den maa være saa omfattende og saa sammenhængende, at
Mennesket kan have Forestillingen om en stor Tingenes Orden,
i Forhold til hvilken de Skuffelser, Modsigelser og Lidelser, Livet
bringer, kunne staa som Noget, der dog ogsaa hører med til Tilværelsen
og har sin Ret til at være, paa samme Tid som det ikke
formaar at bryde Tilliden til den store Sammenhæng, indenfor
hvilken Livets Værdier udfolde sig. Den Modstand, der mødes,
overrasker maaske i Øjeblikket, men indordnes dog som Led i
Erfaringernes Række og medvirker til at bestemme Totalfølelsens
Klangfarve. Luther blev efter Sagnet forstyrret af Djævelen under
sit Arbejde paa Bibeloversættelsen; men da han havde opdaget,
hvem han havde for sig, sagde han blot: „Vær du der længe
nok!“ Djævelen havde nu engang sin Plads i hans Verdensorden
og kunde tages humoristisk. Enhver Livsanskuelse har sin Djævel,
--- Oprører, Frister, Anklager, --- som det gælder om at sætte
paa den Plads, der tilkommer ham. Der er stedse spredte Elementer
tilbage fra det Kaos, ud fra hvilket Karakterudviklingen
foregaar, og der kan under Livets Gang komme nye, uorganiserede
og usammensmeltede Elementer til, der kunne sætte baade
Totalfølelse og Forstaaelse paa Prøve. (Smlgn. § 26.)

Den anden Fordring, Forstaaelsen maa opfylde, er den, at den
skal kunne gøre sig gældende paa umiddelbar Maade og med
% --- p. 113 ---
\setcounter{footnote}{0}%
\opage{113}samlet Kraft. Den maa fremtræde som en praktisk Intuition, som
en Forvisning, der er Frugt af Menneskets Erfaring og Eftertanke.
Der behøver ikke at være nogen udtrykkelig Bevidsthed om de
Kilder, fra hvilke denne Intuition stammer. Hvis der gøres Forsøg
paa at give Grunde for Intuitionen, vil det i Reglen vise sig, at
disse Grunde ere fundne eller konstruerede bagefter og ikke have
Noget at gøre med de Aarsager, en rent objektiv Undersøgelse
opdager. Men hvorledes dette end forholder sig, saa maa For-
staaelsen, hvis den skal være Baggrund for Humor, have Koncentrationens
Karakter, staa som en fast Kyst, mod hvilken Oplevelserne
brydes uden at kunne bryde den.

Til den Forstaaelse, der her er Tale om, kunne Bidrag ydes fra
højst forskellig Side. Den kan snart mere skyldes Hovedet, snart
mere Hjertet. Det kan være Videnskab, der for den væsentligste
Del danner den intellektuelle Baggrund, eller det kan være en
religiøs Forestillingskreds, Spekulationer eller poetiske Syner og
Billeder, eller maaske blot den sunde Menneskeforstands uvilkaarlige
Overblik over Tingene. Det Afgørende er, at der paa en
eller anden Maade har dannet sig en fast Tankesammenhæng.
Spinoza siger med Rette, at Følelser, der ere fast knyttede til en
sammenhængende Tankegang (affectus secundum ordinem ad intellectum
ordinati et concatenati), vanskeligere opløses end vage
og usammenhængende Følelser. Dog er det af Betydning, at Arbejdet
paa at naa en fast Tankesammenhæng stedse kan tages op
paany. Kan Djævelen ikke forstaas indenfor den Sammenhæng,
vi hidtil have fundet, saa kan maaske hans uventede Optræden
give en Anledning til at indse, at vi have arbejdet med for snever
Horizont.

Den intellektuelle Baggrund er kun en enkelt Betingelse for
Humor. Det ses deraf, at selv hvor denne Baggrund er den samme
hos flere forskellige Individer, er det dog ikke sagt, at den store
Humor udvikler sig hos dem alle. I Afhandlingens to sidste Kapitler
vil jeg komme nærmere ind paa dette Spørgsmaal. Her vil
jeg kun bemærke, at der dog vil være et vist Slægtskab mellem den
intellektuelle Baggrund og det Grundlag (Overlegenhed, Sindshøjhed,
Vemod, Længsel, Sympati se Kap. IV), som den store Humor
forudsætter. Baade dette Grundlag og hin Baggrund føre,
% --- p. 114 ---
\setcounter{footnote}{0}%
\opage{114}hver paa sin Maade, Sindet ud over den snevre Horizont, inden
for hvilken Livet saa ofte bevæger sig. (Smlgn. § 13 c og § 23).

33. Hvor forskellige de Tankeverdener end kunne være, som
danne den store Humors intellektuelle Baggrund, saa maa de dog
alle hænge nøje sammen med den Maade, hvorpaa man til given
Tid og paa givet Sted skelner mellem Indbildning og Virkelighed.
Den Maalestok, efter hvilken man foretager denne Skelnen, altsaa
Virkelighedskriteriet, er stedse Udgangspunktet for Tankeverdenens
Udvikling. Og paa alle Standpunkter bestaar dette Kriterium
i den faste Sammenhæng, man mener at kunne finde mellem de
Iagttagelser, der staa til Raadighed, --- i Muligheden af en Totalsammenhæng,
i hvilken de enkelte Oplevelser kunne indordnes.
Det bestaar aldrig i noget aldeles Enkelt og Enestaaende; selv
om en enkelt Begivenhed eller en enkelt Iagttagelse skulde have
gjort et saadant Indtryk, vakt en saadan Følelse, at alt Andet taber
sin Betydning i Forhold dertil, saa vil der dog stedse være en Bestræbelse
for at finde et Forhold mellem hin enkelte Oplevelse og
alle de andre. Disse ville blive tydede efter den, naar den ikke
kan tydes efter dem. De kirkelige Mystikere tydede deres ekstatiske
Oplevelser efter disses Forhold til den Sammenhæng, Kirkens
Dogmatik hævdede\footnote[1]{Smlgn. min \textit{Religionsfilosofi} § 62.}; hvor ekstatiske Tilstande oprinde hos
psykologisk erfarne Mennesker i nyere Tid, vil der derimod
hævdes en bestemt Modsætning mellem, hvad der opleves i saadanne
Tilstande, og de dogmatiske Forestillinger, det gængse religiøse
Liv hviler i.\footnote[2]{\textit{Th. Flournoy: Une mystique moderne.} (Archives de Psychologie XV. Genève 1915), p. 99 ff.; 122 f.; 134 f.} Forholdet mellem Oplevelse og Tydning er
idethele forskelligt paa de forskellige Udviklingstrin, saaledes som
Auguste Comte har paavist i sin Lære om de forskellige Stadier
i Erkendelsens Udvikling. Videnskabelig Tydning bestaar i at lade
den ene Oplevelse belyse den anden, indtil de alle tilsidst vise
sig at staa i en stor indre Sammenhæng, saaledes at man maaske
endog kan udlede de følgende Oplevelser af de foregaaende.

Allerede den Opgave, der stilles ethvert Menneske, den, at
% --- p. 115 ---
\setcounter{footnote}{0}%
\opage{115}orientere sig i Verden, at skelne mellem, hvad der er Iagttagelse,
hvad der er Erindring, og hvad der er Fantasi, vil føre Mennesket
til at se sig selv som Led i en stor Sammenhæng. Hvad han kalder
Virkelighed, bliver da den Helhed, han bevidst eller ubevidst
danner af de Iagttagelser og Erindringer, der staa til hans Raadighed.
At der existerer en Virkelighed for ham, viser han gennem
de Handlinger, han foretager; de forudsætte jo, at han mener sig
staaende overfor Noget, som er Andet og Mere end hans egne
Forestillinger og Fantasier. Her gælder det gamle simple Kendemærke:
Vis mig din Tro af dine Gerninger! --- Hvor stor den
Helhed er, der for den Enkelte udgør Virkeligheden, og hvor fast
den er sammenføjet, kan være højst forskelligt. Derfor ere de
menneskelige Virkelighedsbilleder saa forskellige. Men den Lov,
efter hvilken de dannes, er stedse den samme, lige fra de Vildes
animistiske Naturopfattelse til den moderne Videnskabs lovbundne
Kosmos. I Bestræbelsen for at sammenarbejde saa mange Iagttagelser
og Erindringer saa nøje som muligt, kundgør det store
Sandhedsbegreb sig, der under forskellige Former og i utallige
Grader ligger til Grund paa alle Trin af menneskeligt
Aandsliv.

Humoristen er Realist. Hans Ophøjethed over det Enkelte har
Virkelighedens store Sammenhæng til Baggrund. Derpaa beror
det, at Humor (særligt Shakespeare’s Humor) kan kaldes Troskab
mod Virkeligheden. Troskab hviler paa den Forudsætning, at den
eller det, jeg er tro imod, forbliver Et med sig selv under alle
Omskiftelser, udvikler sig i Kontinuitet med sig selv. Ingen enkelt
Iagttagelse er her nok. Der maa, uvilkaarligt eller med fuld
Bevidsthed, dannes en Forestilling, en Opfattelse, som kan
svare til alle de forskellige Iagttagelser, hvori Troskabens Genstand
træder frem, en saadan Forestilling, som jeg i min Psykologi
har kaldt en typisk Individualforestilling. Troskaben forudsætter,
at denne Forestilling stadigt vil gælde, at den Lov, hvorefter
Genstanden hidtil har udviklet sig, vil være Udtryk for dens
Udvikling fremdeles. Og Troskaben selv bestaar da i en Kontinuitet
i vor egen Færd, der svarer til Kontinuiteten i Genstandens
Væsen. Det er hverken for Genstandens eller for selve Troskabens
Vedkommende Uforanderlighed, der forudsættes; Forud%
% --- p. 116 ---
\setcounter{footnote}{0}%
\opage{116}sætningen er kun, at Forandringerne ske i samme Retning og
efter samme Lov.\footnote[1]{Smlg. \textit{Psykologi.} V B, 9; VI C, 6. --- \textit{Personlighetsprincipen i Filosofin,} p. 18—27.}

Troskab mod Virkeligheden vil da vise sig i, at vi stole paa de
Love, der hidtil have vist sig at gælde for vore Oplevelser, og
ere villige til at tage dem, som de ere, med de Goder og Onder,
de Opgaver og Arbejder, de Midler og Hindringer, som de føre
med sig.

Man kunde indvende mod det nøje Forhold, som her er hævdet
mellem Humor og det, der staar som Virkelighed, at Falstaff dog
maa kaldes Humorist, men at han jo egentligt lever paa en stor
Løgn. Hertil vilde jeg svare, at Falstaff’s komiske Virkning paa
os netop beror paa, at han indenfor den Virkelighed, han konstruerer
sig, en Virkelighed, i hvilken Alt faar den modsatte Betydning
af, hvad det har i den „virkelige“ Verden, --- at han
indenfor denne Verden optræder med Troskab og Konsekvens.
Han forklarer sin Fedme ved, at Sorg og Modgang blæser et
Menneske op; han reducerer Æresbegrebet in absurdum, fordi
Æren ikke kan sætte en Arm eller et Ben paa igen; og som den
gamle Svirebroder, han er, forsvarer han sin Sviren med, at unge
Mennesker dog ogsaa skulle leve! Med fuldt Alvor opbygges hele
det System, ud fra hvilket han fælder sine Domme. Den Verdensorden,
han konstruerer, lader netop Individer som ham staa som
berettigede. Hans Humor er Trofasthed mod det Virkelighedsbillede,
han har dannet sig, og som han ufortrødent fastholder
overfor Vennernes Kritik og ydre Forholds Indgriben.

Naar Falstaff ikke blot selv optræder som Humorist, men ogsaa
tages humoristisk af os, er det, fordi hans Billede faar sin
Plads og sin Betydning gennem Kontrastvirkning indenfor hele
det Billede af den alvorlige Tid og dens Kampe, som Shakespeare
opruller i „Henrik IV“. Han maa ses i hele denne Sammenhæng
for at kunne stemme os humoristisk. Sidney Lee mener, at naar
Falstaff virker humoristisk, er det, fordi Kontrasten mellem hans
Alderdom og hans kaade Levevis bringer det melankolske Element
ind, som ikke kan undværes i Humoren. (Life of Shakespeare
p. 245.) Men en saadan moralsk Betragtning vilde ganske
% --- p. 117 ---
\setcounter{footnote}{0}%
\opage{117}sikkert umuliggøre den humoristiske Virkning, ligesom det i
Goethe’s „Faust“ vilde gaa for Mefistofeles’ Vedkommende. Man
maa stille sig paa Guds Standpunkt i Forspillet og sige, at der nu
engang ogsaa maa gives saadanne Fyre.\footnote[1]{Ligesom Gud betragter Mefistofeles humoristisk, saaledes ogsaa omvendt. Mefistofeles’ Humor overfor Gud („Von Zeit zu Zeit seh’ ich den Alten gern“) betinges ved, at han ser Gud som et Element i sin egen, mefistofeliske Verdensorden; ud fra denne kunde han fra sin Side have sagt om Gud: „Es muss auch solche Käuze geben!“ --- Ogsaa i mindre Dimensioner kan en saadan gensidig Humor gives. En gammel Skolekammerat af mig, som var bleven en ivrig Politiker, sagde i sin Tid til mig, da vi havde talt om min videnskabelige Interesse, at han oprigtigt talt saa humoristisk paa alt Saadant og kun vurderede det efter dets Betydning paa det politiske Marked. Med den bonhomie, som han stedse bevarede, føjede han til, at han godt kunde forstaa, om jeg fra mit Synspunkt ogsaa tog hans politiske Stræben humoristisk. Jeg tror dog, at ingen af Delene slog ganske til.} Det er i Kraft af dette
Princip, at Falstaff gør sin humoristiske Virkning. Derfor virker
det skurrende og skærende, naar Prins Henrik, saasnart han er
bleven Konge, forstøder sin gamle Staldbroder med haarde Ord.
Som Bradley træffende siger i sin Afhandling „The Rejection of
Falstaff“ (Oxford Lectures p. 259): „Shakespeare havde her skabt
saa vidunderligt et Væsen og sat ham saa fast paa hans intellektuelle
Trone, at Forsøget paa at støde ham fra Tronen igen ikke
kunde lykkes.“ Shakespeare var her ikke tro mod det Billede af
det virkelige Liv, han havde skabt, mod den Virkelighed, der ikke
blot frembød Karakterer som den bekymringsfulde Konge, den
fremstormende Hotspur og den heltemodige Prins Henrik, men
ogsaa havde en Plads til gamle Hans, hans Vid og hans Meriter.

34. Vort Verdensbillede kan og maa ændres i Mangt og Meget
under fremskridende Erfaring, og selv om vi stole paa, at Forandringerne
følge den Lov, vi hidtil have sporet, saa er Begrebet
Virkelighed dog et Idealbegreb, og Begrebet Sandhed --- som den
størst mulige Sammenhæng mellem det størst mulige Antal af
Iagttagelser --- indeholder en Opfordring til stedse fortsat Forsken.\footnote[2]{Smlgn. \textit{Psykologi} V D. --- \textit{Den menneskelige Tanke}, p. 257---262; 289 f.}
En absolut fast Baggrund er kun en mulig Antagelse, saa%
% --- p. 118 ---
\setcounter{footnote}{0}%
\opage{118}længe ukritisk Tryghed og dogmatisk Tillidsfuldhed raade. Erfaringsverdenen
er uafsluttelig i Tid og i Rum; Aarsagsrækken
kan ikke fuldstændigt gennemføres; der fremtræder ofte Begivenheder,
der komme saa pludseligt og overraskende og ere saa afvigende
fra tidligere Erfaringer, at der bliver en stor Opgave at
finde Sammenhængen paany. Men den store Humor forudsætter
jo ikke blot en Forstaaelse af enkelte Oplevelser hver for sig,
men ogsaa Muligheden af, at alle Oplevelser kunne indgaa som
Elementer i et Helhedsbillede. Og derved stilles en uendelig Opgave.
Og naar kan den blive fuldstændigt løst? Desuden er jo
det Enkelte og Individuelle selv en lille Verden, hvis indre Sammenhæng
det gælder at udrede. Og hvis det staar som en Torso
eller som et Misfoster, gælder det at finde, hvad der voldte Ødelæggelsen
eller hæmmede Væxten.

Der kan fremdeles, naar Værdibegrebet anvendes, fremtræde
en grel Modsætning mellem den almene Tingenes Orden, der roligt
fuldbyrder sig, og Enkeltvæsenernes Skæbne, en Modsætning, der
snart kan virke tragikomisk, snart komitragisk. Allerede før den
moderne Udviklingslære med dens Ide om Kampen for Tilværelsen
fremkom, skrev Tennyson sine berømte Linier om Naturen
(In Memoriam 56):

\begin{quote}
So careful of the type she seems,\\
So careless of single life.
\end{quote}

Heri er det træffende udtrykt, hvor vanskeligt det er at anvende
Værdibegreber paa et samlet Billede af Tilværelsen.

For det Første er Værdiernes Række uafsluttelig. Det følger
allerede af den stadige Mulighed for Forskydning: hvad der først
har Værdi som blot Middel, kan faa selvstændig Værdi som Formaal,
og hvad der først stod som Formaal, maaske som et sidste
Formaal, der begrænsede Ønskernes Horizont, kan faa Betydning
som Middel for endnu fjernere Formaal. Hvor kunne vi da finde
en absolut Grundværdi, paa hvilken der kunde bygges en afsluttende
Vurdering? Og er det givet, at selv om vi fandt en sidste
Maalestok, alle enkelte Værdier vilde være kommensurable med
den?

For det Andet dække jo Værdi og Virkelighed ikke hinanden.
Værdibegrebet udspringer indenfor den psykiske Verden, men den
% --- p. 119 ---
\setcounter{footnote}{0}%
\opage{119}synes kun at være en forsvindende lille Del af Tilværelsen i Forhold
til den materielle Verden med den uhyre Udfoldelse af Masser
og Kræfter. Værdiernes Verden staar da, synes det, i Tilværelsens
paa en ensom Plads, ligesom Gyvelbusken i Leopardi’s Digt
(La Ginestra), der voxer paa en enkelt Plet i Lavaørkenen. I sin
bitre Stemning holder Digteren dette Billede frem for alle dem,
„der have vænnet sig til at prise Livet“: „lad dem komme her
og se, om de levende Væseners Ve og Vel virkeligt ligger Naturen
paa Hjerte!“

Alle disse Vanskeligheder faa dog ikke nødvendigvis den intellektuelle
Baggrund for den store Humor til at forsvinde. Følgen
bliver kun, hvad der egentligt ligger i Sagens Natur, at ingen
intellektuel Baggrund for menneskeligt Følelsesliv kan opstilles
en Gang for alle. Et Drama behøver, naar man ikke hævder Fordringen
om „Stedets Enhed“, ikke at spilles for samme Bagtæppe
gennem alle Akter, men Bagtæppet kan skifte efter Handlingens
Gang. Naar der blot er en Sammenhæng mellem de skiftende
Scener, er Enheden hævdet. En Fastsættelse en Gang for alle af
de intellektuelle Forudsætninger vilde endog gøre den store Humor
umulig. Ti den store Humor vender sig ikke blot mod de enkelte
Oplevelser og lader dem fortone sig overfor Baggrunden. Den vender
sig ogsaa mod selve Baggrunden som den, der er underlagt Forandringens
Lov. Enhver Baggrund, der vilde gøre Fordring paa
absolut Gyldighed, vilde fremkalde Humoristens Smil. Den moderne
Erkendelsesteori antager heller ikke noget absolut Første
eller Sidste, hverken hvad Tankens Former eller dens Principer
eller dens Resultater angaar. Vi søge at bygge Tankerækker, som
vi kunne gennemføre overfor alle foreliggende og kommende
Iagttagelser. Hvor saadanne Tankerækkers Anvendelighed viser
sig at have en Grænse, undersøge vi, om Grænsen ikke skulde
være en Følge af, at den Tankerække, vi havde indledet, selv
kun var et Led i en mere omfattende Tankerække. Lyset forsvinder
om Aftenen og kommer igen om Morgenen; men denne
Forsvinden og Opstaaen betyder ikke noget Spring i Naturen; vi
maa blot gaa fra den Tankerække, der kun dannes af Iagttagelser
fra Jordens Stade, over til den Tankerække, der kan bygges
op, naar man gaar ud fra Jordens Omdrejning og dens Stilling til
% --- p. 120 ---
\setcounter{footnote}{0}%
\opage{120}Solen. Allerede Galilei sagde: „Tænk Jorden borte, da gives der
ingen Solopgang og ingen Solnedgang, ingen Horizont, ingen
Meridian, ingen Dag og ingen Nat!“ Hvad der synes at være en
Grænse for Tanken, er altsaa kun et Tegn paa, at Tanken for
tidligt har sluttet sit Arbejde. --- Dette gælder for Værdibegrebernes
Rækker, ligesaa vel som for Grundenes og Aarsagernes
Rækker.\footnote[1]{Smlg. \textit{Den nyere Filosofis Historie} I, p. 176. --- \textit{Den menneskelige Tanke.} p. 196. (Kritik af Renouvier’s og Dühring’s Opstilling af „det bestemte Antals Lov“).}

Forstaaelse er et aandeligt Arbejde, der stedse maa optages
paany. Den store Humor vilde bygge paa Sand, hvis den oversaa
dette. Forholdet mellem Forstaaelse og Humor falder paa en Maade
ind under Forholdet mellem Handling og Humor, som vi undersøge
i næste Kapitel. Humoristen maa opfatte Paastanden om at
være naaet til en absolut Baggrund som Udtryk for, at Tanken er
bleven træt og forvexler sit Hvilested med en absolut Afslutning.
Han ved, at med fremskridende Erfaring vil der kunne ske to
Ting: vore hidtil vundne Resultater kunne falde bort, --- eller
vore hidtil staaende Problemer kunne falde bort. To af vore ældre
danske Tænkere greb her hver sin af disse Muligheder. Holberg
mente, at naar vi ikke fandt flere Problemer i Verden, var det,
fordi vor Erfaring var begrænset; med udvidet Erfaring vilde man
se, hvor illusoriske mange af vore Løsninger vare. J. S. Sneedorff
mente derimod, at med udvidet Erfaring vilde mange af vore nuværende
Problemer falde bort, fordi de kun skyldes vor snevre
Horizont.\footnote[2]{\textit{Danske Filosofer}, p. 5 og 18.} De have begge Ret, og den store Humor kan bruge
begge Synsmaader. Den spøger baade med de absolute Løsninger
og med de absolute Gaader. Og denne Spøg grunder sig netop i
den store Alvor, med hvilken den store Humor omfatter Tankearbejdet.
Inderst inde hviler den paa en Tro paa Sandheden, som
den, der kræver stadigt nyt Arbejde paa at udvide Erfaringen og
gennemtænke den vundne Erfaring.

Ligesom Modsætningen mellem Alvor og Spøg, Sorg og Glæde,
Tragik og Komik, Storhed og Ringhed, Fornuft og Urimelighed
opløses i Humor som Totalfølelse, saaledes forholder det sig ogsaa
% --- p. 121 ---
\setcounter{footnote}{0}%
\opage{121}med Modsætningen mellem Forstaaelse og Uvidenhed. Humoristens
Standpunkt er træffende angivet i Udtrykket docta ignorantia,
saaledes som dette i den tidlige Renæssance opstilledes af
Nicolaus Cusanus.\footnote[1]{\textit{Den nyere Filosofis Historie} I, p. 81.} Humoristens Stade betinges intellektuelt ved
Modsætningsforholdet mellem Tilværelsens rige Fylde og den fulde
Forstaaelse, der til hver given Tid staar til Raadighed. Intet
af disse Modsætningsled kan falde bort, uden at det hele Stade
ændres. Kaster man Vrag paa den rationelle Udformning, fordi
den ikke kan rumme Alt, ender man i Mystik, der igen hurtigt
falder tilbage til Dogmatiken.\footnote[2]{Da Nicolaus Cusanus havde skrevet sin mærkelige Bog \textit{De docta ignorantia} [1440) rejstes der en interessant Diskussion om, hvorvidt man behøvede at vedligeholde Indsigten i Erkendelsens Betingelser, naar man engang havde indset, at den var begrænset. Nogle mente, at Konsekvensen maatte blive ren Mystik. \textit{(Autour de la docte ignorance.} Par \textit{Vansteenberghe.} Münster 1915.)} % PRINTED AS IS: „[1440)“

35. Paa lignende Maade staar Humoristen overfor den Kendsgerning,
at Mennesker, selv de skarpest tænkende Mennesker, afvige
i saa høj Grad fra hverandre i deres Verdensanskuelse. Historien
viser, at forskellige Videnskaber til forskellige Tider have
gjort Fordring paa at være den eneste eller den højeste Videnskab.
Saaledes var for Hobbes Fysiken Et og Alt, for Moleschott Kemien,
for Hegel Logiken, for Bergson den vitalistiske Biologi, medens
for Leibniz og Lotze en afsluttende Verdensanskuelse kun kunde
naas gennem Anvendelse af psykologiske Analogier paa Tilværelsen
idethele. Selve de Brydninger, der opstaa mellem de forskellige
Forsøg paa at tyde Tilværelsen ud fra en af dens enkelte
Sider, optager Humoristen i sin Verdensanskuelse som et Vidnesbyrd
om, hvorledes vi Mennesker nu engang ere stillede overfor
Tilværelsen. Han kan heri med Sibbern finde et Exempel paa, at
al Udvikling, ogsaa Tankens, foregaar sporadisk. Han sympatiserer
med enhver grundig og ærlig Stræben efter at gøre et alment
Synspunkt gældende, alt imedens han fastholder Opgavens Storhed
og Uendelighed. Denne Modsætning betinger, som allerede sagt,
en væsentlig Side af hans Totalfølelse.

Jeg har engang for mit eget Vedkommende forsøgt at beskrive
% --- p. 122 ---
\setcounter{footnote}{0}%
\opage{122}en Form for denne Totalfølelse. Da jeg ikke kan give nogen bedre
Beskrivelse, og da den, jeg i sin Tid gav, desuden er den Text,
jeg i nærværende Afhandling har søgt at udvikle nærmere, saa
hidsætter jeg den her (efter „Den menneskelige Tanke“ p. 377 f.)

„Der gives et Standpunkt, paa hvilket der gennem Livets Blanding
af Tragedie og Komedie, af Sejr og Nederlag, er vunden en
Grundstemning, som paa en Gang har Karakteren af Alvor ---
formedelst det Store, der er oplevet, og den Smerte, det ofte har
kostet, --- og af Ironi overfor ethvert Forsøg paa i Billeder at ville
udtrykke, hvad der i sin Fylde og ved sine Modsætninger sprænger
enhver Form, man forsøger at slaa fast. Den Livsanskuelse
eller Livsfølelse, der saaledes opstaar, fastholder Overbevisningen
om en stor Enhed i Tilværelsen, hvad Værdier saavel som alle
andre Emner og Elementer angaar, og kan derfor vende sig spøgende
mod Livets Tilskikkelser, fordi de ere magtesløse overfor
hin Forvisning, ligesaavel som mod Forsøgene paa at udsige det
Uudsigelige. Spøgen faar dog Vemodets Karakter, fordi der stadigt
føles en Modsigelse ved, at det Store saa ofte er knyttet til det
Ringe, og ved, at det ikke lykkes at faa sagt et afsluttende Ord
om Livets Kærne og Livets Kaar. Bag Ironien og Vemodet vil
der ligge en stor Resignation, snart mere praktisk, snart mere
kontemplativ, --- snart bitter, snart frejdig.“ ---

Jeg har for mit Vedkommende lært at betragte den store Humor
som en af de ideale menneskelige Livstyper, og det er den, jeg
mest sympatiserer med. Men der kan jo gives baade en ulykkelig
og en lykkelig Kærlighed. (Smlgn. § 26 og § 32). Og baade ulykkelig
og lykkelig Kærlighed til Noget kan føre til Forstaaelse.

At den store Humor, saaledes som den her er beskreven, i og
for sig skulde være det højeste Livsstade, kan, som tidligere udtalt,
ikke hævdes. Den store praktiske Stræben paa den ene Side,
den tragiske Lidelse paa den anden Side kunne begge være højere
Livsstader end Humor. Som vi senere (Kap. IX) ville se, kan ogsaa
Erkendelsesglæden udvikle sig til Totalfølelse og blive det Centrale
i Sjælen, og det vil ikke være de ringeste Aander, hos hvem
dette sker. Noget Tilsvarende vil ogsaa kunne ske med den kunstneriske
Trang. Det er da her, ligesom ved den praktiske Stræben,
en stor Opgave, der tager Sindet fangent og drager al Energi til
% --- p. 123 ---
\setcounter{footnote}{0}%
\opage{123}sig. Men de store Opgaver og de tragiske Skæbner falde ikke i
Alles Lod. Selv den, der har arbejdet med Alvor, og som har havt
sin Sorg, kan dog være klar over, at hans Stræben var som Fnug
overfor deres, der stedtes i den store Strid, og at hans Sorg var
som Leg overfor deres, i hvis Lod de dybeste Lidelser faldt.

% --- p. 124 ---
\setcounter{footnote}{0}%
\opage{124}\begin{center}VII. HANDLING OG HUMOR\end{center}


36. Kun rent teoretisk kan der skelnes mellem Følelse og Villie.
Enhver Lyst eller Ulystfølelse hænger sammen med en Trang, en
uvilkaarlig Bevægelse i en eller anden Retning. Naar Trangen er
forbunden med Evne til at naa et hensigtsmæssigt Resultat, kalde
vi denne Forbindelse af Trang og Evne Instinkt. Trangen bliver til
Drift, saasnart der med Følelsen af Lyst eller Ulyst forbinder sig
en Forestilling om, hvad der kan bevare og forøge Lystfølelsen
eller afværge og forminske Ulystfølelsen. Den aktive Side af en
Følelsestilstand er ikke Noget, der først kommer til, efter at Følelsen
er opstaaet, men er knyttet til den fra først af. Spørgsmaalet
er fra først af ikke, hvorledes Trang eller Drift kunne opstaa af
Følelsen, men hvorledes de kunne hæmmes saaledes, at Følelsen
kan miste sin aktive Karakter, idetmindste saaledes, at den ikke
fører til øjeblikkelig Handlen. Først under den senere Udvikling
af Følelseslivet opstaar der maaske en Modsætning mellem hvilende
og aktive Følelser.

Hvad der gælder al Følelse, gælder ogsaa en Følelse som Humor.
Den Trang og Drift, som kundgør sig i den, vil kunne gaa
i to Retninger. Der vil røre sig en Trang eller Drift til at fordybe
sig i det Store, hævde det, kæmpe for det; men der vil tillige
røre sig en Stræben efter at sikre selv det Ringeste dets Plads i
Verden, maaske endog holde sin Haand over Dumhed og Ondskab,
hvis de nu engang høre med og bør høre med --- i Kraft
af Principet felix culpa, eller fordi selve det Onde kun ophæves
grundigt ad den frie Udviklings Vej.

Der ligger, som vi tillige have set (§ 8 Slutning) en Etik i
Spiren i enhver Totalfølelse, og det kommer af, at der ligger en
% --- p. 125 ---
\setcounter{footnote}{0}%
\opage{125}Villen (Ordet taget i videste Betydning) bag enhver Totalfølelse,
en Villen, der fører til Vurdering og Handlen. Alvoren i Humoren
beror, som al Alvor, paa, at en Værdi, et Formaal staar som
ledende og bestemmende overfor de enkelte Stemninger og de
enkelte Handlinger.

Mange Psykologer stole i den Grad paa den skematiske Oversigt
over de sjælelige Tilstandes Egenskaber („Elementer“), som en foreløbig
Analyse kan tilvejebringe, at de betragte de konkrete, virkeligt
givne og meget komplicerede sjælelige Tilstande kun som mekaniske
Forbindelser af saadanne Elementer. Ligesom efter en hyppig
populær Opfattelse Grundstofferne i Naturen existere først,
maaske „i Begyndelsen“, og saa de forskellige sammensatte Stoffer
senere opstaa, saaledes vil man opbygge den sjælelige Verden
af skematisk opstillede sjælelige Elementer. Man forvexler virkelig
Tilblivelse med Fremstillingen i en Lærebog, hvor man begynder
med Elementerne. Det kan ikke nægtes, at paa det sjælelige Omraade
ere Fornemmelser, Forestillinger og Følelser, i deres abstrakte
og isolerede Form, det Simpleste, vi kunne tænke os.
Allerede den sædvanlige Sprogbrug har foretaget den Analyse,
der fører til at opstille dem. En barnlig Psykologi opererer nu
med disse Abstraktioner. Men de ere revne ud af den Sammenhæng
med Trang og Stræben, i hvilken de oprindeligt forekom.
Det er først en speciel Udvikling, en Differentiation, der skyldes
bestemte Betingelser, som fører til en vis Isolation af Sjælelivets
forskellige Sider eller Egenskaber. Ligesom Lyst eller Ulyst fra
først af ere knyttede til Trang og Drift, ere vore Fornemmelser
fra først af Betingelser for Udløsning af Reflex- og Instinkthandlinger,
og vore Forestillinger fra først af nøje knyttede til Forventning
og Drift. Selv videnskabelig Forsken begynder som
Stræben efter at finde Midler til bestemte Formaal; først senere
kunne de Tanker, som derved fremkaldes, faa umiddelbar Interesse.
Det er ikke Handling, der er et specielt Tilfælde af Fornemmelse,
Forestilling og Følelse. Det er omvendt disse, der i
den isolerede Karakter, den psykologiske Analyse giver dem, bero
paa en Hæmning af den uvilkaarlige Trang til Handlen. Selv
efterat Differentiationen er indtraadt, bevare alle Erkendelses- og
Følelseselementer et Forhold til Handling. De beholde Præg af
% --- p. 126 ---
\setcounter{footnote}{0}%
\opage{126}at være Overgangsled indenfor en Stræben i en vis Retning, eller
af at være fremkaldte af en saadan Stræben.

Den genetiske Psykologi, der søger at følge Sjælelivet i dets
faktiske Udfoldelse, berigtiger paa en gavnlig Maade de Misforstaaelser,
der kan opstaa ved Forvexling af de ved Analyse af det
differentierede Sjæleliv fundne Elementer med det begyndende
Sjælelivs Former. Hvad vi paa et senere Trin kunne opfatte som
sondret og som til en vis Grad uafhængigt af Trang og Stræben,
have vi ikke Ret til uden videre at lægge ind paa alle Sjælelivets
Trin. Der gør sig her en psykologisk Antinomi gældende, som
skyldes den Kendsgerning, at Eftertanken først kommer til at
virke, efterat Sjælelivet længe har været i frodig Udfoldelse; derfor
har den ondt ved at forstaa sin egen Fortid og tror, at de
Forskelligheder, den nu kan finde, have været tilstede paa alle Stadier
af det Sjæleliv, hvis sene Frugt den selv er.\footnote[1]{Se min \textit{Psykologi} Kap. IV og \textit{Den menneskelige Tanke,} p. 10—12. --- Det genetiske Synspunkts Betydning overfor det rent analytiske er i den nyeste Tid med Eftertryk fremholdt af \textit{Edward Thorndike (Animal Intelligence.} 1911. --- \textit{Elements of Psychology} 1905). Se ovenfor § 3.} --- Man vilde meget
vel kunne fremstille hele Fornemmelsens, Forestillingens og Følelsens
Psykologi som Kapitler indenfor Villiens Psykologi. Endog
Logiken kunde man finde Plads til der, nemlig i Kapitlet om
Overvejelse.

37. Naar vi nu vende tilbage til Humoren, maa det indrømmes, at
den kan gaa over til en ren kontemplativ Tilstand, en Æstetiseren,
i hvilken Sindet svælger i Tanker og Billeder, medens det personlige
Grundlag (se Kap. IV), særligt Sympatielementet, udviskes.
Det kan da ske, at man vilkaarligt tumler med Objekterne og
forsøgsvis anbringer dem i forskellige Kombinationer for at nyde
sin egen Magt og sit egen Lune. Eller man ser paa Virkeligheden % PRINTED AS IS: „sit egen“
som noget Fjernt, ligesom gennem en omvendt Teaterkikkert.
Eller man glæder sig blot over Formernes og Farvernes
Spil, ligegyldig overfor de virkelige Existenser, der ligge bagved.
Saaledes ser den unge Mand i Ibsen’s „Paa Vidderne oppe fra
Fjeldørkenen paa, hvad der foregaar nede i Dalen; Billedets virkelige
Emner vedkommer ham ikke, skønt det er hans egen
% --- p. 127 ---
\setcounter{footnote}{0}%
\opage{127}Moders Hus, der brænder, og hans egen Kæreste, der rider til
Kirke som en Andens Brud:

\begin{quote}
Det Alt tilhobe jeg stod og saa\\
oppe fra Livets Vidder;\\
et højere Lys over Synet laa ---\\
men se, det kan nu Ingen forstaa,\\
som nede i Sværmen sidder.
\end{quote}

Noget Saadant har Goethe sandsynligvis tænkt paa, naar han
har kaldt Humor for Hjerteløshed, og naar han tillige mente, at
Humor ikke har nogen Lov eller noget Hold i sig selv og derfor
let udarter til Vilkaarlighed. (Brev til Zelter. 30. Oktbr. 1808.)\footnote[1]{I \textit{Kierkegaard’s Papirer VI} p. 38 findes et Stykke med Overskrift: „Replik af en humoristisk Individualitet“. Hvad han skildrer her, er en Blaserthed og Indifferentisme, der ikke stemmer med den Maade, hvorpaa han, omtrent samtidigt, i „Uvidenskabelig Efterskrift“ har karakteriseret Humor.}
Han tænker paa mange Romantikeres lunefulde Leg med Alt
mellem Himmel og Jord, som ikke levnede Tid til alvorlig kunstnerisk
Gennemarbejdelse. Og dog er, siger han, „Naturens og
Kunstens højeste, ja, eneste Operation Udformningen (die Gestaltung),
og indenfor Formen igen Individualiseringen (die Spezifikation),
der fører til, at enhver Ting bliver og vedbliver at være
et Ejendommeligt og Betydningsfuldt.“ --- Snarere maatte man
vel sige, at naar Humoren kunstnerisk ikke fører til Udformning
af individuelle Skikkelser, er det enten, fordi Humoristen ikke er
Kunstner, eller fordi han ikke i Virkeligheden er Humorist (i den
Betydning, Goethe selv tager Ordet, naar han taler om Sterne).
Den store Humor, saaledes som den i det Foregaaende er beskreven,
vil, hvor der er kunstnerisk Evne, netop føre til med
Interesse at sysle med det Enkelte, indtil det staar fuldt udpræget
i sin Ejendommelighed, den Ejendommelighed, der netop giver
det Værdi og bestemmer dets Plads indenfor det store Hele. Den
store Humor kan ikke tænkes uden en dobbelt Interesse, Interesse
baade for det Hele og det Enkelte.

38. Men ligger der overhovedet ikke en Fare i en Totalfølelse
som Humor, i hvilken saa mange forskellige, tildels indbyrdes
% --- p. 128 ---
\setcounter{footnote}{0}%
\opage{128}modsatte Elementer ere forenede? Vil en saadan Forening af
Modsætninger ikke let føre til, at der sker en Udjevning eller
Udligning, som svækker, --- at der slaas ind paa en Middelvej,
saa at den frugtbare Energi, Modsætningerne netop i deres extreme
Form kunne udløse, falder bort? --- Denne Bemærkning skyldes
en ung Ven, med hvem jeg ofte har drøftet Emnet for denne Afhandling.
Jeg tror, at enhver Totalfølelses Tilbliven kan medføre
en Fare, idet dens Elementer forbindes saaledes, at de ikke længere
udløses hvert for sig. Det er Sammensmeltningens Klangfarve
eller Organisationens Lov, der nu bestemmer, hvorledes der
reageres overfor en ny Oplevelse. Der er jo idethele en Fare i
ethvert opnaaet Resultat, enhver afsluttet Tilstand, hvad enten den
nærmest har Følelsens eller Erkendelsens Præg. Istedetfor at betragtes
som Stade for ny Orientering og nyt Arbejde bliver en
saadan Tilstand let en Afslutning af Liv og Virken. Men i og for
sig ligger det ikke i Begrebet Totalfølelse, at denne skal være
en definitiv Afslutning. Den kan meget vel kun betyde en Station
paa Vejen, et praktisk Opgør, der er nødvendigt for at kunne fortsætte
Livet saaledes, at Fortidens Oplevelser kunne komme Fremtiden
til Gode. Men der er desuden en Omstændighed af almindelig
psykologisk Interesse, som i denne Sammenhæng er af
Betydning.

En Totalfølelse har væsentligt Karakteren af Sindelag, ikke af
Sindsbevægelse. Den skyldes jo en Sammenfatten af en Række af
Oplevelser, som hver for sig have fremkaldt mere eller mindre
dybtgaaende Sindsbevægelser og derigennem have afsat blivende
Spor i Sindet. Og den betegner en fast Retning, i hvilken Følelseslivet
nu er slaaet ind, selv om der kan opstaa Svingninger
eller Krusninger. Og naar Bevægelsen i denne Retning begunstiges
eller hæmmes, vil der stadigt kunne udløses Sindsbevægelser,
som saa ofte ere af helt anden Natur end Sindelaget (Smlgn.
§ 2 og § 7). Kærlighed til en Sag, man har gjort til sin, kan afføde
Kamplyst, hvor det gælder om at forsvare den. Og det, at denne
specielle Sindsbevægelse er bleven udløst i Sagens Tjeneste, kan
saa igen virke tilbage paa selve Totalfølelsen overfor Sagen, ændre
dens Klangfarve eller bringe et nyt Element ind i den hele Organisation.
Istedetfor den stille Syslen med Følelsens Genstand
% --- p. 129 ---
\setcounter{footnote}{0}%
\opage{129}kan nu træde en Stemning af væbnet Agtpaagivenhed overfor
hvad der kan true. Men det kan ogsaa være, at naar den Situation,
der fremkaldte Sindsbevægelsen, er ophørt, falder Sindet til
Ro i sit gamle Leje.

Der er ved Opfattelsen af Mennesker en Forskel at gøre mellem
Karakteren som Helhed og den særlige Maade, paa hvilken
den i en bestemt Situation lægger sig for Dagen. Allerede en
ældre Æstetiker (Hurd, citeret i Lessing’s „Hamburgische Dramaturgie“
Nr. 73) har bemærket, at selv Shakespeare’s kraftigst
tegnede Karakterer ikke i hver Ytring eller Handling udtrykke
deres væsentlige eller herskende Egenskaber, men kun, naar Situationen
giver en naturlig Anledning dertil. Dog vilde jeg hellere
vende Sætningen om og sige, at naar Situationen giver Anledning
dertil, kan der fremkomme Handlinger eller Ytringer, som ikke
direkte kunne siges at ligge i de herskende Egenskaber, skønt de
ikke ere uforenelige med dem. Det vidner om Shakespeare’s dybe
Forstaaelse af Menneskenaturen. Enhver Karakter er som et Orkester,
der kun i enkelte Episoder afbrydes af Solo’erne. Sindelaget
er et Orkester, Sindsbevægelsen en Solo. Man gør de dramatiske
Digtere Uret, naar man overser denne Forskel. Allerede
Aristoteles (Poetik. Kap. 15) begaar en saadan Uret, naar han
finder en Inkonsekvens i Euripides’ Ifigeneia deri, at hun fremtræder
som en Anden, naar hun trygler sin Fader om Livet, end
da hun senere villig gaar i Døden. Forklaringen er dog ganske
simpel. Det er jo i Mellemtiden gaaet op for hende, at hendes
Offerdød er nødvendig for Grækerhærens Skyld, og at hendes
Vægring vilde udsætte den ridderlige Achilleus for Undergang.\footnote[1]{Se \textit{Gertz’} Oversættelse af \textit{Ifigeneia i Avlis}. Noten til v. 1368.}
En lignende misforstaaende Kritik har jeg ogsaa hørt overfor
Euripides Medeia, fordi hun jamrede sig i Enrum, men optræder
med kongelig Værdighed overfor Stadens Kvinder. De fleste
Mennesker give jo dog mere efter for deres Sindsbevægelser i
Enrum, end naar de optræde overfor andre Mennesker. Denne
Forskel behøver aldeles ikke at have Noget med Karakteren at
gøre, og den indeholder ingen Modsigelse. En lignende Kritik er
rettet mod Sofokles’ Antigone, der trodser Kongens Forbud mod
% --- p. 130 ---
\setcounter{footnote}{0}%
\opage{130}at begrave sin Broder, og herved viser Kraft og Mod, men senere,
da hun skal lide Døden, jamrer over sin tidlige Død. Der er dog
ingen Modsigelse i, at hun nu, efter at have gjort, hvad hun ansaa
for sin Pligt, og derved vist en kæk Trods, nu optræder som den
græske unge Pige, hun i Virkeligheden er. Der er ingen Modsigelse
i, at Orkestret lyder igen, efterat den ved en bestemt Situation
motiverede Solo er endt.

Kun viser Erfaring, at det i Sindelagets mere harmoniske Tider
ikke altid er let at forstaa, hvad man under en pludselig Sindsbevægelse
uvilkaarligt udførte, idet der paa en Gang vaktes Trang
og Evne. Handlingsøjeblikkene frembyde den største Koncentration,
Sjælelivet raader over; her er den inderligste Samling af alle
Sindets Kræfter i én eneste Retning. I Beslutningsøjeblikket kan
Koncentrationen stige til Ekstase. Intet Under da, at man senere
hen ikke forstaar den Iver, den Begejstring, eller den Hidsighed,
eller den Kraft, man i saadanne Øjeblikke har udvist. Man kan i
Erindringen svimle ved det, der under Oplevelsen eller Udførelsen
ikke vakte Svimmelhed. I Handlingsøjeblikket var der
Sindsbevægelser fremme, der traadte i Modsætning til det
ellers raadende Sindelag og for en Stund lagde Beslag paa al
Energi.

En ivrig Fjeldbestiger fortalte Ludvig Feilberg (Samlede
Skrifter p. 736) om en farlig Nedstigen fra en høj, iset Skrænt.
„Naar jeg siden,“ tilføjede han, „tænkte tilbage paa den Gang, da
jeg, skubbende mig forsigtigt nedad paa Bagen, søgte at nærme
mig Kanten af Afgrunden saa meget, at jeg kunde se ned, paakom
der mig en Følelse af Rædsel. Det var egentligt mærkeligt,
ti i selve Øjeblikket var der ikke Tale om nogen Rædsel . . . .
Man havde andet at bestille.“ --- Antropologen G. H. Schneider
beretter et lignende Træk fra egen Erfaring. Under en botanisk
Udflugt paa Kreta blev han angreben af nogle bevæbnede Hyrder
og kunde kun redde sig ved en halsbrækkende Flugt ned over
en stejl Skraaning. „Fra Blok til Blok sprang eller snarere styrtede
jeg, ofte i farlige Hop, i Dødsangst ned over Bjerget og naaede
saa den nærmeste Landsby\dots{}. Senere kunde jeg kun forbavses
over min ufejlbare Sikkerhed ved de farligste, rask paa hinanden
følgende forskellige Spring fra Sten til Sten, og dog kunde der
% --- p. 131 ---
\setcounter{footnote}{0}%
\opage{131}ikke være Tale om Overvejelse.“ (G. H. Schneider: Der thierische
Wille, p. 203.) --- Et kosteligt og tillige rørende Exempel forekommer
i Goldschmidt’s „Min Onkels Tømmerplads“. Den gamle
Købmand var i Spænding og Angst for et af sine Skibe, hvis
brave Kapitain han særligt holdt af. Pludseligt gaar Butiksdøren
op, og Kapitainen træder ind. „Onkel saa op --- næste Øjeblik
satte han over Disken og om Halsen paa den Indtrædende.“ Da
Bevægelsen havde lagt sig, saa man den gamle Mand med betænksomt
Blik maale Højden og Bredden af Disken, han saa uforvarende
havde sat over.

Man kan altsaa ikke fra Totalfølelsens tilsyneladende Ro slutte
til Fraværelsen af Energi i øjeblikkelige Situationer. Naar der
tales om en Følelses Styrke, maa man erindre, at den kan ytre
sig paa to Maader, i Inderlighedens Form i Sindelaget, og i Heftighedens
Form i Sindsbevægelsen, og at disse to Maader kunne
staa i Vexelvirkning med hinanden.

Naar det fra Erindringens Verden kan være vanskeligt at forstaa
den Energi, hvormed der handledes i Fortiden, saa er det
ofte endnu vanskeligere ud fra Nutidens Tilstand at tænke sig
ind i den Tilstand, en fremtidig Handlen kræver. Her kan Afstanden
mellem Opgavens Storhed og det, man selv tror at kunne
udrette, staa som saa stor, at den ikke kan udlignes; enhver Mulighed
for tilnærmelsesvis at løse den blegner. Amiel’s Dagbøger
(Journal intime) have her deres stadige. atter og atter varierede % PRINTED AS IS: „stadige.“
Tema. „Idealet,“ siger han (I, p. 22; 102), „forgifter al ufuldkommen
Besiddelse for mig\dots{}. Det gælder for mig Alt eller Intet.
Ironien tager, paa Grund af en Sammenligning med det skimtede
eller drømte Ideal, hverken sig selv eller Virkeligheden alvorligt.“
Det var Amiel’s Ulykke, at han ikke kunde blive Humorist, ikke
kunde komme til at se det Smaa, --- ogsaa det Smaa, han selv
kunde udrette, --- ikke blot i dets Modsætning til, men ogsaa i
dets Sammenhæng med det Store. Deraf hans Villiesløshed, hans
Mistillid til sig selv; han kunde ikke tage sig selv alvorligt, i det
højeste kunde han føle Medlidenhed med sig selv\footnote[1]{Je n’aime que le sérieux, et je ne puis prendre au sérieux mes circonstances ni moi-rmême\dots{} je me prends perpétuellement en pitié au nom de tout ce qui est beau et admirable. \textit{Journal intime.} Genève 1908. I. p. 239.}. Humoren % PRINTED AS IS: „moi-rmême“
% --- p. 132 ---
\setcounter{footnote}{0}%
\opage{132}forbinder Idealets Storhed og Handlingens Ringhed, Ironien over
egen Begrænsning og Beundringen for det Fuldkomne. Den hævder
den positive Betydning af Enkens Skærv. Den muliggør
anstrengt Arbejde, alt imedens den giver Rum for den Tanke, at
vi dog, naar den højeste Maalestok anlægges, ere unyttige Tjenere.
Ironien kan her træde i Humorens Tjeneste, idet man i Handlingens
Øjeblik tager sig selv alvorligt, skønt man meget godt véd,
at selv det største Arbejde kun er Spøg, naar Idealet stilles højt
nok; --- bag denne Spøg ligger der en Alvor, som netop er den,
der hævder den ideale Maalestok. (Smlgn. § 18).

Det Sindelag, i hvilket Humoren bestaar, er spredt og indirekte
tilstede i Arbejdet, naar dette gøres uden Vigtighed, uden Sentimentalitet,
uden trykkende Tungsind, uden snakkesalig Dvælen.
Forudsætningen er, at selv om Idealet, der stod for os, ikke fuldt
ud realiseres, og selv om under Livets Gang meget Værdifuldt
forsvinder, saa rinder Værdikilden dog stadigt i det Skjulte, og det
gælder blot om at grave dybt nok for at faa den til at sprudle
frem. I selve det Sind, der trøstigt sætter en ufuldkommen Handling
ind i Forandringernes og Illusionernes Verden, er der større
Værdi end i den vege Drømmers aandfulde Reflexioner og Fantasier.
Ligesom den tragiske Helt netop viser sin Højhed i Undergangen,
saaledes viser den, der ikke af Ufuldkommenheden i alt
Endeligt afskrækkes fra at yde, hvad han formaar, at der er
Højhed i hans Sind.

39. I sit dybsindige, om indtrængende pykologisk Erfaring
vidnende Skrift „Sygdommen til Døden“, (Skrifter XI), det Ypperste
af Alt, hvad han har ydet\footnote[1]{Se om dette Skrift \textit{Danske Filosofer}. p. 163.}, har Kierkegaard skildret de Modsætningsforhold,
vi i dette Kapitel have beskæftiget os med. Jeg
nævner hans Udtryk for de vigtigste Modsætninger og Humorens
Forhold til dem.

„Uendelighedens Fortvivlelse ved at mangle Endelighed“ opstaar,
naar man ikke vil finde sig i den Begrænsning, al Virkeliggørelse
% --- p. 133 ---
\setcounter{footnote}{0}%
\opage{133}af noget Stort er underlagt. „Mulighedens Fortvivlelse ved at
mangle Nødvendighed“ opstaar, naar man hildes i tomme Ønsker
og Forhaabninger eller ogsaa i Angst og Tungsind. Fra disse
Former for Fortvivlelse befrier den store Humor. Og det Samme
gælder overfor de to modsatte Former— „Endelighedens Mangel
paa Uendelighed“ og „Nødvendighedens Mangel paa Mulighed“. Ti
Humoren peger ud mod større Horizonter og derved mod flere
Muligheder, end det snevre Blik, for hvilket alle Sunde saa let
synes at være lukkede, kan opdage.

Saa dybsindigt Kierkegaard end har skildret hine Modsætninger
i deres extreme Former, saa vilde han dog bestemt have
nægtet, at Humor nogensinde skulde kunne raade Bod paa dem.
I sit Ungdomsskrift om Begrebet Ironi (XIII p. 379) kalder han
Solgers Standpunkt, som væsentlig falder sammen med, hvad jeg
kalder den store Humor, en Slags kontemplativ Andagt og mener,
at det kun fører til at leve i en uvirkelig Virkelighed, og at det kun
betoner den negative Side ved det Store og ikke naar op til en
højere Virkelighed, som efter hans Opfattelse kun naas ved at
anlægge teologiske Forudsætninger. Og i sit filosofiske Hovedskrift
(„Videnskabelig Efterskrift“) (VI. p. 525) betragter han Humor, % PRINTED AS IS: „Videnskabelig“ (sense: Uvidenskabelig)
som han dog i „Stadier paa Livets Vej“ (VI. p. 393) havde betegnet
som Udtryk for Aandsstyrke, som en Tilbagekaldelse, en
Dragen sig ud af Existens, et baglængs Perspektiv, en Stillestaaen,
et Liv i Erindringens Verden.

Hermed stemmer det dog ikke godt, at det „Stadium“, der af
ham betegnes som Humor, ligger senere i Rækken end det etiske
„Stadium“. Ifølge Kierkegaard skal et følgende Stadium ikke fornegte
det foregaaende, men vise sin Overlegenhed ved i sig at
have optaget alt det Sande og Berettigede i dette. Humor vilde
ikke være et højere Stadium end det etiske, naar den fornegtede
dette. Den vilde da selv falde for en etisk Dom, istedetfor netop at
udvise rigere Betingelser for det Etiske, som den i det foregaaende
beskrevne store Humor gør det. Men Hemmeligheden er, at
alle „Stadier“, der laa mellem Nydelsesstadiet (det „æstetiske“)
og det strengt religiøse mere og mere tabte deres Betydning for
Kierkegaard. Ogsaa „Humor“ blev klemt inde og trykket ihjel
imellem Nydelsens Evangelium og Lidelsens Evangelium, og
% --- p. 134 ---
\setcounter{footnote}{0}%
\opage{134}Kierkegaard tog egentligt tilbage, hvad han havde beskrevet som
en høj, rent human Maade at tage Livet paa. Han har ydet fine
Bidrag til den store Humors Psykologi, men til Slut kunde han
dog ikke lade den komme til sin Ret som en paa Livserfaringens
Grund fremvoxet Livstype.

% --- p. 135 ---
\setcounter{footnote}{0}%
\opage{135}\begin{center}VIII. HISTORISKE BETINGELSER\\ FOR HUMOR\end{center}


40. Den psykiske Proces, hvorved en Totalfølelse bliver til,
kan ikke gaa for sig under alle Forhold. Enhver Totalfølelse betinges
ved sociale og historiske Forhold, faar en historisk Karakter,
og det bliver Genstand for særlig Undersøgelse, hvorvidt
Betingelserne for denne Opstaaen ere tilstede til en vis given
Tid. Det er et af de Spørgsmaal i Psykologien, hvor denne maa
benytte den historiske (sociologiske, sammenlignende) Metode.
Hverken Selviagttagelse, Experiment eller Analyse kan her gøre
Udslaget. Det gaar her med Humor, som det gaar med andre
Totalfølelser, f. Ex. Nationalfølelse, Pligtfølelse, religiøs Følelse,
Følelse for Naturskønhed o. s. v. Man hører undertiden den Mening,
at det er skadeligt for Videnskaben at indrømme, at der
faktisk existerer saadanne Følelser, fordi der derved indrømmes
en vis Gyldighed for det, der er Genstand for dem, eller for de
Forestillinger, hvortil de ere knyttede. Men i Anerkendelsen af,
at en Følelse existerer, ligger der aldeles Intet med Hensyn til
Gyldigheden eller Værdien af dens Genstand. Og hvis den viser
sig at være knyttet til bestemte historiske Betingelser, er det klart,
at hvis disse Betingelser falde bort, vil Følelsens Tid ogsaa være
forbi. Efter den Opfattelse af det religiøse Problem, jeg har udviklet
i min Religionsfilosofi (og senere i Slutningsafsnittet af
„Den menneskelige Tanke“), er det afgørende Spørgsmaal i religionsfilosofisk
Henseende netop, om der stedse vil være Mulighed
for religiøs Trang og Følelse, efterat de religiøse Forestillingers
objektive, videnskabelige Gyldighed er opgiven. Den psykologiskhistoriske
Metode afløser her den spekulative eller metafysiske.
Og denne Metode maa ogsaa anvendes, hvad Humor angaar.

% --- p. 136 ---
\setcounter{footnote}{0}%
\opage{136}Enkeltfølelser af Glæde eller Sorg, af Ophøjet eller Latterligt,
af Tragisk eller Komisk kunne være mulige, uden at dog den
ejendommelige Komplikation bliver mulig, hvorved der af disse
Følelsers Sammenspil udvikler sig en Totalfølelse som Humor.
Denne er i hvert Tilfælde, som vi have set (§ 10) først iagttaget
og beskrevet i nyere Tid, især i de sidste Aarhundreder. Vel
har man deraf ikke uden videre Lov at slutte, at den ikke har
existeret forud. Der er Forskel mellem Betingelsen for en sjælelig
Tilstands Existens og Betingelsen for, at den iagttages og beskrives.
Men vi have derfor dog heller ikke Lov til at paastaa,
at hvad vi nu kunne iagttage eller se beskrevet, ogsaa har været
muligt paa tidligere Udviklingstrin. Der er f. Ex. Forskelle og
Distinktioner, som vi med Lethed gøre, efterat Udviklingen har
naaet et vist Trin, men som ikke uden videre kunne anvendes
paa tidligere Trin. Og noget Tilsvarende gælder for Totaltilstande,
der bero paa et Sammenspil af mangfoldige Elementer og Betingelser,
som ikke altid ere tilstede.

Psykologer og Æstetikere ere ret enige om, at Humor er et
moderne Fænomen. Cervantes og Shakespeare nævnes ofte som
de første Forfattere, i hvis Værker denne Totalfølelse spores som
Grundstemmning. Oldtidens Opfattelse af Karakterer og Skæbner % PRINTED AS IS: „Grundstemmning“
var mere enkel; man holdt sig til store og rene Linier, og meget
sammensatte Personligheder og Begivenheder fremstilles ikke.
De Gamle havde ganske vist en levende Følelse af den Begrænsningens
Lov, alt Menneskeligt, ja, Alt i Verden er underlagt.
Ifølge en „guddommelig Misundelse“, der svarer til, hvad moderne
Romantik har kaldt „Verdensironien“, ere stor Lykke og
brat Undergang nøje sammenknyttede. Her var der forsaavidt
Mulighed for en Totalfølelse, i hvilken modsatte Stemninger kunde
indgaa som Elementer. Men Grækerne synes at være stansede
ved det successive Forhold mellem disse Modsætninger
af Lys og Mørke, Held og Vanheld, uden at Modsætningsleddene
kom til at gøre sig gældende paa én Gang og derved kunde indgaa
som Elementer i en Totalfølelse (Smlg. § 6).

Platon omtaler i Dialogen „Filebos“ flere blandede Følelser,
uden at den Maade, paa hvilken de ere opstaaede, nærmere angives.
Ikke blot som Følelser ved et Drama, men ogsaa i „Li%
% --- p. 137 ---
\setcounter{footnote}{0}%
\opage{137}vets hele Tragedie og Komedie“ opleve vi saadanne blandede
Følelser (p. 50 B). Tragedie og Komedie stode for Platon (i en
Periode af hans Tankes Historie) i nøje Forbindelse. Men desværre
siges der i den berømte Slutning af „Symposion“ kun
det, at det er den samme Mands Sag at digte Komedie og Tragedie.
Hos Platon kunde man maaske forudsætte den store Humor,
naar man ser, hvorledes han, netop i „Symposion“, stiller
Livets Modsætninger sammen og gennem Spøg og Ironi lader de
største Tanker komme til Udfoldelse. Men til fuld Udvikling som
Totalfølelse er den dog ikke kommen, dels vel paa Grund af den
Pessimisme, der prægede hans sidste Aar, dels, og maaske især,
fordi han var mest optagen af at udvikle sine Tanker i deres
systematiske Form. Til dette sidste Punkt ville vi i næste Kapitel
vende tilbage.

At Platon ikke manglede Blik for, hvorledes der gennem forskellige
Følelsers Skiften kan udvikles en Totaltilstand, ses, som
ovenfor (§ 2) omtalt, af hans Beskrivelse af, hvorledes Karakterfasthed
(„Tapperhed“) bliver til. Men for den Totalfølelse, vi
kalde Humor, har han ikke havt Blik. Han har ikke staaet i det
levende Forhold til Livet som hans store Mester. Hans Psykologi
har ikke evnet at fremstille den Totaltilstand, der laa bag ved
dennes Livsvirksomhed, og som netop er det største Exempel
paa den store Humor. Dette strider ikke imod den Sætning, at
medens Ironi er en antik Følelse, er Humor en moderne Følelse.
Ti Sokrates staar netop ved et Gennembrud, hvor det Antike
viser ud over sig selv, og det er netop derved, at han kan staa
som Repræsentant, ikke blot for en enkelt Tid, men for al menneskelig
Stræben idethele. Jeg kommer tilbage ogsaa til dette
Punkt i næste Kapitel, men bemærker foreløbigt Følgende.

Hegel og Kierkegaard have paavist, at Ironi var den naturlige
Form for det Gennembrud i den antike Oldtid, da Personligheden
for første Gang stiller sig i Modsætning til den givne, overleverede
Tingenes Orden og stræber at blive en lille Verden for sig selv.
Foreløbigt gælder det blot om at hævde Retten til at være sig
selv, tænke sine egne Tanker, tage Livet paa sin egen Maade.
Men saalænge der endnu ikke kan udvikles en ny omfattende
Tankeverden og en ny social Orden, saalænge er det naturligt, at
% --- p. 138 ---
\setcounter{footnote}{0}%
\opage{138}den Enkelte tilsyneladende gaar ind paa det Overleverede og forsøger
at stille sig paa de Andres Stade, skønt det i de foreliggende
Tanker og Ordninger Givne ikke for Alvor kan udtrykke
hans virkelige Stade. Han kæmper for en formel Ret; en saadan
udtrykker nemlig Personlighedsprincipet (se § 16); --- men der
mangler foreløbigt et reelt Indhold. Ofte vil da Spøgen stikke
Hovedet frem bag Alvoren paa drillende Maade. Hvad man overser,
naar man betragter Sokrates blot som Ironiker, er, at den
Spøg, der ligger bag den tilsyneladende Alvor, selv igen kan være
Udtryk for en dybere liggende Alvor (se § 18). De nye Tanker
og Ordninger, der kunne afløse den givne Kultur, kunne foreløbigt
kun ytre sig paa rent subjektiv Maade, ved Personlighedens
Tilbagegang til sig selv og ved dens Frigørelse fra det hidtil Udformede.
Alt Objektivt begynder som subjektivt, i Subjektivitetens
Form. Derfor er en Mand som Sokrates ikke blot den store Ironiker,
men (uden at vide det) den store Bebuder. Den Spøg, der
ligger bag den udadvendte Alvor, er forsaavidt ikke blot Befrielsens
Spøg (§ 13), men staar i Tjeneste hos en Alvor, der endnu
ikke kan udformes paa positiv Maade. Ironien bliver et Udtryk
for Humor. Jo mere det subjektive Grundlag for Ironien uddybes,
des mere bliver den til Humor. I en vis Forstand kan Humoren
siges at være mere subjektiv end Ironien, fordi den forudsætter
et dybere Lag i Personligheden, end Ironien i og for sig behøver
at gøre. Saaledes har Schopenhauer udtalt, at Ironien er objektiv,
Humoren subjektiv, og han finder heri Forklaringen af, at Ironiens
Mesterstykke fremkom i Oldtiden, medens Humor er
en moderne Følelse. („Welt als Wille und Vorstellung“. II.
Kap. 8). Men dersom det er Personlighedsprincipets Opgave
at gennemtrænge baade Tankelivet --- ved at lægge Vægt paa,
at al Erkendelse og al Overbevisning skal vindes ved Selvarbejde
--- og det sociale Liv --- ved at hævde den Enkeltes
Frihed overalt, hvor den ikke krænker Andres Frihed, --- saa
bebudes der i den store Humor en Objektivitet af højere Rang
end den, Ironien søger at frigøre sig fra. Den nye Verden,
der laa som Spire i Sokrates’ Personlighedsprincip, behøvede
Aartusinder for at udfolde sig, og den kæmper endnu for
Livet.

% --- p. 139 ---
\setcounter{footnote}{0}%
\opage{139}41. Totalfølelser ere ikke blot mere sammensatte, men ogsaa
mere individuelt nuancerede end Enkeltfølelser. Ti de Enkeltfølelser,
der indgaa som Momenter i en Totalfølelse, kunne indenfor
denne nye Helhed staa i højst forskelligt Forhold til hverandre,
og derved betinges betydningsfulde Forskelligheder i Totalfølelsens
Klangfarve hos forskellige Individer. Forskelle mellem Individer,
hvad Følelseslivet angaar, bero snarere paa saadanne
forskellige Klangfarver, end paa at nogen Enkeltfølelse helt skulde
mangle hos noget Menneske. Omvendt vil naturligvis selve Muligheden
af en Totalfølelses Opstaaen i høj Grad være betinget
ved de oprindelige individuelle Ejendommeligheder.

De forskellige Temperamenter ville, hvad Humor angaar, ikke
frembyde lige god Jordbund for deres Opstaaen. Der er, som jeg
tror,\footnote[1]{Se min \textit{Psykologi} VII C, 1. --- Smlgn. ogsaa \textit{Alfr. Lehmann: Hauptgesetze des menschl. Gefühlslebens.} p. 390.} Grund til at skelne mellem otte Hovedtemperamenter, der
bestemmes dels ved den Hurtighed og Styrke, hvormed man paavirkes
og reagerer, dels ved Lyst- eller Ulystfølelsens Fremhersken.
To af disse Temperamenter ville afgive god Jordbund for
Humor, nemlig det mørke, stærke og langsomme (det „melankolske“)
og det mørke, svage og langsomme (som ikke har noget
gængse Navn), idet de fra først af, hvert paa sin Vis lade Livets
mørke Sider komme til at spille en fremtrædende Rolle. De lyse
Temperamenter, baade i deres stærke og svage, hurtige og langsomme
Former, ville ikke saa let lade de store Skygger fremtræde
i det Billede af Livets Landskab, som uvilkaarligt danner
sig i Bevidstheden. Og de andre mørke Temperamenter, baade
det stærke og hurtige (det „koleriske“) og det svage og hurtige
(som er anonymt), ville ikke let gøre en saadan dvælende
Tilstand mulig, som Humor forudsætter; „Hurtighed“ betyder jo
Tilbøjelighed til at gaa over til Handling eller til andre Tilstande.
--- Sammen med Temperamentet virke de Erfaringer, som gøres,
og den Eftertanke, som øves.

Naturligvis kunne de nævnte mørke Temperamenter gøre Humor
umulig, naar de ikke lade Erfaringerne komme til deres
fulde Ret. Alvoren skal jo i Humoren kunne have Spøgen til Udtryksform
og maa kunne forenes med Mildhed og Forstaaelse af,
% --- p. 140 ---
\setcounter{footnote}{0}%
\opage{140}hvad Livet bringer, selv om det ikke uden videre passer med den
herskende Grundstensmening. De lyse Sider ved Livet maa ogsaa % PRINTED AS IS: „Grundstensmening“ (sense: Grundstemning)
gøre deres Virkning.

Af de to modsatte Livsstemninger, Optimismen og Pessimismen,
synes dog gennemgaaende Optimismen at være en større
Hindring for Opstaaen af den store Humor end Pessimismen.
Optimismen har en Tilbøjelighed til at bevæge sig paa Overfladen,
baade hvad Tanke og Følelse angaar. Den lægger sig Tingene til
Rette paa den nemmest mulige Maade og tager ikke de teoretiske
og praktiske Problemer saa grundigt, som den store Humor forudsætter.
Men Pessimisme kan ogsaa gøre Humor umulig, især
ved sin Ringeagt for Erfaringerne.

I sin Karakteristik af George Eliot bemærker Leslie Stephen,
hvad enhver Læser af hendes Værker vil sande, at hun egentligt
kun fremtræder som Humorist i sine første Bøger (Adam Bede,
The Mill on the Floss, Silas Marner). Senere hen gjorde hendes
tiltagende melankolske Stemning, at Livets lyse Sider trængtes
tilbage i hendes Skildringer. Men Grunden laa tillige og for en
væsentlig Del deri, at hun i sine tidligere Værker byggede paa
egne Oplevelser og Erindringer, medens hun senere henfaldt til
en direkte Konstrueren af Personer og Begivenheder. --- Her
fremtræder Vexelvirkningen mellem Temperament og Erfaring
tydeligt. ---

Vigtigere end Temperamentet er dog maaske den tidligere omtalte
Forskel mellem Enfødte og Tvefødte (se § 2), --- mellem
Naturer, der udvikle sig paa kontinuerlig Maade, og Naturer, hvis
Udvikling fører gennem Brud og Spring. De Tvefødte ville erfare
Livets modsatte Sider stærkere, fordi deres Erfaringer ere successive,
ikke simultane. Alt Andet lige, virke successive Modsætninger
stærkere end simultane, der komme som Dele af samme Oplevelse.
De Tvefødte have Mulighed for at opleve stærkere Kontraster,
end de Enfødte som Regel opleve. Det kan endog for
saadanne Tvefødte, hvis Liv spaltes i Perioder, mellem hvilke
endog Erindringens Baand brister, være saare vanskeligt at naa
til Sammensmeltning eller Organisation af de i høj Grad modsatte
Følelser, der successivt beherske Sindet. Men sker det
alligevel, at Forbindelsen lykkes, ville de Tvefødte blive Humo%
% --- p. 141 ---
\setcounter{footnote}{0}%
\opage{141}rister i højeste Stil. Hos Enfødte trænge Modsætningerne ikke saa
let ind i det inderste Sjæleliv; deres Udvikling har Expansionens,
Selvudfoldelsens Karakter. Dog kunne Oplevelser og Erfaringer
give den i sig selv mest harmoniske Karakter et saadant Indblik
i Livets Dybder, i dets forviklede Væsen, at de faa Betingelser
for en Totalfølelse som den store Humor. Hvis Shakespeare,
som meget kan tyde paa, har hørt til de Enfødte, saa har han
kunnet faa Betingelserne for sin Humor paa denne Maade.

At Humor især er en moderne Følelse, hænger sammen med
dens udpræget individuelle Karakter, der igen er en Følge af
dens sammensatte Karakter, som frembyder Mulighed for en
Mængde forskellige Klangfarver. Etisk og religiøs Følelse ere ogsaa Totalfølelser, men de have en Støtte, en for alle en social
Gruppes Individer fælles Støtte i Samfundsnormer og i Kultusformer,
saa at individuelle Nuancer ikke saa let fremkomme eller
i hvert Tilfælde ikke saa let opdages --- og hvis de opdages, ofte
staa som Noget, der skal fortrænges. Den religiøse Kultus har i
Aartusinder været et vigtigt Grundlag for Totalfølelser. Gennem
den føres de enkelte Individer ind i skiftende Stemninger af Frygt
og Haab, Sorg og Glæde, Lidelse og Herlighed, som Punkt for
Punkt bestemmes ved objektive, anskuelige Billeder og Overleveringer.
„Langfredag var en bitter Dag, men skøn var Paaskemorgen.“
Dette Motto passer paa mange Former af orientalsk og
græsk Kultus saavel som paa den kristne Kirkes største Fest.
Det individuelle Stemningsliv bevæger sig her i foreskrevne
Rytmer. Humor er derimod et Livsresultat, der kun kan naas
paa en for hvert enkelt Individ ejendommelig Maade, idet det
beror paa Individets Forhold til de Oplevelser, der falde i dets
Lod, og maaske ikke falde i nogen Andens Lod. Ogsaa paa det
etiske og det religiøse Omraade træder i nyere Tid Individualiseringen
mere og mere frem, hvorved Følgen bliver en mere
kompliceret Karakter af etiske og religiøse Problemer.

Det slaar ikke ganske til, hvad Jean Paul (Vorschule der
Aesthetik. § 76) siger, at der gives en Alvor for Alle, men en
Humor kun for nogle Faa. Ti mere og mere individualiseres ogsaa
Alvoren. Og de største Sammenstød finde Sted, hvor alvorlige
Mennesker ikke forstaa hverandre. ---

% --- p. 142 ---
\setcounter{footnote}{0}%
\opage{142}Humor forudsætter altsaa en Komplicering, en Subjektivering
og en Individualisering af Følelseslivet, der væsentligt tilhører
den nyere Tid, skønt Sokrates staar som den store Banebryder
for denne hele Udvikling. Som alle store Aander er han paa en
Gang en stor Begynder og et Symptom paa, hvad der var ved
at begynde, og hvad der senere skulde udvikles gennem mangfoldige
historiske og sociale Ændringer. Ved en betydningsfuld
Ændring paa det intellektuelle Omraade maa vi endnu dvæle
lidt.

42. Det græske Aandsliv udviklede sig under en begrænset
Horizont, baade hvad Naturkundskab og hvad historisk Erfaring
angaar, Modsætningen mellem Tilværelsens store Sammenhæng og
de begrænsede Menneskeidealer og Menneskeskæbner kunde ikke
gøre sig gældende for Alvor. Tilværelsen stod som et sluttet Hele,
der kunde opfattes som Hjemsted for Alt, hvad Trang og Tanke
kunde forme. Middelalderen bevarede det begrænsede Verdensbillede.
Græske Tanker anvendtes mere eller mindre tvungent
ved Dogmernes intellektuelle Udformning. Og efterat den urkristelige
Tids levende, ekstatiske Forventning havde veget Pladsen
for mere taalmodige Tanker, blev det Opgaven at indordne
Følelseslivet under den Ramme, som Oldtidens Verdensbillede og
den dogmatisk udformede Troslære i Forening dannede. Indenfor
denne Ramme kunde Kærlighed og Begejstring, Angst og
Jubel udfolde sig trygt og rigt. Undertiden kunde der ud fra
Følelseslivets Trang lægges Grunden til Omformning af Kultus
og derigennem til Udformning af nye Dogmer, som det f. Ex. er
sket ved Madonnadyrkelsen. Men gennemgaaende sørgede de
kirkelige Autoriteter for, at de dogmatiske Grænser ikke overskredes,
selv ikke ved Mystikens dristigste Svingninger. En paa
førstehaands Livserfaring grundet Totalfølelse havde ikke Betingelser
for at udvikle sig.

Her bragte Renaissancen en Forandring, dels ved sin Opdadagelse % PRINTED AS IS: „Opda-dagelse“
af det naturlige Menneske, dels ved sin Udvidelse af
Verdensbilledet. Her kom nye aandelige Livsvilkaar. Menneskelivet
i dets rent naturlige Udfoldelse blev set paa Baggrund af en
uoverskuelig Verdensorden. Denne Modsætning førte dels til en
% --- p. 143 ---
\setcounter{footnote}{0}%
\opage{143}dristig Selvhævdelse som Grundlag for al personlig Udvikling,
dels til som Karakteregenskab og Grundsætning at fremhæve det
høje Sind, der lægger sig for Dagen ved, at Livet med alle dets
Interesser og Opgaver ses som en forsvindende Episode indenfor
en uoverskuelig Tingenes Orden. --- Ordet Højsind (magnanimitas
eller sublimitas) var et Slagord i Renaissancens Literatur. Den
Egenskab, det betegner, stod som en naturlig Udfoldelse af den
frie Kraft, med hvilken man nu var tilsinds at tage fat paa Livet
og dets Opgaver. Under forskellige Former fremtraadte dette
Begreb hos den ældre Renaissances Moralfilosofer\footnote[1]{\textit{Fiorentino: Il Risorgimento Filosofico nel Quatrocento.} Cap. IV. --- \textit{A. Herrlin: Renæssancens Etik.} (Trykt som Manuskript.)}, og fik især
videregaaende Indflydelse paa den følgende Tid gennem Telesio’s
Etik, ifølge hvilken Højsind staar som den højeste Form for
Selvhævdelse.\footnote[2]{\textit{Den nyere Filosofis Historie.} I. p. 96.} Hos Bruno, Descartes og Spinoza forener sig saa
med den rent psykologiske Udredning af Højsindet Indflydelsen
paa Sindet af det nye Verdensbillede i dets Uoverskuelighed.
Hvad der hos en Natur som Pascal vakte Rædsel, Menneskehedens
forsvindende Plads i Universet, det vakte hos Renaissancens
store Tænkere netop en Følelse af Adel og Værdighed.

Det var ikke blot den større Udstrækning i Tid og Rum, der
paavirkede Sindet. Udvidelsen af Verdensbilledet efterfulgtes af
en stedse mere indgaaende Forstaaelse af Lovmæssigheden i Tilværelsen.
Verdensbilledet godtgjorde mere og mere en indre
Sammenhæng mellem Alt i Verden, det Ringeste som det Højeste,
gennem almene Love.\footnote[3]{\textit{Den nyere Filosofis Historie.} Anden Bog.} Støvgranet og Kloden bevæge sig efter
de samme Love. Intet staar isoleret og tilfældigt. Selv det Ubetydeligste
har sin bestemte Plads i Tilværelsen, og det i Kraft af
den samme Lov, som anviser det Store dets fremragende og
iøjnefaldende Plads.

I samme Retning som Naturopfattelsen virkede snart ogsaa den
historiske Forstaaelse, der udviklede sig i det attende og det nittende
Aarhundrede. Indenfor Menneskelivet, der var en Episode
i Verdenslivet, kom igen det enkelte Folks, det enkelte Slægtleds
og det enkelte Individs Liv til at staa som Episoder. Herved dan%
% --- p. 144 ---
\setcounter{footnote}{0}%
\opage{144}nedes en nærliggende Baggrund, som maatte virke bestemmende
ved den Enkeltes Opgør af sine Livserfaringer. Ogsaa her stod
det Store og det Ringe i indre Forbindelse for det forstaaende
Blik, og ogsaa fra dette Synspunkt blev der stillet den Opgave
at forene Sympatien med det Smaa med Sansen og Ærefrygten
for det Store.

En sidste intellektuel Betingelse hidførte den kritiske Filosofi
(Erkendelsesteorien). Den paaviste, hvorledes al vor Videnskab
udvikler sig af Trangen til at skelne mellem Indbildning og
Virkelighed, og at denne Skæbne kun kan gennemføres ved Anvendelse % PRINTED AS IS: „Skæbne“ (sense: Skelnen)
af den menneskelige Tankes egne Love. Verdensbilledet
og Virkelighedsbegrebet ere til hver given Tid den menneskelige
Tankes Værk. Og hverken hint Billede eller dette Begreb kan
nogensinde afsluttes, saalænge Erfaringens Kilde risler, og saalænge
der er Liv i Tanken. Virkelighed er et Ideal, efter hvilket
vi stadigt maa stræbe, og som gør ethvert Forsøg paa absolut
Afslutning illusorisk. (Smlgn. „Den menneskelige Tanke“. IV B).
Men dette maa blive af afgørende Betydning for den Totalfølelse,
der danner sig som Resultat af Oplevelserne. Tankens Kampe og
Sejre, og de Kaar og Grænser, der vise sig at være den satte,
ere Kendsgerninger af stor Betydning ved dette Opgør. (Smlgn.
§ 34—35).

--- Dersom alle de her antydede intellektuelle Forudsætninger
kunne blive af Betydning for Dannelsen af en Totalfølelse som
Humor, forstaas det, at den --- ganske afset fra dens komplexe,
subjektive og individuelle Karakter —, først sent og langsomt
har kunnet udvikle sig i sin hele Fylde, og at den maa kunne
fremtræde under meget forskellig Former. % PRINTED AS IS: „forskellig Former“

De intellektuelle Betingelser ere i sig selv ikke nok. (Smlgn
§ 23). De ville, som det følgende Kapitel vil vise, faa forskellig
Betydning efter de forskellige individuelle Forhold, under
hvilke de optræde.

% --- p. 145 ---
\setcounter{footnote}{0}%
\opage{145}\begin{center}IX. HUMOR OG FILOSOFI\end{center}


43. Den store Humor, saaledes som jeg har forsøgt at beskrive
den, er en Livsanskuelse, eller i hvert Tilfælde en Livsstemning.
Den kan drage Næring af Alt, hvad Mennesket erfarer, lider og
nyder, virker og føler. Derfor vil ogsaa et Menneskes Filosofi,
om han har nogen, øve sin Indflydelse her. Og der er jo, som vi
have set, intellektuelle Betingelser, som kunne gøre sig gældende
ved denne Livsstemnings Opstaaen. Men det er kun medvirkende
Betingelser. Og Filosofi i og for sig selv vil her aldrig være
eneste Aarsag. Filosofi kan føre til at gøre vort aandelige Regnskab
op, klare de Forudsætninger, som Videnskaben stiltiende
bygger paa, og som Livet stiltiende stoler paa. Ligesaalidt som
Filosofien kan træde istedetfor al anden Videnskab, ligesaalidt
kan den træde istedetfor Livet. Filosofien maa ønske, at de andre
Videnskaber udfolde sig rigt og kraftigt, ti derved stilles der
stedse nye Opgaver i erkendelsesteoretisk Henseende; der bliver
nyt Arbejde at gøre med ad Analysens Vej at kaste Lys over
menneskelig Erkendelses Natur og Forudsætninger. Filosofien
maa ogsaa ønske, at Livet føres energisk og fyldigt, ti da bliver
der nye Opgaver for den, der vil undersøge det Grundlag, hvorpaa
Menneskers Livsførelse bygger. Livet og dets Fænomener ere
Emner for Filosofi; kun indirekte og delvis kunne de tillige være
Virkninger af Filosofi. Filosofien er, som Sibbern saa smukt og
rigtigt hævdede, kun et Element i det hele aandelige Liv, og den
øver sin Indflydelse under Brydning med andre Magter i Tankens
og i Livets Verden.\footnote[1]{\textit{F. C. Sibbern: Om Filosofiens Begreb.} 1843. (p. 82—86: „Filosofien som Led og Moment i den hele aandelige Leven“). Smlgn. for mit Vedkommende \textit{Den menneskelige Tanke} (1910). p. 380 og min Afhandling \textit{Aandelig Kultur} (Mindre Arbejder III).} Den personlige Livserfaring og dens
% --- p. 146 ---
\setcounter{footnote}{0}%
\opage{146}Udbytte falder ikke uden videre sammen med, hvad filosofisk
Tænkning ad sin egen Vej har fundet. Der er endnu et Arbejde
at gøre for at bringe Harmoni mellem de to Indflydelser. Det er
ikke blot i Slægtens Udvikling, at de aandelige Magter brydes.
Denne Brydning foregaar i enhver vaagen Personligheds Indre
og paa forskellig Maade i de forskellige Individer. Her er et rigt,
endnu meget lidet opdyrket Omraade for psykologisk Forsken.
Det er Forholdet mellem Erkendelsesarbejdet og det tildels aldeles
uvilkaarlige Arbejde, der øves indenfor Følelses- og Villieslivet.
Spørgsmaalet er, om det, der foregaar paa de to Arbejdsomraader,
hviler paa de samme Forudsætninger, --- om de Forudsætninger,
det -ene Omraade hviler paa, kunne udledes af det
andets Forudsætninger, --- eller om der maaske er Strid mellem
de to Grupper af Forudsætninger.

Her ville vi nu undersøge dette Punkt ad historisk Vej, hvad
den Totalfølelse, vi her have beskæftiget os med, angaar. Vi ville
undersøge, hvorvidt vi hos store Filosofer finde, at deres Tankegang
og Tankeindhold har kunnet afgive Baggrund for Humor
som Totalfølelse. Det vil vise sig, at Filosoferne gennemgaaende
have været mest optagne af at udforme den intellektuelle Baggrund,
hvorimod det Grundlag, den store Humor forudsætter,
kraftigst har kunnet udfolde sig der, hvor den intellektuelle Baggrund
gik i Et med praktisk Livserfaring. Filosofiens Historie har
én stor Humorist at opvise, og det er en Mand, der tilhører
Livet nok saa meget som Videnskaben. Hos Sokrates gik Liv og
Tanke i Et som hos ingen anden Filosof, og hans Livsstemning
kan ikke benævnes bedre end ved Ordet Humor, i den Betydning,
i hvilken det er brugt i denne Afhandling. Andre Filosofer
gaa overvejende op i Tankearbejdet og faa kun indirekte Betydning
for den store Humors Historie. Her ses det netop, at Filosofien
kun er et Element i Livets hele Sammenhæng, og at dens Betydning
indenfor denne Sammenhæng betinges ved Vexelvirkningen
med mange andre Elementer. Der er en Arbejdsdeling,
som strækker sig helt ind i Aandslivets inderste Regioner, og.
som gør et fuldkomment Sammenspil af alle det indre Livs Elementer
saa vanskeligt og saa sjeldent. Derved forklares, at Filofien % PRINTED AS IS: „Filo-fien“ (sense: Filosofien)
kun har én stor Skikkelse at stille ved Siden af den store
% --- p. 147 ---
\setcounter{footnote}{0}%
\opage{147}Digter, der har givet de mest gribende Billeder af Livets Modsætninger,
og som har havt den dybeste Forstaaelse af, hvorledes
den maa være tilmode, der gennemlever og gennemrystes af disse
Modsætninger. Sokrates og Shakespeare ere de to største Humorister.
Den Ene koncentrerede sit Liv om en eneste Opgave;
den Anden havde Livets rige Fylde til sit Grundlag.

Hermed er naturligvis ikke sagt, at Filosoferne gennemgaaende
skulde have manglet Totalfølelse. Men deres Totalfølelse faar
væsentligt Karakteren af Erkendelsesglæde, af intellektuel Begejstring
over at finde tankemæssige Synspunkter for Tilværelsen.
Deres Totalfølelse er altsaa forskellig fra den, vi undersøge i
denne Afhandling. Men denne Forskel er saa karakteristisk, at
det vil kunne tjene til Belysning af vort Emne at dvæle nærmere
ved den.

44. I mit første Studenteraar var jeg tilstede ved en Doktordisputats,
hvor Sibbern med stort Eftertryk hævdede Urigtigheden
af Kierkegaard’s Opfattelse af Sokrates som Ironiker. I Modsætning
dertil slog han fast, at Sokrates var Humorist. De to Ting
udelukkede ifølge Sibbern hinanden. Og overfor Kierkegaard
havde Sibbern utvivlsomt Ret heri. Kierkegaard har jo ogsaa selv,
især i „Uvidenskabelig Efterskrift“, korrigeret den Opfattelse han
havde gjort gældende i sin Afhandling „Om Begrebet Ironi med
særligt Hensyn til Sokrates“. (Se ovenfor § 16 og § 18). Denne
Berigtigelse hænger sammen med den tidligere omtalte udførligere
Undersøgelse af Forholdet mellem Ironi og Humor i
„Uvidenskabelig Efterskrift“. Naar Sokrates her bestemmes som
Etiker med Ironi til Inkognito, men hen til Grænsen af det Religiøse,
saa betyder det, ifølge Kierkegaards egen Lære om „Stadierne“,
at han var Humorist, ti Ironien ligger i „Stadiernes“
Række foran det Etiske, og Grænsen mellem det Etiske og
det Religiøse er netop det Sted, hvor Humor ifølge Kierkegaard
har sin Plads i Rækken. (Skrifter VII p. 73 og 437). Allerede
i „Stadier paa Livets Vej“ (VI p. 342) hedder det, at Sokrates’
Alvor var skjult i Spøg. Men Alvor skjult i Spøg er Humor.
--- I sine „Papirer“ (udg. af Heiberg og Kuhr, VI p. 118) korrigerer
Kierkegaard med særligt Eftertryk Disputatsens Opfattelse,
% --- p. 148 ---
\setcounter{footnote}{0}%
\opage{148}som han forklarer ved Indflydelsen af Hegel. --- Han kunde have
føjet til: og af Teologien.

Men det var kun overfor Kierkegaard’s Opfattelse (i Disputatsen),
at Sibbern havde Ret i sin Paastand, at Ironi og Humor
udelukke hinanden. Ironi kan, som vi have set (§ 18), meget vel
være en Form for Humor. Og en historisk Undersøgelse fører
netop til det Resultat, at dette er den mest træffende Betegnelse
for Sokrates’ Standpunkt.

Den sokratiske Ironi var væsentligt en Metode, den velbekendte
analytiske Metode. Denne bestaar i at gaa ind paa en Antagelse
og fra den at gaa tilbage til dens Forudsætninger, som saa enten
kunne stemme med eller stride mod andre Forudsætninger eller
Antagelser, som hyldes af den, der hævder hin Antagelse. Saaledes
gaar Sokrates med tilsyneladende Alvor ind paa, at den,
han taler med, virkeligt forstaar, hvad han siger sig at vide. Ved
at drage alle Konsekvenser af denne Antagelse kommer han da
ofte til at blotte Uvidenheden eller Selvmodsigelsen hos den fra
først af saa tillidsfulde Deltager i Samtalen. Det ser ud, som om
det var en Sport, Sokrates her øvede, og utvivlsomt har han havt
Tilfredsstillelse ved at øve den Tankevirksomhed, hans Metode
krævede. Men det gjaldt for ham ikke om Spøg og Leg. Den
ironiske Behandling skulde ægge til nyt, alvorligt og selvstændigt
Tankearbejde, og dette var det, Sokrates egentligt vilde. Der
laa Sympati, „Erotik“, bag Ironien. Spøgen var gennemsyret af
den alvorligste Interesse for Mennesker. Dette viser sig især i
nogle Udtalelser i Platon’s „Apologi“ og „Symposion“, der ifølge
nyere Undersøgelser\footnote[1]{\textit{Heinrich Maier: Sokrates.} Sein Werk und seine geschichtliche Stellung, Tübingen 1913. p. 106 ff; 137 ff.} ere at betragte som sokratesbiografiske. I
„Apologien“ (p. 29 C) siger Sokrates: Saa længe jeg er istand
dertil, vil jeg ikke høre op at stræbe efter Sandhed („at filosofere“),
og hvem det saa er, jeg træffer paa, vil jeg sige til ham:
skammer du dig dog ikke over at tragte efter at faa saa mange
Penge, saa megen Anseelse og Ære som muligt, medens du ikke
bekymrer dig om Indsigt og Sandhed og om, at din Sjæl kan
blive saa god som muligt. Og Alkibiades vidner (i „Symposion“ p.
215—216), hvor betagen og fængslet han blev i Samtale med
% --- p. 149 ---
\setcounter{footnote}{0}%
\opage{149}Sokrates. Hans Hjerte bankede, hans Taarer strømmede, og han
følte den største Misfornøjelse med sig selv. Sokrates var det
eneste Menneske, overfor hvem han skammede sig. Og naar Sokrates
i Alvor aabnede sig for ham, havde han gennem den ydre,
satyragtige Skal opdaget et guddommeligt Indre.

Dersom disse Ytringer have historisk Grundlag, har Ironien
været en ydre Form for et stort og herligt Indhold, der vel ikke
blev udtrykt i bestemt og positiv Form, men i hvilket et højt Begreb
om Sandhed, om menneskelig Værdighed og om Sandhedens
Betydning for menneskelig Værdighed har ligget til Grund. Der
har været en ideal Baggrund, hvad vi nu kunne kalde Personlighedsprincipet,
i Forhold til hvilken Sokrates’ hele Færd blandt
Mennesker maa ses. Hans Spøg og hans Tankesport have stedse
havt Alvor til Grundlag. Det er ikke rigtigt, med Lipps\footnote[1]{\textit{Komik und Humor.} p. 238.}, at anvende
paa Sokrates, i Anledning af Aristofanes’ „Skyer“, den
Sætning, at den ler bedst, som ler tilsidst. Ti tilsidst lo Sokrates
ikke, og ingen Humorist gør det. Sokrates har agtet alt Andet
ringe i Forhold til en Selvbesindelse, der kunde bære Livet. Det
gjaldt at blive klar over sin virkelige Trang og sin virkelige Evne,
for at man ikke skulde komme ind paa Veje, der ikke stemmede
med Ens Personlighed. Sokrates er den store Sandhedsstræbens
og den store Selvomsorgs Profet, begge Dele i uadskillelig Forening.
Videnskab og Liv kunne begge med lige Ret beraabe sig
paa ham. Der brugtes allerede i Oldtiden det Udtryk om ham, at
han var den Første, hvis Tænkning angik selve Livsførelsen.\footnote[2]{\textit{Diogenes fra Laërte}. II, 20 (πρῶτος περὶ βίου διελέχθη).}
Han havde høje Tanker om, hvad det var at være Menneske,
særligt at være Græker, allersærligst at være en Atenienser. Og
i Forhold dertil fandt han de sædvanlige Goder at være af ringe
Værdi. I sin egen klare Tanke skulde Mennesket finde det Gode,
der betingede alle andre Goder, og uden hvilket disse i Virkeligheden
ikke ere utvivlsomme og sikre Goder.\footnote[3]{Selv hos \textit{Xenofon,} der ellers giver et saa nøgternt og overfladisk Billede af Sokrates, kommer denne Tanke frem med filosofisk Sikkerhed i \textit{Memorabilia} IV, 6 (Samtalen med Euthydemos). Men \textit{Heinrich Maier} (\textit{Sokrates.} p. 57—62) har da ogsaa paavist, at hele det Parti af Memorabilia, hvortil denne Samtale hører, skyldes Laan eller Efterligning af platoniske Dialoger.}

% --- p. 150 ---
\setcounter{footnote}{0}%
\opage{150}Gennem den sokratiske Forbindelse af Ironi og Humor forstaar
man bedst Sokrates’ Forhold til Religionen. Han er gaaet ind paa
Folketroens Forestillinger og har i dem fundet Udtryk for sin
Tillid til, at hans Livsarbejde ikke var forgæves, uden at han iøvrigt
lagde Vægt paa at udforme eller begrunde denne Tro nærmere.
Kun det holdt han sig til, at „en god Mand kan intet Ondt
vederfares, hverken i Liv eller Død“ („Apologi“ p. 41). Denne
Tro er udsprungen af hans Tro paa den etiske Stræbens Gyldighed.
Man kan af den Begrundelse, hvormed han (i Platon’s
„Fajdros“) afviser rationalistiske Omtydninger af Myter, vistnok
danne sig en Forestilling om, hvorledes han stod overfor Folketroens
Forestillinger. Han afviser ikke saadanne Omtydninger af
Ærefrygt for Guderne, men dels fordi der, naar han begyndte
paa Sligt, vilde strømme saa mange Uhyrer til og fordre Omtydning,
at det vilde blive uoverkommeligt, dels (og det er Hovedsagen)
fordi han endnu har nok at gøre med at finde Rede i sig
selv og afgøre, om Harmoni eller kaotisk Lidenskab er det Væsentlige
i hans Natur. Her have vi Spøg og Alvor, Ironi og Humor
paa én Gang. Overfor Selverkendelsens store, alvorlige Opgave
spøgede han med al Teologiseren. Platon’s „Eutyfron“,
hvor Etikens Afhængighed af Gudetro bestrides, kan tages som
Vidnesbyrd om, at det ikke blot var overfor rationalistisk Teologi
han, i Tillid til „Gudebillederne i sit Indre“ (for at bruge Alkibiades’
Udtryk), stillede sig ironisk.

Overfor Udødelighedstroen har han, om man kan tro Platon’s
„Apologi“, særligt stillet sig humoristisk. Han henstiller det der
som uafgjort, om Døden er en Søvn, eller om den er en Overgang
til et andet Liv, der skulde føres i Samkvem med Guder. Det er
tydeligt nok den første Mulighed, der har ligget ham nærmest.\footnote[1]{Smlgn. \textit{Heinrich Maier: Sokrates.} p. 435 f. --- I Modsætning hertil mener \textit{Burnet (Greek Philosophy from Thales to Platon.} London 1914. p. 182), at „Apologien“ maa fortolkes i Overensstemmelse med „Fajdon“, hvis Tankegang han anser for sokratisk (p. 193 f). Jeg kan ikke finde Overensstemmelse mellem „Apologien“ og „Fajdon“.}
Men en dogmatisk Afgørelse og en stor Vægt paa en saadan Af%
% --- p. 151 ---
\setcounter{footnote}{0}%
\opage{151}gørelse vilde ikke stemme med den store Betydning, han tillægger
klar Erkendelse, ej heller med den centrale Stilling det etiske
Synspunkt indtager for ham i Forhold til det religiøse. Den dobbelte
Mulighed, han drøfter, er ikke Udtryk for Tvivl og Uro i
Sindet, men for en indre Sikkerhed, hvad Livets væsentlige
Opgave angaar. Humoristisk er især den Udtalelse, at hvis der
var et andet Liv, vilde det være herligt for ham at fortsætte sin
Virksomhed der og anvende sin Metode paa de homeriske Helte
for at se, om det skulde forholde sig med dem som med saa
mange af hans kære Landsmænd, at de mente at være vise uden
at være det! ---

Det synes da at være fuldt berettiget, naar Sokrates paa flere
Steder i det Foregaaende er opfattet som Repræsentant for den
store Humor. Hans Ironi har staaet i dens Tjeneste.

45. Det ene Punkt, hvorom Sokrates havde koncentreret sig,
straalede for hans store Discipel ud i en stor Tankeverden. Platon
fandt i de Tanker, der gøre os Tingene forstaaelige, den
sande Virkelighed, og overfor dem stode Erfaringens Fænomener
og Livets Forhold som flygtige Billeder. Han fordrede af Tænkeren
og af Statsmanden en streng videnskabelig Dannelse, der
gennem Matematiken som Mellemled skulde føre til Indsigt i
Erkendelsens og dermed, efter hans Opfattelse, ogsaa i Tilværelsens
sidste Grunde.\footnote[1]{Se om \textit{Platon’s} Betydning i Tænkningens Historie \textit{Den menneskelige Tanke.} p. 117—121; 123; 140.} Og naar det høje Trin er naat, vil Sindet
være renset fra al Lavhed og Smaalighed. „Tror du,“ siger han
(Staten VI p. 486), „at den, som har et højt Sind, og som har
faaet et Overblik over al Tid og al Væren, kan betragte det menneskelige
Liv som noget Stort?“ --- Vel forlanger Platon, at
Tænkeren, efterat han har naat Erkendelsens Højdepunkt, skal
stige ned mellem Menneskene og udføre et Hverv i Statens Tjeneste.
Men det kommer som en ydre Nødvendighed. Det gælder
at hindre, at de Slette komme til at regere. Derfor maa de Gode,
de, der nødigst ville beskæftige sig med Statssager, tage Styret.
Der bliver altsaa ikke noget direkte Forhold mellem den ideale
Tilværelse og Livet i den praktiske Verden. Og skønt Platon,
% --- p. 152 ---
\setcounter{footnote}{0}%
\opage{152}som det rige Menneske og den store Digter, han var, havde en
levende Opfattelse af Livets Modsætninger, og Livet stod for ham
som Komedie og Tragedie paa én Gang (se § 40), saa glemmer
han dog i Ideernes Verden baade Tilværelsens Tragik og dens
Komik. Forandringer og Modsætninger raade kun i den sanselige
Verden; men i Ideverdenen raade Uforanderlighed og Enhed.
Hverken den tragiske Jammer eller den komiske Latter er der
Brug for i Tankens Verden. Og overvættes Latter eller Jammer
er af det Onde, fordi Sindet derved sættes i saa stærk Bevægelse
i en enkelt Retning, at dets Fasthed og Enhed svækkes. Platon
anvender ikke paa Komedie og Tragedie, hverken Livets eller
Kunstens, hvad han saa indtrængende hævder om Tonekunsten,
at der kan vækkes en Stemning, som kan vedblive at holde sig,
selv om der ikke kan gives Grunde for de Tanker eller de Gerninger,
den afføder („Staten“ III p. 401). Der er her en bestemt
Modsætning mellem hans Antydning i Slutningen af „Symposion“,
at det maa ligge til samme Mand at digte Komedie og Tragedie,
og hans Standpunkt i „Staten“. Han tager i dette sidste Skrift
udtrykkeligt Muligheden af at forbinde Komedie og Tragedie tilbage.
(„Staten“ III p. 295 B).\footnote[1]{\textit{Hans Ræder (Platons philosophiscke Entwickelung.} Leipzig 1905 p. 208) finder ingen Modsigelse mellem de to Dialoger paa dette Punkt, idet Enheden af Komedie og Tragedie i „Symposion“ skal opstilles som et Ideal, medens denne Enhed i „Staten“ opfattes som faktisk umulig. Men det er i Virkeligheden to Idealer, der stilles overfor hinanden: overfor det poetiske Ideal i „Symposion“ stilles i „Staten“ det etiske Ideal af en paa Selvbegrænsning og indre Harmoni bygget Personlighed, der ikke kan give Rum for den store Modsætning mellem det Tragiske og det Komiske.} % PRINTED AS IS: „philosophiscke“

En Sammensmeltning eller Organisation, som den, der ligger i
Begrebet Totalfølelse, ansaa Platon for umulig, naar Livets største
Modsætninger kom til at skaffe sig det fulde, lidenskabelige Udtryk.
Intellektuel og etisk Interesse var tilsidst saa stærk hos
ham, at hans rige menneskelfge Erfaring ikke kunde afsætte sin % PRINTED AS IS: „menneskelfge“
fulde Virkning i hans Sind, i hvert Tilfælde ikke en Virkning,
som hans Tanke kunde anerkende og udtrykke. For ham, som
idethele for Oldtidens Tænkning, stod nu engang det Uforanderlige
og det Foranderlige som de to største, uforenelige Modsætninger,
% --- p. 153 ---
\setcounter{footnote}{0}%
\opage{153}mellem hvilke der tilsidst maatte træffes et Valg. Og for Platon
var der ingen Tvivl om, hvor den sande Værdi laa.

Hans Totalfølelse blev intellektuel Begejstring, Erkendelsesglæden,
som vaktes ved Tankens Arbejde paa at naa til de evige
Ideer og derved hæves op over Forandringens Strøm.

Alligevel faar Platon stor Betydning for Udviklingen af den
Totalfølelse, hvormed vi have beskæftiget os, ved sin store Tanke
om Tilværelsen som en stor Tingenes Orden, i Forhold til hvilken
Menneskelivets Begivenheder stode som forsvindende. Dermed
var en intellektuel Baggrund given, som muliggjorde en
Humor af anden Klangfarve end den sokratiske. Men Platon var
saa optagen af Udarbejdelsen af denne Baggrund, og saa begejstret
for den, at der ikke blev Tid og Kraft til at lade den
psykiske Sammenhæng mellem denne Baggrund og Oplevelserne
i Erfaringens Verden komme til dens fulde Ret.
Der manglede desuden det positive Forhold mellem Uforanderlighed
og Foranderlighed, som først den nyere Tid skulde
bringe.

Naar saa godt som al Energi tilflyder en enkelt Side af Sjælelivet,
bliver der naturligvis saa meget des mindre Energi tilovers
til de andre Sider. Derved forstaas, at de, der ere mest optagne
af at udarbejde og gennemarbejde den intellektuelle Baggrund
for Livet, ikke selv komme til at staa som Repræsentanter for
den store Humor. Vi staa her ved en ny Grænse for Humor som
Livsanskuelse eller Livsstemning, forskellig fra de Grænser, som
paa den ene Side Tragiken, paa den anden Side Handlingen
satte. Den Grænse, det intellektuelle Arbejde kan sætte, kommer
ikke udefra. Den store Humor behøver ikke at forudsætte en
dogmatisk, en Gang for alle afsluttet Baggrund, der skulde udelukke
fortsat intellektuelt Arbejde. Tvertimod vil den store Humor,
som vi have set (§ 34—35), drage den faktisk vundne Erkendelse
med ind og tage den „humoristisk“, fordi der stedse vil
være Mulighed for nye Gaader og nye Løsninger. Bag enhver
Baggrund viser sig Muligheden af en endnu mere omfattende
Baggrund. Problemernes stadige Genopstaaen i nye Former staar
som et komisk, maaske komitragisk Fænomen, og der kan fremkaldes
et Smil over det travle Erkendelsesarbejde. Men selve
% --- p. 154 ---
\setcounter{footnote}{0}%
\opage{154}denne Stemning vil ikke saa let opstaa hos dem, der staa midt i
Problemernes Skærsild og føle deres Svien. ---

At Sokrates er den eneste Humorist i stor Stil blandt Filosoferne,
beror paa, at hos ham det intellektuelle Arbejde gik i Et
med det praktisk-pædagogiske Arbejde overfor Mennesker. Gennem
Ironien som Metode søgte han at føre de Enkelte til at besinde
sig paa den store Baggrund, et Menneske kan finde i sit eget
Indre, hvad enten denne Baggrund kan formes i klare Tanker
eller ikke. Spøgen, Ironien var her en Vej til Alvoren.

Den anden store Humorist i det menneskelige Aandslivs Historie
førtes af sine Oplevelser til at frembringe en Billedverden,
indenfor hvilken Sorg og Glæde, Mørke og Lys kunde staa som
Elementer i en Totalerfaring, en Billedverden, der lader baade
det Smaa og det Store komme til dets Ret, uden at det Store
knuser det Smaa, og uden at det Smaa drager det Store ned i
det Lave.

Baade den filosofiske og den poetiske Humorist stode, hver
paa sin Vis, i et levende Forhold til Livet, saa rig og ophøjet
deres indre Verden end var.

46. I Henhold til den tidligere antydede historiske Betragtning
(Kap. VIII), ifølge hvilken vi først i nyere Tid kunne vente
Grundlag og Baggrund for den store Humor, gaa vi nu lige fra
Platon til Renaissancen. I Middelalderen vil det snarest være i
Folketro og Folkeskik, man vil finde Humor, saaledes i Forestillingen
om den dumme Djævel. Et Udtryk for fuldt Opgør med
Livserfaringerne kunde den under Middelalderens Forudsætninger
ikke blive.

Bruno og Spinoza have i store Træk grundlagt den Verdensanskuelse,
som er ejendommelig for den nyere Tid, givet de
første Helhedsbilleder paa Grundlag af, hvad den nye Astronomi
og den nye Fysik havde bragt for Dagen.

Jeg har allerede (§ 9) omtalt Bruno’s Beskrivelse og Forstaaelse
af Totalfølelser. Hans filosofiske og hans poetiske Begavelse virkede
sammen paa dette Punkt, som paa saa mange andre Punkter
i hans Tænkning. Men den Totalfølelse, han kendte, var ikke
Humor, men Livsmod og Højsind, den Følelse, der idethele er
% --- p. 155 ---
\setcounter{footnote}{0}%
\opage{155}karakteristisk for Renaissancen (se § 42). Den kan blive Grundlag
for Humor (se § 20), men behøver ikke at blive det.

Hvad Bruno har ydet for den intellektuelle Baggrund for det
moderne Følelsesliv, var to Ting: Udvidelsen af det fysiske Verdensbillede
i det Uoverskuelige, og Hævdelsen af, at skal der
antages guddommelige Kræfter i Verden, saa maa de søges i
Elementernes og i Sjælenes Indre, ikke et Sted i Rummet. Bruno’s
hele Aandskraft toges i Beslag af Arbejdet med disse store
Tanker. Og den Tapperhed, hvormed han bestod i den Spænding
og Møje, dette Arbejde voldte, laa til Grund for hans Skildring
af Heltesindet eller Højsindet. Men Tankens stadige Higen lagde
Beslag paa alle Sindets Kræfter, og den intellektuelle Følelse
indgik ikke nogen Forbindelse med en forstaaende Sympati
overfor det Endelige og Begrænsede i Menneskelivet og i Naturen.


For Spinoza stod ikke Tilværelsens Uendelighed, men dens
Lovmæssighed i første Række. Han indskærper som det vigtigste
Middel til aandelig Frihed, at vore Følelser knyttes til gyldige Tanker,
i hvilke vi ganske kunne finde Hvile. (Ethica V, 4 Schol).
Følelseslivets Bølgegang skal ophæves ved, at vi føle os Et med
Verdenssubstansen, med det, der bestaar under alle Forandringer,
idet vi indse vor Aands Sammenhæng med hele Naturen (De
emendatione intellectus). Da tabes det Smaa og det Skiftende af
Syne, og Erkendelsesglæden bliver den herskende Stemning.
Spinoza har ganske vist ikke manglet Interesse for det Enkelte
og Særlige, især paa Sjælelivets Omraade. Han har jo ikke
blot skrevet „Etiken“s første og anden Bog, der handle om
Verdenssubstansen og om den menneskelige Erkendelse, men
ogsaa dens tredie og fjerde Bog, der undersøge de menneskelige
Følelser og Lidenskaber. Men de to Tankegange, --- den, der
søger Tilværelsens og Tankens Grundforhold, og den, der udforsker
Erfaringens Enkeltheder, --- komme ikke til at forene sig i en
og samme Række af Oplevelser, saa at de begge komme til at
bestemme Totalfølelsen. Livets Højdepunkt var for Spinoza den
Skuen (Intuition), der med ét Blik overskuede og gennemskuede
den store Sammenhæng, i hvilken Alt føjer sig ind. I sin Beundring
for Naturens uendelige Magt og Fylde, der ikke blot for%
% --- p. 156 ---
\setcounter{footnote}{0}%
\opage{156}maar at frembringe det i vore Øjne Fuldkomne, men ogsaa det i
vore Øjne Ufuldkomne (Ethica I. Slutning), --- en Bemærkning,
han stillede i skarp Modsætning til Menneskenes overtroiske
Frygt eller bornerede Forargelse, --- i denne Beundring hævedes
han over Livets Komedie og Tragedie, og disse Modsætninger
kom ikke til paa positiv Maade at bestemme hans Sindelag overfor
Livet.

Det vilde være urigtigt heraf at slutte, at Spinoza i sit personlige
Liv var uberørt af de Følelser, hvis Sammenspil betinge den
store Humor. Baade Medfølelse og Latterfølelse have spillet en
Rolle hos ham.

Andres Ulykker gik ham, berettes det af hans første Biografer
(Lucas, Colerus), nærmere til Hjertet end hans egne, og det var
ham en hæslig Trøst, man kunde finde for sine egne Lidelser
deri, at Andre lide som vi. Kun den rent passive Medlidenhed,
der fører til blind og tankeløs Indgriben, forkastede han. Efter
Spinoza’s „Omvurdering af alle Værdier“ ere kun de Følelser
i og for sig gode, der ere Udtryk for aktiv Holdning og
Brug af egen Kraft og Eftertanke. Kun Glæde, ikke Sorg kan
derfor i sig selv være god. Dog kunne hos Mennesker, der ikke
ledes af klar Eftertanke, ogsaa Sorg og sorgfulde Stemninger
(Medlidenhed, Frygt og Anger f. Ex.) være af Værdi.

Men han saa ikke Muligheden af, at en Følelse, der i og for
sig ikke har ubetinget Værdi, kan indgaa som Element i en
Totalfølelse, som derved kommer til at staa som Frugt af mangesidig
Livserfaring og gør Sindet vel rustet til at møde, hvad Livet
fremdeles kan bringe.

En sund og frisk Latter regnede Spinoza til de gode Ting i
Livet, med stærk Betoning af sin Modsætning til asketisk og pietistik % PRINTED AS IS: „pietistik“
Tankegang (Eth. IV, 45 Schol.). Han gjorde kun en Forskel
mellem den uvilkaarlige Latter (risus) og Haanlatteren (irrisio).
Saasnart Latteren faar en bestemt Genstand, bliver den for Spinoza
Udtryk for Haan og staar i Modsætning til den Forstaaelse,
der for ham er det Højeste. Det er denne Opfattelse, der ligger
til Grund for hans berømte Udtalelse, at han bestræber sig for
hverken at le, sørge eller harmes over menneskelige Handlinger,
men at forstaa dem (Tractatus politicus. I, 4). Han anvendte dette
% --- p. 157 ---
\setcounter{footnote}{0}%
\opage{157}Princip, ikke blot i en statspolitisk Afhandling, men ogsaa overfor
Tidens Begivenheder. Tidens Forvirring, skriver han i et Brev
(Ep. 30) under en Krig mellem Holland og England, vækker
hverken hans Latter eller hans Graad, men fører ham til at forske
og iagttage den menneskelige Natur; Mennesket er jo, som alt
Andet, en Del af Naturen, hvis Sammenhæng med de andre Dele
og med hele Naturen vi ikke kende. Erkendelsesglæden samlede
al hans Interesse; den var hans Totalfølelse. Men den store Humor
hviler netop paa en Forudsætning, som Spinoza ikke hyldede,
den nemlig, at Forstaaelse og Latterfølelse ikke udelukke hinanden.
Latteren kan blive Udtryk, ikke for Afstand, Haan eller Foragt,
men for en Enhedsfølelse. Selve Forstaaelsen kan føre Latter
eller Smil med sig. ---

For Renaissancens Mænd gjaldt det først og fremmest om
Selvhævdelse. Det var Modsætningen til Middelalderens Autoritetstro
og Askese, der gjorde sig gældende hos dem. De optoges
helt af Arbejdet paa at bygge en ny Tankeverden op som Ramme
for Livet. Gennem dette Arbejde kom Bruno til „den heroiske
Affekt“ og Spinoza til „den intellektuelle Kærlighed“. Det var
deres Totalfølelser, hvorimod den ejendommelige Syntese, som
Humoren forudsætter, ikke kom istand.

47. Medens Renaissancen lagde Vægten paa Selvhævdelsen og
de Følelses- og Villiesformer, der kunne udvikle sig af den, førtes
de engelske Psykologer og Moralfilosofer i det attende Aarhundrede,
--- Rækken Shaftesbury, Hutcheson, Hume og Adam
Smith, --- til at betone de sociale og sympatiske Følelser. Dette
blev af Betydning for den Totalfølelse, vi her have undersøgt,
idet denne kan opstaa af en Sympati, som bevirker en psykologisk
Forbindelse mellem Personligheden og dens Objekt og
muliggøre, at Objektet, var det end nok saa ringe, ses i sin Værdi
og Betydning. Deraf følger dog ikke, at de nævnte Filosofer selv
vare Humorister.

Shaftesbury var en begejstret Skønhedsdyrker. Det var ham
om at gøre at bekæmpe Fanatisme og Overtro, der ville sætte
snævre Grænser for Livets frie Udfoldelse. I denne Kamp anser
han Komiken for en god og retmæssig Forbundsfælle. Dette er
% --- p. 158 ---
\setcounter{footnote}{0}%
\opage{158}Grundtanken i hans lille Skrift „Essay on the Freedom of Wit and
Humour“ (1709). (Ordet „Humour“ er her taget i den ældre,
Jonson’ske Betydning af Ordet). „Jeg for min Del,“ siger han,
„frygter ikke det skeptiske Vid. Ingen Sag, som jeg er i mindste
Grad interesseret for, kan skades af et smagfuldt Vid.“ I denne
Tankegang ligger der en Spire til den store Humor, for hvem det
nu engang er saaledes, at selv det Værdifuldeste har en Side,
fra hvilken det er Begrænsning og Endelighed underlagt, saa at
der kan smiles, om ikke les deraf. Men for Shaftesbury drejede
det sig her kun om et Middel, et Kunstgreb, der ikke paa indre
Maade hænger sammen med en hel Livsstemning.

Hutcheson’s Hovedinteresse er etisk, og Smith’s etisk og nationaløkonomisk.
Deres psykologiske Interesse var afledet og underordnet.
Hos Hume kunde man derimod vente at finde Humor.
Han har et klart Blik for Modsigelserne i Verden og for Misforholdet
mellem hvad menneskelig Følelse kræver, og hvad Virkeligheden
udviser. Men de Erfaringer, han i denne Henseende har
gjort, synes ikke at have ført til nogen Totalfølelse. Han opgiver
tvertimod at stræbe efter et samlet Udtryk for Livserfaringerne.
En saadan Stræben hviler efter hans Mening paa en Overvurdering
af Livets Betydning. Derimod gives der --- saaledes slutter
Hume sine Reflexioner (i Afhandlingen „The Sceptic“) --- Naturer,
for hvem det at tænke over Livet er den mest underholdende
Maade at anvende Livet paa!

48. Den engelske Litteratur i det attende Aarhundrede bevægede
sig, især hvad Romanlitteraturen angik, i en Retning, der
frembyder en interessant Parallel til den filosofiske Bevægelse.
Der blev lagt Vægt paa Følelseslivet, paa Sympatien, paa det gode
Hjerte, og Livets ydre Begivenheder fulgtes med et mildt Blik og
et venligt Smil. Vi have her med den Række af Forfattere at gøre,
som Thackeray (The English Humourists of the Eighteenth Century.
1853) har sammenstillet under Navn af de engelske Humorister.
Thackeray’s Definition af Humor gaar ud paa, at den er en Forbindelse
af Vid og Sympati (wit and love)\footnote[1]{Thackeray haaber, at han ogsaa selv kan regnes til Humoristerne efter denne Definition, skønt han i „Times“ er bleven kaldt en kedelig Misantrop, der ikke ser noget Godt nogetsteds. --- Thackeray selv, før ham Dickens og efter ham George Eliot danne en Række engelske Humorister, hvis Karakteristik i Forhold til det 18. Aarhundredes Forfattere af samme Retning vilde frembyde stor Interesse. For mit Formaal vil Karakteristiken af Carlyle (se § 51) være mere nærliggende.}. I denne specielle
% --- p. 159 ---
\setcounter{footnote}{0}%
\opage{159}Betydning synes det ikke at passe paa alle de Mænd, Thackeray
omtaler: Swift, Addison, Steele, Pope, Fielding. Om Swift har jeg
tidligere (§ 15) søgt at vise, at han er Satiriker, ikke Humorist.
Den af de nævnte Forfattere, den nævnte Definition bedst passer
paa, er Sterne, skønt Thackeray snarest vilde kalde ham sentimental.
Goethe, der ellers ikke holdt af „Humorister“, gjorde en
Undtagelse med Sterne (se ovenfor § 25). I sin „Sentimental
Journey“ („en Rejse for at lære det menneskelige Hjerte at kende“)
siger Sterne, at hans Rejse skyldes „Hjertets Trang til Naturen
og til de Følelser, der udspringes af Naturen og faa os til at elske
hverandre og Verden mere end før.“ Dermed er Grundlaget givet.
Det fører til, snart i Alvor, snart i Spøg, snart med Taarer, snart
med Latter at skildre menneskelige Karakterer og Skæbner, med
Sans for de smaa Ting og for Livet i de smaa Kredse. Hvad
Shakespeare i sine Kæmpeværker havde formaaet at fremstille,
Livets store Rigdom og stærke Kontraster, paa Grundlag af dyb
Livserfaring og trofast Kærlighed til det virkelige Liv, det blev
nu af ringere Aander ført ud i vide Kredse, ned i de smaa Tings
Verden, som Mesteren kun episodisk havde givet en Plads i sin
Verden. Det skortede her ofte paa den kunstneriske Evne til at
lade Enkelthederne forme sig skarpt og individualiseret. Den naive,
overstrømmende Følelse legede ofte med Emnerne paa lunefuld
Maade, og istedetfor at være det dybe, men skjulte Grundlag
trængte det gode Hjerte sig gerne frem og talte direkte med, saa
at Hjerteligheden blev Sentimentalitet. Dertil kom Tidens glatte
Optimisme, som gjorde Følelseslivet overfladisk. Man gik let hen
over Disharmonierne. Sansen for det Uendelige og Hemmelighedsfulde,
der havde været saa stærkt udviklet hos Shakespeare og
Milton, var forsvunden hos denne Tids Digtere\footnote[1]{\textit{L. Stephen: English Thought in the 18. Century.} II. p. 371.}.

Alligevel var det af Betydning, at den aandelige og literære
Strømning opkom, som Eftertiden har kaldt Humor, skønt det her
% --- p. 160 ---
\setcounter{footnote}{0}%
\opage{160}ikke var den store Humor. Den var Tegn paa en for Livets Mangfoldighed
mere aaben Psykologi. Og i Løbet af Aarhundredet vaagnede
igen Sansen for Shakespeare’s Storhed, især gennem en stor
Skuespillers Gengivelse af hans Karakterer\footnote[1]{Se om Betydningen af Garrick’s Skuespilkunst for en Shakespearsk Renaissance: \textit{L. Stephens: English Literature and Society in the 18. Century.} p. 83 o. fl. St.}, og derved kunde igen
Sansen for stor Humor vækkes.

49. Renaissancen havde gjort en stor intellektuel Baggrund
mulig, og Poesien havde, snart i stor Stil, snart i smaa Dimensioner,
vakt Sansen for Oplevelsernes Mangfoldighed og Ejendommelighed.
Ved det attende Aarhundredes Slutning kom saa dertil et Forsøg
paa at opgøre Forholdet mellem Tankearbejdet og de andre Sider
af Sjælelivet. Fra det store Verdensbillede og den store Verdensanskuelse,
som Renaissancens Naturvidenskab og Filosofi havde
gjort mulige, steg Immanuel Kant ned til den Kraft i den menneskelige
Aand, som det skyldes, at den overhovedet kan danne
Billeder og Anskuelser, forme Tanker og opstille Idealer. Han
afslørede, hvad han har kaldt „den skjulte Kunst i Menneskets
Indre“, den uvilkaarlige sammenfattende Virken, der ytrer sig i
al Sansning og Tænkning. Og naar der spurgtes om, hvorvidt der
kunde tillægges denne Virkens Resultater Gyldighed, blev hans
Svar, at denne Gyldighed beroede paa, om de vare nødvendige
for at kunne skelne mellem Erfaring og Indbildning, mellem Virkelighed
og Drøm. Derved blev Grundlaget taget bort for alle spekulative
Systemer, som søgte ud over enhver mulig Erfaring.
Teologi og spiritualistisk Psykologi bleve umulige som Videnskaber.
Ikke uden indre Kamp var Kant naat til dette Resultat. Ti han
havde, som han selv sagde, den Skæbne at være forelsket i Metafysiken,
--- og nu saa han, at den var umulig. I denne delte Stemning
havde allerede hans første Undersøgelser bragt ham. Det var
gennem Humor, at han kom ud over Deltheden. I „Träume
eines Geistersehers, erläutert durch Träume der Metaphysik“.
(1766), hersker der en humoristisk Stemning. Overfor det høje og
strenge Sandhedsbegreb, han nu havde faaet Øje paa, stode spekulative
Systemer for ham i et komisk Lys, uagtet det havde været
% --- p. 161 ---
\setcounter{footnote}{0}%
\opage{161}hans eget inderlige Ønske at kunne frembringe et saadant System.
Ogsaa senere, efter Udarbejdelsen af hans berømte Hovedværk,
kan han se med Humor paa sit Arbejde. I en af sine Optegnelser
spøger han med den pedantiske Karakter, „Kritik der reinen
Vernunft“ unegteligt for en Del frembyder; han trøster sig med,
at Pedanteri skal Pedanteri fordrive. Hans sokratiske Karakter,
der lod ham foretrække grundig og selvstændig Viden for Spekulation
saavel som for blot Ophoben af Kundskab og satte det at
filosofere over det at have Filosofi, gav ham en fast Baggrund for
Betragtningen af frugtesløse, men ofte des mere tillidsfulde
Tankeforsøg.

I Kant’s Følelsesliv spillede det Ophøjede den største Rolle ---
Universets Uendelighed udadtil og den i Samvittigheden fremtrædende
indre Uendelighed var efter en berømt Udtalelse af ham
to Ting, der opfyldte hans Sind med bestandigt ny Ærefrygt. Men
som de faste, fornuftige Reglers Mand fæstnede han stedse mere
sit Liv og sin Tanke i bestemte Former, som det var ham pinligt
og tilsidst umuligt at tænke sig anderledes. Den Humor, der kunde
røre sig hos ham i hans yngre Aar paa det intellektuelle Omraade,
kom ikke til at give hans hele Personlighed og Livsstemning sit Præg.

49. Det er indenfor den tyske Romantik, at Begrebet om den % PRINTED AS IS: „49.“ a second time (sense: 50.)
store Humor først udtrykkeligt dannes. Som overalt gaar ogsaa
her Sagen forud for Begrebet. Vi have henpeget til to store Skikkelser
som Repræsentanter for, hvad Begrebet indeholder --- den
ene i Oldtiden, den anden ved Indgangen til den nyere Tid.

De engelske Humorister i det attende Aarhundrede repræsentere,
som allerede bemærket, ikke den store Humor. Deres Horizont
var for begrænset, og deres Optimisme lod ikke Livets Modsætninger
faa deres fulde Indflydelse. Man bevægede sig paa et
jævnt, borgerligt Niveau, og da Følelsen ingen Modsætning og
Modstand havde at overvinde, blev den til Sentimentalitet, især
hos Efterlignere. Ifølge Jean Paul (Vorschule § 32) fik Sterne i
Tyskland en hel Skare Efterlignere, som i Virkeligheden ikke vare
Andet end „Ausplauderer lustiger Selbstbehaglichkait“. Der manglede
en stor Ide som Baggrund. En saadan Ide raadede den tyske
Romantik over.

% --- p. 162 ---
\setcounter{footnote}{0}%
\opage{162}Ikke som om de tyske Romantikere selv havde skabt denne
Baggrund. Som jeg allerede har bemærket i „Den nyere Filosofis
Historie“ (i Indledningen til ottende Bog), byggede de væsentligt
paa den store tyske Renaissance i sidste Del af det attende Aarhundrede (Lessing og Kant,
Herder og Hamann, Schiller og Goethe),
den Udvidelse af den aandelige Horizont, som den havde bragt, og
den Forestilling om Aandslivets Oprindelighed og Mangfoldighed,
som den havde vakt. Af særlig Betydning blev det desuden, at
man lærte forskellige Aands- og Tankeretninger, mange Tidsaldres
Tros- og Stemningsliv at kende, og fik Sans for og Trang til at
sætte sig ind i mangfoldige Synsmaader og Standpunkter og at se
med Ironi paa hver enkelt Aandsform, naar den vilde være den
eneste gyldige. Hos Fr. Schegel have vi allerede fundet dette Træk % PRINTED AS IS: „Schegel“
(se § 16). Med en Blanding af Ironi og Beundring saa man tilbage
paa tidligere Tiders Tro, paa deres Kunst og deres Tænkning. Man
havde en dyb Overbevisning om, at hvad der saaledes stadigt søgte
nye Former og sprængte de gamle, det rørte sig ogsaa nu hos
Alle, hvis Sans og Sind vare opladte. Og det var ikke Noget, der
blot rørte sig i Menneskenes Sind. Hvad der for Tanke og Trang
var det Højeste, det maatte tillige være det Inderste i Tilværelsen,
skønt ingen enkelt Del af vor Erfaring fuldstændigt kan udtrykke
det. Kun en højere Skuen eller en spekulativ Tanke aabnede
Adgang dertil. Hver enkelt Erfaring, hver Forestilling, der var
hentet fra et specielt Iagttagelsesomraade, kunde kun faa symbolsk
Betydning.

Romantikens Filosofi er en utaalmodig Idealisme. Man stræbte
ubændigt, gennem en Slags Magi, at trænge ind til „det, der holder
Verden sammen i dens Inderste“, og man havde ingen Sans
for Arbejdet i det Enkelte, ud fra specielle Udgangspunkter og
ved specielle Metoder. Sligt stod, ligesom det jevne, borgerlige
Livs Opgaver og Interesser, som en Spøg i Forhold til den ideale
Verden, der afslørede sig i saa mange Former.
Romantiken virkede befriende ved at pege ud over de snevre
Former, hvori Liv, Kunst, Religion og Tænkning saa let stivne.
Og den kunde derved gøre stor Humor mulig. Men den romantiske
Humor led af en stor Ufuldkommenhed. Ti Livets smaa
Forhold og Tankens smaa Skridt kunne ganske vist staa som for%
% --- p. 163 ---
\setcounter{footnote}{0}%
\opage{163}svindende og ubetydelige i Forhold til det Ideale og Uendelige,
vi kunne ane; men det er dog gennem Troskaben i det Smaa, at
de store Ting gennemføres. Romantiken opfattede enhver begrænset
Stræben, enhver ensidig Retning blot som noget Negativt, som
Noget, der gjorde Fordring paa at være uendeligt uden at være
det, Noget, der satte Grænse for det omfattende Blik. Derfor
rettedes Ironiens Medusahoved imod dem. Der var her Mulighed
for en Ironi som den, man med Urette fandt hos Fr. Schlegel
(se § 16), og som den, Kierkegaard beskrev i sin Afhandling
om Begrebet Ironi og senere i sin Karakteristik af „det
æstetiske Stadium“. Men der var ogsaa Mulighed for en Forbindelse
af Ironi og Begejstring, som fører og førte til den store
Humor.

Allerede Novalis hævdede, at Fr. Schlegel egentligt var Humorist.
Men den romantiske Humors Teoretiker var Solger, som
tillige anslog den store Humors Strenge. Hos ham forenes det
attende Aarhundredes Humor, der er Forbindelse af Vid og Sympati,
med Romantikens Ide om det Enkelte som pegende ud over sig
selv til en stor Helhed. I sin Livsanskuelse og sin Kunstopfattelse
forener han Ironi og Begejstring. Gennem Ironien fortæres det
Enkelte i dets Isolation, gennem Begejstringen ses det som Led i
Helheden. I hans „Erwin“ er det den store Længsel (§ 21), der
danner Grundlaget; hans efterladte Skrifter synes at vise, at han
i sine sidste Aar slog ind paa en mere mystisk Retning. Han
er Romantiker ved sin Længsel efter en Form for aandeligt
Liv, i hvilken Forskellen mellem Kunst, Religion og Filosofi er
forsvunden. Filosofien skal bane Vej for et saadant Højdepunkt
ved gennem fri Erkendelse at paavise det Evige og Uendelige
som indeholdt i de Love, der gælder for de endelige Ting. Kunsten
skal gennem Billeder vise det Endelige som gennemstraalet
af det Uendelige. Og Religionen er en umiddelbar Leven i det
Evige og Uendelige. Paa Højdepunktet gælde ingen af disse tre Former
for Aandslivet; Forskellen mellem dem falder bort. Den højeste
Tilstand er Begejstringen; men den afføder Ironi overfor Virkeligheden
og dens Fænomener.

Det er et personligt Livsstandpunkt, Solger her har skildret.
Hans abstrakte Udtryk have gjort, at Kierkegaard har skildret hans
% --- p. 164 ---
\setcounter{footnote}{0}%
\opage{164}Stade som rent formelt og negativt og betragter ham blot som en
ufuldkommen Forgænger for Hegel. Men medens for Hegel alle
aandelige Livsformer vare ufuldkomne Former for filosofiske Spekulationer,
hævdede Solger, at den højeste Livsform maatte være
en saadan, til hvilken Filosofi, Kunst og Religion hver ydede
sit Bidrag.

Mangelen ved den romantiske Humor, en Mangel, der træder
mere frem hos Solgers Forgænger Jean Paul end hos Solger selv,
bestaar i, at der ikke kommer noget positivt Forhold mellem det
Smaa og det Store. Der lægges overvejende Vægt paa Umuligheden
af en fuldkommen Fremstilling af det Højeste. Den store
Humor naar først sin fulde Udfoldelse, naar det Smaa og det Begrænsede
ses som nødvendige Former for det Store og Uendelige og
derigennem faa positiv Betydning. Den romantiske Humor kunde
føre, og førte ogsaa, til Ringeagt overfor Erfaringsvidenskaben og
overfor praktiske Reformbestræbelser. Hos Carlyle ville vi tydeligt
se dette. Det realistiske Element i den store Humor, som udspringer
af dens Troskab mod Virkeligheden, kommer ikke til sin Ret
i Romantiken. Hverken Tragedien (Kap. V) eller Forstaaelsen (Kap.
VI) eller Handlingen (Kap. VII) --- de tre Omraader, overfor
hvilke Humoren skal vise sig paa én Gang i sin positive Betydning
og i sin nødvendige Begrænsning --- anerkendes i deres Selvstændighed.
---

Sin fulde Udvikling har den store Humor ikke naat i tysk
Aandsliv. Den tyske Humor bevarer mere eller mindre en romantisk
spekulativ Karakter. Den engelske Humor er mere jordbunden,
men paa Virkelighedens faste Grund er saa Shakespeare’s Verden
af vældige Skikkelser i dramatisk Vexelvirkning voxet frem, en
Verden, i hvilken den inderligste og tillige den mest plastisk udformede
Humor fremtræder. --- Det er karakteristisk, at det franske
Sprog bruger den engelske Form (humour) for at udtrykke
Humor i den Betydning, i hvilken vi her tage den. Den Sammensmeltning
eller Organisation, den store Humor forudsætter, synes
ikke at have naturlige Betingelser paa fransk Jordbund. Det er i
hvert Tilfælde mere karakteristik for det franske Naturel at svinge % PRINTED AS IS: „karakteristik“
mellem meget forskellige Stemninger, som kunne udfolde sig i en
vis Uafhængighed af hinanden. Derved forklares de Overraskelser,
% --- p. 165 ---
\setcounter{footnote}{0}%
\opage{165}Franskmændene ofte byder Verden. Saaledes hævder en fransk
Filosof, at der ikke for en grundig Iagttager er foregaaet noget
psykologisk Mirakel med hans Folk, men at man allerede tidligere
havde kunnet spore de udmærkede Egenskaber, det har lagt for
Dagen i den nu staaende Krig\footnote[1]{Nous avions pu discerner, sous une frivolité apparente, un esprit sérieux et capable d’application, sous des dehors frondeurs, la résolution de se soumettre librement à une discipline raisonnable, et sous la „blague“ superficielle, le sentiment et le respect de l’ideal. \textit{Terraillon:} Allocution prononcée à la Distribution des prix. Collège de Saint-Claude. 1915, p. 9.}. --- Det er maaske dristigt at
slutte, at psykisk Kemi, nøje Forbindelse og Sammenhæng af forskellige
psykiske Tendenser, skulde ligge mere for Anglosakser
og Germaner end for Franskmænd. Men der er i hvert Tilfælde
her en Opgave for nærmere Undersøgelse, naar engang Racepsykologien
har traadt sine Børnesko. 

51. Carlyle er den betydeligste Personlighed, som repræsenterer
den romantiske Humor. Grundlaget var for ham Livets store
Alvor, som hans tunge Sind aabnede hans Sans for, og som den
kalvinistiske Tro, han var opvokset i, havde indpræget i hans
Tænkemaade. Siden levede han sig ind i tyske Digteres og Filosofers
Tankeverden. Især har (som han selv i sine ældre Dage har
udtalt) Goethe’s „Wilhelm Meister“ paavirket ham, og her igen
Ideen om de tre Arter af Ærefrygt, der indskærpes i „Wanderjahre“:
Ærefrygten for det Store, det, som er over os, --- for det
Jevnbyrdige, det, som er ved vor Side, --- for det Smaa, det, som
staar under os. Ærefrygt danner Grundlaget for Carlyle’s Følelsesliv
og for hans Livsanskuelse. Hvad han tilegnede sig af tysk
Filosofi, var især den store Modsætning mellem Livets indre
Kærne og det ydre Hylster. Kant’s Modsætning mellem „Tingen
i sig selv“ og „Fænomenet“ anvendte han, ikke, som Kant selv,
til at afstikke Videnskabens Omraade og slaa sig ned der, men til
at udtrykke, at videnskabelig Forsken ikke fører til Noget, da den
ikke fører til Alt. Hele Erfaringen blev for ham til en Symbolik,
der forholdt sig til Sagen, som Klæderne til Kroppen. Hans Filosofi
var „Klædernes Filosofi“, med Hovedvægten paa, at Klæderne
ikke gøre Manden. (Sartor Resartus. 1836). Dette betinger hans
% --- p. 166 ---
\setcounter{footnote}{0}%
\opage{166}Forhold til Videnskabens Lære saavel som til Religionens Dogmer,
og til alle praktiske Bestræbelser. Alle disse Ting regner
han til Klæderne, eller til Lapperne, der sættes paa Klæderne.
Han staar overfor saadanne Bestræbelser med en Blanding af
Spot og Medlidenhed. Naar Carlyle saa hyppigt siger „stakkels“
om et Menneske, er det dels Sympati, dels Foragt, der udtrykkes
ved dette Ord. Paa drastisk Maade skildrer han det i hans Øjne
Meningsløse i at arbejde for Sandhed og Retfærdighed i Verden.
Og dog var hans Hjerte varmt. Og han holdt af den sunde, spontane
Latter, kunde selv le hjerteligt, og mente, at det Menneske,
som blot én Gang havde let af sit fulde Hjerte, kunde ikke være
helt slet (se § 13). Men hans Latter gik underligt ved Siden af
den ofte bitre Spøg eller brød igennem som en Kontrastvirkning
til denne. Carlyle er Humorist ved sin Sans for Modsætningen
mellem Kærne og Skal og ved sin Forstaaelse af, at Skallerne
skifte og maa skifte. Men hans Idealisme er, ligesom Romantikens
idethele, af den utaalmodige Art, der krævede Alt eller Intet, og
saa ved denne tilsyneladende modige Fordring i Virkeligheden
spærrede sig inde i Fantasiens og Stemningens Verden. Der manglede
et positivt Baand mellem Hovedtanken og dens specielle
teoretiske og praktiske Anvendelser, hvilke han derfor heller ikke
vilde anerkende som Anvendelser. Hans ubændige Temperament
fik ham til at haane, hvad han i Kraft af sit eget Princip maatte
have anerkendt som nødvendigt, skønt ufuldkomment Udtryk for
Realitet. Der ytrede sig her en hemmelig Tvivl hos ham.

Man har med Rette fremhævet, at det er det moderne
Livs komplicerede Forhold, som forvirre og forstemme ham. Hans
Tanke bevægede sig stedse om de simple, brede Kendsgerninger
i Livet. Da han fra sin Ensomhed i Højskotland blev hensat til
Edinburgh og London, traadte det civiliserede Liv ham imøde med
al sin Mangfoldighed i Forsken og Tragten. Som Rousseau\footnote[1]{Det var efter Carlyle’s eget Sigende Rousseau’s „Confessions“, som først havde vakt ham til dybere Eftertanke. (Ifølge hvad han har fortalt Emerson, der igen fortalte George Eliot det. Se \textit{George Eliot’s Life,} as related in her letters and journals. New edition. p. 104). Dette var mig ikke bekendt, da jeg skrev min Karakteristik af Carlyle i „Den nyere Filosofis Historie.“ (Niende Bog) og i „Mindre Arbejder“ III p. 125.} stod han
uforstaaende og forkastende overfor sin Tids Kultur. Men medens
% --- p. 167 ---
\setcounter{footnote}{0}%
\opage{167}Rousseau dog lærte at forstaa, at der gives en naturlig Kultur,
kom Carlyle aldrig ud over den skarpe Modsætning, som han i
enkelte Øjeblikke, under sit sygelige og heftige Sinds Indflydelse
udtrykte paa den voldsomste Maade. Og dog havde han i sin
Humor og i de Tanker, der dannede dens Baggrund, Betingelser
for at naa ud over denne Modsætning.

Den romantiske Humor naaede, derpaa er Carlyle det største
Exempel, ikke til den fulde Sammensmeltning eller Organisation.
De Modsætninger, der ikke kunde indgaa fuld, indre Forbindelse,
gjorde sig da gældende i successiv Kontrast. Det berettes om
Carlyle, at naar han paa sin barske Vis havde revet ned paa Alt
i Verden, kunde han pludseligt holde inde og gennem en kraftigt
frembrydende Latter vende tilbage til mere jevne Regioner. Hans
egne Profettirader maatte da komme til at staa for ham selv i et
komisk Lys.\footnote[1]{Henry James den Ældre (William James’ og Novelleforfatteren Henry James’ Fader) beskriver en kostelig Scene fra Carlyle’s Tebord. En Amerikaner disputerede med Carlyle om Politik. Carlyle blev voldsommere og voldsommere i sine Repliker, og Mrs. Carlyle traadte derfor Gæsten paa Foden for at faa ham til at holde op med det farlige Tema. Men Amerikaneren raabte højt: „Hvorfor træder De ikke paa Deres Mands Tæer, Mrs. Carlyle? Han er sikkert værre end jeg!“ Hele Selskabet brast i hjertelig Latter, og Carlyle førte an. (\textit{Henry James: Some personal recollections of Carlyle}. Literary Remains. Boston 1897. p. 454).}

52. Ifølge den psykologiske Karakteristik, der i det Foregaaende
er given af den store Humor, har den en afgjort realistisk Karakter
og er derfor bleven kaldt Troskab mod Virkeligheden. Naar
den, som det synes, sjeldent findes hos fremragende Filosofer,
maa Grunden, som ovenfor sagt, især søges i, at det bliver den
intellektuelle Interesse, om hvilken Alt samler sig for dem, og
som derfor komme til at bestemme den raadende Totalfølelse.
Ofte kan Grunden ligge i en utaalmodig Idealisme, der fører til
at betone Modsætningen til Erfaringen mere end Sammenhængen
med den. James Sully (Essay on Laughter. p. 397. 401), som
ogsaa har gjort den Iagttagelse, at Filosofer sjeldent ere Humorister,
finder Grunden i, at Filosofen af sine Studier fjernes for
% --- p. 168 ---
\setcounter{footnote}{0}%
\opage{168}meget fra Scenen for det virkelige Liv. Jeg vilde dog lægge
Hovedvægten paa Modsætningen mellem Tankelivet og de andre
Sider af Sjælelivet. Hos enhver Videnskabsmand vil der her
kunne fremtræde en større eller mindre Modsætning. Og selv om
Filosofen har særlig Anledning til at undersøge Forholdet mellem
Tanke og Personlighed, saa bliver det let selve denne Undersøgelse,
der lægger Beslag paa hele Interessen og Energien. Ogsaa
Tankelivet har sin Skæbne, medfører sine Oplevelser, der
kunne blive bestemmende for den sluttelige Stemning ved Livet.
Selv om det virkelige Liv griber ind med stærk Haand, kan det
da alligevel blive Tankens Arbejde og dette Arbejdes Held eller
Vanskæbne, der bliver det Afgørende. Saa vigtig den tankemæssige,
videnskabelige Baggrund for Livet end er, saa er dens Betydning
for en Totalfølelse som Humor dog kun indirekte. Og
som vi have set (§21) kan en ad rent praktisk Vej vunden
grundig Livserfaring have samme Betydning.

Der er i det Foregaaende skelnet mellem Baggrund og Grundlag.
Grundlaget, som især ligger i Følelses- og Villieslivet, er
vigtigere end den intellektuelle Baggrund. Det viser sig da ogsaa,
at man fra en Forskers almindelige Tankegang ikke kan slutte
til Karakteren af hans Totalfølelse ved Livet. Til at drage en
saadan Slutning udfordres Kendskab til det hele konkrete Grundlag
for Personligheden. Desværre er det saa vanskeligt at kende
det, og man maa i Regelen her bygge paa en Hypotese ud fra
den Totalfølelse, der faktisk har dannet sig, og til hvilken ogsaa
Baggrunden har givet sit Bidrag.

Maaske kunne store og stærke Totalfølelser snarere danne sig
hos Kunstnere end hos Tænkere. Der er et nærmere Forhold
mellem Billede og Følelse end mellem Tanke og Følelse. Uvilkaarligt
faar Følelsen ved det, som opleves, Udtryk i Billeder,
som paa anskuelig Maade fremstille eller symbolisere det Oplevede
i hele dets Ejendommelighed, medens for Tanken den enkelte
Oplevelse let kommer til at staa som et Exempel blandt
flere. Derved staar Kunstneren i et nærmere Forhold til Livet i
dets konkrete, i individuelle Former. Det er vel derfor, Schiller
skriver i et Brev til Goethe, at Digteren er det eneste sande
Menneske, medens den bedste Filosof kun er en Karikatur i
% --- p. 169 ---
\setcounter{footnote}{0}%
\opage{169}Forhold til ham\footnote[1]{Smlgn. min Mindetale om Schiller i \textit{Mindre Arbejder.} Tredie Række. p. 179 f.}. Dog kan ogsaa Kunstneren gøre Livet til blot
Emne for Betragtning (smlgn. § 37), saa at Fantasiens Verden
for ham bliver den eneste Verden. Schiller selv raadede en Ven
til at blive i „Ideernes lyse og stille Verden“ og foreløbigt
lade „den stakkels, uværdige og umodne Menneskehed“ sørge
for sig selv.

Hvor Energien anvendes overvejende til Livet i Tanke eller
Fantasi, vil der mangle Energi til Sammensmeltning eller Organisation
af de Følelser, Livets Gang kan vække. Saa kan der da ofte
i Regionerne nedenunder Tanke- og Billedverdenen tumle sig
mangfoldige Følelser paa kaotisk Maade. Hypokondri eller Letsind,
Flegma eller Opfarenhed, Smaalighed eller Højsindethed,
Bitterhed eller Sentimentalitet, Sanselighed eller Idealitet, ---
disse og mange flere Tendenser kunne danne en sublunarisk
Verden, som den, der kun kender de Tanke- eller Fantasiværker,
der stamme fra Æterverdenen foroven, ikke aner. Det er ikke
altid, at Tanke- eller Fantasilivet, foruden at være det Omraade,
hvor Energien først og fremmest sættes ind, tillige kan danne
Baggrund for Udviklingen af en Totalfølelse i stor Stil. Kun naar
det sker, ere vi paa Menneskelivets Højdepunkter.

Det er iøvrigt ikke blot Tænkere og Digtere, hos hvem denne
Modsætning mellem en Verden, i hvilke de leve deres mest energiske
Liv, og en sublunarisk Verden af helt anden Struktur kan
findes. Ogsaa hos praktiske Naturer, Handlingsmennesker, kan
der være en tilsvarende Modsætning, nemlig mellem den Del af
deres Leven, som gaar op i Arbejde, offentlig Virksomhed eller
Erhvervsvirksomhed, og de Sider i deres Sjæleliv, som ikke tages
i Beslag af denne Virken. ---

Man kunde (smlgn. § 38) ogsaa se Sagen fra den modsatte
Side og finde det betænkeligt, om Følelseslivet kunde forbindes
til fuldkommen Totalitet, da derved derved den Ensidighed bliver % PRINTED AS IS: „derved derved“
umulig, som alvorligt Arbejde og energisk Handlen kræver.

Gustav Frøding har i en lille Afhandling fra Aaret 1890 paa
en smuk Maade hævdet Humorens Betydning og Ret overfor
den ensidige Rolle som Sandhedsstræben og moralsk Tragten
% --- p. 170 ---
\setcounter{footnote}{0}%
\opage{170}efter hans Mening spillede i Tidens Literatur. Han mente, at der
var mere Humor i gamle Dage end nu, uagtet der nuomstunder
prædikes saa meget om Kærlighed. „Man är så kärleksful i ord
och gärning, att man alldeles glömmer att vara det i själ och
sinne, hvilket ändå är det, som verker försonande människor
mellan.“ „Humor är hundarne, som slicka såren på den lidande
Lazarus, och vänligare än den kristliga kärleken aktar den icke
för rof at skänka den rike mannen den vattendroppe, efter hvilkan
hans tunga trådde.“ --- Men da Frøding i Aaret 1910 optrykte
sin Afhandling (i hvilken han nærmest havde skildret Humor paa
Grundlag af Sympati), havde han tabt Troen paa den store Humors
Ret og mente, at han havde overvurderet den. Den bør
efter hans Mening ikke være eneherskende.

Jeg har i det Foregaaende heller ikke været blind overfor de
Betænkeligheder, den kan vække, idet jeg har fremdraget Farerne
ved enhver afsluttende Totalfølelse. Men gennem Undersøgelse
af Humorens Forhold til Forstaaelse, til Tragik og til Handling
har jeg søgt at vise, at der af Humorens egen Natur følger visse
Begrænsninger, der ingenlunde behøve at komme udefra.

Det var en genial Ide af Kierkegaard at opfatte den store Humor
som Højdepunktet af, hvad der kan naas ad rent menneskelig
Vej, i Modsætning til de religiøse Standpunkter, der forudsætte
et Forhold til Noget, som ligger ud over menneskelig
Erfarings og Videns Rækkevidde. Blot det, at der kan danne sig
en saadan Totalfølelse, som rummer Livets modsatte Elementer,
uden at udelukke Muligheden af, at disse Elementer, om nødvendigt,
kunne udløses paa ny i deres Særegenhed, --- blot det
er allerede af afgørende Betydning. Derved fremtræder et typisk
Exempel paa human Livsanskuelse.

53. Tilsidst ville vi betragte Humorens Forhold til en realistisk
Tankegang, hvorved jeg forstaar en Tankegang, der finder Sandheds-
og Virkelighedskriteriet i den faste, lovmæssige Sammenhæng
i Alt, hvad der fremtræder som Genstand for Iagttagelse.


Humor er, som det flere Gange er udtalt i det Foregaaende,
en realistisk Følelse, og forsaavidt kunde den synes at være helt
% --- p. 171 ---
\setcounter{footnote}{0}%
\opage{171}naturlig paa et udpræget realistisk Stade. Men det viser sig ogsaa
her, at Tankegangen kun er en Betingelse, ikke nødvendigvis
afgørende for en Totalfølelses Opstaaen. Baggrunden er ikke nok,
det kommer især an paa Grundlaget. Medens Idealistens Fare er,
at han let gaar op i de rene Ideer, uden at komme i nøje Forhold
til Livets konkrete Realiteter, vil Realistens Fare være, at
han tynges af Livets „Virkeligheder“ og ikke kommer til at gøre
sin Forstaaelse af dem til blot Baggrund for sit Følelsesliv. Baggrunden
kommer da til at hæmme Grundlagets Udfoldelse. Fortoner
Livet sig for Idealisten til en Æterverden, saa kvæles det
for Realisten let i „Virkeligheden“s tætte Taager. Saaledes er nu
engang Menneskets Lod, efter Goethe’s berømte Ord (Grenzen
der Menschheit): hæve vi os højt, kunne vi intet Fodfæste finde,
og staa vi fast paa Jorden, naa vi ikke højt op!

Hvis Realismen skal blive Baggrund for en fri og frisk Udvikling
af Følelseslivet, maa Realisten se kritisk paa sin egen
Realisme. Den samme Betingelse gælder for Idealisten. Kant, der
i sine „Träume eines Geistersehers“ havde vendt sig kritisk og
humoristisk mod en dogmatisk Idealisme, har i sin „Kritik der
reinen Vernunft“ banet Vejen for en kritisk Realisme.

Naar Realisten for Alvor spørger sig selv, hvad Virkelighed
(Realitet) egentligt er, vil han blive opmærksom paa, at naar vi
ere i Tvivl om, hvorvidt Noget er virkeligt eller ikke, gaa vi, i
det daglige Liv saavel som i Videnskaben, den Vej, at vi samle
saa mange Iagttagelser som muligt angaaende vedkommende
Emne og søge at finde saa nøje Sammenhæng som muligt mellem
disse Iagttagelser indbyrdes. Man experimenterer ikke blot i
Videnskaben, men ogsaa i det daglige Liv. Virkelighedskriteriet
opfordrer til stedse nyt Arbejde --- baade for at finde eller frembringe
nye Iagttagelser og for at finde nøjere Sammenhæng mellem
de fundne Iagttagelser. Begrebet Virkelighed staar som et
Ideal, vi kunne nærme os, men som stedse stiller nye Fordringer.
Væbnet med dette Ideal skyder den store Humor al dogmatisk
Realisme tilside. Ingen foreliggende Virkelighed tilfredsstiller
Virkelighedsbegrebet. (Smlgn. § 34).

Indenfor en saadan Betragtning af Tilværelsen, som dette paa
én Gang experimentale og ideale Virkelighedsbegreb gør mulig,
% --- p. 172 ---
\setcounter{footnote}{0}%
\opage{172}kan man forene Interesse for og Forstaaelse af det Enkelte i dets
Ejendommelighed og dets Værdi med Henblikket til en større
Helhed. Saa forsvindende det Enkelte end kan synes at være,
saa bliver dets Værdi dog aldrig Nul. Og selve den større Helhed,
af hvilken det er et Led, viser sig igen forsvindende overfor
mere omfattende Helheder. Det kan maaske forudses, at ogsaa
dens Tid engang vil være forbi, saaledes som Shakespeare’s
Prospero aner det, og som moderne Forskere ofte have antaget
det, hvad vort Verdenssystem angaar. Man behøver overhovedet
ikke at være Optimist for at være Humorist; tvertimod vil Optimisme
(smlgn. § 48) udelukke Humor ved at hindre en tilstrækkeligt
grundig Indtrængen i alle Sider af Livet. Humoristen vil ikke
høre til dem, der glæde sig over, hvor herligt vidt vi have drevet
det; saadan Glæde vil snarest vække hans Ironi. Han véd,
at Kampen mellem det Menneskelige og det Dyriske er lang;
den gaar, ligesom en moderne Krig, fra Skyttegrav til Skyttegrav.
Men saalænge der endnu kæmpes, ere Haab og Tillid baade
mulige og nødvendige. For Charles Darwin var denne Langsomhed
i Udviklingen ikke den største Anstødssten. Hvad der stod
for ham som en Vanskelighed var, at naar Udviklingen var
skreden frem gennem uhyre Tidsperioder, saa skulde Alt igen
opløses i et Kaos. Han skriver til Hooker (9. Febr. 1865): „Jeg
er ganske enig i, hvor ydmygende Menneskets langsomme Fremskridt
er; men Enhver har sin egen Kæphest, hvad Frygt angaar
(his own pet horror), og det langsomme Fremskridt, ja endog
den personlige Tilintetgørelse synker i mine Øjne ned til en Ubetydelighed
i Sammenligning med Forestillingen eller rettere Visheden
om, at Solen engang vil blive kold, og vi alle fryse ihjel.
At tænke sig en Fremskriden gennem Millioner af Aar, med enhver
Verdensdel mylrende af gode og oplyste Mennesker, og Alt
ende saaledes, uden at der sandsynligvis kommer nogen ny Begyndelse,
før vort Planetsystem igen er blevet til rødglødende
Gas! Sic transit gloria mundi --- med en Nemesis (with a vengeance).“
--- Darwin opgav dog ikke sin Tro paa, at Udviklingen
i Verden havde Overvægt over Opløsningen. Han fandt, ligesom
Kant, det faste Punkt i Pligtfølelsen. „Sin Pligt kan Mennesket
gøre.“ Ud fra Muligheden af at virkeliggøre noget Værdifuldt
% --- p. 173 ---
\setcounter{footnote}{0}%
\opage{173}straalede tilsidst det Lys, i hvilket det blev ham muligt at leve
og stræbe\footnote[1]{Se min Afhandling om \textit{Darwin og Filosofien} i det i Cambridge udgivne Festskrift „Darwin and Modern Science“, paa Dansk i „Mindre Arbejder“ Tredie Række. (p. 225—228).}.

Skulde det forholde sig saaledes, at Værdiernes Kilde udtømtes
engang, fordi Alt tilsidst kom i absolut Ligevægt, saa vil, som
allerede sagt (§ 30), Humoristen ikke kunne tage dette humoristisk.
Saa indrømmer han, at Tragiken faar det sidste Ord. Ja,
han vilde ønske, at den kunde faa det sidste Ord. Ti det vil desværre
gaa saaledes, at der forud for det psykiske Livs Udslukning
vil gaa en Svækkelses- og Sløvhedsperiode, hvor der end
ikke bliver Tale om tragisk Lidelse. Det Tragiske forudsætter
Kraft. Men Humoristen arbejder, saalænge han kan, for at finde
nye Værdier, og han stiller sig kritisk overfor alle Teorier, Spekulationer
og Perspektiver, der uden tilstrækkeligt Grundlag dogmatisere
om „de sidste Ting“. Fysiken véd i Grunden ikke Mere
om dem end Teologien. Ligesaalidt som der er Grund til, maaske
ikke engang Mulighed for at tænke sig en Afslutning af Tiden
og Rummet, ligesaalidt kan der nogensinde føres Bevis for, at
alle Muligheder ere udtømte, „alle Sunde lukte“. I Kraft af
Uoverskueligheden af den intellektuelle Baggrund kan Tilværelsen
stedse komme til at staa som ny.

54. Ud fra Realismen kommer man ikke med Nødvendighed
til Humor. Men kritisk Realisme kan maaske bedre end de fleste
idealistiske Retninger afgive Baggrund for Humor.

En nødvendig Bestanddel af en saadan kritisk Realisme er
Anerkendelsen af den Kendsgerning, at enhver Personlighed i
større eller mindre Grad og med større eller mindre Bevidsthed
former og maa forme sin egen Livsanskuelse, som man ikke
uden videre kan udlede af de positive eller negative Dogmer, til
hvilke de Enkelte bekende sig. Netop fra et realistisk Standpunkt
maa det kræves, at Enhver bygger sin Livsanskuelse paa sine
egne Oplevelser og ved Hjælp af den Erkendelse, han har vunden
ved sin egen Tankes Arbejde. Kun saa bygger den paa
Virkelighedens Grund, paa det, der virkeligt er Virkelighed for
% --- p. 174 ---
\setcounter{footnote}{0}%
\opage{174}den Enkelte. Livsanskuelse er kun noget værd, naar den virkeligt
er Udtryk for, hvorledes Livet rører sig paa det bestemte Sted i
Tilværelsen, den Enkelte indtager. Fordringen om personlig Sandhed
er realistisk, skønt den først er fremsat af en Idealist.

Erfaring viser mere og mere, hvor intetsigende dogmatiske
Erklæringer kunne være i Forhold til det indre Livs Realiteter.
Fra hvad der bekendes, synges og forkyndes, kan man ikke drage
nogen sikker Slutning, naar man ikke ad anden Vej faar Adgang
til det indre Liv, der ligger bagved. Bag Teologien eller Antiteologien
søger Realisten, og Humoristen med ham, Psykologien.
Han véd, at der i Sindets indre Verden i al Stilhed kan foregaa
store Ting, og de larmende Stridigheder om Registret og Overskrifterne
til Livets Bog lokke kun et Smil frem hos ham. Personlighedsprincipet
virker ikke blot opløsende overfor alskens
Dogmatisme og alskens overfladisk Fællesskab; men det udtrykker
det indre Livs Uudtømmelighed.\footnote[1]{Se \textit{Religionsfilosofi}\textsuperscript{2}. p. 284—286.} Der aabner sig her en indre
Verden, som hører med til den intellektuelle Baggrund for
Følelseslivet. I hver enkelt Personlighed kan denne indre Verden
kun finde Udtryk paa en mere eller mindre ensidig Maade. Det
gaar med Livets Kunst ligesom med Billedkunsten, der, efter
hvad Julius Lange saa energisk har hævdet, maa være ensidig,
fordi „intet enkelt Kunstværk kan meddele den alsidige Opfattelse
af Menneskeskikkelsen“.\footnote[2]{Se \textit{Religionsfilosofi}\textsuperscript{2}. p. 256.}

Ved at vedkende sig denne Betragtning vedkender den store
Humor sig tillige sin egen Begrænsning. Selv om den skulde
have nok saa stor Betydning som en karakteristisk Totalfølelse,
i hvilken alt Menneskeligt kan komme til sin Ret, saa vil den
dog aldrig kunne forkynde sig selv som den eneste rette Maade
at tage Livet paa. Den forholder sig her humoristisk til sig selv.
Paa Virkelighedens Grund ser den Menneskelivet udfolde sig
fra højst forskellige Udgangspunkter og under højst forskellige
Brydningsforhold, og den forstaar, at al Udvikling foregaar sporadisk.
Derfor undrer den sig ikke over, at Totalfølelserne blive
af forskellig Karakter. Det er ikke underligt, at ikke Alle ere
Humorister. Og selv om de vare det, vilde de være det hver paa
% --- p. 175 ---
\setcounter{footnote}{0}%
\opage{175}sin Maade. Harmoni og Fællesskab ere store Goder, men maa
ikke købes med Opofrelse af Oprindelighed og Ægthed hos de
Enkelte. Udviklingens sporadiske Karakter gør Aandens Verden
paa én Gang stor og lille, tragisk og komisk, --- og dette indeholder
(smlgn. § 35) en Anledning til Opstaaen af den Totalfølelse,
vi her have beskæftiget os med.
% --- p. 177 ---
\opage{177}\begin{center}REGISTER\end{center}
\entry \textbf{A}cedia 96.
\entry Alvor 68 ff. 125.
\entry Amiel 131.
\entry Antigone 100. 129.
\entry Aristoteles 101 f. 129.
\par\smallskip
\entry \textbf{B}aggrund (for Humor) 86. 113 f. 145 ff. 166. 171.
\entry Benn 102.
\entry Beslutning 10. 19. kfr. 123 fi. % PRINTED AS IS: „fi.“ (sense: ff.)
\entry Bradley (A. C.) 98 f. 117.
\entry Brandes (G.) 100.
\entry Bruno 39. 143. 154 f.
\entry Burton 41.
\par\smallskip
\entry \textbf{C}arlyle 34. 55. \textit{165---167}
\entry Cervantes 43.
\entry Comte 114.
\entry Cusanus 121.
\par\smallskip
\entry \textbf{D}ante 96---98.
\entry Darwin 172.
\entry Descartes 9.
\entry Differentiation 125.
\entry Dowden 43. 67.
\par\smallskip
\entry \textbf{E}nfødte 13. 140.
\entry Enkeltfølelse \textit{7---12}. % the print: italic „7—12“ (the layer: „1.—12.“)
\entry Etik. (og Totalfølelse) 36 f. 125 f.
\par\smallskip
\entry \textbf{F}allstaff 116 f.
\entry Feilberg (Ludvig) 10. 12.
\entry Felix culpa 80. 83. 124.
\entry Filosofi og Livsanskuelse 45. \textit{145---175}.
\entry Freud 102. kfr. 13.
\entry Frøding 169 f.
\par\smallskip
\entry \textbf{G}alilei 120.
\entry Genetisk Psykologi 16. 126.
\entry George Eliot 54. 82. 140.
\entry Goethe 55. 57. 92 f. 100. 127. 171.
\entry Goldschmidt 131.
\entry Grundlag (for Humor) 86. 113 f. 166. 171.
\entry Gunnarson (G.) 104.
\entry Fru Gyllembourg 83.
\par\smallskip
\entry \textbf{H}aan 52. 59.
\entry Handling 125 ff.
\entry Hegel 100 f. 137.
\entry Heroisk Lidenskab 39.
\entry Hirn (Yrjo) 68. 73.
\entry Hobbes 48. 55.
\entry Holberg 42. 120.
\entry Home (Henry) 46.
\entry Hume 158.
\entry Humor. Ordets Betydning 40—45. Arter 52 ff. \textit{75---87}. Den store Humor 43 f. 47. 54. \textit{56---58}. Humor og Ironi 68 f. Humor og Virkelighed 70 ff. Humor som Livskunst 92 ff. Humor og Filosofi 145 ff.
\entry Hypokondri 95 ff. 169.
\par\smallskip
\entry \textbf{I}figencia 120.
\entry Ironi 61 ff. 147 ff.
\par\smallskip
\entry \textbf{J}ames (W.) 28.
\entry Jean Paul 47. 141. 161. 164.
\entry Jeg (realt) 11---15. 23. 27. 35. 54.
\entry Jonson (Ben) 41.
\par\smallskip
\entry \textbf{K}aalund 57.
\entry Kamm (Adèle) 90.
\entry Kant 34. 76. \textit{169 f.}
\entry Kierkegaard 21. 34. 47. 64. 74. 89. 93. \textit{104---109}. \textit{132---134}. 147.

% --- p. 178 ---
\opage{178}\entry Kristendom (Urkristendom) 81 f. 105.
\entry Kultus 141.
\par\smallskip
\entry \textbf{L}ange (Julius) 27. 174.
\entry Latter 48 ff.
\entry Lehmann (Alfr.) 22. 29.
\entry Leopardi 119.
\entry Lessing 89.
\entry Lichtenberg 55.
\entry Lotze 58. 65.
\entry Længsel 53. \textit{80---83}.
\par\smallskip
\entry \textbf{M}aier (Heinrich) 148---150.
\entry Medeia 129.
\entry Mefistofeles 117.
\entry Mozart 9.
\par\smallskip
\entry \textbf{N}ietzsche 59.
\entry Novalis 163.
\par\smallskip
\entry \textbf{O}ndt (det Ondes Udspring) 80.
\entry Optimisme 140.
\entry Organisation (af Følelser) 18. 26. \textit{30---33}.
\par\smallskip
\entry \textbf{P}ersonlighed 13---102.
\entry Personlighedsprincipet 64.
\entry Pessimisme 140.
\entry Platon 15. 26. 38. 55. 136 f. \textit{151---153}.
\entry Poincaré (Henri) 37.
\entry Prospero 99. 107. 162.
\entry Psykoanalyse 13. 102.
\entry Psykologi 34. 76.
\par\smallskip
\entry \textbf{R}ealt Jeg, se Jeg.
\entry Religion 108ff. kfr. 104. 150.
\entry Ribot 26. 49.
\entry Rousseau 10. 166.
\entry Ræder (Hans) 152.
\par\smallskip
\entry \textbf{S}ammensmeltning (af Følelser 18. 23. \textit{26---30}. % PRINTED AS IS: the parenthesis not closed
\entry Sandhed 115. 117.
\entry Satire 60 f.
\entry Schlegel (Fr.) 65 ff. 69. 163.
\entry Schlegel (A. W.) 67.
\entry Shaftesbury 157f.
\entry Shakespeare 40. 42 f. 67 f. 85. \textit{97---100}. 116 f. 129. 141. 147. 159 f.
\entry Shand 26. 31. 36.
\entry Sibbern 25. 29. 38. 73. 87. 121. 145. 147.
\entry Sidney Lee 100. 116.
\entry Sindelag og Sindsbevægelse 128 ff.
\entry Sindshøjhed (Højsind) 03. \textit{77---79}. % PRINTED AS IS: „03.“
\entry Smerte 91.
\entry Sneedorff 120.
\entry Sokrates 21. 61 f. 138. \textit{146---151}.
\entry Solger 47. 73—93. 163.
\entry Spinoza 19. 97. 113. \textit{155---157}.
\entry Sporadisk Udvikling 75.
\entry Sterne 93. 189. % PRINTED AS IS: „189“ (the book has 178 pages)
\entry Stoicisme 79.
\entry Stuckenberg 83.
\entry Suddard (Mary) 40.
\entry Sully (James) 49. 51. 54. 167.
\entry Swift 60. 159.
\entry Sympati 58. 84 f.
\par\smallskip
\entry \textbf{T}emperament 139.
\entry Thackeray 54. 158 f.
\entry Thorndike 16. 126.
\entry Totalfølelse \textit{18---26}. 76. 76. 77. 128. 145—148.
\entry Totaltilstand 1—25.
\entry Tragik 89 ff. 173.
\entry Troskab 43. 115f.
\entry Tvefødte 14. 140.
\par\smallskip
\entry \textbf{V}elasquez 27.
\entry Vemod 14. 30. 53. \textit{80---83}.
\par\smallskip
\entry \textbf{W}ilhelm Meister 89.
\par\smallskip
\entry \textbf{V}irkelighedskriterium 114 f. 160. 171.
\entry Værdi \textit{68---73}. 118. Værdiens Bestaaen 99. 106. 109f. Værdi og Virkelighed 80 f. 118 f.

\end{document}
